% !TeX root = ../Book2.tex
% 纲要标题：# 第四编　Morita 等价、中心单代数与非交换环
% 纲要位置：计划纲要.md，第 1560 行。

\BookPart{Morita 等价、中心单代数与非交换环}{b2:part:04}
\BookPartGuide
两个环可以有不同的元素和乘法表，却具有等价的模范畴；一个代数也可以在扩大底域后成为矩阵代数，尽管它在原域上不分裂。第16章比较一般环的模范畴，第17、18章再研究有限维中心单代数的分裂与稳定分类。第19、20章不再以有限维或 Artin 条件作为共同背景，而转向更广的非交换环与代数：第19章分别研究满足 Ore 条件\footnote{奥尔（Øystein Ore，1899-1968），挪威裔美国数学家，研究非交换环及分式理论。}时的分式、对所有代入成立的多项式恒等式，以及生成空间的增长；第20章则用滤过、Ore 扩张和重写系统研究正规形。

\ProseParagraphBreak
\chapref{b2:ch:16}从有限生成投射生成元出发建立 Morita 理论。通过模展示和余核，张量函子与 Hom 函子的互逆关系从自由模扩展到任意模及其同态。矩阵环和满幂等元的角环提供了显式的等价，而中心、简单模、投射性和长度说明这种比较保留哪些结构。向量空间维数及最少模生成元数则可以改变。

\ProseParagraphBreak
\chapref{b2:ch:17}把底域固定为代数的中心，证明中心单性在任意域扩张下保持，建立中心化子与 Skolem--Noether 定理\footnote{斯科伦（Thoralf Albert Skolem，1887-1963），挪威数学家及逻辑学家，研究数理逻辑、集合论和代数。Noether 已在本卷前文介绍。}。中心化子定理对简单子代数不额外假设其中心在底域上可分；构造有限可分分裂域时，则主动在相应的除代数中寻找可分最大子域，再取正规闭包得到有限 Galois 分裂域。

\ProseParagraphBreak
\chapref{b2:ch:18}把矩阵大小略去，组成 Brauer 群。循环代数将分裂问题转成范数方程，约化特征多项式则由分裂模型下降到底域。随后以交叉积和\(2\)-余圈证明相对 Brauer 群的对应，并证明周期整除指数。周期为二、指数为四的具体除代数例子提醒读者，两种数值记录不同的信息。

\begin{center}
\begin{tikzcd}[column sep=large,row sep=large]
&\begin{gathered}\text{第16章}\\\text{模范畴的等价}\end{gathered}&\\
\begin{gathered}\text{已有基础}\\\text{模、张量与环结构}\end{gathered}\arrow[ur]\arrow[r]\arrow[dr]
&\begin{gathered}\text{第17章}\\\text{中心单代数}\end{gathered}\arrow[r]
&\begin{gathered}\text{第18章}\\\text{Brauer群与范数}\end{gathered}\\
&\begin{gathered}\text{第19、20章}\\\text{分式、恒等式与增长}\\\text{滤过与正规形}\end{gathered}&
\end{tikzcd}
\end{center}

第16章直接使用第五章的范畴等价、第六章的伴随及第七章的投射模，第十五章的基本角代数是它的一项应用。第17章主要使用第十一章的双模与中心化子、第十三章的矩阵环分解，以及第十四章的稠密性。第17、18章都需要第一卷的有限扩张、可分性与有限 Galois 理论，第18章进一步使用域范数、本原元素及中国剩余定理。Brauer 稳定等价通过中心除代数代表与 Morita 理论相联系，其群运算则由代数的张量积给出。

\ProseParagraphBreak
第十四章的根与 Nakayama 引理接到第十五章的投射覆盖，再进入第十六章；第十四章的简单模与稠密性则接到第十七章的包络代数判据，再进入第十八章。

\ProseParagraphBreak
\chapref{b2:ch:19}从分母不能随意移动的困难出发，构造右 Ore 分式环，并说明 Goldie 条件\footnote{戈尔迪（Alfred William Goldie，1920-2005），英国数学家，研究非交换环，给出半素环具有半单分式环的判据。}如何保证半单分式环存在。多项式恒等式另从自由代数出发；Amitsur--Levitzki 恒等式\footnote{阿米楚尔（\OriginalHebrewName{שמשון עמיצור}，1921-1994），以色列数学家，研究环与多项式恒等式。列维茨基（\OriginalHebrewName{יעקב לויצקי}，1904-1956），以色列数学家，研究环的根及矩阵环的恒等式。}联系第九章的外代数与 Cayley--Hamilton 定理。增长及 Gelfand--Kirillov 维数\footnote{盖尔范德（\OriginalName{Израиль Моисеевич Гельфанд}，1913-2009），苏联及美国数学家，研究泛函分析和表示论。基里洛夫（\OriginalName{Александр Александрович Кириллов}，1936-），俄罗斯数学家，研究 Lie 群的表示与轨道方法。}则比较生成空间的幂怎样变大，与普通的有限维数有不同含义。

\ProseParagraphBreak
\chapref{b2:ch:20}用具体构造证明 Weyl 代数、量子平面和 Heisenberg 包络代数中心商的正规基，并通过 Ore 扩张继续研究正规式计算。滤过可以把某些非交换问题化为伴随分次中的问题，但每次提升都需要穷尽性和降次归纳的下界。最后的 Diamond 引理说明重写规则何时给出唯一正规字，也给出在指定字序下有限关系却需要无限 Gröbner 基\footnote{格罗布纳（Wolfgang Gröbner，1899-1980），奥地利数学家，研究代数几何；这一计算工具由其学生提出并以他命名。}的例子。一般 Lie 代数\footnote{李（Marius Sophus Lie，1842-1899），挪威数学家，创立连续变换群理论，推动微分方程与几何的研究。}的 PBW 结论在本章用于结构比较，完整证明见第五卷第15章的第15.2节。

\ProseParagraphBreak
第二十章的 Rees 代数\footnote{里斯（David Rees，1918-2013），威尔士数学家，研究半群与交换代数。}把原代数与其伴随分次放进同一个带参数的代数中：令参数为零得到伴随分次，令参数为一得到原代数。这两个特殊纤维的同构直接由构造证明；平坦性则使用第七章的主理想整环上无挠模平坦定理。一般群上同调、路径与表示的进一步应用、Lie 代数包络代数的系统理论，分别在后续各卷展开。

% !TeX root = ../../Book2.tex
% 纲要标题：## 第16章　Morita 理论
% 纲要位置：计划纲要.md，第 1562 行。

\chapter{Morita 理论}
\label{b2:ch:16}
\BookChapterNumber{b2:ch:16}

\begin{chapterabstract}
本章证明模范畴的等价如何由有限生成投射生成元实现，并核对张量函子、Hom 函子与自同态反环的方向。矩阵环和满幂等角环给出具体等价，进一步辨认等价所保持的模论性质及环本身可能改变的性质。
\end{chapterabstract}

\begin{learninggoals}
\begin{itemize}
\item 区分作为对象的生成元、投射性与有限生成性，证明生成元的三种等价刻画。
\item 在左模约定下正确使用 \(S=\operatorname{End}_R(P)^{\mathrm{op}}\)，写出伴随单位、余单位、互逆双模和有限对偶基公式。
\item 从给定范畴等价恢复有限生成投射生成元，并用任意自由展示证明反向的 Morita 等价。
\item 计算矩阵环与满角环的等价，辨认所保留的模结构及可以改变的坐标、维数和生成元数。
\end{itemize}
\end{learninggoals}

设 \(k\) 为域、\(n\ge1\)，令 \(A=M_n(k)\)，并在一个左 \(A\) 模 \(M\) 上取矩阵单位 \(E_{ij}\)。元素 \(m\) 可拆成 \(\sum_iE_{ii}m\)，而 \(E_{ij}\) 把第 \(j\) 个分量送往第 \(i\) 个分量。知道 \(E_{11}M\) 后，其余分量似乎都能由这些矩阵单位重新造出；模同态也由它在这一分量上的作用决定。这说明环的元素和乘法不同，并不妨碍它们的全部模具有相同的组织方式。

\ProseParagraphBreak

从一个矩阵角回收整个模，需要说明重建映射在任意模上都互逆，而且对模同态相容。一般环未必有现成的矩阵单位，这时要问是否存在一个模 \(P\)，从它出发仍能恢复整个模范畴。生成元性质要求每个模都是 \(P\) 某个直和的商；投射性使 \(\operatorname{Hom}_R(P,-)\) 保持正合列；有限生成性使这个函子保持任意直和。三者各自保护一项重建所需的操作，由此构造的 Hom 与张量函子才可能在全部模上互为拟逆。

\ProseParagraphBreak

本章需要第五章的范畴等价、第六章的双模张量与 Hom–张量伴随，以及第七章的投射模和有限对偶基。第十一章的角环自同态公式与第十五章的基本角代数提供具体应用；中心和有限长度的讨论还使用前两编已建立的自然性与合成列。等价始终指全部左模范畴的等价，因而还须构造在一般模上的作用，有限自由模上的检验并不充分。

\ProseParagraphBreak
一般环的结构定理及矩阵环例子可参阅 Weibel 的《The K-book》\cite[Chapter II, Theorem 2.7, Example 2.7.2]{Weibel2013KBook}。该书采用右模，本章固定左模并相应使用自同态反环。有限维代数的版本可参阅 Etingof 等人的《Introduction to Representation Theory》\cite[Theorem 9.6.4, Problem 9.6.5, p. 215; Section 9.7, p. 216]{EtingofEtAl2011RepresentationTheory}；那里的有限维模范畴与本章全部模范畴的范围须加以区分。

% !TeX root = ../../Book2.tex
% 纲要标题：### 16.1 模范畴与生成元
% 纲要位置：计划纲要.md，第 1564 行。

\section{模范畴与生成元}
\label{b2:sec:1601}

上一编把矩阵环 \(M_n(D)\) 的模分解为简单列模的直和。当 \(n>1\) 时，一个列模比正则模小，仍可用它的直和与商构造其他模。对于一般环，也要问哪些模能够这样替代正则模，产生整个模范畴。本节把这种性质与投射性、有限生成性分开说明。始终使用左模；\(R\text{-}\mathbf{Mod}\) 包含全部幺性左 \(R\) 模，而不只包含有限生成模。后面的等价也将在全部模上成立。

% 纲要要点（供目录核对）：
%
% * 投射生成元；区分模忠实、非零对象检测与同态检测，正式证明迹理想判据及对象检测反例
% * 有限生成投射生成元；用两组有限直和因子分别控制有限投射与生成，正式给出三条件各自缺失的反例
% * 矩阵环与原环的模范畴；矩阵单位读取与恢复分量，正则模未必与正则模对应
%

\subsection{投射生成元}
\label{b2:subsec:1601:01}

正则模通过自由模的满射产生所有模。若想用一个模 \(P\) 替代它，最直接的要求就是每个模都能由若干份 \(P\) 的直和取商得到。“生成元”的名称表达了这个要求；下面用 Hom 函子的忠实性定义它，并证明这两种说法等价。

\begin{definition}\label{b2:def:module-generator}
设 \(R\) 为含幺环，\(P\) 为幺左 \(R\) 模。若函子 \(\operatorname{Hom}_R(P,-)\) 忠实，即对任意幺左 \(R\) 模 \(M,N\) 及不同的 \(R\) 线性映射 \(f,g:M\rightarrow N\)，都存在 \(R\) 线性映射 \(h:P\rightarrow M\) 使 \(fh\ne gh\)，则称 \(P\) 为左 \(R\) 模范畴的\term[mo-fanchou-shengchengyuan]{生成元}[generator]。若 \(P\) 还是投射模，则称它为\term[toushe-shengchengyuan]{投射生成元}[projective generator]。
\end{definition}

这里的“生成元”是一个模对象，不是某个模内部的一个元素。条件 \(M\ne0\Rightarrow\operatorname{Hom}_R(P,M)\ne0\) 只检测非零对象，定义中的忠实性则比较不同同态；二者的差别见\cref{b2:prop:object-detection-not-generator}的双数例子。正则模 \(R\) 是生成元：若 \(f(m)\ne g(m)\)，取 \(h(r)=rm\)，在一处求值便可检测两同态不同。

\ProseParagraphBreak
每个模 \(M\) 本来就有自由满射 \(R^{(I)}\to M\)，所以只要 \(P\) 的直和能满射到 \(R\)，便能对自由模的各个分量作同样构造，再满射到 \(M\)。而 \(R\) 由一生成，一的原像又会把这个构造缩到有限份 \(P\)；这就是下面第三个条件的来源。

\begin{theorem}[生成元的等价刻画]\label{b2:thm:generator-characterizations}
设 \(R\) 为含幺环，\(P\) 为幺左 \(R\) 模，所论模均为幺左 \(R\) 模。以下条件等价：
\begin{equivalences}
\item\label{b2:item:generator-faithful} \(\operatorname{Hom}_R(P,-)\) 忠实。
\item\label{b2:item:generator-quotients} 对每个左 \(R\) 模 \(M\)，都存在指标集 \(I\) 及满同态 \(P^{(I)}\rightarrow M\)。
\item\label{b2:item:generator-finite-summand} 存在有限整数 \(n\geq0\)，使正则左模 \({}_RR\) 是 \(P^n\) 的直和因子。
\end{equivalences}
这些等价不要求 \(P\) 已经投射或有限生成。
\end{theorem}
\begin{proof}
先设 Hom 函子忠实。对任意 \(M\)，将全部同态 \(h:P\rightarrow M\) 作为指标，定义求值映射
\[
q:\bigoplus_{h\in\operatorname{Hom}_R(P,M)}P\longrightarrow M,
\qquad (p_h)_h\longmapsto\sum_h h(p_h).
\]
若像 \(T\) 不等于 \(M\)，商映射 \(\pi:M\rightarrow M/T\) 非零，却满足 \(\pi h=0\) 对每个 \(h\) 成立。于是 \(\operatorname{Hom}_R(P,\pi)\) 与零同态诱导的映射相同，违反忠实性。因此 \(q\) 满射，得到第二项。

将第二项用于 \(M=R\)。一的某个原像只有有限个非零坐标，故限制到相应的 \(P^n\) 后，映射 \(q_0:P^n\rightarrow R\) 的像仍包含一，因而包含每个 \(r=r1\)，所以仍满射。这里有限的 \(n\) 来自直和元素的有限支集，不要求 \(P\) 有限生成。取 \(v\in P^n\) 使 \(q_0(v)=1\)，则 \(s(r)=rv\) 为左线性截面，\(q_0s=\operatorname{id}_R\)。所以 \(R\) 是 \(P^n\) 的直和因子。

若第三项成立，任意模 \(M\) 的自由满射 \(R^{(I)}\rightarrow M\) 可以与各分量上的分裂满射 \(P^n\rightarrow R\) 复合，得到某个 \(P\) 直和到 \(M\) 的满射，所以第二项成立。最后，若 \(f\ne g:M\rightarrow N\)，而每个 \(h:P\rightarrow M\) 都满足 \(fh=gh\)，则对第二项给出的满射 \(P^{(I)}\rightarrow M\)，两复合在每个直和分量上相同，因而整体相同。满射性迫使 \(f=g\)，矛盾。故 Hom 函子忠实。
\end{proof}

生成元必为忠实模，即 \(\operatorname{Ann}_R(P)=0\)：若 \(rP=0\)，第三项使 \(r\) 也零化直和因子 \(R\)，故 \(r=r1=0\)。但模的零化子为零，弱于 Hom 函子忠实。例如 \(\mathbb Q\) 作为 \(\mathbb Z\) 模忠实，但 \(\operatorname{Hom}_{\mathbb Z}(\mathbb Q,\mathbb Z)=0\)：任意同态的每个值都被所有正整数整除，只能为零。因此任何 \(\mathbb Q\) 的直和都不能满射到 \(\mathbb Z\)，它不是生成元。

\subsection{有限生成投射生成元}
\label{b2:subsec:1601:02}

\begin{definition}\label{b2:def:progenerator}
设 \(R\) 为含幺环，\(P\) 为幺左 \(R\) 模。若 \(P\) 有限生成、投射，且是左 \(R\) 模范畴的生成元，则称 \(P\) 为\term[youxian-toushe-shengchengyuan]{有限生成投射生成元}[progenerator]。投射性指沿任意满同态的提升性质，生成元性质采用\cref{b2:thm:generator-characterizations}的任一等价条件。
\end{definition}

\cref{b2:prop:finite-projective-summand}与生成元判据合在一起说明，这等价于同时存在有限整数 \(m,n\) 及模 \(Q,Q'\)，使
\[
P\oplus Q\cong R^m,\qquad R\oplus Q'\cong P^n.
\]
第一个分解把 \(P\) 放在有限自由模中作为直和因子，控制其有限生成与投射性；第二个分解把正则模放在有限份 \(P\) 中，保证它能够产生全部模。这两类分解不能互相替代，\cref{b2:prop:progenerator-hypothesis-counterexamples}将分别检验三个条件缺失时的情形。

\ProseParagraphBreak
有限生成性用一组元素来表述，但范畴等价并不指定两个对象底集合的元素如何对应。要在等价下追踪投射模的有限生成性，需要把它改写成能由 Hom 和直和识别的条件：从这个模到任意直和的同态，是否总能只用有限个目标分量表达？

\begin{definition}\label{b2:def:compact-module-object}
设 \(R\) 为含幺环，\(P\) 为幺左 \(R\) 模。对任意指标集 \(I\) 及幺左 \(R\) 模族 \((M_i)_{i\in I}\)，若由各分量包含映射诱导的自然映射
\[
\bigoplus_{i\in I}\operatorname{Hom}_R(P,M_i)
\longrightarrow\operatorname{Hom}_R\left(P,\bigoplus_{i\in I}M_i\right)
\]
均为双射，则称 \(P\) 为左 \(R\) 模范畴中的\term[jin-duixiang]{紧对象}[compact object]。此处专指 Hom 保持任意直和的范畴条件，不指拓扑上的紧性，也不同于以 Hom 保持滤过余极限刻画的有限展示性。
\end{definition}

\begin{proposition}\label{b2:prop:compact-projective-finite}
设 \(R\) 为含幺环，所论模均为幺左 \(R\) 模。每个有限生成左 \(R\) 模都是左 \(R\) 模范畴中的紧对象。若左 \(R\) 模 \(P\) 还投射，则 \(P\) 有限生成，当且仅当 \(P\) 是紧对象。
\end{proposition}
\begin{proof}
自然映射始终单射：若有限支集族 \((f_i)_i\) 合成到直和中的同态为零，再接上第 \(j\) 个坐标投影，便得到 \(f_j=0\)，对每个 \(j\) 都成立。若 \(P\) 有有限生成组，任意 \(f:P\rightarrow\bigoplus_iM_i\) 在这些生成元上的像共同只有有限支集，故全像落在一个有限子直和中。于是 \(f\) 是有限多个分量同态之和，证明满射。

反过来，设 \(P\) 投射且紧。由\cref{b2:thm:projective-characterizations}，可取分裂自由满射 \(q:R^{(I)}\rightarrow P\) 及截面 \(s:P\rightarrow R^{(I)}\)。紧性使 \(s\) 的像落在某个有限子直和 \(R^J\) 中。限制 \(q\) 到 \(R^J\) 后仍满足 \(qs=\operatorname{id}_P\)，所以 \(P\) 为有限自由模的直和因子，因而有限生成。
\end{proof}

对投射对象，这个判据把有限生成性转成了由 Hom 与直和描述的条件，后面的等价定理将使用这一形式。

\subsection{矩阵环与原环的模范畴}
\label{b2:subsec:1601:03}

\begin{proposition}\label{b2:prop:matrix-morita-equivalence}
设 \(R\) 为含幺环、\(n\ge1\) 为整数，所论模均为幺性左模。令 \(A=M_n(R)\)、\(e=E_{11}\)，将 \(eAe\) 按 \(r\mapsto rE_{11}\) 识别为 \(R\)。以下函子互为拟逆：
\[
\mathcal L:R\text{-}\mathbf{Mod}\longrightarrow A\text{-}\mathbf{Mod},
\quad N\longmapsto N^n,
\qquad
\mathcal K:A\text{-}\mathbf{Mod}\longrightarrow R\text{-}\mathbf{Mod},
\quad M\longmapsto eM.
\]
\(\mathcal L\) 使用矩阵乘列的左作用，\(\mathcal K\) 使用角环作用；在同态上分别逐坐标作用和限制到 \(eM\)。
\end{proposition}
\begin{proof}
\(eN^n\) 恰为只有第一坐标可能非零的列，取该坐标自然同构到 \(N\)。另一方向用 \(E_{1j}\) 将第 \(j\) 个位置的数据读到第一角，再用 \(E_{i1}\) 写回第 \(i\) 个位置。相应地，定义
\[
\Theta_M:(eM)^n\longrightarrow M,\qquad
(m_1,\ldots,m_n)\longmapsto\sum_iE_{i1}m_i,
\]
并令 \(\Psi_M(m)=(E_{11}m,E_{12}m,\ldots,E_{1n}m)\)。每个坐标都在 \(eM\)。由 \(E_{1j}E_{i1}=\delta_{ij}e\) 与 \(\sum_iE_{i1}E_{1i}=1\)，两映射互逆。

对 \(a=(a_{ij})\in A\)，恒等式 \(aE_{j1}=\sum_iE_{i1}(a_{ij}E_{11})\) 说明 \(\Theta_M\) 把列上的矩阵作用变为 \(M\) 的原作用，故 \(A\) 线性。两映射只使用矩阵单位的作用，与任意 \(A\) 模同态交换，所以构成自然同构。第一坐标同构也显然与 \(R\) 模同态相容，因此两函子确为范畴等价。
\end{proof}

这个等价不要求 \(R\) 交换。例如 \(M_n(R)\) 的正则左模被 \(\mathcal K\) 送到首行模 \(eA\cong R^n\)。取 \(n=2\)，矩阵 \(m=\begin{psmallmatrix}a&b\\c&d\end{psmallmatrix}\) 对应于
\[
\Psi_A(m)=\left(
\begin{pmatrix}a&b\\0&0\end{pmatrix},
\begin{pmatrix}c&d\\0&0\end{pmatrix}
\right).
\]
逆映射把第一个首行留在第一行，再由 \(E_{21}\) 把第二个首行移到第二行，便恢复 \(m\)。矩阵的全部条目因此可以从首行模的两份副本恢复，而 \(eA\) 本身是秩二的 \(R\) 模。这里正则模 \({}_AA\) 的对应对象是 \({}_RR^n\)，等价不要求两边指定的正则模互相对应。取出 \(eM\)，再由矩阵单位恢复整个 \(M\)，给出了后面用双模张量积重构模的具体模型。

\begin{proposition}\label{b2:prop:trace-generator-criterion}
设 \(R\) 为含幺环、\(P\) 为幺左 \(R\) 模。子集
\[
\operatorname{tr}_R(P)=\sum_{f\in\operatorname{Hom}_R(P,R)}f(P)
\]
是 \(R\) 的双边理想，称为 \(P\) 的\term[ji-lixiang]{迹理想}[trace ideal]。以下条件等价：\(P\) 是全部幺左 \(R\) 模范畴的生成元；\(\operatorname{tr}_R(P)=R\)；存在非负整数 \(n\)、元素 \(p_1,\ldots,p_n\in P\) 和左 \(R\) 线性映射 \(f_1,\ldots,f_n:P\to R\)，使
\[
\sum_{i=1}^n f_i(p_i)=1_R.
\]
这里 \(n=0\) 时按空和为零理解，不要求 \(P\) 有限生成或投射。
\end{proposition}
\symidx{\(\operatorname{tr}_R(P)\)：模的迹理想}
\begin{proof}
每个 \(f(P)\) 都是正则左模的子模，所以它们的和是左理想。对任意 \(r\in R\)，映射 \(f_r:P\to R\)、\(f_r(p)=f(p)r\) 仍左线性，因为 \(f_r(ap)=af(p)r=af_r(p)\)。故迹理想也对右乘封闭，是双边理想。

迹理想为全环，当且仅当其中含有一；而这个和中的每个元素都由有限多个像相加得到，所以这又等价于所述有限和等式。若等式成立，定义
\[
q:P^n\longrightarrow R,\qquad (x_i)_i\longmapsto\sum_i f_i(x_i),
\qquad
s:R\longrightarrow P^n,\qquad r\longmapsto(rp_i)_i.
\]
两映射都左线性，且 \(q(s(r))=r\sum_i f_i(p_i)=r\)，故 \(R\) 是 \(P^n\) 的直和因子，由\cref{b2:thm:generator-characterizations}，\(P\) 为生成元。反过来，生成元判据给出有限份 \(P\) 到 \(R\) 的分裂满射。取一的原像 \((p_i)_i\)，并将满射限制到第 \(i\) 个分量得到 \(f_i\)，就有 \(\sum_i f_i(p_i)=1_R\)。
\end{proof}

迹判据中的有限和控制的是 \(1_R\)。相比之下，\cref{b2:prop:finite-dual-basis}中的有限对偶基控制的是每个 \(p\in P\) 的重构式，因而控制 \(\operatorname{id}_P\)。这两种有限数据分别检验生成元性质与有限生成投射性。

\begin{proposition}\label{b2:prop:object-detection-not-generator}
设 \(k\) 为域、\(R=k[\varepsilon]/(\varepsilon^2)\)，\(\varepsilon\) 的剩余类仍记为 \(\varepsilon\)，并令 \(P=R/(\varepsilon)\) 为幺左 \(R\) 模。对每个非零幺左 \(R\) 模 \(M\)，都有 \(\operatorname{Hom}_R(P,M)\ne0\)；但 \(\operatorname{Hom}_R(P,-)\) 不忠实，\(P\) 不是生成元。
\end{proposition}
\begin{proof}
在 \(M\) 中取 \(m\ne0\)。若 \(\varepsilon m=0\)，令 \(v=m\)；否则令 \(v=\varepsilon m\)，由 \(\varepsilon^2=0\) 仍有 \(v\ne0\) 且 \(\varepsilon v=0\)。于是 \(r+(\varepsilon)\mapsto rv\) 良定义，给出 \(P\to M\) 的非零左模同态。

考虑非零左线性映射 \(u:R\to R\)、\(u(r)=r\varepsilon\)。每个 \(h:P\to R\) 的像都被 \(\varepsilon\) 零化；写 \(r=a+b\varepsilon\) 可知 \(\varepsilon r=a\varepsilon=0\) 当且仅当 \(r\in k\varepsilon\)，故 \(h(P)\subseteq k\varepsilon\)，从而 \(uh=0\)。于是 \(u\ne0\) 和零同态经 Hom 函子后成为同一个映射，忠实性失败，\(P\) 不是生成元。
\end{proof}

\begin{proposition}\label{b2:prop:progenerator-hypothesis-counterexamples}
设 \(k\) 为域，所论模均为幺性左模。有限生成投射生成元定义中的三个条件，均不能由另外两个推出：
\begin{enumerate}
\item 在 \(R=k\times k\) 上，\(P=k\times0\) 有限生成且投射，但不是生成元。
\item 在 \(R=\mathbb Z\) 上，\(P=\mathbb Z\oplus\mathbb Z/2\mathbb Z\) 有限生成且为生成元，但不投射；\(\operatorname{Hom}_{\mathbb Z}(P,-)\) 不保持所有满同态。
\item 在 \(R=k\) 上，\(P=\bigoplus_{i\ge1}ku_i\) 投射且为生成元，但不有限生成，也不是紧对象；\(\operatorname{Hom}_k(P,-)\) 不保持任意直和。
\end{enumerate}
\end{proposition}
\begin{proof}
在第一例中，令 \(e=(1,0)\)。\(P=Re\) 是 \(R=Re\oplus R(1-e)\) 的直和因子，又由 \(e\) 生成，所以有限生成且投射。非零模 \(Q=0\times k\) 满足 \(eQ=0\)。若 \(h:P\to Q\)，则 \(h(e)=h(e^2)=eh(e)=0\)，所以 \(h=0\)。故 \(Q\) 上的恒等映射与零同态无法由 \(P\) 检测，\(P\) 不是生成元。

第二例的两个直和分量各由一个元素生成，而 \(\mathbb Z\) 是 \(P\) 的直和因子，所以由生成元判据，\(P\) 是有限生成的生成元。取自然满射 \(q:\mathbb Z\to\mathbb Z/2\mathbb Z\) 及投影 \(f:P\to\mathbb Z/2\mathbb Z\)。若有 \(h:P\to\mathbb Z\) 满足 \(qh=f\)，则 \(h\) 在 \(\mathbb Z/2\mathbb Z\) 分量上的限制应是从二阶元素到整数的同态，只能为零；但 \(f\) 在此分量上为恒等，矛盾。因此 \(P\) 不投射，\(\operatorname{Hom}_{\mathbb Z}(P,q)\) 也不满射。

第三例中，\(P\) 自由，因而投射；它含 \(ku_1\cong k\) 为直和因子，故是生成元。有限个向量的支集之并仍有限，所以不能生成全部 \(u_i\)，\(P\) 不有限生成。将恒等映射看成 \(P\to\bigoplus_{i\ge1}ku_i\)。如果它来自 \(\bigoplus_i\operatorname{Hom}_k(P,ku_i)\)，其像就会落在有限个坐标的直和中，不能包含所有 \(u_i\)，矛盾。所以 \(P\) 不是紧对象，Hom 函子不能保持这个直和。
\end{proof}

% !TeX root = ../../Book2.tex
% 纲要标题：### 16.2 双模实现的函子
% 纲要位置：计划纲要.md，第 1570 行。

\section{双模实现的函子}
\label{b2:sec:1602}

为了把两个模范畴联系起来，需要一个同时接受两侧标量作用的模。先从任意 \((R,S)\) 双模 \(P\) 出发，写出张量函子、Hom 函子及它们的伴随映射；然后再把 \(S\) 选为 \(P\) 的自同态环的反环。这样可以分别检查构造本身的合法性，以及何种条件使它成为等价。

% 纲要要点（供目录核对）：
%
% * 张量函子与 Hom 函子；调用第7.1节的有限投射 Hom–张量同构
% * 自同态环和反环约定；逐式核对预复合和右作用，正式证明任意环上的行列互逆双模及一般模上的余单位
% * 伴随的单位及余单位
%

\subsection{张量函子与 Hom 函子}
\label{b2:subsec:1602:01}

设 \(R,S\) 为含幺环，给定幺性双模 \({}_RP_S\)，定义
\[
F(M)=\operatorname{Hom}_R(P,M),\qquad
G(N)=P\otimes_SN.
\]
这里 \(M\) 为幺左 \(R\) 模，\(N\) 为幺左 \(S\) 模，因此 \(F:R\text{-}\mathbf{Mod}\to S\text{-}\mathbf{Mod}\)、\(G:S\text{-}\mathbf{Mod}\to R\text{-}\mathbf{Mod}\)。\(G(N)\) 的左 \(R\) 作用留在第一个因子；\(F(M)\) 的左 \(S\) 作用由
\begin{equation}\label{b2:eq:morita-hom-action}
(sf)(p)=f(ps)
\end{equation}
给出：\(s\) 先在输入端作用，再由 \(f\) 求值，因而这是预复合。它仍为左 \(R\) 线性映射，因为 \(P\) 的两侧作用相容；而
\[
\bigl(s_1(s_2f)\bigr)(p)=f\bigl((ps_1)s_2\bigr)
=\bigl((s_1s_2)f\bigr)(p),
\]
所以确为左 \(S\) 作用。

\ProseParagraphBreak
对 \(a:M\rightarrow M'\)，令 \(F(a)(f)=af\)；对 \(b:N\rightarrow N'\)，令 \(G(b)=\operatorname{id}_P\otimes b\)。复合与恒等关系由同态复合和张量函子性得到。它们都是加性函子。由\cref{b2:thm:hom-tensor-adjunction}，有关于两变量自然的对应
\[
\operatorname{Hom}_R(P\otimes_SN,M)
\cong\operatorname{Hom}_S\bigl(N,\operatorname{Hom}_R(P,M)\bigr),
\]
把 \(h\) 送到 \(n\mapsto\bigl(p\mapsto h(p\otimes n)\bigr)\)。因此 \(G\) 左伴随于 \(F\)。一般地，\(F\) 左正合，\(G\) 右正合；\(P\) 左投射时，\cref{b2:cor:projective-hom-exact}才使 \(F\) 正合。

\ProseParagraphBreak
若这里的双模 \({}_RP_S\) 作为左 \(R\) 模有限生成并投射，\cref{b2:prop:finite-projective-hom-tensor}便把已经定义的 \(F(M)=\operatorname{Hom}_R(P,M)\) 写成张量函子。具体地，令 \(P^*=\operatorname{Hom}_R(P,R)\)，它接受 \((S,R)\) 双模作用
\[
(s\lambda)(p)=\lambda(ps),\qquad
(\lambda r)(p)=\lambda(p)r.
\]
该命题给出关于每个左 \(R\) 模 \(M\) 自然的左 \(S\) 模同构
\[
P^*\otimes_RM\longrightarrow F(M),\qquad
\lambda\otimes m\longmapsto\bigl(p\mapsto\lambda(p)m\bigr).
\]
目标的左 \(S\) 作用正是式\eqref{b2:eq:morita-hom-action}中的作用，因此这个同构保留当前的双模侧别。公式中的标量 \(\lambda(p)\) 作用在 \(m\) 左侧；交换这两个符号在一般环上并没有意义。

\subsection{自同态环和反环约定}
\label{b2:subsec:1602:02}

现在固定一个含幺环 \(R\) 及幺左 \(R\) 模 \(P\)，取
\[
\Delta=\operatorname{End}_R(P),\qquad S=\Delta^{\mathrm{op}}.
\]
\symidx{\(S=\operatorname{End}_R(P)^{\mathrm{op}}\)：投射生成元对应的自同态反环}
若 \(s\in S\) 对应自同态 \(\alpha_s\in\Delta\)，定义 \(ps=\alpha_s(p)\)。因为
\[
\alpha_{st}=\alpha_t\circ\alpha_s,
\qquad (ps)t=\alpha_t\bigl(\alpha_s(p)\bigr)=p(st),
\]
这是右 \(S\) 作用；各 \(\alpha_s\) 左 \(R\) 线性，保证它与左作用相容。因此 \(P\) 成为 \((R,S)\) 双模，上一个小节的函子与伴随立即适用。

两侧作用可记为
\[
{}_RP_S,\qquad {}_SP^*_R,\qquad
P^*=\operatorname{Hom}_R(P,R),\qquad
S=\operatorname{End}_R(P)^{\mathrm{op}}.
\]
自同态的复合在右 \(S\) 作用中反序一次，因而这里必须取反环。

\ProseParagraphBreak
在此约定下，\(F(P)=\operatorname{End}_R(P)\) 作为左 \(S\) 模自然同构于正则左模 \(S\)：将 \(\alpha_s\) 对应于 \(s\)。事实上，\(t\alpha_s\) 的作用是先做 \(\alpha_t\) 再做 \(\alpha_s\)，即 \(\alpha_s\alpha_t=\alpha_{ts}\)，正对应左乘 \(t\)。所以在将来由这些函子给出的等价中，与正则左 \(S\) 模对应的是选定的 \(P\)，未必是正则左 \(R\) 模。这是一个左模识别；作为环，\(\Delta=S^{\mathrm{op}}\)，不能因此省去反环。

\ProseParagraphBreak
最基本的例子是 \(P={}_RR\)。每个左线性自同态都是右乘某个 \(a\in R\)，所以 \(\Delta\cong R^{\mathrm{op}}\)、\(S\cong R\)，Hom 在一处求值后恢复原模。若取左自由模 \(P=R^{1\times n}\)、\(n\ge1\)，写成行向量方便记录右作用；非交换系数不能随意搬到另一侧，所以不能把以下公式直接改成左乘列：每个左线性自同态唯一为右乘矩阵 \(B\in M_n(R)\)，两次右乘的算子复合满足
\[
D_BD_C(p)=(pC)B=p(CB).
\]
所以 \(S\cong M_n(R)\)。映射 \(f:P\rightarrow M\) 对应其在各标准行上的值组成的列
\[
\bigl(f(e_1),\ldots,f(e_n)\bigr)^{\mathsf T}\in M^n.
\]
由 \(e_iB=\sum_j b_{ij}e_j\)，左 \(S\) 作用恰为矩阵乘这个列。这便与\cref{b2:prop:matrix-morita-equivalence}的显式等价相合。

\subsection{伴随的单位及余单位}
\label{b2:subsec:1602:03}

对任意 \((R,S)\) 双模，伴随的单位和余单位由\cref{b2:thm:hom-tensor-adjunction}的两种求值公式给出：
\begin{align}
\eta_N:N&\longrightarrow FG(N),
&\eta_N(n)(p)&=p\otimes n,\label{b2:eq:morita-unit}\\
\varepsilon_M:GF(M)&\longrightarrow M,
&\varepsilon_M(p\otimes f)&=f(p).\label{b2:eq:morita-counit}
\end{align}
第一式比较先取张量再取 Hom 所得对象与原来的 \(N\)，第二式则将一个输入 \(p\) 与同态 \(f\) 配对，尝试恢复原来的 \(M\)。第二式的平衡性为 \(f(ps)=(sf)(p)\)，第一式的左 \(S\) 线性为 \(p\otimes sn=ps\otimes n\)。相应的左 \(R\) 线性分别来自 \(P\) 的左作用和 \(f\) 的线性。因此两者都是所在模范畴中的同态，不只是集合映射。

\begin{proposition}\label{b2:prop:morita-adjunction-maps}
设 \(R,S\) 为含幺环，\({}_RP_S\) 为幺性 \((R,S)\) 双模，令 \(F=\operatorname{Hom}_R(P,-)\)、\(G=P\otimes_S-\)。对左 \(S\) 模 \(N\) 和左 \(R\) 模 \(M\)，定义
\[
\eta_N:N\longrightarrow FG(N),\quad n\longmapsto\bigl(p\longmapsto p\otimes n\bigr),
\qquad
\varepsilon_M:GF(M)\longrightarrow M,\quad p\otimes f\longmapsto f(p).
\]
这些映射关于 \(N,M\) 自然，并满足
\begin{equation}\label{b2:eq:morita-triangles}
\varepsilon_{G(N)}G(\eta_N)=\operatorname{id}_{G(N)},\qquad
F(\varepsilon_M)\eta_{F(M)}=\operatorname{id}_{F(M)}.
\end{equation}
若 \(S=\operatorname{End}_R(P)^{\mathrm{op}}\)，且 \(P\) 的右 \(S\) 作用取自其左 \(R\) 线性自同态，则 \(\eta_S\) 与 \(\varepsilon_P\) 都是同构，不要求 \(P\) 已是生成元。
\end{proposition}
\begin{proof}
若 \(b:N\rightarrow N'\)，两种从 \(N\) 到 \(FG(N')\) 的路径在 \(n,p\) 处都为 \(p\otimes b(n)\)；若 \(a:M\rightarrow M'\)，余单位的两种路径在 \(p\otimes f\) 处都为 \(a\bigl(f(p)\bigr)\)，所以自然。第一条三角恒等式把 \(p\otimes n\) 送到 \(\eta_N(n)(p)=p\otimes n\)；第二条把 \(f\) 送到函数 \(p\mapsto\varepsilon_M(p\otimes f)=f(p)\)，故都成立。

在所给自同态环约定下，识别 \(P\otimes_SS\cong P\) 及 \(F(P)\cong S\)。\(\eta_S(s)\) 对应 \(p\mapsto ps=\alpha_s(p)\)，再对应到 \(s\)，所以是恒等映射。\(\varepsilon_P\) 在这些识别下为 \(P\otimes_SS\rightarrow P\)、\(p\otimes s\mapsto ps\)，是张量积的单位同构。
\end{proof}

这些特殊分量为同构，并不说明其他模上的分量也是同构。为了在全部模上得到同构，需要把模写成适当直和之间映射的余核，并证明函子保留相应的直和与正合尾段。有限生成、投射和生成元三个条件将在下一节分别承担这些任务。

\begin{proposition}[行列双模的互逆配对]\label{b2:prop:matrix-row-column-bimodules}
设 \(R\) 为含幺环、\(n\ge1\)，令 \(S=M_n(R)\)、\({}_RP_S=R^{1\times n}\)、\({}_SQ_R=R^{n\times1}\)，各作用由矩阵乘法给出。令 \(u_i\)、\(v_i\) 分别为第 \(i\) 个标准行、标准列，则矩阵乘法诱导双模同构
\[
P\otimes_SQ\longrightarrow R,\quad p\otimes q\longmapsto pq,
\qquad Q\otimes_RP\longrightarrow S,\quad q\otimes p\longmapsto qp.
\]
它们的逆分别为 \(r\mapsto ru_1\otimes v_1\) 和
\[
(b_{ij})\longmapsto\sum_{i,j}v_i b_{ij}\otimes u_j.
\]
对任意幺左 \(R\) 模 \(M\)，将 \(\operatorname{Hom}_R(P,M)\) 识别为左 \(S\) 模 \(M^n\) 后，余单位也是同构：
\[
P\otimes_SM^n\longrightarrow M,\quad
(p_i)_i\otimes(m_i)_i^{\mathsf T}\longmapsto\sum_i p_i m_i,
\]
其逆为 \(m\mapsto u_1\otimes(m,0,\ldots,0)^{\mathsf T}\)。
\end{proposition}
\begin{proof}
结合律给出两种矩阵配对的平衡性与外侧作用相容性。为核对第一同构的逆，将行列按标准基展开，得到
\[
(p_i u_i)\otimes(v_jq_j)
=u_1\otimes (p_iE_{1i})(v_jq_j)
=\delta_{ij}(p_iq_j)u_1\otimes v_1.
\]
第一步在 \(S\) 上移动 \(p_iE_{1i}\)，最后一步移动 \((p_iq_j)E_{11}\)。因此全张量等于 \((pq)u_1\otimes v_1\)，两复合均为恒等。第二同构的逆把 \(E_{ij}\) 送到 \(v_i\otimes u_j\)；对任意纯张量，\(R\) 上的平衡关系给出
\[
q\otimes p=\sum_{i,j}v_i(q_i p_j)\otimes u_j,
\]
恰为外积矩阵 \(qp\) 的逆像。

最后，任意行 \(p\) 可写成 \(u_1B\)，其中 \(B\) 的首行为 \(p\)、其余行为零。在张量中把 \(B\) 移到 \(M^n\) 上，得到
\[
p\otimes(m_i)_i^{\mathsf T}
=u_1\otimes\left(\sum_i p_i m_i,0,\ldots,0\right)^{\mathsf T}.
\]
这就验证了所列余单位及其逆在任意 \(M\) 上的两个复合。所有系数保持原次序，证明不要求 \(R\) 交换。
\end{proof}

% !TeX root = ../../Book2.tex
% 纲要标题：### 16.3 Morita 等价定理
% 纲要位置：计划纲要.md，第 1576 行。

\section{Morita 等价定理}
\label{b2:sec:1603}

第五章把 Morita 等价定义为模范畴的等价，却没有说明怎样从环本身判断它。本节证明：选择一个有限生成投射生成元，再取其自同态环的反环，恰好得到全部这样的等价。先从一个已给的范畴等价恢复这个模，再证明投射生成元确实使上一节的单位、余单位在每个对象上成为同构。

% 纲要要点（供目录核对）：
%
% * 从范畴等价恢复投射生成元；由双积恢复加法，先保持紧性和投射性，再恢复有限生成性
% * 由投射生成元构造等价；将单位与余单位从生成对象经任意直和、展示及余核推广到全部模
% * 互逆双模及右模版本
% * 满幂等元与角环的 Morita 等价；正式纳入非零对象检测判据和求值逆，并证明双边理想及商环对应
%

\subsection{从范畴等价恢复投射生成元}
\label{b2:subsec:1603:01}

\begin{lemma}\label{b2:lem:module-equivalence-structure}
设 \(R,S\) 为含幺环，\(H:R\text{-}\mathbf{Mod}\rightarrow S\text{-}\mathbf{Mod}\) 为幺左模范畴之间的等价，则 \(H\) 为加性函子，保持并反映核、余核、正合列、投射性与内射性，也保持任意直和，并将生成元送到生成元。
\end{lemma}
\begin{proof}
核、余核及直和由泛性质刻画，投射性、内射性与生成元也可用同态的提升、延拓和商对象刻画，因此这些性质都能随等价转移。需要先说明的是，同态的加法本身也能由范畴结构恢复。由\cref{b2:prop:equivalence-preserves-universal-objects}，等价保持零对象及任意积、余积，特别保持任意直和。零同态是经过零对象的复合，因此也被保持。对同态 \(f,g:M\rightarrow N\)，加法可写成
\[
f+g=\nabla_N(f\oplus g)\Delta_M,
\]
其中 \(\Delta_M:M\rightarrow M\oplus M\) 的两个投影均为恒等，\(\nabla_N:N\oplus N\rightarrow N\) 的两个限制均为恒等。直和同时是积与余积，故等价在相应的自然直和识别下保持这两个映射及此式，所以保持加法。无需在模范畴的等价之外另加一个加性要求。

等价保持并反映单态射和满态射的消去性质；模范畴中的单态射就是单射，因为非零核元素可以通过从 \(R\) 出发的同态检测。满态射就是满射，因为若像不为全体，非零商映射与零映射在其像上相同，会破坏右消去性。因此单射与满射也被保持并反映。

以核为例，设 \(j:K\rightarrow M\) 是 \(f:M\rightarrow N\) 的核，等价记为 \(H\)。若 \(a:Y\rightarrow H(M)\) 满足 \(H(f)a=0\)，取同构 \(u:H(X)\rightarrow Y\)，再由全忠实性把 \(au\) 唯一提升为 \(b:X\rightarrow M\)。忠实性给出 \(fb=0\)，故 \(b=jc\) 且 \(c\) 唯一。于是 \(a=H(j)H(c)u^{-1}\)，而任何别的分解提升后都等于 \(c\)。所以 \(H(j)\) 是核。

对余核 \(q:N\rightarrow C\)，给定 \(a:H(N)\rightarrow Y\) 且 \(aH(f)=0\)，选 \(u:H(X)\rightarrow Y\)，把 \(u^{-1}a\) 提升为 \(b:N\rightarrow X\)，则 \(bf=0\)，由余核性质唯一写成 \(b=cq\)，所以 \(a=uH(c)H(q)\) 且分解唯一。故余核也保持。像可表为余核的核，因此也保持。对拟逆使用同样论证，得到反映性，正合列因核与像的相等而双向保持。

投射性由沿满射提升的性质刻画。将一个提升问题经拟逆送回，满射仍满，再作提升并经等价送出，就保持投射性；反向相同。对偶地，内射性由沿单射延拓的性质刻画；等价及拟逆都保持单射，故也保持并反映内射性。若 \(P\) 是生成元，每个源范畴对象都为 \(P\) 某个直和的商。等价保持直和与满射，且每个目标对象都同构于某个 \(H(X)\)，故每个目标对象都为 \(H(P)\) 某个直和的商。因此 \(H(P)\) 也是生成元。
\end{proof}

目标范畴的正则模 \({}_SS\) 是投射生成元，而且 \(\operatorname{Hom}_S(S,Y)\) 在 \(1\) 处求值就恢复 \(Y\)。因此把它经拟逆送回得到 \(P=\Psi(S)\)，应当能用 \(\operatorname{Hom}_R(P,-)\) 恢复原等价。投射性与生成元性质已有范畴刻画；有限生成性则不能直接用元素和生成元的个数转移，需要先把正则模的紧性转移给 \(P\)，再利用 \(P\) 的投射性。

\begin{theorem}[由范畴等价恢复投射生成元]\label{b2:thm:morita-converse}
设 \(R,S\) 为含幺环，\(\Phi:R\text{-}\mathbf{Mod}\rightarrow S\text{-}\mathbf{Mod}\) 为幺左模范畴的等价，取其拟逆 \(\Psi\)，令 \(P=\Psi({}_SS)\)，则幺左 \(R\) 模 \(P\) 为有限生成投射生成元，并有环同构
\[
S\cong\operatorname{End}_R(P)^{\mathrm{op}}.
\]
通过这个同构赋予 Hom 左 \(S\) 作用后，\(\Phi\) 自然同构于 \(\operatorname{Hom}_R(P,-)\)。
\end{theorem}
\begin{proof}
正则模 \(S\) 投射且为生成元，\cref{b2:lem:module-equivalence-structure}使 \(P=\Psi(S)\) 也投射且为生成元。固定拟逆提供的自然同构 \(\nu:\operatorname{Id}_{S\text{-}\mathbf{Mod}}\rightarrow\Phi\Psi\)。为证明有限生成，先证明 \(P\) 紧。对任意幺左 \(R\) 模族 \((M_i)\)，有自然同构
\[
\begin{aligned}
\operatorname{Hom}_R\left(P,\bigoplus_iM_i\right)
&\cong\operatorname{Hom}_S\left(S,\Phi\left(\bigoplus_iM_i\right)\right)\\
&\cong\operatorname{Hom}_S\left(S,\bigoplus_i\Phi(M_i)\right)\\
&\cong\bigoplus_i\operatorname{Hom}_S\bigl(S,\Phi(M_i)\bigr)\\
&\cong\bigoplus_i\operatorname{Hom}_R(P,M_i).
\end{aligned}
\]
第一项同构来自 \(\Phi\) 的全忠实性及 \(\nu_S:S\cong\Phi(P)\)，第二项来自等价保持直和。第三项使用正则模 \(S\) 由 \(1\) 生成：从 \(S\) 到直和的同态由 \(1\) 的像决定，而这个像只有有限支集。最后一项再用全忠实性及 \(\nu_S\)。这些同构与各直和嵌入相容，所以其复合正是紧性中的自然映射之逆。\cref{b2:prop:compact-projective-finite}于是说明投射模 \(P\) 有限生成。

等价加性且全忠实，故 \(\Psi\) 在自同态环上给出环同构
\[
\operatorname{End}_S({}_SS)\cong\operatorname{End}_R(P).
\]
由\cref{b2:prop:regular-hom-endomorphism}，正则左模的自同态是右乘映射 \(r_s:x\mapsto xs\)，并有 \(r_s\circ r_t=r_{ts}\)。因此 \(s\mapsto\Psi(r_s)\) 给出 \(S\cong\operatorname{End}_R(P)^{\mathrm{op}}\)，相应的右作用为 \(ps=\Psi(r_s)(p)\)。

对每个 \(M\)，再用正则模的求值对应
\[
\operatorname{Hom}_R(P,M)
\cong\operatorname{Hom}_S(S,\Phi(M))
\cong\Phi(M).
\]
具体地，\(h:P\rightarrow M\) 先对应 \(a=\Phi(h)\nu_S:S\rightarrow\Phi(M)\)，再对应 \(a(1)\)。后复合 \(M\rightarrow M'\) 的同态与这两步相容，故所得同构对 \(M\) 自然。它也保持左 \(S\) 作用：Hom 中的 \(s\cdot h\) 为 \(h\circ\Psi(r_s)\)，而 \(\nu\) 的自然性给出
\[
\Phi\bigl(h\circ\Psi(r_s)\bigr)\nu_S
=\Phi(h)\nu_S r_s=a r_s.
\]
在 \(1\) 处求值，右边为 \(a(s)=s\cdot a(1)\)，所以该自然同构为左 \(S\) 模同构。
\end{proof}

\subsection{投射生成元与等价构造}
\label{b2:subsec:1603:02}

\begin{theorem}[Morita 等价]\label{b2:thm:morita-equivalence}
设 \(R\) 为含幺环，\(P\) 为幺左 \(R\) 模范畴的有限生成投射生成元，令 \(S=\operatorname{End}_R(P)^{\mathrm{op}}\)，并使 \(s\in S\) 在 \(P\) 上的右作用为 \(ps=\alpha_s(p)\)，其中 \(\alpha_s\) 是 \(s\) 对应的左 \(R\) 线性自同态，则
\[
F=\operatorname{Hom}_R(P,-),\qquad G=P\otimes_S-
\]
分别为从幺左 \(R\) 模到幺左 \(S\) 模、从幺左 \(S\) 模到幺左 \(R\) 模的互逆等价，而且式\eqref{b2:eq:morita-unit}与式\eqref{b2:eq:morita-counit}的单位、余单位在每个模上都是同构。结合\cref{b2:thm:morita-converse}，任意两个含幺环 \(R,S\) 在 Morita 意义下等价，当且仅当存在幺左 \(R\) 模的有限生成投射生成元 \(P\) 使 \(S\cong\operatorname{End}_R(P)^{\mathrm{op}}\)。
\end{theorem}
\begin{proof}
上一节已知 \(\varepsilon_P\) 与 \(\eta_S\) 为同构。要把它们扩展到所有对象，可以先扩展到 \(P\) 或 \(S\) 的任意直和，再把任意模写成这些直和之间某个映射的余核。这里有限生成性与\cref{b2:prop:compact-projective-finite}保证 \(F\) 保持任意直和；投射性与\cref{b2:cor:projective-hom-exact}保证 \(F\) 正合。\cref{b2:prop:tensor-direct-sums}与\cref{b2:thm:tensor-right-exact}说明 \(G\) 保持直和且右正合。因此 \(FG,GF\) 都保持直和与正合尾段。

由\cref{b2:prop:morita-adjunction-maps}，\(\varepsilon_P\) 为同构；因两边函子都保持直和，且求值公式与各嵌入相容，每个 \(\varepsilon_{P^{(I)}}\) 也是同构。对任意幺左 \(R\) 模 \(M\)，生成元性质给出满射 \(q:P^{(J)}\rightarrow M\)。记其核为 \(K\)，再对 \(K\) 使用生成元性质，得到满射 \(P^{(I)}\rightarrow K\)。将它与核的嵌入复合，得到 \(d:P^{(I)}\rightarrow P^{(J)}\) 及正合列
\[
P^{(I)}\longrightarrow P^{(J)}\longrightarrow M\longrightarrow0.
\]
生成元性质的这两次使用，分别保证 \(M\) 有由 \(P\) 给出的生成元和关系。应用右正合函子 \(GF\)，再用余单位的自然性，得到交换的正合尾段
\[
\begin{tikzcd}[column sep=2.6em,row sep=2.5em]
GF(P^{(I)}) \arrow[r,"GF(d)"] \arrow[d,"\varepsilon_{P^{(I)}}"'] &
GF(P^{(J)}) \arrow[r,"GF(q)"] \arrow[d,"\varepsilon_{P^{(J)}}"'] &
GF(M) \arrow[r] \arrow[d,"\varepsilon_M"'] & 0\\
P^{(I)} \arrow[r,"d"'] &
P^{(J)} \arrow[r,"q"'] &
M \arrow[r] & 0
\end{tikzcd}
\]
前两列的竖直映射为同构，并把 \(\operatorname{im}GF(d)\) 对应到 \(\operatorname{im}d\)，因而在对这两个像取商后诱导同构。两行的最后一个对象都是相应第一映射的余核，故由余核泛性质，这个同构正是 \(\varepsilon_M\)。

已知 \(\eta_S\) 为同构，两边保持直和，使每个 \(\eta_{S^{(I)}}\) 也为同构。任意幺左 \(S\) 模 \(N\) 本来就有自由展示
\[
S^{(I)}\longrightarrow S^{(J)}\longrightarrow N\longrightarrow0,
\]
这里指标集允许无限；这一侧用的是正则模 \(S\) 的生成元性质，不再需要 \(P\) 生成 \(R\) 模。把前两个映射记为 \(d':S^{(I)}\to S^{(J)}\)、\(q':S^{(J)}\to N\)。应用右正合函子 \(FG\)，由单位的自然性得到
\[
\begin{tikzcd}[column sep=2.6em,row sep=2.5em]
S^{(I)} \arrow[r,"d'"] \arrow[d,"\eta_{S^{(I)}}"'] &
S^{(J)} \arrow[r,"q'"] \arrow[d,"\eta_{S^{(J)}}"'] &
N \arrow[r] \arrow[d,"\eta_N"'] & 0\\
FG(S^{(I)}) \arrow[r,"FG(d')"'] &
FG(S^{(J)}) \arrow[r,"FG(q')"'] &
FG(N) \arrow[r] & 0
\end{tikzcd}
\]
前两列的竖直映射为同构，第三列是它们在余核上诱导的映射，故 \(\eta_N\) 也为同构。

所以 \(\eta:\operatorname{Id}\rightarrow FG\)、\(\varepsilon:GF\rightarrow\operatorname{Id}\) 都是自然同构，正好给出两函子互为拟逆。任意等价的反方向已经由\cref{b2:thm:morita-converse}证明，得到最后的充要条件。
\end{proof}

这就是\term[Morita-dengjia-dingli]{Morita 等价定理}[Morita equivalence theorem]。仅在正则模或有限自由模上检验单位、余单位，不足以替代证明中通过任意直和和余核所作的扩展。

\ProseParagraphBreak
同态也可显式恢复。若 \(a:F(M)\rightarrow F(M')\) 左 \(S\) 线性，则它唯一来自
\[
\varepsilon_{M'}\,G(a)\,\varepsilon_M^{-1}:M\longrightarrow M'.
\]
确实，式\eqref{b2:eq:morita-triangles}给出 \(F(\varepsilon_M)=\eta_{F(M)}^{-1}\)，单位的自然性便说明将上式施加 \(F\) 后得到 \(a\)。唯一性来自等价的忠实性。这使对象对应与同态对应同时可计算。

\begin{corollary}[互逆双模]\label{b2:cor:morita-inverse-bimodules}
设 \(R\) 为含幺环，\({}_RP\) 为幺左模范畴的有限生成投射生成元，\(S=\operatorname{End}_R(P)^{\mathrm{op}}\)，并令 \(ps=\alpha_s(p)\)，其中 \(\alpha_s\in\operatorname{End}_R(P)\) 对应 \(s\in S\)。令 \(Q=P^*=\operatorname{Hom}_R(P,R)\)，其幺 \((S,R)\) 双模作用为 \((s\lambda)(p)=\lambda(ps)\)、\((\lambda r)(p)=\lambda(p)r\)。则有双模同构
\[
{}_RP\otimes_S Q_R\cong {}_RR_R,\qquad
{}_SQ\otimes_RP_S\cong {}_SS_S.
\]
第一式把 \(p\otimes\lambda\) 送到 \(\lambda(p)\)；第二式把 \(\lambda\otimes p\) 送到 \(S\) 中对应于 \(x\mapsto\lambda(x)p\) 的元素。
\end{corollary}
\begin{proof}
第一式是\eqref{b2:eq:morita-counit}在 \(M=R\) 时的余单位，由\cref{b2:thm:morita-equivalence}为左 \(R\) 模同构。它也保持右 \(R\) 作用，因为 \((\lambda r)(p)=\lambda(p)r\)。

由\cref{b2:prop:finite-projective-hom-tensor}取 \(M=P\)，映射
\[
Q\otimes_RP\longrightarrow\operatorname{End}_R(P),\qquad
\lambda\otimes p\longmapsto T_{\lambda,p},\quad
T_{\lambda,p}(x)=\lambda(x)p
\]
为同构。有限对偶基解释了为什么这些自同态足以构成所有自同态。由\cref{b2:prop:finite-dual-basis}取 \((p_i,\lambda_i)\)，则 \(x=\sum_i\lambda_i(x)p_i\)，因而每个 \(R\) 线性自同态 \(T\) 都满足
\[
T(x)=\sum_i\lambda_i(x)T(p_i).
\]
所以任意自同态都是上述形式的有限和，且已证同构的逆将 \(T\) 送到 \(\sum_i\lambda_i\otimes T(p_i)\)。

将自同态 \(T\) 看作反环 \(S\) 的元素时，左、右 \(S\) 作用分别对应于两种复合次序：
\[
T_{s\lambda,p}=T_{\lambda,p}\circ\alpha_s,\qquad
T_{\lambda,ps}=\alpha_s\circ T_{\lambda,p}.
\]
由于 \(S\) 的乘法与自同态复合反向，这两式分别是 \(S\) 中的左乘 \(s\) 与右乘 \(s\)。故第二式是 \((S,S)\) 双模同构；这里 \(\operatorname{End}_R(P)\) 与 \(S\) 的识别用的是其加性群及这两个作用，而不是把两者的环乘法认作相同。
\end{proof}

这两个同构说明 \({}_RP_S\) 与 \({}_SQ_R\) 在双模张量意义下互为逆：它们的两种复合分别恢复正则双模 \({}_RR_R\) 与 \({}_SS_S\)。同一对双模因此也能从右模一侧构造等价。

\begin{corollary}\label{b2:cor:morita-left-right}
对任意含幺环 \(R,S\)，全部幺左模范畴等价当且仅当全部幺右模范畴等价：
\[
R\text{-}\mathbf{Mod}\simeq S\text{-}\mathbf{Mod}
\quad\Longleftrightarrow\quad
\mathbf{Mod}\text{-}R\simeq\mathbf{Mod}\text{-}S.
\]
\end{corollary}
\begin{proof}
若左模范畴等价，由\cref{b2:thm:morita-converse,b2:cor:morita-inverse-bimodules}取互逆双模 \({}_RP_S\) 与 \({}_SQ_R\)。函子 \(-\otimes_RP\) 把幺右 \(R\) 模送到幺右 \(S\) 模，\(-\otimes_SQ\) 把幺右 \(S\) 模送到幺右 \(R\) 模。两次复合经张量结合律分别同构于 \(-\otimes_R(P\otimes_SQ)\cong-\otimes_RR\) 和 \(-\otimes_S(Q\otimes_RP)\cong-\otimes_SS\)，故互为拟逆。反过来，把右 \(R,S\) 模分别视为左 \(R^{\mathrm{op}},S^{\mathrm{op}}\) 模，对刚证的方向再用一次，便得到左模范畴的等价。
\end{proof}

\subsection{满幂等元与角环的 Morita 等价}
\label{b2:subsec:1603:03}

\begin{proposition}[满幂等角环]\label{b2:prop:full-idempotent-morita}
设 \(R\) 为含幺环，\(e\in R\) 满足 \(e^2=e\)，令 \(S=eRe\)，其单位元为 \(e\)。条件 \(ReR=R\)、幺左 \(R\) 模 \(Re\) 为生成元，以及对每个幺左 \(R\) 模 \(M\) 都有 \(eM=0\Longrightarrow M=0\)，三者等价。在这些条件下，\(Re\) 为有限生成投射生成元，且
\[
R\text{-}\mathbf{Mod}\simeq S\text{-}\mathbf{Mod}
\]
由 \(M\mapsto eM\) 与 \(N\mapsto Re\otimes_SN\) 实现，其中 \(eM\) 取自然的幺左 \(S\) 模结构。对每个幺左 \(R\) 模 \(M\)，求值映射
\[
\varepsilon_M:Re\otimes_SeM\longrightarrow M,\qquad
re\otimes m\longmapsto rm
\]
为自然的左 \(R\) 模同构。
\end{proposition}
\begin{proof}
由 \(R=Re\oplus R(1-e)\)，\(Re\) 是正则左模的直和因子，故有限生成且投射。任意左模同态 \(Re\rightarrow R\) 由 \(e\) 的像 \(b\in eR\) 决定，公式为 \(re\mapsto rb\)。因此若 \(Re\) 为生成元，就有某个满射 \((Re)^{(L)}\rightarrow R\)。给 \(1\) 取一个原像，这个原像仅有有限支集；结合各分量上的同态公式便得到 \(1=\sum_{i=1}^n a_i e b_i\)，所以 \(ReR=R\)。此处和式有限来自这个原像的有限支集。

反过来，若 \(1=\sum_{i=1}^n a_i e b_i\)，映射 \((Re)^n\rightarrow R\)、\((p_i)\mapsto\sum_i p_i b_i\) 的像包含一，因而满射。其截面可取 \(r\mapsto(ra_i e)_i\)，所以 \(R\) 是 \((Re)^n\) 的直和因子。\cref{b2:thm:generator-characterizations}使 \(Re\) 为生成元。

若 \(ReR=R\) 且 \(eM=0\)，对每个 \(m\in M\) 有 \(m=\sum_i a_i e b_i m=0\)，所以 \(M=0\)。若 \(ReR\ne R\)，则 \(M=R/ReR\) 为非零幺左模，而 \(eM=0\)。这就证明第三个条件与满性等价，也说明非满的取角会把某个非零模送到零。

由\cref{b2:prop:corner-endomorphism}，\(\operatorname{End}_R(Re)^{\mathrm{op}}\cong S\)。与该命题中的同态识别一样，一个 \(f:Re\rightarrow M\) 完全由 \(f(e)\) 决定，因为 \(f(re)=rf(e)\)，且 \(f(e)=ef(e)\in eM\)。因此有自然同构
\[
\operatorname{Hom}_R(Re,M)\longrightarrow eM,\qquad f\longmapsto f(e),
\]
逆为 \(m\mapsto(re\mapsto rm)\)。若 \(re=r'e\)，因 \(m=em\) 有 \(rm=r'm\)，故逆公式良定义。对 \(s\in S\)，右乘 \(s\) 是 \(Re\) 的自同态，而 Hom 中的左作用满足 \((s\cdot f)(e)=f(es)=sf(e)\)，所以此同构保留左 \(S\) 作用。这里 \(eM\) 一般不是 \(M\) 的左 \(R\) 子模，所用的标量环是 \(S\)。当 \(e\) 满时，\(Re\) 已是有限生成投射生成元，\cref{b2:thm:morita-equivalence}给出题述等价。

求值映射 \(\varepsilon_M\) 也有直接的逆公式。固定一次表达 \(1=\sum_i a_i e b_i\)，则 \(m=\sum_i a_i e b_i m\)；要使求值恢复 \(m\)，就把每一项拆成 \(a_i e\otimes e b_i m\)。由此得到加性映射
\begin{equation}\label{b2:eq:full-idempotent-evaluation-inverse}
\delta_M(m)=\sum_i a_i e\otimes e b_i m.
\end{equation}
于是 \(\varepsilon_M\delta_M(m)=\sum_i a_i e b_i m=m\)。反向复合在 \(re\otimes m\)、\(m=em\) 上得到
\[
\sum_i a_i e\otimes e b_i r e m
=\sum_i a_i e b_i r e\otimes m=re\otimes m,
\]
其中中间一步移动的平衡标量是 \(e b_i r e\in S\)。故 \(\delta_M\) 与左 \(R\) 线性映射 \(\varepsilon_M\) 互为逆，\(\delta_M\) 也左 \(R\) 线性。求值与模同态相容，因而此同构自然。另一方面，\(e(Re\otimes_SN)\) 自然识别为 \(S\otimes_SN\cong N\)：右 \(S\) 模分裂投影 \(Re\rightarrow S\)、\(p\mapsto ep\) 与其包含映射张量后仍互为投影与截面，故它的像正是 \(e(Re\otimes_SN)\)，最后使用单位张量同构。
\end{proof}

对含幺环 \(B\) 及 \(n\ge1\)，在矩阵环 \(R=M_n(B)\) 中，\(e=E_{11}\) 满，因为 \(1=\sum_iE_{i1}eE_{1i}\)。上述公式正好恢复\cref{b2:prop:matrix-morita-equivalence}。当 \(B=k\) 为域且 \(n>1\) 时，令 \(A=M_n(k)\)，则幺左 \(A\) 模 \(Ae\) 是有限生成投射生成元，却不是自由模：\(\dim_kAe=n\)，而任意非零自由左 \(A\) 模都包含一个正则模 \(A\) 的直和分量，故其 \(k\) 维数至少为 \(n^2>n\)，无论自由基有限还是无限都不可能同构于 \(Ae\)。

\ProseParagraphBreak
若 \(R=B\times C\)、\(e=(1,0)\)，且 \(C\ne0\)，则 \(ReR=B\times0\ne R\)，取角只保留第一个坐标；第二个坐标上的非零模被送到零，所以不能给出等价。

\ProseParagraphBreak
对有限维代数 \(A\)，\cref{b2:prop:basic-corner}从每种有限生成不可分解投射模的同构类型中选取一份，构造了满足 \(AeA=A\) 的幂等元，并证明 \(eAe\) 为基本代数。\cref{b2:prop:full-idempotent-morita}进一步说明 \(A\) 与这个基本代数的全部幺左模范畴等价。因此可先在基本代数上研究模，再经上述张量函子返回原代数。

\begin{proposition}[满角环的双边理想对应]\label{b2:prop:full-corner-ideal-correspondence}
设 \(R\) 为含幺环，\(e\in R\) 为满幂等元，即 \(e^2=e\) 且 \(ReR=R\)，令 \(S=eRe\)。对双边理想 \(I\subseteq R\) 与 \(K\subseteq S\)，映射
\[
I\longmapsto eIe,\qquad K\longmapsto RKR
\]
给出 \(R\) 与 \(S\) 的双边理想格之间互逆且保持包含关系的对应，其中 \(RKR\) 表示元素 \(rkt\)、\(r,t\in R\)、\(k\in K\) 的有限和组成的双边理想。对 \(R\) 的每个双边理想 \(I\)，记 \(\bar e=e+I\)，则 \(\bar e\) 是 \(R/I\) 的满幂等元，并有环同构
\[
\bar e(R/I)\bar e\cong S/eIe.
\]
因此 \(R/I\) 与 \(S/eIe\) 的全部幺左模范畴等价。
\end{proposition}
\begin{proof}
若 \(I\) 为 \(R\) 的双边理想，则 \(eIe\) 为 \(S\) 的双边理想；若 \(K\) 为 \(S\) 的双边理想，则上述有限和组成的 \(RKR\) 为 \(R\) 的双边理想。固定 \(1=\sum_{i=1}^n a_i e b_i\)。对 \(x\in I\)，两次使用这个等式得到
\[
x=\sum_{i,j}a_i e(b_i x a_j)e b_j\in R(eIe)R.
\]
反向包含 \(R(eIe)R\subseteq I\) 来自 \(I\) 的双边理想性质，故 \(R(eIe)R=I\)。

对 \(k\in K\) 有 \(k=eke\)，所以 \(K\subseteq e(RKR)e\)。反向包含可逐项检查：若 \(r,t\in R\)、\(k\in K\)，则
\[
e(rkt)e=(ere)k(ete)\in K.
\]
故 \(e(RKR)e=K\)。两种映射显然都保持包含关系，因此给出互逆的理想格对应。

商映射限制到 \(S\) 给出满环同态
\[
S\longrightarrow\bar e(R/I)\bar e,\qquad
ere\longmapsto\bar e(r+I)\bar e.
\]
其核为 \(S\cap I=eIe\)，因为 \(s\in S\cap I\) 时 \(s=ese\)。商环同构定理便给出所述角环同构。\(1=\sum_i a_i e b_i\) 在 \(R/I\) 中仍成立，故 \(\bar e\) 满；由\cref{b2:prop:full-idempotent-morita}应用于 \(R/I\) 即得模范畴等价。此论证也包含 \(I=R\) 的情形，此时两个商环均为零环。
\end{proof}

% !TeX root = ../../Book2.tex
% 纲要标题：### 16.4 Morita 不变量
% 纲要位置：计划纲要.md，第 1583 行。

\section{Morita 不变量}
\label{b2:sec:1604}

模范畴的等价保留对象、同态及其相互关系，却不把一个模的底集合或坐标照搬过去。若一种性质还使用了指定的正则模、到向量空间的忘却函子或额外的型，就须另行检查这些数据是否相容。例如，相应模的合成长度数值确实保持，而它们作为同一个底域上的向量空间，维数可以改变；两个正则模也未必彼此对应。中心和子模结构则可以从模范畴本身恢复。

% 纲要要点（供目录核对）：
%
% * 中心、简单模、投射模和内射模；中心作为恒等函子自然自变换环，子模作为单射子对象及其因子关系
% * 半单性、不可分解性与有限长度；正式证明有限展示与环的左 Noether、左 Artin 性保持，以及与除环等价的单 Artin 判据
% * 维数、生成元数及附加型不是 Morita 不变量
% * 在同一实向量空间上比较正定与不定型及其不同伴随运算，说明普通 Morita 数据没有指定型
% ---
%

\subsection{中心、简单模与投射性}
\label{b2:subsec:1604:01}

\begin{proposition}\label{b2:prop:morita-center}
设 \(R,S\) 为含幺环，幺性左模范畴 \(\mathcal C_R=R\text{-}\mathbf{Mod}\) 与 \(\mathcal C_S=S\text{-}\mathbf{Mod}\) 等价。对 \(T=R,S\)，恒等函子的自然自变换以逐分量加法、复合作为乘法、恒等自然变换作为单位，组成环
\(\operatorname{Nat}(\operatorname{Id}_{\mathcal C_T},\operatorname{Id}_{\mathcal C_T})\)。映射
\[
Z(T)\longrightarrow
\operatorname{Nat}(\operatorname{Id}_{\mathcal C_T},\operatorname{Id}_{\mathcal C_T}),
\qquad
c\longmapsto\bigl(\theta^c_M:m\longmapsto cm\bigr)_{M\in\mathcal C_T}
\]
是含幺环同构。因此 \(Z(R)\cong Z(S)\)。
\end{proposition}
\begin{proof}
中心与自然自变换环的同构已由\cref{b2:prop:natural-center}证明。其逆映射是在正则模的一处求值：\(c=\theta_T(1)\)。任意模元素 \(m\in M\) 都给出同态 \(f_m:T\to M\)、\(r\mapsto rm\)，自然性强迫 \(\theta_M(m)=cm\)，因此这一求值决定所有分量；加法、复合与单位分别对应 \(c+d\)、\(cd\) 与 \(1\)。

接下来将一个范畴中对所有对象自然定义的自同态，传到另一个范畴。设 \(H:\mathcal C_R\rightarrow\mathcal C_S\) 为等价，固定拟逆 \(K\) 及自然同构 \(\beta:HK\Rightarrow\operatorname{Id}_{\mathcal C_S}\)。给定 \(\theta\in\operatorname{Nat}(\operatorname{Id}_{\mathcal C_R},\operatorname{Id}_{\mathcal C_R})\)，对每个幺性左 \(S\) 模 \(Y\) 令
\[
\theta'_Y=\beta_Y\,H(\theta_{K(Y)})\,\beta_Y^{-1}.
\]
\(\beta\) 与 \(\theta\) 的自然性保证这些分量构成自然自变换。由\cref{b2:lem:module-equivalence-structure}，\(H\) 加性，故此对应保持逐分量加法；共轭保持复合与单位。对 \(H(X)\)，全忠实性使 \(\beta_{H(X)}=H(b_X)\)，其中 \(b_X:KH(X)\to X\) 为同构。\(\theta\) 对 \(b_X\) 的自然性给出
\[
\theta'_{H(X)}
=H\bigl(b_X\theta_{KH(X)}b_X^{-1}\bigr)
=H(\theta_X).
\]
因此这个对应单射。反过来，目标范畴上的任一自然自变换 \(\tau\) 在各 \(H(X)\) 上的分量可由全忠实性唯一提升为 \(\theta_X\)。\(\tau\) 对 \(H(f)\) 的自然性与 \(H\) 的忠实性保证 \(\theta\) 自然。它的像与 \(\tau\) 在所有 \(H(X)\) 上相同，再用同构 \(\beta_Y:HK(Y)\to Y\) 的自然性，就知在每个 \(Y\) 上也相同。因此两个自然自变换环同构，进而 \(Z(R)\cong Z(S)\)。
\end{proof}

\begin{proposition}\label{b2:prop:morita-submodule-lattices}
设 \(R,S\) 为含幺环，\(H:R\text{-}\mathbf{Mod}\rightarrow S\text{-}\mathbf{Mod}\) 为幺性左模范畴的等价，\(M\) 为幺性左 \(R\) 模，则给定的 \(H\) 诱导 \(M\) 与 \(H(M)\) 的子对象格之间的序同构。把子对象视为子模时，\(N\subseteq M\) 对应于 \(H(N\hookrightarrow M)\) 的像；这一序同构保留零子模、全模、任意子模族的交与和。
\end{proposition}
\begin{proof}
子对象由单射 \(i:N\to M\) 表示，两个单射在定义域之间存在与它们相容的同构时，表示同一子对象。在模范畴中，取 \(i(N)\subseteq M\) 就给出它的子模代表。\cref{b2:lem:module-equivalence-structure}保证 \(H(i)\) 仍单射，所以其像同构于 \(H(N)\)；这里对应的是嵌入的像，而不是对底集合逐元素应用 \(H\)。反过来，给定 \(L\subseteq H(M)\)，选同构 \(u:H(X)\rightarrow L\)，把复合 \(H(X)\rightarrow L\hookrightarrow H(M)\) 唯一提升为 \(a:X\rightarrow M\)。等价反映单射，故 \(a\) 单射，令 \(N=a(X)\)，便得到所需原子模。若两个原子模给出相同像，像之间保持嵌入的同构提升后仍与原嵌入相容，因而两者在 \(M\) 中的像相同，故子模相等。

子模包含等价于两个嵌入之间存在相容的因子映射。全忠实性双向保持这一条件，所以所得双射保持并反映包含关系。零与全模是最小、最大元；交为最大下界，子模之和为最小上界，因此它们也由这个序同构保留。
\end{proof}

非零模简单，当且仅当子模序只有零和全模两个元素，所以简单模在等价下对应。\cref{b2:lem:module-equivalence-structure}已经从满射提升证明投射性也保持。并不要求简单模对应同一个底集合，或投射模仍表现为相同秩的自由模。

\subsection{半单性与有限长度}
\label{b2:subsec:1604:02}

\begin{proposition}\label{b2:prop:morita-invariants}
设 \(R,S\) 为含幺环，\(H:R\text{-}\mathbf{Mod}\rightarrow S\text{-}\mathbf{Mod}\) 为幺性左模范畴的等价，则 \(H\) 保持并反映模的简单性、投射性、内射性、不可分解性、半单性、有限生成性、有限展示性、Noether 性、Artin 性和有限长度。若幺性左 \(R\) 模 \(M\) 有限长度，则
\[
\ell_R(M)=\ell_S\bigl(H(M)\bigr).
\]
相应简单因子的合成重数也保持。此外，\(R\) 为左 Noether 环，当且仅当 \(S\) 为左 Noether 环；左 Artin 性也满足同样的等价。两个环还同时为半单环或同时不是半单环。
\end{proposition}
\begin{proof}
简单性和投射性已经证明，内射性由\cref{b2:lem:module-equivalence-structure}得到。一个模分解为两个非零模的直和，等价于存在相应的嵌入、投影及直和恒等式；这也是\cref{b2:prop:idempotent-direct-sum}中非平凡幂等自同态的投影刻画。等价双向保持这些映射关系和非零对象，故保持并反映不可分解性。半单模是简单模的某个直和；等价保持简单模与任意直和，所以保持半单性，对拟逆使用同样结论得到反映性。

由\cref{b2:prop:morita-submodule-lattices}，子模链的包含、严格性和稳定性双向保持，因此 Noether 性与 Artin 性保持。对一条合成列，等价的正合性使相邻商对应原简单因子的像；每个像仍简单，故得到相同层数的合成列。\cref{b2:thm:module-jordan-holder}使这个层数正是两边的长度，不同简单类型也由全忠实性一一对应，所以重数相同。对拟逆应用同样论证，有限长度也被反映。

有限生成性不能通过比较两边的生成元个数来证明，但可以从子模序直接判定：一个模 \(M\) 有限生成，当且仅当每个非空、向上有向且并为 \(M\) 的子模族中，都有一个成员等于 \(M\)。向上有向是指任意两个成员都包含在某个共同成员中。正向时，每个有限生成元属于族中某个成员，反复使用有向性便找到包含全部生成元的共同成员，它只能等于 \(M\)；若 \(M=0\)，任取一个成员即可。反向取 \(M\) 的所有有限生成子模，它们向上有向，因为两组有限生成元合在一起仍有限，且并为 \(M\)，因为每个元素都属于自己的循环子模。于是其中一个成员等于 \(M\)。子模序同构保持有向性和最小上界；对有向子模族，这个最小上界正是它们的并，因此保持这一判据。

有限展示还须检查全部关系能否有限生成。令 \(P=H({}_RR)\)，已证的有限生成性与投射性的保持使 \(P\) 为有限生成投射左 \(S\) 模。若 \(M\) 有有限自由展示 \(R^a\xrightarrow{d}R^b\twoheadrightarrow M\to0\)，其中 \(a,b\ge0\) 为整数，正合性与有限直和的保持给出
\[
P^a\xrightarrow{H(d)}P^b\twoheadrightarrow H(M)\longrightarrow0.
\]
由\cref{b2:prop:finite-projective-summand}，可取有限生成投射左 \(S\) 模 \(Q\)，使 \(P^b\oplus Q\cong S^t\)，其中 \(t\ge0\) 为整数。把 \(P^b\) 上的满射与 \(Q\) 上的零映射合起来，得到有限自由满射 \(S^t\to H(M)\)，其核在这个分解下为
\[
\operatorname{im}H(d)\oplus Q.
\]
第一项是有限生成模 \(P^a\) 的像，第二项也有限生成，故这个核有限生成。为核取一个有限自由满射，就得到 \(H(M)\) 的有限自由展示。对拟逆应用同样论证，有限展示性也被反映。

环的链条件还需要从 \(P\) 传到正则模 \({}_SS\)，因为两者未必同构。由\cref{b2:lem:module-equivalence-structure}，\(P\) 也是生成元，因而\cref{b2:prop:finite-projective-summand,b2:thm:generator-characterizations}分别给出有限直和分解
\[
P\oplus Q_1\cong S^u,\qquad
S\oplus Q_2\cong P^v,
\]
其中 \(u,v\ge0\) 为整数，\(Q_1,Q_2\) 为左 \(S\) 模。\cref{b2:thm:chain-conditions-exact}保证 Noether 性和 Artin 性对有限直和及子模传递。第一个分解使 \({}_SS\) 的链条件传到 \(P\)，第二个分解使 \(P\) 的链条件传到 \({}_SS\)。而子模格对应已说明 \(P=H(R)\) 与 \({}_RR\) 的链条件相同，所以两环同时为左 Noether 环，也同时为左 Artin 环。

最后，\cref{b2:thm:semisimple-ring-characterizations}说明环半单当且仅当它的每个左模半单。每个目标对象都同构于等价像中的对象，而半单性又被双向保持，所以两环的这一性质同时成立。这里不能只比较正则模的像，因为等价未必把一个正则模送到另一个正则模。
\end{proof}

有限长度并不要求两边的环都有限维；它由子模链定义，正是等价能保留的结构。例如，若 \(M\) 的两个不可分解直和分量合成长度分别为三和五，它们在等价后仍不可分解，长度仍分别为三和五：分解与非零性都可由投影和零对象辨认，不能因坐标数改变而改变这些长度。

\ProseParagraphBreak
中心同构也给出一个实用限制：两个交换环若 Morita 等价，则它们作为环同构，因为交换环就是自己的中心。反向，环同构当然给出改写标量的模范畴同构。因此对交换环，Morita 等价不会把不同的环合并；其更广的作用主要体现在非交换结构中。

\subsection{维数、生成元数与附加型}
\label{b2:subsec:1604:03}

设 \(k\) 为域、\(n>1\) 为整数，取 \(R=k\)、\(S=M_n(k)\)。它们 Morita 等价，但作为 \(k\) 代数的维数分别为一和 \(n^2\)，前者交换，后者不交换。标准等价把 \(k\) 模 \(V\) 送到列模 \(V^n\)，在同态上逐坐标作用，是 \(k\) 线性的。即使如此，若 \(\dim_kV=m<\infty\)，仍有
\[
\dim_kV=m,\qquad \dim_k(V^n)=nm.
\]
维数变化并没有违反长度不变性：\(V\) 是 \(m\) 个一维简单 \(k\) 模的直和，\(V^n\) 是 \(m\) 个简单 \(M_n(k)\) 列模的直和，故两边的模长度都为 \(m\)。

\ProseParagraphBreak
最少生成元的数值也可能改变。对 \(V=k^m\)，有 \(\mu_k(V)=m\)；在矩阵环一侧，\cref{b2:prop:semisimple-top-generators}给出
\[
\mu_{M_n(k)}(V^n)=\left\lceil\frac mn\right\rceil.
\]
因此等价保持“有限生成”这个性质，却不保持相对于正则模计算的最少生成元数。事实上，
\[
{}_{M_n(k)}M_n(k)\cong(k^n)^{\oplus n},
\]
所以一个正则模包含 \(n\) 个简单列模。一个生成元给出的正则模像至多覆盖 \(n\) 份简单列模；要覆盖 \(V^n\) 中的 \(m\) 份，最少需要 \(\lceil m/n\rceil\) 个生成元。

\ProseParagraphBreak
若 \(k\) 为 \(q\) 元域，则两个环的元素数分别为 \(q\) 和 \(q^{n^2}\)。模的元素从向量换成向量列，环元素从标量换成矩阵，都允许底集合发生变化。使用 Morita 等价时，应说明所比较的是环同构、模范畴结构、模长度还是指定底域上的维数；这些问题不能仅凭“等价”二字互相替换。

\begin{example}\label{b2:exm:morita-forgets-forms}
普通模范畴不记录一个模上额外指定的型。取 \(V=\mathbb R^2\)，在同一个实向量空间上分别指定矩阵为
\[
H_+=I_2,\qquad H_-=\operatorname{diag}(1,-1)
\]
的对称双线性型。它们的底层模完全相同，自同态代数也都是 \(M_2(\mathbb R)\)，但两型不等距：第二个型有非零各向同性向量 \((1,1)\)，第一个没有。
\end{example}

\cref{b2:prop:form-adjoint-operator}将这些型分别变为自同态代数上的伴随运算
\[
\sigma_+(X)=X^{\mathsf T},\qquad
\sigma_-(X)=H_-^{-1}X^{\mathsf T}H_-.
\]
它们对 \(E_{12}\) 的取值分别为 \(E_{21}\) 与 \(-E_{21}\)，故是不同运算。Morita 等价中的对象和态射没有包括 \(H_+\)、\(H_-\) 或这些伴随运算，因而不能仅凭模范畴等价，就声称还给出了带型对象之间的等价。若要比较后者，必须额外指定如何传送型、对偶及等距映射；普通 Morita 定理并未给出这些附加数据。

\ProseParagraphBreak
回到半单环，Artin–Wedderburn 分解中的矩阵大小可以在模范畴等价下改变，而简单模的同构类型仍一一对应。这给出了与除环等价的准确范围。

\begin{corollary}\label{b2:cor:morita-division-ring}
设 \(R\) 为非零含幺环。\(R\) Morita 等价于某个除环，当且仅当 \(R\) 是单 Artin 环，即 \(R\) 是单环且为左 Artin 环。
\end{corollary}
\begin{proof}
若 \(R\) 为单 Artin 环，\cref{b2:prop:simple-artinian}给出 \(R\cong M_n(D)\)，其中 \(D\) 为除环、\(n\ge1\)。由\cref{b2:prop:matrix-morita-equivalence}，这个矩阵环与 \(D\) Morita 等价。

反过来，设 \(R\text{-}\mathbf{Mod}\simeq D\text{-}\mathbf{Mod}\)，其中 \(D\) 为除环。每个幺性左 \(D\) 模都有基，因而是简单模 \({}_DD\) 的直和，且非零简单左 \(D\) 模只有这一种同构类型。\cref{b2:prop:morita-invariants}使每个幺性左 \(R\) 模也半单，并使简单左 \(R\) 模只有一种同构类型。由\cref{b2:thm:semisimple-ring-characterizations}，\(R\) 为半单环。再用\cref{b2:thm:artin-wedderburn}，将 \(R\) 写成有限个除环上矩阵环的乘积。不同矩阵因子的简单列模不同构，因为相应的中心幂等元在一个上作用为恒等，在另一个上作用为零，所以这里只有一个因子。由于 \(R\ne0\)，这个因子存在，由\cref{b2:prop:simple-artinian}即得 \(R\) 为单 Artin 环。
\end{proof}

只有一种简单模类型并不足够。设 \(k\) 为域，\(A=k[\varepsilon]/(\varepsilon^2)\)。对任意简单左 \(A\) 模 \(U\)，\(\varepsilon U\) 是子模；若它等于 \(U\)，则 \(U=\varepsilon U=\varepsilon^2U=0\)，矛盾，故 \(\varepsilon U=0\)。于是 \(U\) 是 \(k\) 上的简单模，同构于 \(k\)。但正则模 \(A\) 并不半单：商映射 \(A\to A/(\varepsilon)\) 若有模截面，一的像应为某个 \(1+c\varepsilon\)，并被 \(\varepsilon\) 零化；实际上 \(\varepsilon(1+c\varepsilon)=\varepsilon\ne0\)。所以这个商映射不能分裂，\(A\) 不与除环 Morita 等价。


\begin{chaptersummary}
生成元使每个模都成为其某个直和的商，有限生成使 Hom 保持任意直和，投射性使 Hom 正合。三者合起来，允许从正则对象上的伴随同构出发，经由展示和余核覆盖全部模。反方向，等价将正则模送到投射生成元，再由紧性恢复有限生成，所以这一构造恰好描述全部 Morita 等价。

\ProseParagraphBreak
自同态环的反环记录左模与右作用之间的乘法次序。互逆双模还给出相应的右模范畴等价。矩阵环等价和满幂等角环等价都能写成明确的矩阵单位及求值公式；满性保证没有一部分模在取角时被丢弃。基本角代数因而保留原有限维代数的整个模范畴。

\ProseParagraphBreak
中心可由恒等函子的自然自变换恢复，子模序则由单态射及其因子关系恢复。这些描述解释了简单性、投射性、内射性、不可分解性、半单性、有限生成性、链条件和长度为什么保持。有限展示性还须检查关系模的有限生成，环的链条件则须把正则模的像与另一侧正则模联系起来；这些结论已在正文中分别证明。与此同时，正则模未必对应正则模，底域维数、元素数和最少生成元数都可能改变。额外指定的型及伴随运算也不是普通模范畴所记录的结构，若要传送这些数据，必须另加相容条件。后续比较中心单代数时，将继续使用这种区分，而不把 Morita 等价误当作环同构。
\end{chaptersummary}

\BookExercises

\begin{exercise}\label{b2:ex:16-trace-generator}
设 \(k\) 为域、\(R=k[t]/(t^3)\)，以 \(t\) 表示剩余类。分别计算 \(R/(t)\)、\(R/(t^2)\) 及 \(P=R/(t)\oplus R/(t^2)\) 的迹理想，判断 \(P\) 是否为生成元。再判断 \(P\oplus R\) 是否为生成元。
\end{exercise}
\begin{solution}
对任意理想 \((a)\)，从 \(R/(a)\) 到 \(R\) 的同态由 \(1+(a)\) 的像 \(b\) 决定，条件恰为 \(ab=0\)。因而在当前交换环中，迹理想为 \(\operatorname{Ann}_R(a)\)。这里
\[
\operatorname{tr}_R(R/(t))=(t^2),\qquad
\operatorname{tr}_R(R/(t^2))=(t).
\]
每个从有限直和出发的同态由各分量同态相加，所以 \(\operatorname{tr}_R(P)=(t)\ne R\)。由\cref{b2:prop:trace-generator-criterion}，\(P\) 不是生成元。另一方面，\(P\oplus R\) 含正则模为直和因子，其迹包含 \(\operatorname{id}_R(R)=R\)，所以是生成元。
\end{solution}

\begin{exercise}\label{b2:ex:16-object-detection}
设 \(k\) 为域、\(n\ge2\)，令 \(R=k[t]/(t^n)\)、\(P=R/(t)\)。对任意左 \(R\) 模 \(M\)，识别 \(\operatorname{Hom}_R(P,M)\)，并求从全部 \(P\) 副本到 \(M\) 的求值映射之像。由此说明 \(P\) 能检测非零对象却不能检测所有非零同态，并刻画哪些 \(M\) 是 \(P\) 某个直和的商。
\end{exercise}
\begin{solution}
在 \(1+(t)\) 处求值得到
\[
\operatorname{Hom}_R(P,M)\cong\{m\in M:tm=0\}.
\]
这也是所有求值像的和，因为每个被 \(t\) 消去的元素都给出一个同态。若 \(M\ne0\)，取 \(x\ne0\)，并取最小正整数 \(j\) 使 \(t^jx=0\)；它存在，因为 \(t^nM=0\)。则 \(t^{j-1}x\ne0\) 被 \(t\) 消去，故 Hom 非零。然而 \(R\) 上乘 \(t\) 是非零同态，却在所有来自 \(P\) 的像上为零。这是\cref{b2:prop:object-detection-not-generator}在较高幂零次数中的同一现象。

\(P\) 的任意直和都被 \(t\) 消去，其商也如此。反过来，若 \(tM=0\)，则 \(M\) 是一个 \(k\) 向量空间；取一组基便得到若干 \(P\cong k\) 的直和与 \(M\) 的同构。因此题述商恰为被 \(t\) 消去的模。
\end{solution}

\begin{exercise}\label{b2:ex:16-projective-not-generator}
设 \(k\) 为域、\(R=k\times k\times k\)，令 \(P=k^2\oplus k\oplus0\)，三个因子分别在对应分量上作用。取第一分量的行坐标，将 \(\operatorname{End}_R(P)^{\mathrm{op}}\) 识别为 \(S=M_2(k)\times k\)。对 \(M=M_1\oplus M_2\oplus M_3\)，写出 \(\operatorname{Hom}_R(P,M)\) 的两个 \(S\) 分量，并求余单位 \(P\otimes_S\operatorname{Hom}_R(P,M)\to M\) 的像。判断哪些模被 Hom 函子送到零。
\end{exercise}
\begin{solution}
三个中心幂等元迫使同态分别作用于对应分量，故
\[
\operatorname{Hom}_R(P,M)\cong M_1^2\oplus M_2
\]
作为左 \(M_2(k)\times k\) 模，第一项为矩阵乘列，第二项为通常的标量作用。第一项的行列求值由\cref{b2:prop:matrix-row-column-bimodules}恢复 \(M_1\)，第二项恢复 \(M_2\)，所以余单位的像恰为 \(M_1\oplus M_2\oplus0\)。第三项不能由 \(P\) 产生；所有 \(0\oplus0\oplus M_3\) 都被送到零，且只有这些模被送到零。\(P\) 是有限生成投射模，但不是生成元，说明增加第一分量的副本并不能弥补缺失的第三分量。
\end{solution}

\begin{exercise}\label{b2:ex:16-generator-not-projective}
设 \(k\) 为域、\(R=k[t]/(t^2)\)、\(P=R\oplus R/(t)\)，令 \(q:R\to R/(t)\) 为商映射。证明 \(P\) 有限生成且为生成元，计算
\[
\operatorname{Hom}_R(P,R)\longrightarrow\operatorname{Hom}_R(P,R/(t))
\]
的像及两个 Hom 空间的 \(k\) 维数，并由此判断投射性。
\end{exercise}
\begin{solution}
\(P\) 有限生成，且含 \(R\) 为直和因子，所以是生成元。第一分量的同态由任意 \(r\in R\) 决定，第二分量的像必须被 \(t\) 消去，故
\[
\operatorname{Hom}_R(P,R)\cong R\oplus kt,\qquad
\operatorname{Hom}_R(P,R/(t))\cong k\oplus k.
\]
维数分别为三和二。后复合 \(q\) 把 \((r,ct)\) 送到 \((r+(t),0)\)，其像只有第一项，维数为一。特别地，从 \(P\) 到第二项 \(R/(t)\) 的投影不能沿 \(q\) 提升。因此 Hom 不保持这一满射，\(P\) 不投射。有限生成和生成元性质并不能替代提升条件。
\end{solution}

\begin{exercise}\label{b2:ex:16-projective-generator-not-compact}
设 \(k\) 为域、\(P=\bigoplus_{j\ge1}ke_j\)。用无限矩阵描述
\[
\bigoplus_{i\ge1}\operatorname{Hom}_k(P,k)
\longrightarrow\operatorname{Hom}_k\left(P,\bigoplus_{i\ge1}k\right)
\]
的像，说明“每列只有有限个非零条目”为什么不足以保证来自左边。
\end{exercise}
\begin{solution}
目标中的同态由各 \(e_j\) 的像确定，故对应每列只有有限个非零条目的矩阵 \((a_{ij})\)。左边是有限多个目标分量上的同态之和，所以对应的矩阵只有有限个非零行；每一行可以有无限多个非零条目。反过来，只有有限个非零行的矩阵确实给出左边的元素。

单位矩阵每列恰有一个非零条目，却有无限多个非零行，因而不在像中。这里有限支集可以随输入的 \(e_j\) 变化，紧性要求的是全像共有的有限目标支集。\(P\) 自由而且含一份 \(k\)，所以投射且为生成元；失败的正是有限生成性。
\end{solution}

\begin{exercise}\label{b2:ex:16-full-idempotent-detects}
设 \(k\) 为域、\(R=T_2(k)\)，分别取 \(e=E_{11}\)、\(f=E_{22}\)。计算 \(eRe\)、\(fRf\)、\(ReR\)、\(RfR\)，并求各自正则模上的求值映射之像。对每个角给出一个被送到零的非零简单模。
\end{exercise}
\begin{solution}
两个角都同构于 \(k\)，但
\[
ReR=kE_{11}\oplus kE_{12},\qquad
RfR=kE_{12}\oplus kE_{22},
\]
都是真理想。求值映射 \(Re\otimes_{eRe}eR\to R\) 的像是 \(ReR\)：纯张量的像为乘积 \(reb\)，其和恰为这个理想；另一角同理。这两张量源各有 \(k\) 维数二，两个矩阵单位的乘积分别给出所示理想的基，故这两个求值映射均单射，却不满射。

两个一维简单模 \(S_1,S_2\) 分别由第一、第二对角条目作用。\(eS_2=0\)，而 \(fS_1=0\)。所以即使角环与标量域相同、求值映射还单射，未满的角仍会遗漏非零对象。
\end{solution}

\begin{exercise}\label{b2:ex:16-rectangular-morita}
取 \(R=M_3(k)\)、\(e=E_{11}+E_{22}\)。证明 \(eRe\cong M_2(k)\)，给出满性的一次有限和表达，并写出对任意 \(M\) 的求值映射 \(Re\otimes_{eRe}eM\rightarrow M\) 的逆。在自然列模 \(k^3\) 上，这个等价得到哪个 \(M_2(k)\) 模？
\end{exercise}
\begin{solution}
角环是左上二阶矩阵块，单位为 \(e\)，所以同构于 \(M_2(k)\)。等式
\[
1=e+E_{31}eE_{13}
\]
证明 \(e\) 满。应用\cref{b2:prop:full-idempotent-morita}及其逆式\eqref{b2:eq:full-idempotent-evaluation-inverse}，代入 \((a_1,b_1)=(1,1)\)、\((a_2,b_2)=(E_{31},E_{13})\)，求值映射的逆为
\[
m\longmapsto e\otimes em+E_{31}\otimes E_{13}m.
\]
这里 \(E_{31}\in Re\)、\(E_{13}m\in eM\)；正向求值后得到 \(em+E_{33}m=m\)，反向复合由角环平衡关系核验。对 \(M=k^3\)，\(eM\) 为前两个坐标组成的自然列模 \(k^2\)。两个简单模的基域维数分别为三和二，但模长度都为一。
\end{solution}

\begin{exercise}\label{b2:ex:16-matrix-hom}
设 \(R\) 为含幺环、\(n\ge1\)，\(M,N\) 为幺左 \(R\) 模，令 \(A=M_n(R)\)。证明每个左 \(A\) 模同态 \(u:M^n\rightarrow N^n\) 都唯一地由一个左 \(R\) 模同态 \(h:M\rightarrow N\) 逐坐标作用得到。
\end{exercise}
\begin{solution}
把 \(m\in M\) 放到第一坐标，因 \(u\) 与 \(E_{11}\) 交换，其像只可能在第一坐标非零，记该坐标为 \(h(m)\)。与 \(rE_{11}\) 交换说明 \(h(rm)=r\cdot h(m)\)，加法相容也直接成立，所以 \(h\) 左线性。第 \(i\) 坐标向量由 \(E_{i1}\) 作用于第一坐标向量得到，因此
\[
u(0,\ldots,m,\ldots,0)=\bigl(0,\ldots,h(m),\ldots,0\bigr).
\]
有限加法得到 \(u(m_i)_i=\bigl(h(m_i)\bigr)_i\)。反向逐坐标的 \(h\) 保持矩阵作用，故确为 \(A\) 模同态；在第一坐标求值又保证唯一性。
\end{solution}

\begin{exercise}\label{b2:ex:16-row-column-evaluation}
设 \(B\) 为任意环，令 \(A=M_2(B)\)、\(S=M_3(B)\)，取矩形矩阵双模 \({}_AP_S=M_{2\times3}(B)\)、\({}_SQ_A=M_{3\times2}(B)\)。证明矩阵乘法给出双模同构
\[
P\otimes_SQ\cong A,\qquad Q\otimes_AP\cong S,
\]
并在矩阵单位上写出逆。
\end{exercise}
\begin{solution}
用 \(p_{i\ell}(b)\) 表示仅第 \((i,\ell)\) 项为 \(b\) 的 \(P\) 元素，\(q_{\ell j}(b)\) 类似。矩阵乘法结合律给出所需平衡性。第一映射的逆为
\[
(a_{ij})\longmapsto\sum_{i,j=1}^2p_{i1}(a_{ij})\otimes q_{1j}(1).
\]
确实，在 \(S\) 上移过 \(E_{1\ell}\)，再移过 \(cE_{11}\)，得到
\[
p_{i\ell}(b)\otimes q_{kj}(c)
=\delta_{\ell k}\,p_{i1}(bc)\otimes q_{1j}(1).
\]
每个纯张量可由这类矩阵单位张量相加，所以这个公式同时验证两个复合为恒等。第二映射的逆为
\[
(s_{\ell k})\longmapsto\sum_{\ell,k=1}^3q_{\ell1}(s_{\ell k})\otimes p_{1k}(1),
\]
在 \(A\) 上使用同样的平衡计算即可。所列映射保持外侧矩阵作用；系数始终按 \(bc\) 的次序相乘，不要求 \(B\) 交换。此题把\cref{b2:prop:matrix-row-column-bimodules}的行列模型推广到两个不同的矩阵大小。
\end{solution}

\begin{exercise}\label{b2:ex:16-dual-bimodule-pair}
设 \(R\) 为任意环、\(u\in R\)，在行列双模模型 \(P=R^{1\times2}\)、\(Q=R^{2\times1}\)、\(S=M_2(R)\) 中取
\[
C=\begin{pmatrix}1&u\\0&1\end{pmatrix}.
\]
若行坐标改为 \(\widetilde p=pC\)，求保持 \(pq\) 的列坐标变换及与右作用相容的矩阵坐标变换。对 \(B=\begin{psmallmatrix}a&b\\c&d\end{psmallmatrix}\) 写出新矩阵，核对两个双模配对。
\end{exercise}
\begin{solution}
因为 \(C^{-1}=\begin{psmallmatrix}1&-u\\0&1\end{psmallmatrix}\)，取
\[
\widetilde q=C^{-1}q,\qquad
\widetilde B=C^{-1}BC
=\begin{pmatrix}
a-uc&au+b-ucu-ud\\
c&cu+d
\end{pmatrix}.
\]
则 \((pB)C=(pC)\widetilde B\)，所以右作用被正确搬到新坐标。两种配对分别满足
\[
\widetilde p\,\widetilde q=pq,\qquad
\widetilde q\,\widetilde p=C^{-1}(qp)C.
\]
前者在 \(R\) 中保持原值，后者在 \(S\) 中随内自同构改变坐标。与外侧作用相容后，这正是\cref{b2:prop:matrix-row-column-bimodules}的两个张量同构的坐标变换。若 \(R\) 非交换，\(au\)、\(ucu\)、\(ud\) 等因子的次序均不可交换。
\end{solution}

\begin{exercise}\label{b2:ex:16-matrix-center}
设 \(R\) 为含幺环、\(n\ge1\)。直接用矩阵单位证明 \(Z\bigl(M_n(R)\bigr)=\bigl\{\,cI_n:c\in Z(R)\,\bigr\}\)。把这个显式同构与 Morita 中心不变量比较。
\end{exercise}
\begin{solution}
与各 \(E_{ii}\) 交换使中心矩阵的非对角条目全为零；再与 \(E_{ij}\) 交换，使全部对角元相同，记为 \(c\)。与在某个对角位置放任意 \(r\in R\) 的矩阵交换，得到 \(cr=rc\)，所以 \(c\in Z(R)\)。反向，中心元素乘单位矩阵与所有矩阵交换。因此 \(c\mapsto cI_n\) 是中心之间的环同构，正体现矩阵环与原环虽可能维数不同，中心仍相同。
\end{solution}

\begin{exercise}\label{b2:ex:16-center-obstruction}
设 \(k\) 为域。证明 \(M_2\bigl(k[\varepsilon]/(\varepsilon^2)\bigr)\) 与 \(M_2(k\times k)\) 不 Morita 等价，尽管它们作为 \(k\) 向量空间的维数相同。
\end{exercise}
\begin{solution}
两个环的 \(k\) 维数都为八。其中心分别为 \(k[\varepsilon]/(\varepsilon^2)\) 与 \(k\times k\)。前者有非零幂零元 \(\varepsilon\)，后者任何幂零元素的两个域坐标都必须为零，所以没有非零幂零元。因此中心不同构，由\cref{b2:prop:morita-center}，两个环不 Morita 等价。单独比较代数维数不能判断等价。
\end{solution}

\begin{exercise}\label{b2:ex:16-dimension-length-generators}
设 \(k\) 为域。在 \(k\) 与 \(S=M_3(k)\) 的标准等价中，求 \(V=k^7\) 的像的 \(k\) 维数、\(S\) 长度与最少 \(S\) 生成元数。该像是否为自由 \(S\) 模？
\end{exercise}
\begin{solution}
像为 \(V^3\cong(k^3)^7\)，其 \(k\) 维数为二十一。自然列模 \(k^3\) 简单，所以 \(S\) 长度为七，与 \(V\) 的 \(k\) 模长度相同。在\cref{b2:prop:semisimple-top-generators}中代入矩阵块大小三及简单模重数七，最少生成元数为 \(\lceil7/3\rceil=3\)，而原 \(k\) 模需要七个生成元。

它不是自由 \(S\) 模。有限秩自由 \(S\) 模的 \(k\) 维数为九的整数倍，二十一不是；无限秩自由模的 \(k\) 维数无限，也不可能。它仍然投射，因为半单环上的每个模投射，或因为等价保持原自由 \(k\) 模的投射性。
\end{solution}

\begin{exercise}\label{b2:ex:16-indecomposable-transport}
设 \(k\) 为域，令 \(R=k[t]\)、\(A=M_2(R)\)，把 \(M=R/(t^2)\oplus R/(t^3)\) 通过标准等价送到 \(A\) 模。求所得模的不可分解分量个数、长度及 \(k\) 维数，并解释为何不能从底 \(R\) 模的直和个数判断 \(A\) 不可分解性。
\end{exercise}
\begin{solution}
由\cref{b2:thm:module-correspondence}，\(R/(t^j)\) 的子模对应包含 \((t^j)\) 的理想，恰为 \((t^\ell)/(t^j)\)、\(0\le\ell\le j\)。这些子模构成一条链，相邻商均为 \(k\)，所以原循环模不可分解，长度分别为二和三。\cref{b2:prop:matrix-morita-equivalence}给出所用等价；全忠实性和\cref{b2:lem:module-equivalence-structure}的加性给出自同态环同构。由\cref{b2:prop:idempotent-direct-sum}，非平凡幂等元的不存在保证两者的像仍不可分解。再由\cref{b2:prop:morita-invariants}保持长度，像为两个不可分解 \(A\) 模的直和，长度为五。

底 \(k\) 维数乘二，故为十。每个像在忘掉矩阵作用后形如 \(\bigl(R/(t^j)\bigr)^2\)，但两个坐标会被矩阵单位相互送入；单个坐标不是 \(A\) 子模，不能将其当作 \(A\) 模的两个直和因子。
\end{solution}

\begin{exercise}\label{b2:ex:16-ring-chain-invariants}
设 \(k\) 为域、\(R=k[t]/(t^4)\)、\(A=M_3(R)\)。证明两环都是左 Noether 环和左 Artin 环，计算两个正则左模的合成长度。说明这是否与 Morita 等价保持模长度矛盾。
\end{exercise}
\begin{solution}
\(R\) 的理想为 \((t^j)\)、\(0\le j\le4\)，相邻商为 \(k\)，所以 \(\ell_R(R)=4\)。它满足两种链条件，由\cref{b2:prop:morita-invariants}及矩阵环等价，\(A\) 也满足两种链条件。

令 \(J=M_3((t))\)。\(J^j=M_3((t^j))\)，故
\[
A\supset J\supset J^2\supset J^3\supset0
\]
的四个层商各为 \(M_3(k)\) 的正则左模，分别是三个简单列模的直和。长度可加给出 \(\ell_A(A)=12\)。这没有矛盾：\cref{b2:prop:matrix-morita-equivalence}把 \(A\) 送到 \(R^3\)，其 \(R\) 长度正为十二；它把 \(R\) 送到自然列模 \(R^3\)，该模的 \(A\) 长度为四。等价保持对应对象的长度，两侧正则模却不是对应对象。
\end{solution}

\begin{exercise}\label{b2:ex:16-finite-presentation}
在 \(\mathbb Z\) 与 \(S=M_2(\mathbb Z)\) 的标准等价中，取 \(M=\mathbb Z/2\mathbb Z\oplus\mathbb Z/3\mathbb Z\)。将其像识别为 \((\mathbb Z/6\mathbb Z)^2\)，并写出一个只用有限自由 \(S\) 模的展示，明确检查全部关系。
\end{exercise}
\begin{solution}
中国剩余同构把 \((\bar a,\bar b)\) 送到 \(3a+4b\pmod6\)，逐坐标使用后得到题述识别。令 \(e=E_{11}\)，定义
\[
q:S\longrightarrow(\mathbb Z/6\mathbb Z)^2,\qquad
q(B)=\text{\(B\) 的第一列模六}.
\]
这是满的左 \(S\) 同态。其核中的第一列每项被六整除，第二列任意，因此
\[
\ker q=S(6e)+S(1-e).
\]
所以完整的有限自由展示为
\[
S^2\xrightarrow{\ d\ }S\xrightarrow{\ q\ }
(\mathbb Z/6\mathbb Z)^2\longrightarrow0,\qquad
d(B,C)=6Be+C(1-e).
\]
这里等式 \(\operatorname{im}d=\ker q\) 检查了全部关系，而不只是列出若干已知关系。\cref{b2:prop:morita-invariants}保证有限展示性保持，本题给出原坐标中的具体实现。
\end{solution}

\begin{exercise}\label{b2:ex:16-compose-bimodule-equivalences}
设 \(R,S,T\) 为含幺环，幺性双模 \({}_RP_S\)、\({}_SQ_T\) 分别通过张量实现从全部幺左 \(S\) 模到全部幺左 \(R\) 模、从全部幺左 \(T\) 模到全部幺左 \(S\) 模的等价。证明复合等价由 \({}_R(P\otimes_SQ)_T\) 实现，并说明该左 \(R\) 模为何为有限生成投射生成元。
\end{exercise}
\begin{solution}
对左 \(T\) 模 \(N\)，\cref{b2:thm:tensor-associativity-unit}给出自然同构
\[
P\otimes_S(Q\otimes_TN)\cong(P\otimes_SQ)\otimes_TN,
\qquad p\otimes(q\otimes n)\longmapsto(p\otimes q)\otimes n.
\]
因此复合确实由题述双模张量实现，作用次序不能颠倒。复合仍为范畴等价，而它将正则模 \(T\) 送到 \((P\otimes_SQ)\otimes_TT\cong P\otimes_SQ\)。由\cref{b2:thm:morita-converse}，正则模在等价下的这个像必为有限生成投射生成元，并有自同态反环同构于 \(T\)。
\end{solution}

\begin{exercise}\label{b2:ex:16-equivalent-to-division-ring}
设 \(D,E\) 为除环，\(m,n\ge1\)。证明 \(M_n(D)\) 与 \(M_m(E)\) Morita 等价，当且仅当 \(D\cong E\)。说明如何从唯一简单模的自同态环恢复除环，而不靠比较矩阵大小。
\end{exercise}
\begin{solution}
每个矩阵环只有一种简单左模类型，分别为自然列模 \(D^n\)、\(E^m\)。任意模范畴等价保持简单对象及其自同态环，所以给出
\[
\operatorname{End}_{M_n(D)}(D^n)
\cong\operatorname{End}_{M_m(E)}(E^m).
\]
与矩阵单位交换迫使自同态逐坐标相同；与任意 \(dE_{11}\) 交换，又使这一坐标映射为右乘某个除环元素。令 \(r_a\) 表示逐坐标右乘 \(a\)，则 \(r_a\circ r_b=r_{ba}\)，故这两个自同态环分别为 \(D^{\mathrm{op}}\)、\(E^{\mathrm{op}}\)。取反环得到 \(D\cong E\)。反向由矩阵环与原环的等价及改写标量得到。

因此\cref{b2:cor:morita-division-ring}中的除环同构类型可由简单对象恢复，而矩阵大小并不能由普通模范畴恢复。
\end{solution}

\begin{exercise}\label{b2:ex:16-full-corner-ideals}
设 \(k\) 为域、\(B=k[t]\)、\(R=M_3(B)\)、\(e=E_{11}+E_{22}\)。利用满角环的理想对应，求 \(R\) 的全部双边理想；计算 \(I=M_3((t^2))\) 对应的角理想及两个商环。
\end{exercise}
\begin{solution}
先对任意 \(n\ge1\) 考察 \(M_n(B)\) 的双边理想 \(L\)。若 \(X=(x_{ij})\in L\)，则
\[
E_{ai}XE_{jb}=x_{ij}E_{ab}\in L.
\]
因而各条目构成同一个 \(B\) 理想 \(\mathfrak a\)，且 \(L=M_n(\mathfrak a)\)；反向这类矩阵确为双边理想。\(B\) 是主理想整环，所以 \(R\) 的全部双边理想为 \(M_3((g))\)，其中 \(g\in k[t]\)，取零或首一多项式可消除单位倍的不唯一性。

这里 \(1=e+E_{31}eE_{13}\)，故 \(e\) 满，角环为 \(M_2(B)\)。\cref{b2:prop:full-corner-ideal-correspondence}把 \(M_3((g))\) 送到 \(M_2((g))\)，反向由 \(RJR\) 恢复。对题设 \(I\)，两个商环为
\[
R/I\cong M_3(k[t]/(t^2)),\qquad
eRe/eIe\cong M_2(k[t]/(t^2)).
\]
它们的矩阵大小不同，但满角在取商后仍满，故两商环 Morita 等价。
\end{solution}

\begin{exercise}\label{b2:ex:16-symmetric-alternating-forgotten}
在 \(V=\mathbb R^2\) 上，分别取矩阵为 \(I_2\) 的对称型与矩阵为 \(J=\begin{psmallmatrix}0&1\\-1&0\end{psmallmatrix}\) 的交替型。求两种伴随运算在 \(M_2(\mathbb R)\) 中的不动子空间维数，并说明为什么普通模范畴没有记录这一区别。
\end{exercise}
\begin{solution}
第一种不动条件是 \(X^{\mathsf T}=X\)，对称矩阵由两个对角条目和一个非对角条目确定，维数为三。第二种条件为 \(J^{-1}X^{\mathsf T}J=X\)。若 \(X=\begin{psmallmatrix}a&b\\c&d\end{psmallmatrix}\)，左边为 \(\begin{psmallmatrix}d&-b\\-c&a\end{psmallmatrix}\)，故不动条件迫使 \(a=d,b=c=0\)，只有标量矩阵，维数为一。底层实模和自同态环没有改变，改变的是另行指定的反转乘法运算；普通模范畴没有把它作为结构的一部分。
\end{solution}

\begin{exercise}\label{b2:ex:16-ideal-autoequivalence}
沿用\cref{b2:ex:krull-schmidt-nonunique}，令 \(s=\sqrt{-5}\)、\(R=\mathbb Z[s]\)、\(I=(2,1+s)\)。证明 \(I\) 是生成元且 \(\operatorname{End}_R(I)\cong R\)，从而给出 \(R\) 模范畴的一个 Morita 自等价。不得因 \(I\) 非自由便认为它不能作为生成元。
\end{exercise}
\begin{solution}
\cref{b2:ex:krull-schmidt-nonunique}已证明 \(I^2=2R\) 及 \(I\oplus I\cong R^2\)，所以由\cref{b2:prop:finite-projective-summand}，\(I\) 有限生成且投射。两个 \(R\) 线性映射 \(f_1(x)=2x\)、\(f_2(x)=(s-1)x/2\) 都从 \(I\) 取值于 \(R\)：第二个的分子属于 \(I^2=2R\)。并且
\[
f_1(2)+f_2(1+s)=4+\frac{s^2-1}{2}=1.
\]
所以 \((x,y)\mapsto f_1(x)+f_2(y)\) 给出 \(I\oplus I\) 到 \(R\) 的满同态，其截面为 \(r\mapsto\bigl(2r,(1+s)r\bigr)\)。因此 \(R\) 是 \(I\oplus I\) 的直和因子，\cref{b2:thm:generator-characterizations}说明 \(I\) 为生成元。

对任意 \(R\) 线性 \(f:I\rightarrow I\)，等式 \(2f(x)=f(2x)=x\cdot f(2)\) 说明它在分式域中为乘 \(c=f(2)/2\)。因 \(cI\subseteq I\)，乘以 \(I\) 得 \(cI^2\subseteq I^2\)，即 \(2cR\subseteq2R\)，消去非零整数二得 \(c\in R\)。反向每个 \(c\in R\) 都给出自同态，且这个对应保持加法、乘法与单位，所以 \(\operatorname{End}_R(I)\cong R\)。

\(R\) 交换，反环仍为 \(R\)。因为 \(I\) 是有限生成投射生成元，由\cref{b2:thm:morita-equivalence}，\(\operatorname{Hom}_R(I,-)\) 与 \(I\otimes_R-\) 通过上述环识别互为等价。这个生成元非自由，正说明自等价可以由非自由投射生成元实现。
\end{solution}

% !TeX root = ../../Book2.tex
% 纲要标题：## 第17章　中心单代数
% 纲要位置：计划纲要.md，第 1594 行。

\chapter{中心单代数}
\label{b2:ch:17}
\BookChapterNumber{b2:ch:17}

\begin{chapterabstract}
中心单代数在适当扩张标量后成为矩阵代数，而它在原域上可能保留非交换除代数的结构。本章以正则双模和稠密性定理建立包络代数判据，证明任意标量扩张和张量积保持中心单性。中心化子定理给出简单子代数的维数关系，Skolem–Noether 定理比较它的不同嵌入。通过在除代数中寻找可分最大子域，建立有限 Galois 分裂域的存在；四元数的显式范数与分裂方程则为下一章的 Brauer 群提供可计算的例子。
\end{chapterabstract}

\begin{learninggoals}
\begin{itemize}
\item 区分单性、中心性与分裂性，证明中心单代数的包络代数判据及次数平方公式。
\item 证明标量扩张、下降与张量积结论，说明中心化子定理在子代数中心不可分时仍然成立。
\item 分清包含极大子域、最大交换子代数与分裂域，构造除代数的可分最大子域并证明有限 Galois 分裂域存在。
\item 用两个模结构证明 Skolem–Noether 定理，辨认其假设，并将四元数的分裂问题化为范数方程。
\end{itemize}
\end{learninggoals}

实四元数代数 \(\mathbb H\) 中，每个非零元素都有逆，因此不可能同构于 \(M_2(\mathbb R)\)：后者有非零而不可逆的矩阵。可是把标量扩到复数，\(\mathbb H\otimes_{\mathbb R}\mathbb C\) 却成为 \(M_2(\mathbb C)\)。同一个代数在原域上保留除法结构，在较大域上呈现为矩阵代数。由扩域后的形状倒推原代数，必须记录底域中的哪些信息仍可恢复。

\ProseParagraphBreak

没有非零真双边理想是单性，中心恰为指定底域是中心性，同构于底域上的全矩阵代数则是分裂性。三者要分开判断：非平凡有限域扩张作为基域代数是单的，中心却更大；实四元数在实数上中心单，仍不在实数上分裂。若其中嵌入一个简单子代数，它与自己的中心化子共同限制了整个代数的维数。进一步若要找到使代数分裂的扩域，就要比较子域、不同嵌入和标量扩张的作用；四元数提供了可以手算的检验场所。

\ProseParagraphBreak

本章需要第十一章的包络代数与双中心化子、第十三章的矩阵环分类和简单模重数，以及第十四章 Jacobson 稠密性定理的有限维结论。有关分裂域的论证使用第一卷定理15.2.3的代数闭包存在性，以及可分多项式和有限 Galois 扩张的基本性质。任意标量扩张保持中心单性不要求可分；有限可分分裂域的存在则需要另行构造可分最大子域。

\ProseParagraphBreak
中心单代数、中心化子与分裂域的系统论述可对照 Gille–Szamuely 的《Central Simple Algebras and Galois Cohomology》\cite[Chapter 2, pp. 19--69]{GilleSzamuely2017CentralSimpleAlgebras}。除代数的基础及实四元数可对照 Roman 的《Advanced Linear Algebra》\cite[Chapter 18, Division Algebras and The Quaternions, pp. 462--463; Theorem 18.10, p. 463]{Roman2008AdvancedLinearAlgebra}。关于有限维实除代数的进一步分类，该书\cite[Theorem 18.12, pp. 466--469]{Roman2008AdvancedLinearAlgebra}给出 Frobenius 定理。本章先建立适用于任意底域的中心单结构与嵌入理论，再在特征不为二时展开四元数的具体计算。

% !TeX root = ../../Book2.tex
% 纲要标题：### 17.1 单代数与中心
% 纲要位置：计划纲要.md，第 1596 行。

\section{单代数与中心}
\label{b2:sec:1701}

半单结构定理把有限维代数中的单块写成除环上的矩阵环，但底域未必就是这些块的中心。中心中的元素同全部乘法相容，所以每个单块也自然成为其中心上的代数。研究一个单块扩大标量后的结构时，就要先确定这个标量域：中心单代数要求它恰为全部中心，接下来将证明这一条件使任意扩域后的代数仍然单。张量积和中心化子的结构也将以此为基础。本章的代数结合且含单位元，子代数与嵌入保持单位元；中心单代数一律非零且有限维。

% 纲要要点（供目录核对）：
%
% * 有限维中心单代数；固定底域与实际中心，包络判据核对稠密性所用右除环及全线性算子目标
% * 分裂代数与分裂域；区别给定映射的同构检测与扩域后新同构的存在，核对有限维自同态的标量扩张、基比较下降及次数平方，正式证明纯不可分扩域后的非单性反例
% * 张量积与反代数；用两个包络作用证明张量积中心单，区别一般单代数的非单张量积
%

\subsection{有限维中心单代数}
\label{b2:subsec:1701:01}

\begin{definition}\label{b2:def:central-simple-algebra}
设 \(k\) 为域，\(A\) 为非零有限维含幺结合 \(k\) 代数。若 \(A\) 没有非零真双边理想，且其中心满足 \(Z(A)=k1_A\)，则称 \(A\) 为 \(k\) 上的\term[zhongxin-dan-daishu]{中心单代数}[central simple algebra]。下文用结构嵌入把 \(k1_A\) 识别为 \(k\)。
\end{definition}

单性限制的是双边理想，中心性还要求指定的底域就是全部中心。例如，有限域扩张 \(L/k\) 作为 \(k\) 代数是单的，但中心为 \(L\)，只有 \(L=k\) 时才中心单。对 \(n\ge1\)，矩阵代数 \(M_n(k)\) 则中心单：单性由\cref{b2:prop:simple-artinian}证明，中心为标量矩阵是\cref{b2:prop:matrix-centers-central-idempotents}在除环为 \(k\) 时的结论。

\ProseParagraphBreak
底域的选择确实会改变这个判断。例如 \(M_2(\mathbb C)\) 作为实代数是单的，中心却是 \(\mathbb C I_2\)，大于 \(\mathbb R I_2\)；它的实维数为 \(2\cdot2^2=8\)。改以 \(\mathbb C\) 为底域后，同一个乘法代数成为中心单代数，维数为 \(2^2=4\)。因此“中心单”要连同底域一起说明；后面得到的维数平方公式也只对这个指定的中心域成立。

\ProseParagraphBreak
一般非零含幺单代数 \(B\) 的中心是域，这不需要有限维假设。若 \(0\ne z\in Z(B)\)，则 \(Bz\) 是非零双边理想，故等于 \(B\)，从而 \(bz=1\) 对某个 \(b\) 成立。由 \(z\) 中心可知 \(zb=1\)，且从 \(zx=xz\) 可推出 \(z^{-1}x=xz^{-1}\)。所以 \(z\) 在中心内可逆。若 \(B\) 还是有限维 \(k\) 代数，\(Z(B)\) 作为 \(B\) 的 \(k\) 子空间也有限维，因而是 \(k\) 的有限扩张。

\ProseParagraphBreak
Artin–Wedderburn 定理给出 \(A\cong M_r(D)\) 时，除环 \(D\) 的中心未必就是原底域 \(k\)。中心单条件进一步固定这个参数，要求 \(Z(D)=k\)。

\begin{proposition}\label{b2:prop:csa-division-description}
设 \(k\) 为域，\(A\) 为中心单 \(k\) 代数，则存在整数 \(r\geq1\) 与有限维中心除 \(k\) 代数 \(D\)，使 \(A\cong M_r(D)\)。反过来，对任意这样的 \(r,D\)，矩阵代数 \(M_r(D)\) 中心单。整数 \(r\) 与 \(D\) 的 \(k\) 代数同构类型由 \(A\) 唯一确定。
\end{proposition}
\begin{proof}
有限维使左理想降链稳定，单性与\cref{b2:prop:simple-artinian}给出 \(A\cong M_r(D)\)。这个构造保持中心中的 \(k\) 标量，因此是 \(k\) 代数同构，\(D\) 有限维。由\cref{b2:prop:matrix-centers-central-idempotents}，\(Z\bigl(M_r(D)\bigr)=Z(D)I_r\)，而 \(A\) 的中心恰为 \(k\)，于是 \(Z(D)=k\)。反向的单性由\cref{b2:prop:simple-artinian}中矩阵环的结论得到，中心由同一中心公式给出。唯一性是\cref{b2:prop:semisimple-multiplicity}对正则模及其自同态除环的唯一性，构造同时保留 \(k\) 作用。
\end{proof}

\begin{theorem}[包络代数判据]\label{b2:thm:csa-enveloping-isomorphism}
设 \(k\) 为域，\(A\ne0\) 为有限维含幺结合 \(k\) 代数，则 \(A\) 中心单，当且仅当左右乘作用给出同构
\begin{equation}\label{b2:eq:csa-enveloping-isomorphism}
\Theta_A:A\otimes_kA^{\mathrm{op}}\longrightarrow\operatorname{End}_k(A),
\qquad a\otimes b^{\mathrm{op}}\longmapsto(x\mapsto axb).
\end{equation}
\end{theorem}
\begin{proof}
\cref{b2:thm:enveloping-bimodules}已证明 \(\Theta_A\) 是代数同态。若 \(A\) 中心单，要证明它满射，就需让左右乘的有限和实现任意 \(k\) 线性算子。把正则双模 \(A\) 看成左 \(A^e\) 模，其子模正是双边理想，故这个模简单。由\cref{b2:prop:regular-bimodule-endomorphism}，稠密性定理所用的右标量除环为
\[
D=\operatorname{End}_{A^e}(A)^{\mathrm{op}}=Z(A)^{\mathrm{op}}=k,
\]
所以目标 \(\operatorname{End}_D(A_D)\) 恰为 \(\operatorname{End}_k(A)\)。\(A\) 对这个右除环有限维，\cref{b2:cor:density-finite-dimensional}用于环 \(A^e\) 及简单模 \(A\)，便使 \(\Theta_A\) 满射。来源与目标的 \(k\) 维数同为 \((\dim_kA)^2\)，故它还是单射。

反过来，若 \(\Theta_A\) 同构，则每个双边理想 \(I\) 都是正则双模 \(A\) 的 \(A^e\) 子模，因而被 \(\Theta_A\) 的全部作用算子保持。满射性使它被每个 \(T\in\operatorname{End}_k(A)\) 保持。若 \(0\ne x\in I\)，对任意 \(y\in A\)，把 \(x\) 扩充为一组基并指定 \(T(x)=y\)，就得到 \(y=T(x)\in I\)。所以 \(I=A\)，证明单性。由正则双模自同态公式，\(Z(A)\) 同构于 \(\operatorname{End}_{A^e}(A)\)；后者现在是全部线性算子的中心化子。在\cref{b2:prop:tensor-factor-centralizers}中取 \(U=A\)、\(W=k\)，便知这个中心化子为 \(k\operatorname{id}_A\)，故 \(Z(A)=k\)。
\end{proof}

张量积中的元素是有限和，所以这个同构具体断言：对每个 \(T\in\operatorname{End}_k(A)\)，都能找到有限组 \(a_\nu,b_\nu\in A\)，使
\[
T(x)=\sum_{\nu=1}^{r}a_\nu x b_\nu\qquad(x\in A).
\]
纯张量 \(a\otimes b^{\mathrm{op}}\) 对应一次左右乘；一般线性算子由这样的操作的有限和实现。

\ProseParagraphBreak
在 \(A=M_n(k)\) 中，若 \(X=(x_{pq})\)，矩阵单位给出
\[
E_{ij}XE_{\ell m}=x_{j\ell}E_{im}.
\]
这项操作读取第 \((j,\ell)\) 个条目，再把它放到第 \((i,m)\) 个位置，其余条目清零。让两组位置遍历所有可能，就得到矩阵空间的全部线性自同态的矩阵单位；它们的线性组合正是任意线性操作。扩大底域后，这些读取、放置公式依然成立，这也显示了包络代数判据适合处理标量扩张的原因。

\subsection{分裂代数与分裂域}
\label{b2:subsec:1701:02}

对域扩张 \(K/k\)，记 \(A_K=K\otimes_kA\)。其乘法逐因子定义，单位为 \(1\otimes1\)。
\begin{definition}\label{b2:def:csa-splitting-field}
设 \(k\) 为域，\(A\) 为中心单 \(k\) 代数。若存在整数 \(n\geq1\) 使 \(A\cong M_n(k)\)，则称 \(A\) 为\term[fenlie-daishu]{分裂代数}[split algebra]，或称它在 \(k\) 上分裂。对域扩张 \(K/k\)，记 \(A_K=K\otimes_kA\)；若 \(A_K\cong M_n(K)\)，则称 \(K\) 为 \(A\) 的\term[fenlie-yu-daishu]{分裂域}[splitting field of an algebra]。
\end{definition}
\symidx{\(A_K=K\otimes_kA\)：代数的标量扩张}

扩张标量能检测一个给定映射是否为同构，但这与扩张后才找到一个同构不同。若 \(f:V\to W\) 是给定的 \(k\) 线性映射，且 \(K\otimes_k f\) 为同构，则 \(f\) 本身就是同构：向量空间的短正合列分裂，张量后仍正合，所以 \(f\) 的核与余核扩张后都为零；非零 \(k\) 向量空间的一组基扩张后仍非空，因而核与余核原本就为零。对于给定的 \(k\) 代数同态，这同样检测它是否为代数同构。

\ProseParagraphBreak
若只知道 \(A_K\cong B_K\)，这个同构却未必来自某个 \(k\) 代数同态 \(A\to B\)。例如实四元数代数 \(\mathbb H\) 与 \(M_2(\mathbb R)\) 扩张到 \(\mathbb C\) 后都同构于 \(M_2(\mathbb C)\)，其中四元数的分裂将在\cref{b2:prop:quaternion-basis-splitting}中具体给出；但 \(\mathbb H\) 是除代数，\(M_2(\mathbb R)\) 有非零不可逆矩阵，所以两者在 \(\mathbb R\) 上不同构。\cref{b2:cor:similarity-descent}中矩阵相似能够下降，是因为首一不变因子在扩张标量后保留，这个额外的分类信息不能直接推广到任意代数的同构问题。

\begin{theorem}[标量扩张]\label{b2:thm:csa-base-change}
设 \(k\) 为域，\(A\) 为中心单 \(k\) 代数，则对任意域扩张 \(K/k\)，标量扩张 \(A_K=K\otimes_kA\) 为中心单 \(K\) 代数，不要求扩张有限或可分。
\end{theorem}
\begin{proof}
要保留中心单性，可以把包络代数判据中的同构扩张标量。由于 \(A\) 有限维，取它的一组有限基后，矩阵单位给出自然同构
\[
K\otimes_k\operatorname{End}_k(A)\cong\operatorname{End}_K(A_K),
\]
将 \(c\otimes T\) 送到 \(d\otimes a\mapsto cd\otimes T(a)\)。它把两侧相应的有限矩阵单位基一一对应，所以是代数同构；这个与自同态代数交换的公式不能无条件用于无限维空间。另有
\[
A_K\otimes_K(A_K)^{\mathrm{op}}
\cong K\otimes_k(A\otimes_kA^{\mathrm{op}}).
\]
将\eqref{b2:eq:csa-enveloping-isomorphism}扩张标量并作这些识别，所得映射正是 \(A_K\) 的左右乘作用。它是同构，且 \(A_K\ne0\)，因为原 \(k\) 基扩张后仍为 \(K\) 基。因此包络代数判据给出所需中心单性。
\end{proof}

若只假设 \(A\) 单，标量扩张后却可能出现非零真双边理想。\cref{b2:prop:inseparable-base-change-counterexample}将给出一个有限纯不可分域扩张作为 \(k\) 代数的反例；它的中心大于 \(k\)，因而不满足上述定理的中心性条件。

\begin{proposition}[中心单性的下降]\label{b2:prop:csa-descent-from-extension}
设 \(k\) 为域，\(A\ne0\) 为有限维含幺结合 \(k\) 代数，\(K/k\) 为域扩张。若 \(A_K=K\otimes_kA\) 为中心单 \(K\) 代数，则 \(A\) 为中心单 \(k\) 代数。
\end{proposition}
\begin{proof}
若 \(I\) 是 \(A\) 的非零真双边理想，则 \(K\otimes_kI\) 是 \(A_K\) 的双边理想。取 \(I\) 的一组 \(k\) 基并扩充为 \(A\) 的基，便知扩张后它仍非零且仍为真子空间，这与 \(A_K\) 单矛盾。因此 \(A\) 单。

再把 \(1_A\) 扩充为 \(A\) 的 \(k\) 基。若 \(z\in Z(A)\)，则 \(1\otimes z\) 与 \(K\) 标量及全部 \(1\otimes a\) 交换，所以在 \(Z(A_K)=K(1\otimes1_A)\) 中。于是 \(1\otimes z=c\otimes1_A\) 对某个 \(c\in K\) 成立。用扩张后的基比较系数，\(z\) 的所有非单位基坐标都为零，其单位坐标原本在 \(k\) 中，故 \(z\in k1_A\)。这证明中心性。
\end{proof}

\begin{corollary}\label{b2:cor:csa-algebraic-closure-splits}
设 \(k\) 为域，\(\Omega\) 为其代数闭包，\(A\) 为中心单 \(k\) 代数，则 \(A_\Omega=\Omega\otimes_kA\cong M_n(\Omega)\)，其中整数 \(n\geq1\) 唯一，且 \(\dim_kA=n^2\)。
\end{corollary}
\begin{proof}
\cref{b2:thm:csa-base-change}使 \(A_\Omega\) 有限维且中心单，因而简单、半单。\cref{b2:prop:finite-dimensional-semisimple-algebra}在代数闭域上把它写成矩阵环的有限乘积；单性排除多个非零因子，所以只剩 \(M_n(\Omega)\)。扩张标量不改变有限基的大小，故 \(\dim_kA=n^2\)，这也唯一确定 \(n\)。
\end{proof}

称 \(n=\sqrt{\dim_kA}\) 为 \(A\) 的\term[zhongxin-dan-daishu-cishu]{次数}[degree of a central simple algebra]，记为 \(\deg A\)。次数为 \(n\)，整个代数的向量空间维数则为 \(n^2\)；因此中心单代数的维数必须是完全平方，例如三维或六维的代数不能中心单。任意分裂域给出的矩阵大小都为这个 \(n\)，因为维数在标量扩张下不变。\secref{b2:sec:1702}还将证明可以选取有限 Galois 分裂域；仅知道在代数闭包上分裂，尚未说明有限可分扩张的存在。
\symidx{\(\deg A=\sqrt{\dim_kA}\)：中心单代数的次数}

\subsection{张量积与反代数}
\label{b2:subsec:1701:03}

反代数 \(A^{\mathrm{op}}\) 与 \(A\) 有相同的中心，双边理想也相同，所以中心单性在取反代数后保持。它的次数等于 \(\deg A\)，因为维数不变。与自身的反代数作张量则有更强的结论：由\cref{b2:thm:csa-enveloping-isomorphism}，
\[
A\otimes_kA^{\mathrm{op}}\cong M_{\dim_kA}(k)
=M_{(\deg A)^2}(k).
\]
无论 \(A\) 本身是否分裂，这个张量积总是分裂的。右侧矩阵的阶数是 \(\dim_kA=(\deg A)^2\)，不能误记为 \(\deg A\)。在\secref{b2:sec:1801}的 Brauer 群中，分裂矩阵代数代表单位元，这个同构将使 \(A^{\mathrm{op}}\) 的类成为 \(A\) 的类的逆元。

\begin{theorem}[中心单代数的张量积]\label{b2:thm:csa-tensor-product}
设 \(k\) 为域，\(A,B\) 为中心单 \(k\) 代数，则 \(A\otimes_kB\) 也是中心单 \(k\) 代数，且
\(\deg(A\otimes_kB)=\deg A\deg B\)。
\end{theorem}
\begin{proof}
令 \(C=A\otimes_kB\)。\(C\) 的左右乘可以分别在 \(A,B\) 两个因子上实现，因此要把它的包络作用与两个因子的包络作用相比较。交换不同张量因子的次序并取结合识别，得到
\[
C\otimes_kC^{\mathrm{op}}
\cong(A\otimes_kA^{\mathrm{op}})\otimes_k(B\otimes_kB^{\mathrm{op}}).
\]
分别应用两个包络代数同构，再用\cref{b2:prop:tensor-factor-centralizers}中已经通过矩阵单位证明的同构，得到
\[
C\otimes_kC^{\mathrm{op}}
\cong\operatorname{End}_k(A)\otimes_k\operatorname{End}_k(B)
\cong\operatorname{End}_k(C).
\]
在纯张量上，这个复合把 \((a\otimes b)\otimes(a'\otimes b')^{\mathrm{op}}\) 送到
\(x\otimes y\mapsto axa'\otimes byb'\)，正是 \(C\) 的左右乘作用。因此包络代数判据保证 \(C\) 中心单。维数乘法给出次数乘法。
\end{proof}

特别地，\(M_r(A)\cong M_r(k)\otimes_kA\) 是中心单代数，次数为 \(r\deg A\)。同构把 \(E_{ij}\otimes a\) 送到第 \((i,j)\) 个条目为 \(a\)、其余为零的矩阵；基与乘法同时对应。另一方面，两个一般单 \(k\) 代数的张量积不必单，中心性在定理中不能删除。例如
\[
\mathbb C\otimes_{\mathbb R}\mathbb C\longrightarrow\mathbb C\times\mathbb C,
\qquad z\otimes w\longmapsto(zw,z\bar w)
\]
是实代数同构。确实，按第一个 \(\mathbb C\) 因子看，来源以 \(1\otimes1,1\otimes i\) 为基，两个像为 \((1,1),(i,-i)\)，线性无关且生成目标；乘法直接保持。恒等嵌入和复共轭使第二因子分别产生两个坐标，而目标中的 \(\mathbb C\times0\) 是非零真双边理想，故张量积不是单环。这里每个因子作为实代数都单，但其中心为 \(\mathbb C\)，大于底域 \(\mathbb R\)。

\ProseParagraphBreak
把底域换成交换环时，有限维性应改为忠实有限投射性；左右乘作用的同构则仍有意义。\appref{b2:app:C}的\cref{b2:def:azumaya-algebra}据此定义 Azumaya 代数\footnote{東屋五郎（假名：あずまや ごろう，罗马音：Gorō Azumaya，1920-2010），日本数学家，研究环论及代数的结构；Azumaya 代数以他命名。}，并由\cref{b2:thm:azumaya-fiber-criterion}证明其剩余域纤维判据。这项推广不将每个纤维预设为已经分裂的矩阵代数。

\begin{proposition}[不可分标量扩张后的非单性]\label{b2:prop:inseparable-base-change-counterexample}
设 \(p\) 为素数、\(t\) 为 \(\mathbb F_p\) 上的超越元，令 \(k=\mathbb F_p(t)\)。在 \(k\) 的代数闭包中取 \(\alpha\) 满足 \(\alpha^p=t\)，并令 \(L=k(\alpha)\)。则 \(X^p-t\) 在 \(k[X]\) 中不可约，\([L:k]=p\)，且以第一个张量因子为 \(L\) 标量时有含幺 \(L\) 代数同构
\[
L\otimes_kL\cong L[\varepsilon]/(\varepsilon^p).
\]
其中 \(1\otimes\alpha-\alpha\otimes1\) 对应 \(\varepsilon\) 的剩余类。它是非零幂零元，生成一个非零真理想，故 \(L\otimes_kL\) 不是单代数。原代数 \(L\) 虽为单 \(k\) 代数，但其中心为 \(L\ne k\)。
\end{proposition}
\begin{proof}
先证明 \(t\) 在 \(k\) 中不是 \(p\) 次方。若 \((f/g)^p=t\)，其中 \(f,g\in\mathbb F_p[t]\)、\(g\ne0\)，约去公因子后有 \(f^p=tg^p\)。左边不可约因子 \(t\) 的指数为 \(p\) 的倍数，右边的指数则模 \(p\) 余一，矛盾。

令 \(\mu_{\alpha,k}\) 为 \(\alpha\) 在 \(k\) 上的首一最小多项式。它整除 \(X^p-t\)，而在代数闭包中 \(X^p-t=(X-\alpha)^p\)，故 \(\mu_{\alpha,k}=(X-\alpha)^d\)，其中 \(1\le d\le p\)。若 \(d<p\)，其 \(X^{d-1}\) 系数为 \(-d\alpha\in k\)。此时 \(d\) 在特征 \(p\) 的域中非零，便有 \(\alpha\in k\)，与 \(t\) 不是 \(p\) 次方矛盾。因此 \(d=p\)，\(\mu_{\alpha,k}=X^p-t\)，并且 \(1,\alpha,\ldots,\alpha^{p-1}\) 是 \(L/k\) 的基。

把 \(L[X]\) 中的 \(X\) 送到 \(1\otimes\alpha\)，把系数 \(c\in L\) 送到 \(c\otimes1\)，得到含幺 \(L\) 代数同态。关系 \(X^p-t\) 在像中为零，故它诱导
\[
L[X]/(X^p-t)\longrightarrow L\otimes_kL.
\]
两侧分别以 \(X^j\) 的剩余类和 \(1\otimes\alpha^j\)（\(0\le j<p\)）为 \(L\) 基，且这些基向量一一对应，所以此映射为同构。在 \(L[X]\) 中令 \(\varepsilon=X-\alpha\)，便有 \(X^p-t=\varepsilon^p\)，得到所述截断多项式环。它以 \(1,\varepsilon,\ldots,\varepsilon^{p-1}\) 的剩余类为 \(L\) 基，故 \(\varepsilon\ne0\) 而 \(\varepsilon^p=0\)；理想 \((\varepsilon)\) 非零，且商为非零域 \(L\)，所以是真理想。最后，\(L\) 是域，故作为 \(k\) 代数单；但 \([L:k]=p>1\)，且 \(Z(L)=L\)，因而它不中心单，也就不满足\cref{b2:thm:csa-base-change}的假设。
\end{proof}

% !TeX root = ../../Book2.tex
% 纲要标题：### 17.2 中心化子、子域与分裂域
% 纲要位置：计划纲要.md，第 1313 行。

\section{中心化子、子域与分裂域}
\label{b2:sec:1702}

一个单子代数 \(B\subseteq A\) 的中心化子记录与它交换的另一部分乘法。维数公式将精确说明这两部分各占多大；但只有 \(B\) 的中心也是 \(k\) 时，二者才直接张量成整个 \(A\)。本节通过简单模及其重数证明这些结论，子代数中心的扩张 \(Z(B)/k\) 可以不可分。

% 纲要要点（供目录核对）：
%
% * 中心化子的维数关系
% * 双中心化子定理
% * 子域与最大交换子代数
% * 可分最大子域与有限 Galois 分裂域
%

\subsection{中心化子的维数关系}
\label{b2:subsec:1702:01}

记
\[
C_A(B)=\{\,a\in A\mid ab=ba\text{ 对所有 }b\in B\,\}.
\]
它是含单位元的 \(k\) 子代数。本节的中心化子在 \(A\) 内取，与第十一章在整个线性自同态代数内取中心化子应分别理解。
\symidx{\(C_A(B)\)：子代数 \(B\) 在 \(A\) 中的中心化子}

\begin{theorem}[中心化子定理]\label{b2:thm:csa-centralizer}
设 \(k\) 为域，\(A\) 为中心单 \(k\) 代数，\(B\subseteq A\) 为非零简单含幺 \(k\) 子代数，\(L=Z(B)\)。令 \(C=C_A(B)=\{\,a\in A\mid ab=ba\text{ 对所有 }b\in B\,\}\)，则
\[
C\text{ 为简单代数},\qquad Z(C)=L,
\qquad \dim_kA=\dim_kB\,\dim_kC,
\qquad C_A(C)=B.
\]
其中两个中心均按它们在 \(A\) 中的实际嵌入识别。结论不要求 \(L/k\) 可分。
\end{theorem}
\begin{proof}
先构造有限维代数 \(E=B\otimes_kA^{\mathrm{op}}\)。\(B\) 以其中心 \(L\) 为底域是中心单代数，而
\[
E\cong B\otimes_L(L\otimes_kA^{\mathrm{op}}).
\]
具体同构将 \(b\otimes(l\otimes a^{\mathrm{op}})\) 送到 \(bl\otimes a^{\mathrm{op}}\)，逆映射在 \(b\otimes a^{\mathrm{op}}\) 上取 \(b\otimes(1\otimes a^{\mathrm{op}})\)。\cref{b2:thm:csa-base-change}说明第二个因子在 \(L\) 上中心单，再由\cref{b2:thm:csa-tensor-product}，\(E\) 为中心单 \(L\) 代数，特别是单 Artin 环。

令 \(E\) 在 \(A\) 上作用为 \((b\otimes a^{\mathrm{op}})x=bxa\)。一个 \(E\) 线性自同态首先与全部右乘交换，所以由\cref{b2:thm:regular-double-centralizer}，它形如左乘 \(c\in A\)。再要求它与左 \(B\) 作用交换，恰好得到 \(c\in C\)。左乘保留乘法次序，因此
\[
C\cong\operatorname{End}_E(A).
\]
写 \(E\cong M_r(D)\)，其中 \(D\) 是除环。\cref{b2:thm:artin-wedderburn}及\cref{b2:prop:semisimple-multiplicity}的矩阵块分类说明，\(E\) 半单，且简单左模只有列模 \(D^r\) 一种类型。\(A\) 是有限维 \(k\) 空间，所以作为 \(E\) 模为有限个副本之和，设 \(A\cong(D^r)^m\)、\(m\ge1\)。因此
\[
\operatorname{End}_E(A)\cong M_m(D^{\mathrm{op}}),
\]
从而 \(C\) 简单。

记 \(d=\dim_kD\)。分别比较上述代数与模的 \(k\) 维数，有
\[
\dim_kB\,\dim_kA=r^2d,\qquad
\dim_kA=mrd,\qquad\dim_kC=m^2d.
\]
这三式给出 \(\dim_kB\,\dim_kC=\dim_kA\)。至此已独立证明：对任意满足定理假设的简单子代数，其中心化子简单，且有上述维数公式；双中心化子和中心的识别尚未用到。现在 \(C\) 也是 \(A\) 的简单含幺子代数，将这个已证的中间结论用于 \(C\subseteq A\)，得到
\(\dim_k C_A(C)=\dim_kA/\dim_kC=\dim_kB\)。而 \(B\subseteq C_A(C)\) 由交换性定义直接成立，有限维与等维使二者相等。最后
\[
Z(C)=C\cap C_A(C)=C\cap B=Z(B)=L,
\]
得到中心的实际识别。
\end{proof}

证明中 \(L\otimes_kA^{\mathrm{op}}\) 的中心单性来自对任意域扩张均成立的标量扩张定理。即使 \(L/k\) 是纯不可分扩张，这一步也仍然有效；真正需要简单性的，是 \(B\) 以其中心为底域的结构。

\subsection{双中心化子定理}
\label{b2:subsec:1702:02}

\begin{corollary}\label{b2:cor:csa-centralizer-factorization}
设 \(k\) 为域，\(A\) 为中心单 \(k\) 代数，\(B\subseteq A\) 为非零简单含幺 \(k\) 子代数，记 \(C_A(B)\) 为 \(B\) 在 \(A\) 中的中心化子。若 \(Z(B)=k\)，则乘法诱导中心单 \(k\) 代数同构
\[
B\otimes_kC_A(B)\longrightarrow A,\qquad b\otimes c\longmapsto bc.
\]
一般地，令 \(L=Z(B)\)，则有 \(L\) 代数同构
\[
B\otimes_LC_A(B)\cong C_A(L).
\]
\end{corollary}
\begin{proof}
两个因子的元素互相交换，所以乘法公式保持张量关系与代数乘法。若 \(L=k\)，来源由\cref{b2:thm:csa-tensor-product}为单代数；映射保持单位元，故核不能为全代数，只能为零。中心化子维数公式说明来源与目标同维，于是满射。

一般情形中，\(B,C_A(B)\) 都是中心单 \(L\) 代数，故它们在 \(L\) 上的张量积单，乘法仍单射，且其像包含于 \(C_A(L)\)。这里仍以 \(k\) 为底域，将简单 \(k\) 子代数 \(L\subseteq A\) 代入中心化子定理，得到 \(\dim_k C_A(L)=\dim_kA/[L:k]\)。张量积来源的 \(k\) 维数为
\(\dim_kB\,\dim_kC_A(B)/[L:k]\)，也等于这个数，因此像就是全部 \(C_A(L)\)。
\end{proof}

例如，把有限扩张 \(L/k\) 通过正则作用嵌入 \(A=\operatorname{End}_k(L)\)。与左乘 \(L\) 交换的映射恰为右乘 \(L\)，因为 \(T(x)=T(x\cdot1)=x\cdot T(1)\)。\(L\) 交换，所以 \(C_A(L)=L\)，维数公式给出 \([L:k]^2=[L:k]\,[L:k]\)。当 \(L\ne k\) 时，\(L\otimes_L L=L\) 远小于 \(A\)；只有在正确的中心上取张量，且把目标写成 \(C_A(L)\)，才能保持一般公式成立。

\ProseParagraphBreak
双中心化子等式还说明，简单子代数可由它在 \(A\) 中的全部交换关系恢复。若删除简单性，这一性质可能失败：第十一章的上三角代数 \(T_2(k)\subseteq M_2(k)\) 的中心化子为 \(kI\)，再取一次得到整个 \(M_2(k)\)，并非原上三角代数。

这些关系可以在 \(A\) 内按包含画成两支：
\[
\begin{tikzcd}[column sep=2.2em,row sep=1.5em]
& & B \arrow[dr,hook] & & \\
k \arrow[r,hook] & L=Z(B)=Z(C) \arrow[ur,hook] \arrow[dr,hook]
& & C_A(L) \arrow[r,hook] & A\\
& & C=C_A(B) \arrow[ur,hook] & &
\end{tikzcd}
\]
两支相乘得到 \(B\otimes_LC\cong C_A(L)\)，而 \(C_A(C)=B\)、\(\dim_kA=\dim_kB\,\dim_kC\)。当 \(L=k\) 时，图中的 \(C_A(L)\) 才是整个 \(A\)。

\subsection{子域与最大交换子代数}
\label{b2:subsec:1702:03}

设 \(A\) 为中心单 \(k\) 代数。若含 \(k\) 的子域 \(L\subseteq A\) 不包含于更大的 \(k\) 子域，称 \(L\) 为\term[baohan-jida-ziyu]{包含极大子域}[inclusion-maximal subfield]。若一个含幺交换子代数不包含于更大的含幺交换子代数，则称它为\term[zuida-jiaohuan-zidaishu]{最大交换子代数}[maximal commutative subalgebra]。前者只在域中比较，后者也允许零因子或幂零元。文献中的 maximal subfield 还可能专指满足 \([L:k]=\deg A\) 的子域；本书在一般中心单代数中分别写明这两个条件，只在除代数语境把包含极大子域简称\term[zuida-ziyu]{最大子域}[maximal subfield]。

\cref{b2:tab:csa-subfield-conditions}按对象所在位置列出容易混淆的条件。最后一行的 \(K\) 只须是 \(k\) 的扩域，可以完全不嵌入 \(A\)。
\begin{table}[htbp]
\centering
\caption{中心单代数中的子域、交换子代数与分裂域}\label{b2:tab:csa-subfield-conditions}
\begin{tabularx}{\linewidth}{@{}p{0.31\linewidth}X@{}}
\toprule
对象 & 条件 \\
\midrule
包含极大子域 \(L\subseteq A\) & 不包含于 \(A\) 中更大的含 \(k\) 子域。\\
最大交换子代数 \(E\subseteq A\) & 不包含于 \(A\) 中更大的含幺交换 \(k\) 子代数。\\
自中心化子域 \(L\subseteq A\) & \(C_A(L)=L\)。\\
次数达到 \(\deg A=n\) 的子域 \(L\subseteq A\) & \([L:k]=n\)。\\
分裂域 \(K/k\) & \(A_K\cong M_n(K)\)；不要求 \(K\subseteq A\)。\\
\bottomrule
\end{tabularx}
\end{table}

\begin{proposition}\label{b2:prop:csa-self-centralizing-field}
设 \(k\) 为域，\(A\) 为中心单 \(k\) 代数，\(n=\deg A\)，\(L\subseteq A\) 为含 \(k\) 的子域，则以下条件等价：
\[
[L:k]=n;\qquad C_A(L)=L;\qquad
L\text{ 是 }A\text{ 的最大交换子代数}.
\]
若 \(A\) 本身为除代数，则每个含 \(k\) 的最大子域都满足这些条件。
\end{proposition}
\begin{proof}
将\cref{b2:thm:csa-centralizer}用于简单子代数 \(L\)，得到 \(n^2=[L:k]\dim_k C_A(L)\)。由于 \(L\subseteq C_A(L)\)，当 \([L:k]=n\) 时二者同维，所以相等；反向相等也立即给出次数条件。

若 \(C_A(L)=L\)，任何包含 \(L\) 的交换子代数都包含于该中心化子，因此只能是 \(L\)。反过来，若 \(c\in C_A(L)\)，则 \(L[c]\) 是包含 \(L\) 的交换子代数，最大交换性使 \(c\in L\)。最后，若环境为除代数 \(D\)，则 \(L[c]\) 是有限维交换整环，因而为域：非零元的乘法单射，有限维使其满射，从而有逆。最大子域的性质就迫使 \(L[c]=L\)，也得到中心化子等于 \(L\)。
\end{proof}

在矩阵代数中，包含极大不保证次数达到 \(\deg A\)。例如 \(M_2(\mathbb C)\) 中每个有限维的 \(\mathbb C\) 子域都等于 \(\mathbb C\)，所以标量域是包含极大子域，却不是最大交换子代数，也只有次数一。对角代数 \(\mathbb C\times\mathbb C\) 是最大交换子代数，但不是域；\(\mathbb C[N]\)、\(N=\begin{pmatrix}0&1\\0&0\end{pmatrix}\)，也是最大交换子代数，因为与 \(N\) 交换的矩阵恰为 \(aI+bN\)，而它含有非零幂零元。

\begin{proposition}\label{b2:prop:maximal-field-splits-csa}
设 \(k\) 为域，\(A\) 为中心单 \(k\) 代数，\(L\subseteq A\) 为含 \(k\) 的子域。若 \([L:k]=\deg A=n\)，则 \(L\) 分裂 \(A\)。特别地，每个有限维中心除 \(k\) 代数的最大子域都是它的分裂域。
\end{proposition}
\begin{proof}
把 \(A\) 看成右 \(L\) 向量空间，其维数为 \(\dim_kA/[L:k]=n\)。作用
\[
L\otimes_kA\longrightarrow\operatorname{End}_L(A),
\qquad l\otimes a\longmapsto(x\mapsto axl)
\]
是 \(L\) 代数同态；右 \(L\) 线性用到 \(L\) 内元素交换，乘法相容性由结合律与相同交换性给出。目标中的 \(A\) 始终是原来的右 \(L\) 空间。来源由\cref{b2:thm:csa-base-change}为单代数，所以保持单位元的映射单射。两侧的 \(L\) 维数均为 \(n^2\)，故为同构，目标为 \(M_n(L)\)。除代数的最大子域由\cref{b2:prop:csa-self-centralizing-field}满足 \([L:k]=\deg A\)，所以也适用。
\end{proof}

\subsection{可分最大子域与有限 Galois 分裂域}
\label{b2:subsec:1702:04}

有限维除代数确有最大子域：子域维数是有界正整数，取具有最大维数的一个即可。不过，这样得到的子域未必可分。为得到有限 Galois 分裂域，下面另作一次可分元素的选择。

\begin{theorem}[有限可分分裂域的存在]\label{b2:thm:csa-separable-splitting}
设 \(k\) 为域。每个中心单 \(k\) 代数都有有限可分分裂域，因而也有有限 Galois 分裂域。更具体地，每个有限维中心除 \(k\) 代数 \(D\) 都含有一个可分最大子域 \(L\)，满足 \([L:k]=\deg D\)，且 \(L\) 分裂 \(D\)。
\end{theorem}
\begin{proof}
由\cref{b2:prop:csa-division-description}，只须对中心除代数 \(D\) 找到所述 \(L\)：一旦 \(L\) 分裂 \(D\)，它也分裂 \(M_r(D)\)。若 \(k\) 有限，则 \(D\) 的元素总数也有限，由\cref{b2:thm:wedderburn-little}，\(D\) 交换；中心为 \(k\) 便迫使 \(D=k\)，取 \(L=k\) 即可。以下设 \(k\) 无限，令 \(n=\deg D\)，\(\Omega\) 为代数闭包。

取同构 \(D_\Omega\cong M_n(\Omega)\) 及 \(D\) 的 \(k\) 基 \(e_1,\ldots,e_{n^2}\)，其像为矩阵 \(E_1,\ldots,E_{n^2}\)。这些矩阵是 \(M_n(\Omega)\) 的 \(\Omega\) 基。要寻找 \(a=\sum_i c_ie_i\)、\(c_i\in k\)，使对应矩阵具有 \(n\) 个互异特征值。

先建立一个明确的多项式检验。对任意 \(H\in M_n(\Omega)\)，令 \(f_H(t)=\det(tI-H)\)，并令 \(U_j\) 表示次数小于 \(j\) 的 \(\Omega\) 多项式空间，\(U_0=0\)。考虑等维空间之间的映射
\[
U_{n-1}\oplus U_n\longrightarrow U_{2n-1},
\qquad(u,v)\longmapsto u f_H+v f_H'.
\]
在单项式基下取其行列式 \(\Delta(H)\)。这是首一多项式 \(f_H\) 与 \(f_H'\) 的结式 \(\operatorname{Res}(f_H,f_H')\) 的一种归一化，也可视为判别式；更换这些空间的固定基只会使行列式乘以一个非零常数。矩阵条目来自 \(f_H\) 及其导数的系数，所以 \(\Delta(H)\) 是 \(H\) 条目的多项式。

该行列式非零当且仅当 \(f_H,f_H'\) 互素。如果互素且 \(u f_H+v f_H'=0\)，则 \(f_H\mid v\)，但 \(\deg v<n\)，故 \(v=0\)，继而 \(u=0\)，映射单射而可逆。若有正次数公因子 \(h\)，则 \((u,v)=(f_H'/h,-f_H/h)\) 是符合次数限制的非零核向量，所以行列式为零。互素又等价于 \(f_H\) 没有重根：根 \(\lambda\) 同时使 \(f_H,f_H'\) 为零，恰好是它在因式分解中出现至少两次。

于是
\[
P(X_1,\ldots,X_{n^2})=\Delta\left(\sum_iX_iE_i\right)
\in\Omega[X_1,\ldots,X_{n^2}]
\]
是非零多项式。因为 \(E_i\) 张成全部矩阵，可在 \(\Omega\) 中选择系数使其和为具有 \(n\) 个互异对角元的对角矩阵，此处 \(\Delta\ne0\)。

即使 \(P\) 的系数在 \(\Omega\) 中，它也不能在无限子域 \(k\) 的全部点上消失。证明可对变量数归纳：一元非零多项式只有有限多个根；多元时先将它按最后一个变量展开，选一个非零系数多项式，用归纳假设为前面的变量取 \(k\) 中的值，使该系数非零，再用一元结论选最后一个值。因此可取 \(c_i\in k\) 使 \(P(c_1,\ldots,c_{n^2})\ne0\)。相应的 \(a=\sum_i c_ie_i\in D\) 在分裂矩阵模型中有 \(n\) 个互异特征值。任意零化矩阵的多项式必须在这 \(n\) 个特征值处为零，故次数至少为 \(n\)；其 \(n\) 次特征多项式又由\cref{b2:thm:exterior-cayley-hamilton}零化该矩阵。因此矩阵最小多项式就是这个没有重根的特征多项式。

\(a\) 在 \(k\) 上的最小多项式也正是上述矩阵的最小多项式。事实上，对左乘算子 \(L_a:D\to D\)，有 \(p(L_a)=0\) 当且仅当 \(p(a)=0\)：必要性只需将算子作用于 \(1\)。\cref{b2:prop:scalar-extension-invariant-factors}说明 \(L_a\) 扩张到 \(\Omega\) 后最小多项式原样不变；在 \(D_\Omega\cong M_n(\Omega)\) 上，这个扩张算子就是矩阵 \(H\) 的左乘，其最小多项式又等于 \(H\) 的最小多项式。因此 \(a\) 的首一最小多项式次数为 \(n\)，且在 \(\Omega\) 中无重根，所以这个最小多项式可分。这里没有预先假定 \(H\) 的特征多项式系数属于 \(k\)。

\(D\) 为除环，使交换子代数 \(k[a]\) 是有限维整环，因而为域。于是 \(L=k[a]\) 是次数 \(n\) 的可分子域，由\cref{b2:prop:csa-self-centralizing-field}及\cref{b2:prop:maximal-field-splits-csa}，它是最大子域并分裂 \(D\)。最后把 \(a\) 送到其最小多项式在 \(\Omega\) 中的一个根，将 \(L\) 嵌入 \(\Omega\)，并取该多项式在 \(\Omega\) 中的分裂域 \(F\)。第一卷命题15.1.2保证其有限性；它由这个可分多项式的全部根生成，按该卷定义19.1.1后的分裂域刻画，是有限 Galois 扩张，并且含有所选的 \(L\) 的像。继续扩张分裂矩阵环仍为矩阵环，所以 \(F\) 也分裂原来的 \(A\)。
\end{proof}

% !TeX root = ../../Book2.tex
% 纲要标题：### 17.3 Skolem–Noether 定理
% 纲要位置：计划纲要.md，第 1609 行。

\section{Skolem–Noether 定理}
\label{b2:sec:1703}

中心化子定理比较一个子代数与同它交换的部分。现在比较同一个简单代数嵌入 \(A\) 的两种方式：结论是，它们只相差 \(A\) 中一个可逆元的内共轭。证明不从方程组中猜这个可逆元，而是把两种嵌入看成两个模结构，再由简单 Artin 环的模分类构造所需同构。

% 纲要要点（供目录核对）：
%
% * 单子代数的嵌入及内自同构；两种嵌入只改变左作用，单 Artin 环的一种简单模类型使维数决定重数
% * 中心单代数的自同构；从右正则模同构在单位元的像构造共轭元，逐个自同构可提升不保证存在群截面
% * 适用假设与非单代数的反例；非单源的不同分量重数及非中心环境的外自同构分别说明失败原因
%

\subsection{单子代数的嵌入与内自同构}
\label{b2:subsec:1703:01}

\begin{theorem}[Skolem–Noether]\label{b2:thm:skolem-noether}
设 \(k\) 为域，\(A\) 为中心单 \(k\) 代数，\(B\) 为非零有限维简单含幺 \(k\) 代数。对任意两个保幺 \(k\) 代数同态 \(\varphi,\psi:B\to A\)，存在 \(u\in A^\times\)，使
\[
\psi(b)=u\cdot\varphi(b)u^{-1}\qquad(b\in B).
\]
这些同态因 \(B\) 简单而自动单射；\(Z(B)/k\) 不必可分。这个版本也见\cite[Main Theorem 7.7.1]{Voight2021QuaternionAlgebras}。
\end{theorem}
\begin{proof}
把 \(B\) 的左作用与 \(A\) 的右正则作用放在同一个代数 \(E=B\otimes_kA^{\mathrm{op}}\) 中，便能只改变嵌入给出的左作用而保留右作用。令 \(L=Z(B)\)。与\cref{b2:thm:csa-centralizer}的证明相同，
\[
E\cong B\otimes_L(L\otimes_kA^{\mathrm{op}})
\]
是两个中心单 \(L\) 代数的张量积，因此为单 Artin 环。这一步只用任意标量扩张的中心单性，不要求 \(L/k\) 可分。

在同一个 \(k\) 向量空间 \(A\) 上定义两种左 \(E\) 模结构：
\[
(b\otimes a^{\mathrm{op}})\cdot_\varphi x=\varphi(b)xa,
\qquad
(b\otimes a^{\mathrm{op}})\cdot_\psi x=\psi(b)xa.
\]
左右乘法交换，所以模公理成立。由\cref{b2:prop:simple-artinian}写 \(E\cong M_r(D)\)，再由\cref{b2:thm:artin-wedderburn}及\cref{b2:prop:semisimple-multiplicity}的矩阵块分类，简单左模只有 \(S=D^r\) 一种类型。两种模都是有限个 \(S\) 的直和，且底层 \(k\) 空间都是 \(A\)，因此
\[
A_\varphi\cong S^{\oplus m}\cong A_\psi,\qquad
m=\frac{\dim_kA}{\dim_kS}.
\]
这里可以只凭维数确定重数，正因为 \(E\) 是单 Artin 环，只有一种简单左模。一般半单环可能有多种简单类型，它们的重数不能只由总维数恢复。于是存在 \(E\) 模同构 \(F:A_\varphi\to A_\psi\)。

\(E\) 线性首先保证 \(F(xa)=F(x)a\)，即 \(F\) 是右正则模之间的同构。右正则模由 \(1\) 生成，因此令 \(u=F(1)\)，便有 \(F(x)=F(1\cdot x)=ux\)。逆同构也由一的像决定，故为左乘某个 \(v\in A\)；两次复合为恒等在 \(1\) 处给出 \(vu=uv=1\)，所以 \(u\) 可逆。再由 \(B\) 的作用相容性，
\[
u\cdot\varphi(b)=F\bigl(\varphi(b)\bigr)=\psi(b)F(1)=\psi(b)u.
\]
右乘 \(u^{-1}\) 得到所需公式。最后，同态的核是 \(B\) 的双边理想，且因保持单位元不能为 \(B\)，所以核为零。
\end{proof}

这称为\term[Skolem-Noether-dingli]{Skolem–Noether 定理}[Skolem--Noether theorem]。若 \(B_1,B_2\subseteq A\) 为简单 \(k\) 子代数，且给定同构 \(\alpha:B_1\to B_2\)，把 \(B_1\) 的包含与 \(\alpha\) 后接 \(B_2\) 的包含作为两种嵌入，就得到 \(\alpha\) 可延伸为 \(A\) 的一个内自同构。所选 \(u\) 一般不唯一；若 \(u,v\) 实现同一嵌入之间的共轭，则 \(v^{-1}u\) 恰好与 \(\varphi(B)\) 的全部元素交换。

\subsection{中心单代数的自同构}
\label{b2:subsec:1703:02}

\begin{corollary}\label{b2:cor:csa-automorphisms-inner}
设 \(k\) 为域，\(A\) 为中心单 \(k\) 代数，则 \(A\) 的每个 \(k\) 代数自同构都是内自同构，并有群同构
\[
\operatorname{Aut}_{k\text{-alg}}(A)\cong A^\times/k^\times.
\]
特别地，对每个整数 \(n\ge1\)，\(\operatorname{Aut}_{k\text{-alg}}\bigl(M_n(k)\bigr)\cong\operatorname{PGL}_n(k)\)，其中
\(\operatorname{PGL}_n(k)=\operatorname{GL}_n(k)/k^\times I_n\)。
\end{corollary}
\begin{proof}
在\cref{b2:thm:skolem-noether}中取 \(B=A\)，两个嵌入分别为恒等映射和给定自同构，得到满群同态
\(A^\times\to\operatorname{Aut}_{k\text{-alg}}(A)\)、\(u\mapsto(a\mapsto uau^{-1})\)。它的核由与全部 \(A\) 交换的可逆元组成，即 \(Z(A)^\times=k^\times\)。因此把商类 \(uk^\times\) 送到该共轭自同构，给出所述商群同构。矩阵代数情形直接代入。
\end{proof}
\symidx{\(\operatorname{PGL}_n(k)\)：一般线性群模去非零标量矩阵所得的商群}

例如，若 \(\operatorname{char}k\ne2\)，\(d\in k^\times\) 非平方，矩阵
\[
T=\begin{pmatrix}0&d\\1&0\end{pmatrix}
\]
给出二次域 \(k(\sqrt d)\) 的嵌入。其非平凡自同构 \(\sqrt d\mapsto-\sqrt d\) 由 \(P=\operatorname{diag}(1,-1)\) 实现，因为 \(PTP^{-1}=-T\)。这个式子直接显示了域自同构如何在更大的矩阵代数内成为共轭。

\begin{corollary}\label{b2:cor:csa-field-normalizer}
设 \(k\) 为域，\(A\) 为中心单 \(k\) 代数，\(L\subseteq A\) 为含 \(k\) 的子域。令
\[
N_{A^\times}(L)=\{\,u\in A^\times\mid uLu^{-1}=L\,\}.
\]
共轭给出正合列
\[
1\longrightarrow C_A(L)^\times\longrightarrow N_{A^\times}(L)
\longrightarrow\operatorname{Aut}_k(L)\longrightarrow1.
\]
若 \([L:k]=\deg A\)，则核为 \(L^\times\)。
\end{corollary}
\begin{proof}
正规化条件保证共轭限制为 \(L\) 的 \(k\) 自同构。其核恰为在 \(L\) 上逐点恒等的那些可逆元，也就是 \(C_A(L)^\times\)。对任意 \(\sigma\in\operatorname{Aut}_k(L)\)，把包含 \(\iota:L\hookrightarrow A\) 与 \(\iota\circ\sigma\) 作为两种简单代数嵌入，\cref{b2:thm:skolem-noether}给出实现它的可逆元，因此该同态满射。最后的核识别由\cref{b2:prop:csa-self-centralizing-field}。
\end{proof}

满射保证对每个 \(\sigma\in\operatorname{Aut}_k(L)\)，存在 \(u_\sigma\in N_{A^\times}(L)\) 实现它，却不保证能同时选择这些元素，使 \(u_{\sigma\tau}=u_\sigma u_\tau\) 对全部 \(\sigma,\tau\) 成立。两边虽然实现同一个域自同构，比值仍可能是中心化子中的非平凡单位。\cref{b2:prop:quaternion-normalizer-nonsplit}将在实四元数中给出没有群截面的具体正合列；下一章的循环代数还将记录实现元相乘时的这种差异。

\subsection{适用条件与非单代数反例}
\label{b2:subsec:1703:03}

先看简单子代数假设。取 \(A=M_3(k)\)、\(B=k\times k\)，定义
\[
\varphi(x,y)=\operatorname{diag}(x,y,y),\qquad
\psi(x,y)=\operatorname{diag}(x,x,y).
\]
两者都是含幺的单射代数同态，但 \(\varphi(1,0)\) 的秩为一，\(\psi(1,0)\) 的秩为二。矩阵内共轭保持矩阵的秩，所以这两个嵌入不可能由内共轭相连。此例的环境 \(A\) 中心单，失败来自 \(B\) 有两个非零简单因子。对 \(\varphi\) 的表示，两种简单分量的重数为 \(1,2\)；对 \(\psi\) 则为 \(2,1\)。总维数同为三，分量重数却不同，正是前述证明只对单 Artin 环成立的原因。

\ProseParagraphBreak
中心性也不能任意删去。把 \(\mathbb C\) 看成实代数，它是单代数，但中心大于 \(\mathbb R\)。复共轭是非恒等实代数自同构，而交换代数中的每个内自同构都是恒等映射，所以复共轭并非内自同构。

\ProseParagraphBreak
非单代数甚至可以通过交换其不同因子产生外自同构。例如 \(k\times k\) 的映射 \((x,y)\mapsto(y,x)\) 保持代数结构，却不是内自同构，因为整个代数交换。因而“自同构都是内的”不是任意有限维代数的性质；它来自中心单性及前述简单模比较的共同作用。

% !TeX root = ../../Book2.tex
% 纲要标题：### 17.4 四元数代数、范数与分裂判据
% 本轮补充的纲要细目：
% * 以四元数左右乘构造范数相似变换，证明乘子及对角标量核；不由核的计算断言满射
% 纲要位置：计划纲要.md，第 1325 行。

\section{四元数代数、范数与分裂判据}
\label{b2:sec:1704}

四元数代数是次数二的中心单代数模型。它的乘法只需两个参数就能写下，是否分裂却与一个范数方程的可解性有关。本节假设 \(\operatorname{char}k\ne2\)，并固定 \(a,b\in k^\times\)；这一假设用于基的验证、共轭及二次扩张的描述。

% 纲要要点（供目录核对）：
%
% * 四元数代数
% * 矩阵代数及非分裂例子
% * 与二次型及范数方程的联系
%
% ---
%

\subsection{四元数代数}
\label{b2:subsec:1704:01}

\begin{definition}\label{b2:def:quaternion-algebra}
设 \(k\) 为特征不等于 \(2\) 的域，\(a,b\in k^\times\)。将自由结合代数 \(k\langle i,j\rangle\) 对下列关系生成的双边理想取商：
\[
i^2=a,\qquad j^2=b,\qquad ji=-ij
\]
所得含幺 \(k\) 代数称为\term[siyuanshu-daishu]{四元数代数}[quaternion algebra]，记为 \((a,b)_k\)。
\end{definition}
\symidx{\((a,b)_k\)：参数为 \(a,b\) 的四元数代数}

关系使每个字都能先把 \(i\) 移到 \(j\) 左边，再将两个指数降到零或一，所以 \(1,i,j,ij\) 生成这个代数。它们还线性无关，这一点不能仅由关系写法直接断言。

\begin{proposition}\label{b2:prop:quaternion-basis-splitting}
设 \(k\) 为特征不等于 \(2\) 的域，\(a,b\in k^\times\)，\(Q=(a,b)_k\)，则 \(Q\) 以 \(1,i,j,ij\) 为 \(k\) 基，并且是次数二的中心单 \(k\) 代数。若域扩张 \(K/k\) 含有 \(\alpha\) 满足 \(\alpha^2=a\)，则有分裂同构
\[
Q_K=K\otimes_kQ\longrightarrow M_2(K),\qquad
i\longmapsto I=\begin{pmatrix}\alpha&0\\0&-\alpha\end{pmatrix},
\quad j\longmapsto J=\begin{pmatrix}0&b\\1&0\end{pmatrix}.
\]
\end{proposition}
\begin{proof}
这两个矩阵满足 \(I^2=aI_2\)、\(J^2=bI_2\)、\(JI=-IJ\)，所以给出由商代数到矩阵代数的同态。在对角位置，\(I_2,I\) 线性无关，因为 \(2\alpha\ne0\)；在非对角位置，\(J,IJ\) 线性无关，因为其两个坐标组成的行列式为 \(-2\alpha b\ne0\)。故四个矩阵为 \(M_2(K)\) 的基。原四个生成元若有 \(k\) 线性关系，映射后也有同一关系，因此系数全零。这证明 \(\dim_kQ=4\)，扩张后同态便是基之间的同构。

现在验证 \(Q\) 本身中心单。若 \(U\) 为非零双边理想，则 \(K\otimes_kU\) 是 \(Q_K\cong M_2(K)\) 的非零理想，因矩阵环单而等于整个 \(Q_K\)。维数比较得 \(\dim_kU=4\)，所以 \(U=Q\)。若 \(z\in Z(Q)\)，其像在 \(M_2(K)\) 中为标量。以 \(I_2,I,J,IJ\) 为基比较系数，得到 \(z\) 只能是原 \(1\) 的 \(k\) 倍数，故中心为 \(k\)。维数为四使次数为二。
\end{proof}

对 \(q=x_0+x_1i+x_2j+x_3ij\)，定义它的\term[siyuanshu-gonge]{四元数共轭}[quaternion conjugation]为
\[
\bar q=x_0-x_1i-x_2j-x_3ij.
\]
生成元的像 \(i\mapsto-i,j\mapsto-j\) 在反代数中仍满足原关系，所以该映射反转乘法；它在上述基上作用两次为恒等，因此
\begin{equation}\label{b2:eq:quaternion-conjugation}
\overline{qr}=\bar r\bar q,\qquad\overline{\bar q}=q.
\end{equation}
直接按乘法关系展开，交叉项成对消去，得到
\begin{equation}\label{b2:eq:quaternion-norm}
q\bar q=\bar q q
=x_0^2-a x_1^2-b x_2^2+ab x_3^2.
\end{equation}

\begin{proposition}\label{b2:prop:quaternion-norm}
设 \(k\) 为特征不等于 \(2\) 的域，\(a,b\in k^\times\)，\(Q=(a,b)_k\)。对 \(q=x_0+x_1i+x_2j+x_3ij\in Q\)，令 \(\bar q=x_0-x_1i-x_2j-x_3ij\)。由于
\[
q\bar q=\bar q q=x_0^2-a x_1^2-b x_2^2+ab x_3^2\in k,
\]
可将这个标量记为 \(N(q)\)，称为\term[siyuanshu-fanshu]{四元数范数}[quaternion norm]。对任意 \(q,r\in Q\)，有
\[
N(qr)=N(q)N(r),\qquad
q\in Q^\times\Longleftrightarrow N(q)\ne0.
\]
可逆时 \(q^{-1}=N(q)^{-1}\bar q\)。在\cref{b2:prop:quaternion-basis-splitting}的分裂矩阵模型中，\(N(q)\) 就是矩阵行列式。
\end{proposition}
\begin{proof}
因 \(N(r)\in k\) 在中心，
\[
N(qr)=qr\bar r\bar q=q\cdot N(r)\bar q=N(q)N(r).
\]
若范数非零，\eqref{b2:eq:quaternion-norm}给出所述双侧逆。若 \(q\) 可逆，则乘法性给出 \(N(q)N(q^{-1})=N(1)=1\)，所以范数非零。

在分裂模型中，\(q\) 的矩阵为
\[
\begin{pmatrix}
x_0+\alpha x_1&b(x_2+\alpha x_3)\\
x_2-\alpha x_3&x_0-\alpha x_1
\end{pmatrix}.
\]
其行列式为 \(x_0^2-a x_1^2-bx_2^2+abx_3^2\)，正好是 \(N(q)\)。
\end{proof}
\symidx{\(N(q)=q\bar q\)：四元数范数}

这里的 \(N\) 是中心单代数约化范数在次数二情形的具体形式；\secref{b2:sec:1803}将从一般分裂矩阵模型重新得到这一识别。

\begin{proposition}\label{b2:prop:quaternion-norm-similitudes}
设 \(k\) 为特征不为二的域，\(Q=(a,b)_k\)、\(a,b\in k^\times\)，\(N(x)=x\bar x\) 为四元数范数。对 \(u,v\in Q^\times\)，定义 \(k\) 线性映射
\[
L_{u,v}:Q\longrightarrow Q,\qquad x\longmapsto uxv^{-1}.
\]
它是范数二次型的相似变换，乘子为 \(N(u)/N(v)\)。映射 \((u,v)\mapsto L_{u,v}\) 是从 \(Q^\times\times Q^\times\) 到该相似变换群的群同态，其核为 \(\{(c,c)\mid c\in k^\times\}\)。特别地，内自同构 \(x\mapsto uxu^{-1}\) 保持 \(N\)。
\end{proposition}
\begin{proof}
范数的对角系数为 \(1,-a,-b,ab\)，均非零，所以二次型非退化。\cref{b2:prop:quaternion-norm}的乘法性给出
\[
N(uxv^{-1})=\frac{N(u)}{N(v)}N(x).
\]
由特征不为二，极化同样使相应双线性型乘以这个标量。又有
\[
L_{u,v}L_{u',v'}(x)=uu'xv'^{-1}v^{-1}
=L_{uu',vv'}(x),
\]
故是群同态。若 \(L_{u,v}\) 为恒等，在 \(x=1\) 处即得 \(u=v\)；再对全部 \(x\) 使用等式，得到 \(u\) 属于 \(Q\) 的中心，也就是 \(k^\times\)。反过来，对角标量显然作用为恒等。这证明核的描述，而取 \(u=v\) 就得到等距结论。
\end{proof}

这给出一族由代数乘法产生的几何变换。仅从核的计算不能断言它们穷尽全部相似变换；这里只使用上述明确构造，也不预先调用正交群的进一步分类。

\subsection{矩阵代数与非分裂例子}
\label{b2:subsec:1704:02}

\begin{proposition}\label{b2:prop:quaternion-division-or-split}
设 \(k\) 为特征不等于 \(2\) 的域，\(a,b\in k^\times\)，\(Q=(a,b)_k\)，\(N:Q\rightarrow k\) 为其四元数范数，则 \(Q\) 或者是除代数，或者同构于 \(M_2(k)\)。它分裂，当且仅当二次型 \(N\) 有非平凡零点，即存在 \(q\in Q\setminus\{0\}\) 使 \(N(q)=0\)。
\end{proposition}
\begin{proof}
由\cref{b2:prop:csa-division-description}，\(Q\cong M_r(D)\)，其中 \(D\) 是除代数，且 \(4=r^2\dim_kD\)。若 \(r=1\)，\(Q\) 是除代数；否则只能 \(r=2,\dim_kD=1\)，即 \(D=k\)，所以 \(Q\cong M_2(k)\)。

由\cref{b2:prop:quaternion-norm}，每个非零元素可逆当且仅当每个非零元素的范数都非零。因此有非平凡零点恰好等价于不是除代数，按前一分类又等价于分裂。
\end{proof}

若 \(a=\alpha^2\in k^{\times2}\)，\cref{b2:prop:quaternion-basis-splitting}的两个矩阵已经定义在 \(k\) 中，因此 \((a,b)_k\) 直接分裂。同样，交换生成元 \(i,j\) 给出 \((a,b)_k\cong(b,a)_k\)，所以 \(b\) 为平方时也分裂。

\ProseParagraphBreak
实四元数代数
\[
\mathbb H=(-1,-1)_{\mathbb R}
\]
则不分裂，因为其范数为 \(x_0^2+x_1^2+x_2^2+x_3^2\)，对任意非零实向量严格为正。因而每个非零四元数都可逆；矩阵代数 \(M_2(\mathbb R)\) 却有非零奇异矩阵，二者不可能同构。扩张到 \(\mathbb C\) 后，参数 \(-1\) 成为平方，所以 \(\mathbb H\otimes_{\mathbb R}\mathbb C\cong M_2(\mathbb C)\)。

\ProseParagraphBreak
更一般地，对非零实数 \(a,b\)，只要其中一个为正，就因该参数为平方而分裂；若二者都为负，分别把 \(i,j\) 除以 \(\sqrt{-a},\sqrt{-b}\)，得到平方均为 \(-1\) 的生成元，故代数同构于 \(\mathbb H\)。因此实四元数代数恰有分裂矩阵代数与 \(\mathbb H\) 两种类型。这里分类的是四维四元数代数，没有在此断言所有实中心单代数的进一步分类。

\subsection{二次型与范数方程}
\label{b2:subsec:1704:03}

\begin{theorem}[四元数分裂的等价条件]\label{b2:thm:quaternion-splitting}
设 \(k\) 为特征不等于 \(2\) 的域，\(a,b\in k^\times\)，\(Q=(a,b)_k\)，其四元数范数为 \(N(q)=q\bar q\)。则
\[
Q\cong M_2(k)
\quad\Longleftrightarrow\quad Q\text{ 不是除代数}
\quad\Longleftrightarrow\quad
N\text{ 有非平凡零点}.
\]
若 \(a\) 不是 \(k\) 中的平方，令 \(L=k(\alpha)\)、\(\alpha^2=a\)，则上述条件还等价于
\[
b\in N_{L/k}(L^\times)
\quad\Longleftrightarrow\quad
b=x^2-a y^2\text{ 对某些 }x,y\in k.
\]
若 \(a\) 为非零平方，代数总分裂，最后一个方程也总有解。
\end{theorem}
\begin{proof}
\cref{b2:prop:quaternion-division-or-split}已经给出前三个条件的等价。以下处理 \(a\) 非平方时的域范数条件。

\(a\) 非平方时，\(L=k[i]\subseteq Q\) 是二次域，且 \(Q=L\oplus Lj\)。记 \(L\) 的非平凡自同构为 \(u\mapsto\bar u\)，则 \(ju=\bar u j\)。由前述范数公式，对 \(u,v\in L\) 有
\[
N(u+vj)=N_{L/k}(u)-b\cdot N_{L/k}(v).
\]
这里 \(N_{L/k}(x+y\alpha)=(x+y\alpha)(x-y\alpha)=x^2-a y^2\)，可由二次扩张的乘法矩阵或两个共轭直接计算。

若 \(Q\) 分裂，\cref{b2:prop:quaternion-division-or-split}给出非零 \(u+vj\) 的范数为零。不能有 \(v=0\)，否则域范数零使 \(u=0\)，矛盾。因此 \(v\ne0\)，得到
\(b=N_{L/k}(u)/N_{L/k}(v)=N_{L/k}(u/v)\)。反过来，若 \(b=N_{L/k}(z)\)，则 \(z+j\ne0\) 且其范数为零，所以代数分裂。域范数的显式公式又给出最后一个等价。

若 \(a=\alpha^2\) 且 \(\alpha\in k^\times\)，分裂性已证明，而
\[
x=\frac{b+1}{2},\qquad y=\frac{b-1}{2\alpha}
\]
满足 \(x^2-a y^2=b\)，故方程仍可解。
\end{proof}

已知范数方程的一个解，还可以直接构造分裂同构。若 \(a\) 非平方，取 \(z=x+y\alpha\in L\) 满足 \(N_{L/k}(z)=b\)。在二维 \(k\) 空间 \(L\) 上，让 \(i\) 作用为乘 \(\alpha\)，让 \(j\) 作用为 \(v\mapsto z\bar v\)。两者满足平方关系，且
\(z\overline{\alpha v}=-\alpha z\bar v\)，所以反交换，给出 \(Q\to\operatorname{End}_k(L)\)。它保持单位元，由单性为单射，两端维数同为四，故同构。在基 \(1,\alpha\) 下，具体矩阵为
\[
i\longmapsto\begin{pmatrix}0&a\\1&0\end{pmatrix},\qquad
j\longmapsto\begin{pmatrix}x&-ay\\y&-x\end{pmatrix}.
\]
后者的平方为 \((x^2-a y^2)I_2=bI_2\)，可作为公式的直接核验。

\ProseParagraphBreak
范数判据也等价于齐次方程
\[
X^2-aY^2-bZ^2=0
\]
有一个非零解。若 \(Z\ne0\)，除以 \(Z^2\) 就得到 \(b=(X/Z)^2-a(Y/Z)^2\)；若 \(Z=0\)，非零解必有 \(Y\ne0\)，于是 \(a=(X/Y)^2\)，同样推出分裂。反过来，范数方程解 \(b=x^2-a y^2\) 给出点 \((x,y,1)\)，平方参数时也可直接取 \((\alpha,1,0)\)。这样，四元数代数是否为矩阵代数，已经化为一个具体二次方程的可解性问题。

这里的范数既是二次型，也是除代数或矩阵代数的判别工具。若要比较它与 Clifford 代数、Witt 环中的不变量以及矩阵代数上的伴随对合，可接着读\secref{b2:sec:B02}和\secref{b2:sec:B03}；后者需要本章的中心单代数与 Skolem--Noether 定理。


\begin{chaptersummary}
中心单代数的正则双模是简单模，其双模自同态只有底域标量。稠密性因此把左右乘作用识别为全部线性自同态，得到包络代数同构。这个判据同时解释了任意标量扩张和张量积为何保持中心单性；在代数闭包上，结构定理又把代数化为全矩阵环，从而确定其次数。

\ProseParagraphBreak
一般简单子代数的中心可能大于底域，中心化子的中心恰好随之扩大。双中心化子和维数公式仍然成立；在共同中心 \(L\) 上取张量积，所得的是 \(C_A(L)\)，未必是整个 \(A\)。一般中心单代数中的包含极大子域未必有次数 \(\deg A\)；自中心化子域以及除代数中的最大子域则达到这个次数，并可分裂原代数。有限可分分裂域的证明通过非零多项式选择可分元素；约化特征多项式及其系数下降将在下一章建立。

\ProseParagraphBreak
Skolem–Noether 定理把简单代数的两种嵌入改写成两个同维模，由模的同构产生实现内共轭的可逆元。四元数代数则把这些抽象结构落实为四个基向量、四元数共轭和一个乘法范数：范数二次型有非平凡零点恰好等价于分裂，非平方参数情形又等价于二次扩张中的范数方程。左右乘以可逆四元数又给出范数的相似变换，内共轭则给出等距变换。下一章将用稳定等价、循环代数与约化范数进一步组织这些现象。
\end{chaptersummary}

\BookExercises

\begin{exercise}\label{b2:ex:17-simple-central-distinction}
设 \(k\) 为域、\(n\ge1\)，\(L/k\) 为非平凡有限扩张。分别说明 \(k[t]/(t^2)\)、\(L\)、\(M_n(k)\) 是否单、是否以 \(k\) 为中心。解释这些例子为何不能把定义中的两个要求合并为一个。
\end{exercise}
\begin{solution}
双数代数交换，中心是整个二维代数而非 \(k\)，且含非零真双边理想 \((\bar t)\)，所以既不单也不中心。域 \(L\) 是单环，但其中心为 \(L\ne k\)，所以不是中心单 \(k\) 代数。\(M_n(k)\) 单且中心为 \(kI_n\)，因此中心单。还可用 \(T_2(k)\) 分离另一方向：与 \(E_{11},E_{12}\) 交换的上三角矩阵只有标量矩阵，所以它的中心是 \(k\)，但严格上三角矩阵组成非零真双边理想，故它中心而不单。
\end{solution}

\begin{exercise}\label{b2:ex:17-inseparable-tensor}
设 \(k\) 为特征不等于二的域，\(d\in k^\times\) 非平方，\(L=k(\alpha)\)、\(\alpha^2=d\)。求 \(L\otimes_kL\) 的两个非平凡中心幂等元、全部理想及 Jacobson 根，说明它的非单性与\cref{b2:prop:inseparable-base-change-counterexample}中的非单性有何区别。
\end{exercise}
\begin{solution}
在 \(L\otimes_kL\) 中用第一个因子赋予左 \(L\) 结构，令 \(w=\alpha^{-1}\otimes\alpha\)，则 \(w^2=1\)，且 \(1,w\) 为 \(L\) 基。因此
\[
e_+=\frac{1+w}{2},\qquad e_-=\frac{1-w}{2}
\]
是非零正交幂等元，和为一，分别生成一维 \(L\) 因子。具体同构为
\[
L\otimes_kL\longrightarrow L\times L,\qquad
u\otimes v\longmapsto(uv,u\sigma(v)),
\]
其中 \(\sigma(\alpha)=-\alpha\)。它把 \(w\) 送到 \((1,-1)\)，故把 \(e_+,e_-\) 送到两个标准幂等元，保乘法并在基上为同构。

全部理想为 \(0,Le_+,Le_-\) 及全环，两极大理想的交为零，所以 Jacobson 根为零。这次非单来自分成两个域因子，没有非零幂零元；纯不可分反例则出现非零幂零元。两种情形的原代数 \(L\) 都不以 \(k\) 为中心。
\end{solution}

\begin{exercise}\label{b2:ex:17-quaternion-tensor-square}
设 \(k\) 为特征不等于二的域，\(a,b\in k^\times\)，\(Q=(a,b)_k\)。在 \(Q\otimes_kQ\) 中令
\[
s=-a^{-1}(i\otimes i),\qquad t=-b^{-1}(j\otimes j).
\]
证明 \(s,t\) 为交换的对合元，并求四个元素 \(e_{\epsilon,\delta}=(1+\epsilon s)(1+\delta t)/4\)、\(\epsilon,\delta\in\{1,-1\}\)，在四维空间 \(Q\) 上的作用。它们是否为本原幂等元？
\end{exercise}
\begin{solution}
由平方关系，\(s^2=t^2=1\)。两因子中的反交换同时产生两个负号，故 \(st=ts=(ab)^{-1}(ij\otimes ij)\)。由\cref{b2:prop:quaternion-tensor-square}，张量积作用为 \((q\otimes r)x=qx\bar r\)。于是 \(s\) 作用为 \(x\mapsto ixi^{-1}\)，在基 \(1,i,j,ij\) 上的符号为 \((+,+,-,-)\)；\(t\) 的符号为 \((+,-,+,-)\)。

因此 \(e_{+,+},e_{+,-},e_{-,+},e_{-,-}\) 分别投影到 \(k1,ki,kj,kij\)。它们非零、两两正交且和为一，矩阵模型中各有秩一。其角代数都是相应一维空间的自同态代数 \(k\)，没有非平凡幂等元，故四者都是本原幂等元。原四元数是否分裂不影响这个张量平方中的分解。
\end{solution}

\begin{exercise}\label{b2:ex:17-dimension-sixteen}
设 \(k\) 为域，\(A\) 为十六维中心单 \(k\) 代数。对 \(A\cong M_r(D)\)，列出可能的矩阵大小 \(r\) 及 \(\dim_kD\)。若 \(k\) 代数闭，或 \(k\) 有限，分别还能得到什么结论？
\end{exercise}
\begin{solution}
\(\deg A=4\)，而 \(\dim_kD\) 为平方，记 \(\deg D=d\)。维数关系为 \(16=r^2d^2\)，故 \(rd=4\)。可能的 \((r,\dim_kD)\) 为 \((1,16),(2,4),(4,1)\)。这只是参数的可能形式，不保证每个底域都有各类除代数。代数闭时有限维除代数只能是 \(k\)，所以 \(A\cong M_4(k)\)。有限底域时 \(D\) 本身也是有限除环，由\cref{b2:thm:wedderburn-little}交换，再由中心为 \(k\) 得 \(D=k\)，结论相同。
\end{solution}

\begin{exercise}\label{b2:ex:17-finite-field-maximal-subfield}
设 \(q\) 为素数幂、\(n\ge1\)。在 \(A=M_n(\mathbb F_q)\) 中构造一个次数为 \(n\) 的有限可分子域，并证明它的中心化子恰为自身。说明这个构造不需要在有限域上使用“非零多项式总有非零取值”的错误断言。
\end{exercise}
\begin{solution}
第一卷定理17.1.1保证存在有限域 \(L=\mathbb F_{q^n}\)，并可通过该卷定理17.2.3的子域描述包含 \(\mathbb F_q\)。选取 \(L\) 的 \(\mathbb F_q\) 基，左乘作用给出嵌入 \(L\hookrightarrow\operatorname{End}_{\mathbb F_q}(L)\cong M_n(\mathbb F_q)\)。若 \(T\) 与全部左乘交换，则 \(T(x)=x\cdot T(1)\)，故它也是乘以 \(L\) 中某个元素，中心化子就是 \(L\)。有限域扩张可分，因为每个元素满足无重根多项式 \(X^{q^n}-X\)。这里使用现成的有限域扩张与正则作用；例如 \(X^q-X\) 虽是非零多项式，却在 \(\mathbb F_q\) 每一点取零，不能作为无限域取点论证的类比。
\end{solution}

\begin{exercise}\label{b2:ex:17-tensor-factor-centralizer}
设 \(k\) 为域，\(A,B\) 为中心单 \(k\) 代数，在 \(C=A\otimes_kB\) 中把 \(A,B\) 识别为 \(A\otimes1,1\otimes B\)。证明它们互为在 \(C\) 内的中心化子。
\end{exercise}
\begin{solution}
这两个嵌入单射，因为 \(1\) 可以分别扩充为两个代数的基，张量基保证没有新线性关系。两因子逐元素交换，所以 \(1\otimes B\subseteq C_C(A\otimes1)\)。由\cref{b2:thm:csa-tensor-product}，环境 \(C\) 中心单；由\cref{b2:thm:csa-centralizer}，右端维数为 \(\dim_kC/\dim_kA=\dim_kB\)。故包含是相等，反方向交换两因子的角色即可。这里的中心化子在张量代数 \(C\) 内计算，而不是在其底向量空间的全部自同态中计算。
\end{solution}

\begin{exercise}\label{b2:ex:17-field-centralizer}
设 \(k\) 为域，\(L/k\) 为有限扩张、\([L:k]=d\)，\(V=L^r\)、\(r\ge1\)。将 \(L\) 通过标量乘法嵌入 \(A=\operatorname{End}_k(V)\)。求其中心化子和双中心化子，核对维数公式，且不假设 \(L/k\) 可分。
\end{exercise}
\begin{solution}
与全部 \(L\) 标量交换的 \(k\) 线性映射恰是 \(L\) 线性映射，因此中心化子为 \(\operatorname{End}_L(L^r)\cong M_r(L)\)。\(A\) 中心单，\(L\) 为简单子代数，\cref{b2:thm:csa-centralizer}给出双中心化子为原来的 \(L\)。维数分别是 \(\dim_kA=r^2d^2\)、\(\dim_kL=d\)、\(\dim_kM_r(L)=r^2d\)，确有 \(r^2d^2=d\cdot r^2d\)。
\end{solution}

\begin{exercise}\label{b2:ex:17-large-commutative-subalgebra}
设 \(k\) 为域，在 \(M_5(k)\) 中按 \(2+3\) 分块，令
\[
B=\left\{cI_5+\begin{pmatrix}0&X\\0&0\end{pmatrix}:c\in k, X\in M_{2\times3}(k)\right\}.
\]
计算 \(C_{M_5(k)}(B)\)、\(J(B)\) 及 \(B^\times\)，判断 \(B\) 是否为最大交换子代数，并比较其维数与环境代数的次数。
\end{exercise}
\begin{solution}
右上块的两个乘积为零，故 \(B\) 交换，维数为七。将一般矩阵写成 \(T=\begin{psmallmatrix}P&Q\\U&V\end{psmallmatrix}\)。与全部右上块交换等价于
\[
XU=0,\qquad UX=0,\qquad PX=XV
\quad\text{对所有 }X\in M_{2\times3}(k).
\]
用单条目矩阵 \(X=E_{ij}\) 检查，前两式迫使 \(U=0\)；第三式先迫使 \(P,V\) 都对角，再迫使所有对角条目为同一个 \(c\)。因此 \(T\in B\)，而 \(Q\) 任意，故中心化子就是 \(B\)。它因而最大交换。

令 \(I\) 为严格右上块组成的理想，则 \(I^2=0\)，且 \(B/I\cong k\)。幂零理想包含于根，商为域又使根不可能更大，故 \(J(B)=I\)。元素 \(cI_5+N_X\) 可逆当且仅当 \(c\ne0\)，此时逆为 \(c^{-1}I_5-c^{-2}N_X\)。其维数七大于 \(\deg M_5(k)=5\)，含非零平方零元，所以不是域。这是\cref{b2:prop:maximal-commutative-not-field}的另一个具体模型。
\end{solution}

\begin{exercise}\label{b2:ex:17-subfield-degree-divides}
令 \(D=(-1,-1)_{\mathbb Q}\)。比较二次域 \(K_+=\mathbb Q(\sqrt2)\) 与 \(K_-=\mathbb Q(\sqrt{-1})\)：它们能否嵌入 \(D\)，能否分裂 \(D\)？由此说明 \([K:\mathbb Q]\mid\deg D\) 能否作为嵌入的充分条件。
\end{exercise}
\begin{solution}
\(D\) 的范数为四个有理坐标的平方和，所以非零元素范数非零，\(D\) 是次数二的除代数。\(K_-\) 通过 \(\sqrt{-1}\mapsto i\) 嵌入 \(D\)，由\cref{b2:prop:quaternion-basis-splitting}也分裂 \(D\)。

若 \(K_+\) 嵌入 \(D\)，则某个 \(q=x_0+v\)、\(v=x_1i+x_2j+x_3ij\)，满足 \(q^2=2\)。其非标量部分为 \(2x_0v\)。若 \(v=0\)，就有有理数 \(x_0^2=2\)，不可能；若 \(v\ne0\)，则 \(x_0=0\)，而
\[
q^2=-(x_1^2+x_2^2+x_3^2)\le0,
\]
也不可能等于二。因此 \(K_+\) 不嵌入。将 \(K_+\) 视为实数子域后，同一个平方和范数在 \(K_+\) 上仍无非平凡零点，所以 \(D_{K_+}\) 仍为除代数，不分裂。两个域的次数都是二，都满足整除约束；\cref{b2:prop:csa-self-centralizing-field}给出的整除关系只是嵌入的必要条件。
\end{solution}

\begin{exercise}\label{b2:ex:17-inseparable-skolem}
设 \(k=\mathbb F_2(t)\)、\(L=k(\alpha)\)、\(\alpha^2=t\)。在 \(M_2(k)\) 中分别将 \(\alpha\) 送到
\[
I=\begin{pmatrix}0&t\\1&0\end{pmatrix},\qquad
J=\begin{pmatrix}0&1\\t&0\end{pmatrix}.
\]
证明这给出两个 \(k\) 代数嵌入，并找出实现共轭的可逆矩阵。
\end{exercise}
\begin{solution}
两矩阵平方都为 \(tI_2\)。\(t\) 不是 \(k\) 中的平方，故 \(X^2-t\) 不可约，两种代入都给出域 \(L\) 的含幺同态，因而单射。取 \(P=\operatorname{diag}(1,t)\)，直接计算得 \(PIP^{-1}=J\)。此扩张纯不可分，而两个嵌入仍共轭；这核验了 Skolem–Noether 定理不需要源代数中心可分的一个具体情形。
\end{solution}

\begin{exercise}\label{b2:ex:17-normalizer-nonsplit}
设 \(k\) 为特征不等于二的域，\(d\in k^\times\) 非平方，\(L=k(\sqrt d)\)。通过乘法作用把 \(L\) 嵌入 \(A=M_2(k)\)，其中
\[
\sqrt d\longmapsto T=\begin{pmatrix}0&d\\1&0\end{pmatrix}.
\]
求 \(N_{A^\times}(L)\)，并证明其到 \(\operatorname{Gal}(L/k)\) 的正合列具有群截面。与实四元数中的情形比较。
\end{exercise}
\begin{solution}
正则作用的中心化子为 \(L\)，或直接解与 \(T\) 交换的方程，得到矩阵 \(\begin{psmallmatrix}x&dy\\y&x\end{psmallmatrix}\)。其可逆元就是 \(L^\times\)。令 \(C=\operatorname{diag}(1,-1)\)，则 \(CTC^{-1}=-T\)、\(C^2=1\)，所以 \(C\) 实现二次域的非平凡自同构 \(\sigma\)。

任意正规化元限制到 \(L\) 后为恒等或 \(\sigma\)；前者属于 \(L^\times\)，后者与 \(C\) 的比值属于 \(L^\times\)。故
\[
N_{A^\times}(L)=L^\times\cup L^\times C
\cong L^\times\rtimes\operatorname{Gal}(L/k),
\]
其中 \(\sigma(z)=\bar z\)。映射 \(1\mapsto I_2\)、\(\sigma\mapsto C\) 是群截面。\cref{b2:prop:quaternion-normalizer-nonsplit}中所有实现复共轭的 \(zj\) 平方都为负实数；这里则确有平方为一的实现元，故两种正规化子扩张表现不同。
\end{solution}

\begin{exercise}\label{b2:ex:17-quaternion-parameter-change}
在 \(Q=(2,3)_{\mathbb Q}\) 中改取 \(i\) 与 \(j'=(1+i)j\) 为生成元。求新参数，并把得到的同构在两组四元数基中写成坐标变换。直接核对这个变换保留四元数范数。再判断 \((2,-6)_{\mathbb Q}\) 是否同构于 \(Q\)。
\end{exercise}
\begin{solution}
\((1+i)(1-i)=-1\)，由扭曲关系得
\[
(j')^2=(1+i)(1-i)j^2=-3,\qquad j'i=-ij'.
\]
又 \((1+i)^{-1}=i-1\)，所以 \(j=(i-1)j'\) 可恢复，给出 \((2,-3)_{\mathbb Q}\cong Q\)。把新基向量写回原基，有 \(j'=j+ij\)、\(ij'=2j+ij\)，故从新基坐标到原基坐标的变换为
\[
(x_0,x_1,x_2,x_3)\longmapsto
(x_0,x_1,x_2+2x_3,x_2+x_3).
\]
其矩阵行列式为负一，确为可逆变换。代入原范数，得到
\[
x_0^2-2x_1^2-3(x_2+2x_3)^2+6(x_2+x_3)^2
=x_0^2-2x_1^2+3x_2^2-6x_3^2,
\]
正是新参数 \((2,-3)\) 的范数。再把生成元 \(j\) 换成 \(ij\)，其平方为 \(-2\cdot3=-6\)，所以 \((2,-6)_{\mathbb Q}\) 也同构于 \(Q\)，符合\cref{b2:prop:quaternion-parameter-changes}。
\end{solution}

\begin{exercise}\label{b2:ex:17-rational-quaternions}
判断 \((2,3)_{\mathbb Q}\) 与 \((2,7)_{\mathbb Q}\) 哪一个分裂。对分裂者给出两生成元的 \(2\times2\) 有理矩阵；对不分裂者证明对应范数方程无有理解。
\end{exercise}
\begin{solution}
因为 \(7=3^2-2\cdot1^2\)，由\cref{b2:thm:quaternion-splitting}，\((2,7)_{\mathbb Q}\) 分裂，可取
\[
i\longmapsto\begin{pmatrix}0&2\\1&0\end{pmatrix},\qquad
j\longmapsto\begin{pmatrix}3&-2\\1&-3\end{pmatrix}.
\]
两矩阵平方分别为 \(2I_2,7I_2\)，并反交换，所以\cref{b2:prop:algebra-relations-universal}给出含幺同态。由\cref{b2:prop:quaternion-basis-splitting}，来源单且为四维，故这个非零同态单射，两边同维又使它成为同构。

若 \(x^2-2y^2=3\) 有有理解，清分母得到互素整数 \(u,v,w\)、\(w\ne0\)，满足 \(u^2-2v^2=3w^2\)。模三后为 \(u^2+v^2=0\)，只能 \(3\mid u,v\)，因为模三平方只有零和一。原等式继而使 \(3\mid w\)，与互素矛盾。因此该范数方程无解，由\cref{b2:thm:quaternion-splitting}，\((2,3)_{\mathbb Q}\) 不分裂；再由\cref{b2:prop:quaternion-division-or-split}，它为除代数。
\end{solution}

\begin{exercise}\label{b2:ex:17-quaternion-quadratic-centralizer}
设 \(k\) 为特征不等于二的域、\(a,b\in k^\times\)，在 \(Q=(a,b)_k\) 中求 \(C_Q\bigl(k[i]\bigr)\)，并证明 \(k[i]\) 是最大交换子代数。分别说明 \(a\) 非平方及 \(a=\alpha^2\) 时这个子代数的结构。
\end{exercise}
\begin{solution}
写 \(q=x_0+x_1i+x_2j+x_3ij\)，则
\(qi-iq=-2x_2ij-2a x_3j\)。基线性无关且 \(2a\ne0\)，所以交换当且仅当 \(x_2=x_3=0\)，即中心化子为 \(k+ki=k[i]\)。任何更大的交换子代数都落在这个中心化子内，故它最大交换。\(a\) 非平方时它是二次域；若 \(a=\alpha^2\)，则两个幂等元 \((1+i/\alpha)/2\)、\((1-i/\alpha)/2\) 非零、正交且和为一，给出 \(k[i]\cong k\times k\)。后者仍最大交换，却不是子域。
\end{solution}

\begin{exercise}\label{b2:ex:17-finite-quaternion-norm-count}
设 \(k=\mathbb F_q\)、\(q\) 为奇数，\(a,b\in k^\times\)。当 \(a\) 非平方时，证明方程 \(x^2-a y^2=b\) 恰有 \(q+1\) 组解；当 \(a\) 为平方时，恰有 \(q-1\) 组解。由此再证所有 \(k\) 上的四元数代数都分裂。
\end{exercise}
\begin{solution}
非平方时 \(L=\mathbb F_{q^2}=k(\sqrt a)\)，该方程是 \(N_{L/k}(x+y\sqrt a)=b\)。Frobenius 映射\footnote{弗罗贝尼乌斯（Ferdinand Georg Frobenius，1849-1917），德国数学家，研究群表示、矩阵和有限域。} \(z\mapsto z^q\) 固定 \(k\)，却不在全部 \(L\) 上为恒等，因为 \(X^q-X\) 至多有 \(q\) 个根。因此它是这个二次扩张的非平凡自同构，范数为 \(z\mapsto z\cdot z^q=z^{q+1}\)。第一卷推论17.2.1说明 \(L^\times\) 为阶 \(q^2-1\) 的循环群，因此这个映射核大小为 \(q+1\)，像大小为 \(q-1\)，正好满射到 \(k^\times\)。每个非零范数纤维都有 \(q+1\) 个元素，对应唯一的坐标对 \((x,y)\)。

若 \(a=\alpha^2\)，可逆线性换元 \(u=x+\alpha y,v=x-\alpha y\) 把方程化为 \(uv=b\)。每个 \(u\in k^\times\) 唯一确定 \(v=b/u\)，所以有 \(q-1\) 组解。两种情形都可解，\cref{b2:thm:quaternion-splitting}于是给出结论。
\end{solution}

\begin{exercise}\label{b2:ex:17-quaternion-automorphisms-norm}
设 \(k\) 为特征不等于二的域、\(a,b\in k^\times\)、\(Q=(a,b)_k\)。证明 \(Q\) 的每个 \(k\) 代数自同构都保留范数及共轭。说明这里不是单纯从某一组生成元的名字看出共轭不变性。
\end{exercise}
\begin{solution}
由\cref{b2:cor:csa-automorphisms-inner}，任一自同构为 \(q\mapsto uqu^{-1}\)。范数乘法性给出
\[
N(uqu^{-1})=N(u)N(q)N(u)^{-1}=N(q),
\]
所以范数被保持。记 \(T(q)=q+\bar q=2x_0\)。范数公式给出 \(T(q)=N(q+1)-N(q)-1\)，这是由代数的单位元及已经证明不变的范数决定的。因此 \(T\) 也不变，进而由 \(\bar q=T(q)-q\) 得共轭与自同构交换。证明没有要求自同构固定原来的 \(i,j\)。
\end{solution}

\begin{exercise}\label{b2:ex:17-quaternion-conjugation-rotation}
在实四元数 \(\mathbb H\) 中记 \(k_0=ij\)，并取 \(u=i+j\)。计算 \(x\mapsto uxu^{-1}\) 在基 \(1,i,j,k_0\) 上的作用，求其在三维实子空间 \(\mathbb Ri\oplus\mathbb Rj\oplus\mathbb Rk_0\) 上的行列式。
\end{exercise}
\begin{solution}
由 \(u^2=-2\)，有 \(u^{-1}=-(i+j)/2\)。按 \(i^2=j^2=-1\)、\(ji=-ij\) 相乘，得到 \(uiu^{-1}=j\)、\(uju^{-1}=i\)；该内自同构保持乘法，所以 \(uk_0u^{-1}=ji=-k_0\)，并且固定 \(1\)。三维子空间上的矩阵为
\[
\begin{pmatrix}0&1&0\\1&0&0\\0&0&-1\end{pmatrix},
\]
其行列式为一。\cref{b2:prop:quaternion-norm-similitudes}保证它保持范数的限制 \(x_1^2+x_2^2+x_3^2\)，这里可直接看出它固定 \(i+j\) 的方向，并把两个正交方向 \(i-j,k_0\) 变号。
\end{solution}

\begin{exercise}\label{b2:ex:17-explicit-quadratic-conjugator}
设 \(k\) 为特征不等于二的域、\(a,b\in k^\times\)，\(a\) 非平方、\(L=k(i)\subseteq(a,b)_k\)，\(i^2=a\)，并设 \(u\in L^\times\) 满足 \(N_{L/k}(u)=1\)。显式构造 \(z\in L^\times\)，使共轭 \(q\mapsto zqz^{-1}\) 固定 \(L\)，并将 \(j\) 送到 \(uj\)。
\end{exercise}
\begin{solution}
对 \(z\in L^\times\)，有 \(zjz^{-1}=z\bar z^{-1}j\)，所以只需求 \(z/\bar z=u\)。若 \(u\ne-1\)，取 \(z=1+u\ne0\)。由于 \(\bar u=u^{-1}\)，有 \(\bar z=1+u^{-1}=u^{-1}(1+u)\)，故 \(z/\bar z=u\)。若 \(u=-1\)，取 \(z=i\)，则 \(\bar z=-i\)，也有 \(z/\bar z=-1\)。\(L\) 交换，故上述共轭在 \(L\) 上恒等。这个构造给出 Skolem–Noether 共轭在二次域情形中的一个明确实现元。
\end{solution}

\begin{exercise}\label{b2:ex:17-nilpotent-maximal-commutative}
设 \(k\) 为域，在 \(A=M_4(k)\) 中令 \(N=\operatorname{diag}(J_2(0),J_2(0))\)。求 \(k[N]\) 的维数和中心化子。它是否最大交换？再求
\[
B=\operatorname{diag}\bigl(k[J_2(0)],k[J_2(0)]\bigr)
\]
的中心化子，解释单个 Jordan 块的条件为何重要。
\end{exercise}
\begin{solution}
\(N\ne0\)、\(N^2=0\)，故 \(k[N]\) 维数为二。把 \(k^4\) 视为 \((k[t]/(t^2))^2\)，\(N\) 就是乘 \(t\) 的算子。与它交换的算子恰为这个模的自同态，因而
\[
C_A(k[N])\cong M_2(k[t]/(t^2)),
\]
其 \(k\) 维数为八。例如 \(\operatorname{diag}(J_2(0),0)\) 与 \(N\) 交换，却不在 \(k[N]\) 中，添入它产生更大交换代数，所以 \(k[N]\) 不最大交换。

\(B\) 含两个块投影，故其中心化子必须分块对角；每个对角块又必须与 \(J_2(0)\) 交换，所以各块为 \(k[J_2(0)]\)。因此 \(C_A(B)=B\)，\(B\) 最大交换，维数为四，仍不是域。\cref{b2:prop:maximal-commutative-not-field}中的单个 Jordan 块有一个循环向量，此处两块的模有两份循环分量，允许额外的分量同态，故中心化子不再只是算子的多项式。
\end{solution}

% !TeX root = ../../Book2.tex
% 纲要标题：## 第18章　Brauer 群与循环代数
% 纲要位置：计划纲要.md，第 1626 行。

\chapter{Brauer 群与循环代数}
\label{b2:ch:18}
\BookChapterNumber{b2:ch:18}

\begin{chapterabstract}
矩阵大小改变中心单代数的维数，却不改变其除代数代表。Brauer 群把这部分稳定信息组织成张量积下的交换群。循环代数使其中的一些类能够由一个 Galois 自同构与一个范数参数计算；约化迹和约化范数则把分裂矩阵中的运算准确地降回底域。对有限 Galois 扩张 \(K/k\)，交叉积与\(2\)-余圈描述相对 Brauer 群 \(\operatorname{Br}(K/k)\)；最后证明周期整除指数，并给出周期二、指数四的具体除代数。
\end{chapterabstract}

\begin{learninggoals}
\begin{itemize}
\item 用除代数代表判断 Brauer 等价，区分代数的次数、指数与 Brauer 类的周期。
\item 写出循环代数的乘法及分裂矩阵，把分裂性转化为范数方程。
\item 计算约化特征多项式、迹和范数，说明它们与正则作用及子域运算的关系。
\item 从半线性矩阵记录余圈，核对换基、余边及反代数的方向。
\item 运用有限分裂域次数的整除性，辨认周期与指数可能出现的差别。
\end{itemize}
\end{learninggoals}

若 \(D\) 是域 \(k\) 上的中心除代数，\(D\) 与 \(M_r(D)\) 的维数不同，二者的模却保留了相同的除代数部分。把矩阵大小当作分类数据，便会把这种稳定信息遮住；Brauer 等价正是将矩阵因子略去后比较中心单代数。张量积给出这些类的加法，但要证明它确实形成群，还须给出逆元并核对选择矩阵代表时运算不变。

\ProseParagraphBreak

仅有抽象群运算还难以计算具体类。若一个循环 Galois 扩张已经使代数分裂，就可以把自同构和一个底域元素写进乘法关系，得到循环代数；它何时已经是矩阵代数，又归结到范数方程。另一方面，分裂后矩阵的迹与行列式若要回到底域，必须证明它们不依赖分裂同构。这些可算对象把稳定分类与原代数中的运算连接起来。

\ProseParagraphBreak

本章需要第十七章的中心单代数、可分分裂域、中心化子及 Skolem--Noether 定理，也需要第一卷的有限 Galois 扩张与域范数。\secref{b2:sec:1801}建立 Brauer 等价与群运算，\secref{b2:sec:1802}给出循环代数的直接构造，\secref{b2:sec:1803}处理分裂模型中迹与行列式的下降，\secref{b2:sec:1804}由余圈方程研究相对 Brauer 群，\secref{b2:sec:1805}比较周期与指数。

\ProseParagraphBreak
\(\mathbb C/\mathbb R\)、有限域及二次有理扩张上的循环代数，为一般交叉积提供了可计算的例子。中心单代数、循环代数、约化运算与周期和指数的进一步论述，可对照 Gille–Szamuely 的《Central Simple Algebras and Galois Cohomology》\cite[Sections 2.4--2.8; Chapter 3]{GilleSzamuely2017CentralSimpleAlgebras}。

% !TeX root = ../../Book2.tex
% 纲要标题：### 18.1 Brauer 等价
% 纲要位置：计划纲要.md，第 1628 行。

\section{Brauer 等价}
\label{b2:sec:1801}

第十六章的矩阵环 Morita 等价给出 \(D\text{-Mod}\simeq M_r(D)\text{-Mod}\)，而第十七章把每个中心单代数写成 \(M_r(D)\)。矩阵大小可以改变代数的维数，却不改变它背后的除代数。固定中心域 \(k\) 及其标量作用，Brauer 等价正是记录这种模范畴的稳定类型；张量积将在这些类上给出可以取逆的运算。本章固定底域 \(k\)，所有中心单代数有限维、非零，代数同构均保持 \(k\) 标量。

% 纲要要点（供目录核对）：
%
% * 中心单代数的 Brauer 等价；以固定底域的线性 Morita 等价解释稳定分类，正式证明非零角代数保持 Brauer 类
% * 张量积给出的群结构；核对加矩阵块的相容性，以包络同构解释反代数真正给出逆元
% * 除代数代表、指数与周期的区分；指数等于最小有限分裂次数，正式证明扩域后指数的两项整除约束
% * 利用反代数解释第一类对合使 Brauer 类的二倍为零，区分周期二与指数二，保留次数、指数及周期各自的意义
%

\subsection{除代数代表与 Brauer 等价}
\label{b2:subsec:1801:01}

若 \(A=M_p(D)\)、\(B=M_q(D)\)，希望略去原来的矩阵重数，就应允许两侧再加不同大小的矩阵块，因为
\[
M_q(A)\cong M_{pq}(D)\cong M_p(B).
\]
要求两侧都加同一个矩阵大小，会保留原来的重数差异，不能表达这种稳定关系。

\begin{definition}\label{b2:def:brauer-equivalence}
设 \(k\) 为域，\(A,B\) 为中心单 \(k\) 代数。若存在正整数 \(r,s\)，使
\[
M_r(A)\cong M_s(B)
\]
为 \(k\) 代数同构，则称 \(A,B\) \term[Brauer-dengjia]{Brauer 等价}[Brauer equivalent]。
\end{definition}

\begin{proposition}\label{b2:prop:brauer-division-representative}
设 \(k\) 为域。每个中心单 \(k\) 代数都同构于某个 \(M_r(D)\)，其中 \(r\geq1\)，\(D\) 为中心是 \(k\) 的有限维除代数。若 \(A\cong M_r(D)\)、\(B\cong M_s(E)\) 为两个中心单 \(k\) 代数，则 \(A,B\) 在 Brauer 意义下等价，当且仅当 \(D\cong E\) 为 \(k\) 代数同构。上述条件也等价于幺左模范畴 \(A\text{-Mod}\) 与 \(B\text{-Mod}\) 之间存在 \(k\) 线性等价，即该等价在每个 Hom 空间上诱导的映射均为 \(k\) 线性。因此 Brauer 等价是等价关系，每个等价类具有唯一到 \(k\) 代数同构的除代数代表。对任意非零幂等元 \(e\in A\)，角代数 \(eAe\) 以 \(e\) 为单位，也是中心单 \(k\) 代数，并与 \(A\) 在 Brauer 意义下等价；其 \(k\) 标量为 \(ce\)、\(c\in k\)。
\end{proposition}
\begin{proof}
中心单代数是有限维单 Artin 环，由\cref{b2:prop:simple-artinian}可写成 \(M_r(D)\)。矩阵环的中心恰为 \(Z(D)\) 中元素组成的标量矩阵，故 \(Z(D)=k\)。这个结构同构保持中心中的标量，故为 \(k\) 代数同构。若 \(A=M_r(D),B=M_s(D)\)，则 \(M_s(A)\cong M_{rs}(D)\cong M_r(B)\)，所以它们 Brauer 等价。

反过来，定义中的矩阵同构给出同一个单环的两种矩阵分解。\cref{b2:thm:artin-wedderburn}及其唯一性说明除环代表同构。该同构可由一个简单模的自同态反环恢复；矩阵同构保持 \(k\) 的中心作用，所以恢复的同构也保持 \(k\)。于是上述充要条件成立，自反、对称和传递均随之得到。

矩阵 Morita 等价由 \(k\) 双模的张量积构造，保持 Hom 空间的 \(k\) 标量，所以同一个 \(D\) 给出上述 \(k\) 线性范畴等价。反过来，等价保持简单对象，并给出其自同态环的 \(k\) 代数同构。两侧各只有一种简单模类型，自同态环分别为 \(D^{\mathrm{op}}\)、\(E^{\mathrm{op}}\)，故 \(D\cong E\) 也保持 \(k\) 标量。

最后，在 \(A=M_r(D)=\operatorname{End}_D(D^r)\) 中，幂等算子 \(e\) 给出右 \(D\) 空间分解 \(D^r=\operatorname{im}e\oplus\ker e\)。取适应该分解的基，将 \(e\) 写成 \(\operatorname{diag}(I_s,0)\)、\(s\ge1\)。限制到像空间给出 \(eAe\cong\operatorname{End}_D(\operatorname{im}e)\cong M_s(D)\)，保持 \(k\) 作用和单位。因此角代数中心单，且仍有除代数代表 \(D\)。
\end{proof}

因此每个 Brauer 类可以用唯一的中心除代数同构类型表示。例如，所有 \(M_n(k)\) 都与 \(k\) 在 Brauer 意义下等价。实四元数除代数 \(\mathbb H\) 却不与 \(\mathbb R\) 等价，因为这两个除代数不同构。Brauer 等价不表示原代数同构：矩阵大小仍然决定实际的维数。

\subsection{张量积与 Brauer 群结构}
\label{b2:subsec:1801:02}

\begin{theorem}\label{b2:thm:brauer-group}
设 \(k\) 为域。中心单 \(k\) 代数的 Brauer 等价类\footnote{每个有限维代数可由某个 \(k^n\) 上的有限乘法表表示，全部这些乘法表组成集合，取同构及稳定等价类后仍为集合。}组成交换群，运算、单位及逆分别为
\[
[A]+[B]=[A\otimes_k B],\qquad 0=[k],\qquad -[A]=[A^{\mathrm{op}}].
\]
这个群称为 \(k\) 的\term[Brauer-qun]{Brauer 群}[Brauer group]，记为 \(\operatorname{Br}(k)\)。
\end{theorem}
\symidx{\(\operatorname{Br}(k)\)：域 \(k\) 的 Brauer 群}
\begin{proof}
要使张量积下降到已经略去矩阵大小的类上，须检查加矩阵块与张量积相容。张量积仍中心单，由\cref{b2:thm:csa-tensor-product}已知。对矩阵单位逐项相乘，得到
\[
M_r(A)\otimes_k M_s(B)\cong M_{rs}(A\otimes_k B).
\]
其映射把 \((E_{ij}a)\otimes(E_{uv}b)\) 送到以 \((i,u),(j,v)\) 为行列指标、条目为 \(a\otimes b\) 的矩阵单位。它保持乘法并把相应的向量空间基一一对应，故为同构。这说明张量运算不依赖稳定代表元的选择。

张量积的结合、交换因子同构与单位同构给出交换群运算所需的结合律、交换律和单位。\cref{b2:thm:csa-enveloping-isomorphism}给出
\[
A\otimes_k A^{\mathrm{op}}\cong\operatorname{End}_k(A),
\]
右边为 \(M_{\dim_kA}(k)\)，代表零类，所以反代数确实通过张量积将 \(A\) 抵消为分裂代数。这正是它给出群逆元的原因。
\end{proof}

若 \(K/k\) 为域扩张，标量扩张保持中心单性，并与张量积相容，所以诱导群同态
\[
\operatorname{Br}(k)\longrightarrow\operatorname{Br}(K),\qquad
[A]\longmapsto[A\otimes_kK].
\]
其核记为 \(\operatorname{Br}(K/k)\)，称为\term[xiangdui-Brauer-qun]{相对 Brauer 群}[relative Brauer group]；其中的类恰由被 \(K\) 分裂的代数代表，因为零类等价于除代数代表为底域，亦即代数本身为全矩阵代数。
\symidx{\(\operatorname{Br}(K/k)\)：由 \(K\) 分裂的 Brauer 类}

\ProseParagraphBreak
任意有限域的 Brauer 群为零。事实上，有限维除代数在有限底域上只有有限多个元素，由\cref{b2:thm:wedderburn-little}为交换域；若中心正是底域，它只能等于底域。代数闭域的 Brauer 群也为零，因为每个中心单代数都已分裂。

\subsection{除代数代表、指数与周期}
\label{b2:subsec:1801:03}

\begin{definition}\label{b2:def:brauer-index-period}
设 \(k\) 为域，\(A\) 为中心单 \(k\) 代数，写 \(A\cong M_r(D)\)，其中 \(r\geq1\)，\(D\) 为中心除 \(k\) 代数。记 \(\dim_kD=d^2\)，定义
\[
\operatorname{ind}(A)=d
\]
并称它为 \(A\) 的\term[zhongxindan-daishu-zhishu]{指数}[index of a central simple algebra]。把 Brauer 类 \([A]\in\operatorname{Br}(k)\) 的加法阶称为 \(A\) 的\term[Brauer-zhouqi]{周期}[period]，记为 \(\operatorname{per}(A)\)；此处按群元素阶的通常约定，暂允许其取值为 \(\infty\)。\cref{b2:thm:period-divides-index}将证明周期总是有限的，且整除指数。这里 \(\deg A=\sqrt{\dim_k A}=rd\)，零 Brauer 类的周期和指数均为一。
\end{definition}
\symidx{\(\operatorname{ind}(A)\)：中心单代数的指数}
\symidx{\(\operatorname{per}(A)\)：中心单代数的周期}

指数和周期都不随矩阵大小改变，次数却随之改变。例如 \(M_3(\mathbb H)\) 的次数为六、指数为二。\eqref{b2:eq:quaternion-conjugation}中的四元数共轭给出 \(\mathbb H\cong\mathbb H^{\mathrm{op}}\)，所以其 Brauer 类的二倍为零；它又不是零类，故周期恰为二。这个例子区分了次数与另两个数，但不表示一般代数的周期都等于指数。

\ProseParagraphBreak
设 \(k\) 为任意域，\(A\) 为中心单 \(k\) 代数。若一个 \(k\) 线性映射 \(\sigma:A\to A\) 满足 \(\sigma(1)=1\)、\(\sigma(xy)=\sigma(y)\sigma(x)\) 及 \(\sigma^2=\operatorname{id}\)，则称它为\term[diyilei-duihe]{第一类对合}[involution of the first kind]。它作为 \(A\to A^{\mathrm{op}}\) 的 \(k\) 代数同构，给出 \([A]=[A^{\mathrm{op}}]=-[A]\)，故 \(2[A]=0\)。这个条件限制的是周期，不能据此断言指数至多为二；\cref{b2:exm:period-two-index-four}给出周期二、指数四的例子。

\begin{proposition}\label{b2:prop:index-splitting-degree}
设 \(k\) 为域，\(A\) 为中心单 \(k\) 代数。若有限域扩张 \(L/k\) 分裂 \(A\)，则 \(\operatorname{ind}(A)\mid[L:k]\)。而且存在次数恰为 \(\operatorname{ind}(A)\) 的有限可分分裂域。因此
\[
\operatorname{ind}(A)
=\min\{\,[L:k]\mid L/k\;\text{为分裂}\;A\;\text{的有限域扩张}\,\}.
\]
\end{proposition}
\begin{proof}
令 \(A=M_r(D)\)、\(\deg D=d\)。\(A_L=M_r(D_L)\) 分裂，当且仅当 \(D_L\) 分裂，这由单环矩阵分解的除代数唯一性得到。此时 \(D\) 通过 \(D_L\cong M_d(L)\) 作用在列空间 \(L^d\) 上。把它看成左 \(D\) 向量空间，设维数为 \(t\)。比较 \(k\) 维数得
\[
d[L:k]=t\dim_kD=td^2,
\]
所以 \(d\mid[L:k]\)。\cref{b2:thm:csa-separable-splitting}还为除代数 \(D\) 构造了一个次数为 \(d\) 的可分子域并证明它分裂 \(D\)，同一个域也分裂 \(A\)，得到存在性。
\end{proof}

这说明指数记录最小可能的有限分裂域次数；周期记录在 Brauer 群中反复张量多少次才能得到零类。一个是分裂扩张的大小，另一个是群运算的阶数，定义上不应混同。


\begin{proposition}\label{b2:prop:index-after-base-change}
设 \(k\) 为域，\(A\) 为中心单 \(k\) 代数，\(L/k\) 为域扩张，令 \(d=\operatorname{ind}(A)\)、\(e=\operatorname{ind}(A_L)\)，其中 \(A_L=A\otimes_kL\)。则 \(e\mid d\)。若 \(L/k\) 有限，则还有
\[
\dfrac de\mid[L:k].
\]
特别地，若 \([L:k]\) 与 \(d\) 互素，则 \(e=d\)。
\end{proposition}
\begin{proof}
以除代数代表 \(D\) 代替 \(A\) 不改变扩域后的 Brauer 类。写 \(D_L\cong M_s(E)\)，其中 \(E\) 为中心除 \(L\) 代数，次数为 \(e\)。比较次数得 \(d=se\)，故 \(e\mid d\)。

若 \(L/k\) 有限，通过含幺嵌入 \(D\hookrightarrow D_L\cong M_s(E)\)，列空间 \(E^s\) 成为左 \(D\) 向量空间。记其 \(D\) 维数为 \(t\)，比较底层 \(k\) 维数，得到
\[
s e^2[L:k]=t d^2=t s^2e^2.
\]
这些是正整数之间的等式，消去 \(se^2\) 得 \([L:k]=ts\)，即 \(d/e=s\) 整除扩域次数。若该次数与 \(d\) 互素，\(s\) 同时整除二者，只能等于一，所以 \(e=d\)。
\end{proof}

% !TeX root = ../../Book2.tex
% 纲要标题：### 18.2 循环代数
% 纲要位置：计划纲要.md，第 1637 行。

\section{循环代数}
\label{b2:sec:1802}

一个循环 Galois 扩张不仅能提供标量，还能提供乘法的扭转规则。把扩张的生成自同构写成一个新元素的共轭作用，再规定该元素的某次幂，就得到循环代数。它是否分裂，最后化为原扩张中的一个范数方程。

% 纲要要点（供目录核对）：
%
% * 循环 Galois 扩张及扭曲乘法；先由生成关系产生进位公式，再以两次进位恒等式证明结合律，回扣代数展示的泛性质
% * 循环代数的中心单性；分别核对中心与最少非零系数证明，分裂映射使用原代数的右子域空间
% * 范数与分裂条件；一维子域空间的半线性算子产生范数方程，换生成元直接计算参数及恢复原生成元
%

\subsection{循环 Galois 扩张及扭曲乘法}
\label{b2:subsec:1802:01}

设 \(K/k\) 为次数 \(n\ge2\) 的循环 Galois 扩张，\(\operatorname{Gal}(K/k)=\langle\sigma\rangle\)，并取 \(a\in k^\times\)。希望在包含 \(K\) 的代数中加入一个元素 \(u\)，要求 \(ub=\sigma(b)u\)（\(b\in K\)）及 \(u^n=a\)。前一关系反复使用便给出 \(u^ic=\sigma^i(c)u^i\)；若 \(i+j=qn+r\)、\(0\le r<n\)，后一关系将 \(u^{i+j}\) 化成 \(a^qu^r\)。这两个规则决定了乘积的系数和指数。据此，在左 \(K\) 向量空间
\[
C=K\oplus Ku\oplus\cdots\oplus Ku^{n-1}
\]
上，用下式定义乘法，并按 \(k\) 双线性扩张：
\begin{equation}\label{b2:eq:cyclic-algebra-product}
(bu^i)(cu^j)=b\sigma^i(c)a^{\left\lfloor(i+j)/n\right\rfloor}u^{\,i+j-n\left\lfloor(i+j)/n\right\rfloor},
\qquad 0\le i,j<n.
\end{equation}
这里 \(q=\lfloor(i+j)/n\rfloor\) 记录进位次数，每次进位产生一个 \(a\)。乘法一般不是 \(K\) 双线性的，但 \(\sigma\) 固定 \(k\)，所以仍为 \(k\) 双线性。

\begin{proposition}\label{b2:prop:cyclic-algebra-construction}
设 \(K/k\) 为次数 \(n\geq2\) 的循环 Galois 扩张，\(\operatorname{Gal}(K/k)=\langle\sigma\rangle\)，\(a\in k^\times\)。在左 \(K\) 空间 \(C=\bigoplus_{i=0}^{n-1}Ku^i\) 上，按 \(k\) 双线性延拓乘法
\[
(bu^i)(cu^j)=b\sigma^i(c)a^{\lfloor(i+j)/n\rfloor}
u^{\,i+j-n\lfloor(i+j)/n\rfloor}\qquad(b,c\in K,\ 0\leq i,j<n).
\]
则 \(C\) 是维数为 \(n^2\) 的结合含幺 \(k\) 代数。它由子域 \(K\) 和一个可逆元 \(u\) 生成，满足
\[
ub=\sigma(b)u\quad(b\in K),\qquad u^n=a.
\]
更确切地，给定含幺 \(k\) 代数 \(B\)、指定的含幺 \(k\) 代数嵌入 \(\iota:K\hookrightarrow B\)，以及 \(v\in B\) 满足
\[
v\iota(b)=\iota\bigl(\sigma(b)\bigr)v\quad(b\in K),\qquad v^n=\iota(a),
\]
则存在唯一含幺 \(k\) 代数同态 \(\Phi:C\to B\)，使 \(\Phi|_K=\iota\) 且 \(\Phi(u)=v\)。
\end{proposition}
\begin{proof}
比较三个基型元素 \(bu^i,cu^j,du^\ell\) 的乘积，令 \(q_{i,j}=\lfloor(i+j)/n\rfloor\)、\(r_{i,j}=i+j-nq_{i,j}\)。两种结合次序先分别将 \(i+j\) 或 \(j+\ell\) 化到 \(0,\ldots,n-1\)，总进位数满足
\[
q_{i,j}+q_{r_{i,j},\ell}
=\left\lfloor\frac{i+j+\ell}{n}\right\rfloor
=q_{j,\ell}+q_{i,r_{j,\ell}}.
\]
这是将 \(i+j=nq_{i,j}+r_{i,j}\) 或 \(j+\ell=nq_{j,\ell}+r_{j,\ell}\) 代入中间式所得的恒等式。两种乘积因而含有相同的 \(a\) 次幂，剩下的 \(u\) 指数也都是 \(i+j+\ell-n\lfloor(i+j+\ell)/n\rfloor\)。其余系数都为 \(b\sigma^i(c)\sigma^{i+j}(d)\)，因为 \(\sigma^n=1\)，且 \(a\in k\) 被 \(\sigma\) 固定。因此两种乘积相同，双线性再给出全部元素的结合律。\(1\in K\) 为单位，\(k\) 标量位于中心。底向量空间事先已指定为 \(n\) 个 \(K\) 的直和，所以维数确为 \(n^2\)，无需再假设这些形式单项式无关。

关系由定义直接得到，\(u^{-1}=a^{-1}u^{n-1}\)。给定 \(\iota\) 及 \(v\)，唯一可能的映射把 \(\sum_i b_i u^i\) 送到 \(\sum_i\iota(b_i)v^i\)。由 \(v^i\iota(c)=\iota(\sigma^i(c))v^i\) 与 \(v^n=\iota(a)\)，它保持\eqref{b2:eq:cyclic-algebra-product}，所以确为代数同态。
\end{proof}

该代数称为由这些数据给出的\term[xunhuan-daishu]{循环代数}[cyclic algebra]，记为 \((K/k,\sigma,a)\)。上述泛性质是\cref{b2:prop:algebra-relations-universal}中代数展示原则的具体应用：\(K\) 中的加法与乘法已由 \(\iota\) 保持，新生成元的像只需满足所列的两条扭曲关系，就能唯一确定整个代数同态。
\symidx{\((K/k,\sigma,a)\)：循环扩张及参数确定的循环代数}
当 \(n=1\) 时，另约定 \((k/k,\operatorname{id},a)=k\)；这个退化情形不另引入形式生成元。

\subsection{循环代数的中心单性}
\label{b2:subsec:1802:02}

\begin{theorem}\label{b2:thm:cyclic-algebra-central-simple}
设 \(K/k\) 为次数 \(n\geq1\) 的循环 Galois 扩张，\(\operatorname{Gal}(K/k)=\langle\sigma\rangle\)，\(a\in k^\times\)，则循环代数 \(C=(K/k,\sigma,a)\) 是次数 \(n\) 的中心单 \(k\) 代数，且由 \(K\) 分裂。
\end{theorem}
\begin{proof}
当 \(n=1\) 时 \(C=k\)，结论直接成立。以下设 \(n\ge2\)。

要确定中心，先找出哪些元素与子域 \(K\) 交换，再检查它们是否与生成元 \(u\) 交换。设 \(z=\sum_i b_i u^i\) 与 \(K\) 中所有元素交换。比较 \(zc\) 与 \(cz\)，得 \(b_i\bigl(\sigma^i(c)-c\bigr)=0\) 对所有 \(c\in K\) 成立。对 \(i\ne0\)，\(\sigma^i\ne1\)，可以选 \(c\) 使差非零，故 \(b_i=0\)。所以 \(C\) 中 \(K\) 的中心化子就是 \(K\)。若 \(z=b_0\) 再与 \(u\) 交换，则 \(\sigma(b_0)=b_0\)，故 \(z\in k\)。于是中心恰为 \(k\)。

为证单性，需要证明每个非零双边理想都含有 \(1\)。取这样的理想 \(I\)，并在其中选取非零系数个数最少的元素。右乘 \(u^r\)（\(0\le r<n\)）使指数 \(i\) 变为 \((i+r)\bmod n\)，只是循环搬移各分量，并将进位项的系数乘以非零标量 \(a\)。因此可以将任意一个非零系数搬到常数位置，且非零系数个数不变。再左乘这个常数系数在 \(K\) 中的逆，只将所有系数同时乘以非零数，便得到 \(x=1+\sum_{i>0}b_i u^i\in I\)，仍有最少的非零系数。对 \(c\in K\)，交换子 \(xc-cx\in I\) 的常数项为零，其余非零项只能出现在 \(x\) 原有的非零位置，故非零项数严格小于 \(x\)。若它非零，与极小性矛盾；故它对所有 \(c\) 都为零。由前一段的中心化子计算，\(x\in K\)，进而 \(x=1\)。因此 \(I=C\)，代数单。

构造给出 \(\dim_kC=n^2\)，所以 \([K:k]=n=\deg C\)，恰满足\cref{b2:prop:maximal-field-splits-csa}中子域次数等于代数次数的条件，故 \(K\) 分裂 \(C\)。该命题的实际映射为 \(C\otimes_kK\to\operatorname{End}_K(C)\)，\(c\otimes b\) 作用为 \(x\mapsto cxb\)；目标中的空间是原代数 \(C\)，并未换成 \(C_K\)。把它看成右 \(K\) 空间，其维数为 \(n\)，所以目标是 \(M_n(K)\)。
\end{proof}

在次数二、特征不为二时，若 \(K=k(\sqrt d)\)，则 \(u\sqrt d=-\sqrt d\,u\)，\(u^2=a\)。因而循环代数就是四元数代数 \((d,a)_k\)。这里指定自同构是重要数据；较高次数时，换一个生成元需同时跟踪相应参数，不能只把记号中的 \(\sigma\) 删去。

\subsection{范数与分裂条件}
\label{b2:subsec:1802:03}

循环扩张的范数为
\[
\operatorname{N}_{K/k}(b)=\prod_{j=0}^{n-1}\sigma^j(b).
\]
这是第一卷定理18.2.3的共轭元公式在循环 Galois 扩张上的形式。它给出判断 \(C\) 是否为矩阵代数的条件。

\begin{theorem}[循环代数的分裂判据]\label{b2:thm:cyclic-algebra-norm-splitting}
设 \(K/k\) 为有限循环 Galois 扩张，\(\operatorname{Gal}(K/k)=\langle\sigma\rangle\)，\(a\in k^\times\)，则循环代数 \((K/k,\sigma,a)\) 分裂，当且仅当 \(a\in\operatorname{N}_{K/k}(K^\times)\)。
\end{theorem}
\begin{proof}
当 \([K:k]=1\) 时循环代数为 \(k\)，而域范数为恒等映射，结论成立。以下设 \(n=[K:k]\ge2\)。

若 \(a=\operatorname{N}_{K/k}(b)\)，其中 \(b\in K^\times\)，要构造分裂表示，可以让 \(K\) 在自身上左乘，再寻找与扭曲关系相容的算子 \(T\) 来表示 \(u\)。由 \(uc=\sigma(c)u\) 所要求的相容条件是 \(T(cz)=\sigma(c)T(z)\)。\(K\) 是一维 \(K\) 空间，因此任何满足此条件的算子都由 \(T(1)\) 决定，具体为 \(T(z)=\sigma(z)T(1)\)。取 \(T(1)=b\)，就得到 \(k\) 线性算子 \(T(z)=b\sigma(z)\)，并且
\[
T^n(z)=b\sigma(b)\cdots\sigma^{n-1}(b)\sigma^n(z)
=\operatorname{N}_{K/k}(b)z=az.
\]
因此\cref{b2:prop:cyclic-algebra-construction}的泛性质给出含幺 \(k\) 代数同态 \(C\to\operatorname{End}_k(K)\)，也就是\cref{b2:prop:algebra-modules-representations}所说的 \(C\) 在这个空间上的作用。来源单而映射非零，所以单射；两端维数都是 \(n^2\)，故为同构。

反过来，设 \(C\cong\operatorname{End}_k(V)\)，其中 \(\dim_kV=n\)。这个分裂模型给出 \(C\) 的作用，限制到子域 \(K\) 使 \(V\) 成为左 \(K\) 向量空间。比较 \(\dim_kV=[K:k]\dim_KV\)，得到 \(\dim_KV=1\)。选一个 \(K\) 基向量把 \(V\) 识别为 \(K\)。此时 \(u\) 的作用 \(T\) 必须满足 \(T(cz)=\sigma(c)T(z)\)，所以仍由 \(b=T(1)\) 唯一决定，\(T(z)=b\sigma(z)\)。\(u\) 可逆使 \(b\ne0\)，而 \(u^n=a\) 给出 \(\operatorname{N}_{K/k}(b)=a\)：分裂模型中的一维 \(K\) 空间强制产生了所需的范数解。
\end{proof}

当 \(n\ge2\)、\(b\in K^\times\) 时，把生成元改为 \(v=bu\)。每次将 \(u\) 移过下一个 \(b\) 都施加一次 \(\sigma\)，所以
\[
v^n=b\sigma(b)\cdots\sigma^{n-1}(b)u^n
=\operatorname{N}_{K/k}(b)a.
\]
并且 \(vc=\sigma(c)v\) 对每个 \(c\in K\) 成立。新生成元还能恢复 \(u=b^{-1}v\)，故由构造的泛性质得到从参数为 \(\operatorname{N}_{K/k}(b)a\) 的循环代数到 \(C\) 的满同态；两侧维数都为 \(n^2\)，它是同构。因此参数乘上一个范数不改变代数的同构类型。

\ProseParagraphBreak
例如，取 \(\sigma\) 为复共轭，\(\mathbb C/\mathbb R\) 的范数为 \(z\overline z\)，非零范数恰为正实数。故 \((\mathbb C/\mathbb R,\sigma,a)\) 在 \(a>0\) 时分裂，在 \(a<0\) 时同构于 \(\mathbb H\)。前一种情况可以缩放 \(u\) 使平方为一，后一种可以缩放使平方为负一；范数判据同时解释为什么两种情况不能同构。

% !TeX root = ../../Book2.tex
% 纲要标题：### 18.3 约化迹与约化范数
% 纲要位置：计划纲要.md，第 1643 行。

\section{约化迹与约化范数}
\label{b2:sec:1803}

把中心单代数在一个分裂域中写成矩阵以后，可以取通常的迹与行列式。问题是这些数是否仍属于原来的底域，以及换一个分裂域、换一个矩阵同构后是否改变。本节先解决这两个问题，再把所得运算与域迹、域范数比较。始终令 \(\deg A=n\)。本章通过自然交换同构 \(K\otimes_kA\cong A\otimes_kK\) 采用右侧的标量扩张模型，仍将它记为 \(A_K\)。

% 纲要要点（供目录核对）：
%
% * 在分裂域中的定义；解释抵消半线性、Galois 降系数及任意分裂域比较，Cayley–Hamilton 的系数在中心
% * 与域迹、域范数的比较；同一左乘算子的两种坐标计算及首一因子重数比较，区别四元数标量的最小多项式与约化特征多项式
% * 定义对分裂域选择的独立性与标量扩张相容性；泛型坐标处理扩域后的新元素，约化迹可加、约化范数乘法而不加法，正则迹不能在坏特征除以次数
% * 正式证明一般约化迹配对对称非退化及迹映射满射；纯不可分子域的限制配对可以全退化；对合与转置复合证明保持约化特征多项式和非退化对称迹型
%

\subsection{约化迹与约化范数的定义}
\label{b2:subsec:1803:01}

由\cref{b2:thm:csa-separable-splitting}，可取一个固定的有限 Galois 分裂域 \(K_0/k\) 和同构 \(\rho:A\otimes_kK_0\to M_n(K_0)\)。对 \(a\in A\)，暂记
\[
p_a(T)=\det\bigl(TI_n-\rho(a\otimes1)\bigr).
\]
这是 \(K_0[T]\) 中的首一 \(n\) 次多项式。

\begin{theorem}\label{b2:thm:reduced-characteristic-descent}
设 \(k\) 为域，\(A\) 为次数 \(n\) 的中心单 \(k\) 代数，\(a\in A\)。对任意分裂域 \(K/k\) 及 \(K\) 代数同构 \(\rho:A\otimes_kK\rightarrow M_n(K)\)，令
\[
p_a(T)=\det\bigl(TI_n-\rho(a\otimes1)\bigr).
\]
则 \(p_a(T)\) 的系数属于 \(k\)，并且与分裂域 \(K\) 及同构 \(\rho\) 的选择无关。
\end{theorem}
\begin{proof}
对 \(g\in\operatorname{Gal}(K_0/k)\)，记 \(g_M\) 为对矩阵条目逐项施加 \(g\) 的映射。直接复合 \(g_M\circ\rho\) 得到的是 \(g\)-半线性映射，不能与 \(\rho\) 作为 \(K_0\) 代数同构比较。为抵消这一半线性，定义
\[
\rho_g=g_M\circ\rho\circ(\operatorname{id}_A\otimes g^{-1}):
A\otimes_kK_0\longrightarrow M_n(K_0).
\]
首尾两个映射分别为 \(g\)-半线性及 \(g^{-1}\)-半线性，复合后为 \(K_0\) 线性；三个映射均保持乘法且双射，所以 \(\rho_g\) 是 \(K_0\) 代数同构。特别地，对任意有限和有
\[
\rho_g\left(\sum_i a_i\otimes\lambda_i\right)
=\sum_i\lambda_i\,g_M\bigl(\rho(a_i\otimes1)\bigr).
\]
于是 \(\rho_g\rho^{-1}\) 是全矩阵代数的 \(K_0\) 自同构，由\cref{b2:thm:skolem-noether}为内自同构。记 \(B=\rho(a\otimes1)\)，则存在 \(U_g\in\operatorname{GL}_n(K_0)\) 使 \(g_M(B)=U_gBU_g^{-1}\)。令 \(g\) 固定不定元 \(T\)，便有
\[
g\bigl(p_a(T)\bigr)
=\det\bigl(TI_n-g_M(B)\bigr)
=\det\bigl(TI_n-U_gBU_g^{-1}\bigr)
=p_a(T).
\]
因此全部系数属于不动域 \(K_0^{\operatorname{Gal}(K_0/k)}=k\)。

固定 \(K_0\) 而改变 \(\rho\)，同样由内自同构的结论，矩阵只差共轭，故多项式不变。再给定任意分裂域 \(L/k\) 及其分裂同构。取 \(L\) 的代数闭包 \(\Omega=\overline L\)。由于 \(K_0/k\) 是有限代数扩张，可逐次为其生成元在 \(\Omega\) 中选择最小多项式的根，将 \(k\hookrightarrow\Omega\) 延拓为 \(K_0\hookrightarrow\Omega\)。这样就得到同时容纳 \(K_0\) 与 \(L\) 的域。将两个分裂模型都扩张到 \(\Omega\)，它们成为 \(A_\Omega\) 到 \(M_n(\Omega)\) 的两个同构，所以仍只差内自同构。因此两个特征多项式在 \(\Omega[T]\) 中相同。第一模型的系数已属于 \(k\)，于是第二模型也给出同一多项式。
\end{proof}

\begin{definition}\label{b2:def:reduced-trace-norm}
设 \(k\) 为域，\(A\) 为次数 \(n\) 的中心单 \(k\) 代数，\(a\in A\)。取 \(A\) 的任意分裂域 \(K/k\) 和 \(K\) 代数同构 \(\rho:A\otimes_kK\rightarrow M_n(K)\)。由\cref{b2:thm:reduced-characteristic-descent}，多项式 \(\det\bigl(TI_n-\rho(a\otimes1)\bigr)\in k[T]\) 与这些选择无关；称它为 \(a\) 的\term[yuehua-tezheng-duoxiangshi]{约化特征多项式}[reduced characteristic polynomial]，记为 \(\chi^{\mathrm{red}}_{A,a}(T)\)。写成
\[
\chi^{\mathrm{red}}_{A,a}(T)=T^n-c_1(a)T^{n-1}+\cdots+(-1)^n c_n(a),
\]
定义\term[yuehua-ji]{约化迹}[reduced trace]与\term[yuehua-fanshu]{约化范数}[reduced norm]为
\[
\operatorname{Trd}_{A/k}(a)=c_1(a),\qquad
\operatorname{Nrd}_{A/k}(a)=c_n(a).
\]
\end{definition}
\symidx{\(\chi^{\mathrm{red}}_{A,a}\)：约化特征多项式}
\symidx{\(\operatorname{Trd}_{A/k}\)：约化迹}
\symidx{\(\operatorname{Nrd}_{A/k}\)：约化范数}

分裂情形 \(A=M_n(k)\) 中，这就是通常的特征多项式、矩阵迹与行列式。约化特征多项式的系数属于中心 \(k\)，所以可以在非交换代数 \(A\) 中将 \(T\) 代为 \(a\)，不必为系数与 \(a\) 的乘法另作次序选择。\cref{b2:thm:exterior-cayley-hamilton}在分裂模型中给出 \(\chi^{\mathrm{red}}_{A,a}(a)=0\)：其像为零，而 \(A\to A_K\) 单射，故原代数中的值也为零。

\subsection{与域迹、域范数的比较}
\label{b2:subsec:1803:02}

令 \(\lambda_a:A\to A\) 表示左乘 \(a\)，将它看成 \(n^2\) 维 \(k\) 空间上的线性算子。其普通迹和行列式与约化运算并不相同。

\begin{proposition}\label{b2:prop:regular-reduced-characteristic}
设 \(k\) 为域，\(A\) 为次数 \(n\) 的中心单 \(k\) 代数，\(a\in A\)。令 \(\lambda_a:A\rightarrow A\) 为左乘 \(a\) 的 \(k\) 线性算子，则其特征多项式满足
\[
\chi_{\lambda_a}(T)=\bigl(\chi^{\mathrm{red}}_{A,a}(T)\bigr)^n.
\]
因此 \(\operatorname{tr}_k\lambda_a=n\operatorname{Trd}_{A/k}(a)\)，且 \(\det_k\lambda_a=\operatorname{Nrd}_{A/k}(a)^n\)。
\end{proposition}
\begin{proof}
扩张到分裂域后，左乘矩阵 \(B=\rho(a\otimes1)\) 逐列作用在 \(M_n(K)\) 上：每一列都是一份 \(K^n\)，都受到同一个 \(B\) 的作用，列与列之间互不混合。因此这个 \(n^2\) 维线性作用是 \(n\) 份 \(B\) 的直和，特征多项式为 \(\det(TI-B)^n\)。标量扩张不改变一个原系数矩阵的特征多项式，因而等式在 \(k[T]\) 中成立。比较次高次系数和常数项，就得到两个公式。
\end{proof}

若底域特征整除 \(n\)，不能把正则迹除以 \(n\) 来定义约化迹。例如在 \(M_p(k)\)、\(\operatorname{char}k=p\) 中，\(E_{11}\) 的约化迹为一，而它的左乘算子迹为零。定义约化运算所需的下降论证正好避开了这种除法。

\begin{theorem}\label{b2:thm:reduced-field-comparison}
设 \(k\) 为域，\(A\) 为次数 \(n\) 的中心单 \(k\) 代数，\(E\subseteq A\) 为含 \(k\) 的子域，\([E:k]=e\)，并按左乘将 \(A\) 视为左 \(E\) 向量空间。则 \(e\mid n\)。对 \(a\in E\)，令 \(m_a:E\rightarrow E\) 为乘 \(a\) 的 \(k\) 线性算子，则有
\[
\chi^{\mathrm{red}}_{A,a}(T)=\chi_{m_a:E\to E}(T)^{n/e},
\]
并且
\[
\operatorname{Trd}_{A/k}(a)=\frac ne\operatorname{Tr}_{E/k}(a),\qquad
\operatorname{Nrd}_{A/k}(a)=\operatorname{N}_{E/k}(a)^{n/e}.
\]
此处 \(E/k\) 不必可分。特别地，次数为 \(n\) 的子域上的约化迹、范数恰为域迹、范数。
\end{theorem}
\begin{proof}
\cref{b2:prop:csa-self-centralizing-field}已证明任意含 \(k\) 的子域 \(E\subseteq A\) 的次数 \(e\) 整除 \(n\)，这里不要求 \(E/k\) 可分。

\(A\) 作为左 \(E\) 空间的维数为 \(n^2/e\)。选一组左 \(E\) 基 \(v_1,\ldots,v_{n^2/e}\)，则左乘 \(a\in E\) 把 \(\sum_j c_jv_j\) 送到 \(\sum_j (ac_j)v_j\)。每个系数空间都是一份 \(E\)，其底层 \(k\) 作用为 \(m_a\)，故同一个算子 \(\lambda_a\) 的特征多项式为 \(\chi_{m_a:E\to E}^{n^2/e}\)。另一方面，\cref{b2:prop:regular-reduced-characteristic}通过分裂矩阵的列坐标给出它的特征多项式为 \((\chi^{\mathrm{red}}_{A,a})^n\)。比较这两种坐标计算，得到
\[
\left(\chi^{\mathrm{red}}_{A,a}\right)^n
=\left(\chi_{m_a:E\to E}^{n/e}\right)^n.
\]
两个括号中的多项式都首一。在 \(k[T]\) 的唯一分解中，每个首一不可约因子在两边的指数分别为 \(n u\) 和 \(n v\)，所以 \(u=v\)，从而两个括号相等。这是在整数中比较因子指数，并未在 \(k\) 中除以 \(n\)。域迹和域范数按乘法算子的迹与行列式定义，见第一卷定义18.2.1，故比较系数得到所述公式。
\end{proof}

非可分子域上的限制可能丢失迹配对的信息。例如取 \(k=\mathbb F_p(t)\) 和 \(E=k(\alpha)\)，其中 \(p\) 为素数，\(\alpha^p=t\)。多项式 \(T^p-t\in\mathbb F_p[t][T]\) 对素元 \(t\) 满足第一卷定理12.3.1的 Eisenstein 判据，所以它在 \(k[T]\) 中不可约，\([E:k]=p\)，且 \(1,\alpha,\ldots,\alpha^{p-1}\) 为一组基。在这组基下，正则作用 \(a\mapsto m_a\) 将 \(E\) 嵌入 \(A=\operatorname{End}_k(E)\cong M_p(k)\)。对 \(1\le i<p\)，乘 \(\alpha^i\) 将每个基向量送到另一个基向量的标量倍数，没有对角项；乘标量 \(c\) 的迹则为 \(pc=0\)。由迹的线性，\(\operatorname{Tr}_{E/k}(a)=0\) 对所有 \(a\in E\) 成立。比较定理中的 \(e=n=p\)，所以 \(\operatorname{Trd}_{A/k}\) 在 \(E\) 上恒为零。由于 \(xy\in E\) 对 \(x,y\in E\) 成立，配对 \((x,y)\mapsto\operatorname{Trd}_{A/k}(xy)\) 在 \(E\times E\) 上也恒为零；\cref{b2:prop:involution-reduced-trace-form}将证明它在全代数 \(A\) 上非退化。

\ProseParagraphBreak
在特征不为二的域 \(k\) 上，取 \(a,b\in k^\times\) 及四元数代数 \(Q=(a,b)_k\)。\cref{b2:prop:quaternion-norm}中的共轭与范数给出每个 \(x\in Q\) 的关系
\[
x^2-(x+\bar x)x+x\bar x=0.
\]
其中 \(x+\bar x\) 与 \(x\bar x\) 都属于 \(k\)。若 \(x\) 非标量，它不能满足非零的一次 \(k\) 系数关系，否则就属于 \(k\)；所以最小多项式 \(\mu_x(T)\) 恰为这个首一二次多项式。约化 Cayley--Hamilton 等式又使 \(\mu_x\) 整除首一二次多项式 \(\chi^{\mathrm{red}}_{Q,x}\)，故两者相等。若 \(x=c\in k\)，则 \(\mu_c(T)=T-c\)，而分裂矩阵为 \(cI_2\)，所以 \(\chi^{\mathrm{red}}_{Q,c}(T)=(T-c)^2\)。由这两种情形比较系数，约化迹与范数分别为 \(x+\bar x\) 与 \(x\bar x\)，标量情形即为 \(2c\) 与 \(c^2\)。

\subsection{定义的良定性与基本运算性质}
\label{b2:subsec:1803:03}

\begin{proposition}\label{b2:prop:reduced-trace-norm-properties}
设 \(k\) 为域，\(A\) 为次数 \(n\) 的中心单 \(k\) 代数，则约化迹 \(\operatorname{Trd}_{A/k}:A\rightarrow k\) 为 \(k\) 线性映射；对任意 \(a,b\in A\)，约化范数 \(\operatorname{Nrd}_{A/k}:A\rightarrow k\) 满足
\[
\operatorname{Nrd}(ab)=\operatorname{Nrd}(a)\operatorname{Nrd}(b),\qquad
\operatorname{Nrd}(1)=1.
\]
对 \(c\in k\)，有 \(\operatorname{Trd}(c1_A)=nc\)、\(\operatorname{Nrd}(c1_A)=c^n\)。元素 \(a\) 可逆，当且仅当 \(\operatorname{Nrd}(a)\ne0\)。
\end{proposition}
\begin{proof}
在一个分裂域中计算，迹的线性、行列式的乘法性与标量矩阵公式给出前面各项；值已由\cref{b2:thm:reduced-characteristic-descent}属于 \(k\)，故等式就在底域中成立。

若 \(a\) 可逆，乘法性给出 \(\operatorname{Nrd}(a)\operatorname{Nrd}(a^{-1})=1\)。反过来，范数非零使分裂矩阵可逆，故逐列作用的 \(\lambda_a\) 扩域后可逆。它的原 \(k\) 矩阵的行列式在扩域中非零，因而在 \(k\) 中也非零，这就把矩阵的可逆性降回了底域。于是存在 \(b\in A\) 使 \(ab=1\)；再由 \(a(ba-1)=0\) 及 \(\lambda_a\) 单射得 \(ba=1\)。因此 \(a\) 是单位。
\end{proof}

约化范数的乘法性不包含可加性。取 \(A=M_2(k)\)，则 \(E_{11}\) 与 \(E_{22}\) 的行列式均为零，而 \(E_{11}+E_{22}=I_2\) 的行列式为一，所以 \(\operatorname{Nrd}(E_{11}+E_{22})\ne\operatorname{Nrd}(E_{11})+\operatorname{Nrd}(E_{22})\)。

\ProseParagraphBreak
对原元素 \(a\otimes1\)，扩域后的约化系数可以在共同的分裂模型中比较。但 \(A_F\) 还包含坐标不属于 \(k\) 的新元素 \(\sum_i a_i\otimes x_i\)。要同时处理这些元素，需要证明约化系数本身是坐标的 \(k\) 系数多项式，而不是只比较原元素的取值。

\begin{proposition}\label{b2:prop:reduced-base-change}
设 \(k\) 为域，\(A\) 为次数 \(n\) 的中心单 \(k\) 代数，\(F/k\) 为域扩张，则约化特征多项式及其各系数与标量扩张 \(A_F=A\otimes_kF\) 相容。更具体地，取 \(A\) 的 \(k\) 基 \(a_1,\ldots,a_{n^2}\)，则元素 \(\sum_i x_i a_i\in A\) 的约化特征多项式的每个系数都是坐标 \(x_1,\ldots,x_{n^2}\) 的 \(k\) 系数多项式；对 \(A_F\) 中的 \(\sum_i a_i\otimes x_i\)、\(x_i\in F\)，仍由同一个多项式计算。
\end{proposition}
\begin{proof}
选有限 Galois 分裂域 \(K/k\) 及分裂同构 \(\rho\)，令
\[
H(\mathbf X)=\sum_i X_i\rho(a_i\otimes1).
\]
其特征多项式的系数属于 \(K[X_1,\ldots,X_{n^2}]\)。在\cref{b2:thm:reduced-characteristic-descent}的证明中，\(\rho_g\rho^{-1}\) 是整个矩阵代数的一个内自同构，所以比较矩阵 \(U_g\) 与输入元素无关。令 \(g\) 固定所有不定元 \(X_i\)，便有
\[
g\bigl(H(\mathbf X)\bigr)=U_gH(\mathbf X)U_g^{-1}.
\]
因此每个特征系数作为多项式都被 \(\operatorname{Gal}(K/k)\) 固定。逐项比较单项式的标量系数，得到这些系数属于 \(k\)，即特征系数多项式属于 \(k[X_1,\ldots,X_{n^2}]\)。

给定 \(F/k\)，如下降定理的证明那样，将有限代数扩张 \(K/k\) 嵌入 \(F\) 的代数闭包 \(\Omega\)。上述分裂同构扩张后，也是 \((A_F)_\Omega\cong A_\Omega\) 的分裂模型。在此模型中将 \(X_i\) 代为任意 \(x_i\in F\)，就得到 \(\sum_i a_i\otimes x_i\) 的特征多项式。所有系数都由刚得到的 \(k\) 系数多项式取值，故属于 \(F\)；分裂模型选择无关性保证它们就是 \(A_F\) 的约化系数。
\end{proof}

特别地，对原来的 \(a\in A\)，在任何扩张中都有
\[
\operatorname{Trd}_{A_F/F}(a\otimes1)=\operatorname{Trd}_{A/k}(a),\qquad
\operatorname{Nrd}_{A_F/F}(a\otimes1)=\operatorname{Nrd}_{A/k}(a).
\]
泛型矩阵的行列式每一项都恰含 \(n\) 个矩阵条目，所以约化范数的坐标多项式是 \(n\) 次齐次多项式。

\begin{proposition}\label{b2:prop:involution-reduced-trace-form}
设 \(k\) 为域，\(A\) 为次数 \(n\) 的中心单 \(k\) 代数。则
\[
\beta(x,y)\coloneqq\operatorname{Trd}_{A/k}(xy)
\]
是 \(A\) 上非退化的对称 \(k\) 双线性型，且 \(\operatorname{Trd}_{A/k}:A\to k\) 是满射。

若 \(\sigma:A\to A\) 为一个 \(k\) 线性映射，满足 \(\sigma(1)=1\)、\(\sigma^2=\operatorname{id}\)，且对所有 \(x,y\in A\) 有 \(\sigma(xy)=\sigma(y)\sigma(x)\)，则 \(\sigma(x)\) 与 \(x\) 有相同的约化特征多项式，因而约化迹与约化范数均相同。并且
\[
t_\sigma(x,y)\coloneqq\operatorname{Trd}_{A/k}\bigl(\sigma(x)y\bigr)
\]
是 \(A\) 上非退化的对称 \(k\) 双线性型，称为 \(\sigma\) 的\term[duihe-ji-xing]{迹型}[trace form of an involution]。
\end{proposition}
\begin{proof}
约化迹的线性给出 \(\beta\) 的双线性。选分裂域 \(K/k\) 及分裂同构 \(\rho\)，则 \(\beta\) 在分裂模型中由 \((X,Y)\mapsto\operatorname{tr}(XY)\) 给出。矩阵迹满足 \(\operatorname{tr}(XY)=\operatorname{tr}(YX)\)，故 \(\beta(x,y)=\beta(y,x)\)。这个矩阵配对还非退化：若 \(X_{ij}\ne0\)，取 \(Y=E_{ji}\)，便有 \(\operatorname{tr}(XY)=X_{ij}\ne0\)。

取 \(A\) 的一组 \(k\) 基 \(a_1,\ldots,a_{n^2}\)，其像 \(\rho(a_i\otimes1)\) 为 \(M_n(K)\) 的 \(K\) 基。原配对的 Gram 矩阵为 \(G=(\operatorname{Trd}(a_i a_j))_{i,j}\in M_{n^2}(k)\)，它在 \(K\) 中就是上述矩阵配对的 Gram 矩阵，因此 \(\det G\ne0\)。这个行列式原属 \(k\)，故 \(\beta\) 在 \(k\) 上也非退化。由于 \(1_A\ne0\)，存在 \(y\in A\) 使 \(\beta(1_A,y)=\operatorname{Trd}(y)\ne0\)。约化迹是到一维空间 \(k\) 的非零线性映射，因而满射。这里不需要 \(\operatorname{Trd}(1_A)=n\) 非零。

若存在所述 \(\sigma\)，将 \(\sigma\otimes\operatorname{id}_K\) 通过 \(\rho\) 转为 \(M_n(K)\) 上的反自同构 \(\tau\)。转置 \(t:X\mapsto X^{\mathsf T}\) 也反转乘法次序，所以 \(\tau\circ t\) 是 \(K\) 代数自同构。由\cref{b2:thm:skolem-noether}，存在 \(P\in\operatorname{GL}_n(K)\) 使 \(\tau(X)=PX^{\mathsf T}P^{-1}\)。因此 \(\sigma(x)\) 的分裂矩阵与 \(x\) 的分裂矩阵之转置相似，特征多项式不变。\cref{b2:thm:reduced-characteristic-descent}把这个等式降回 \(k\)，比较系数就得到约化迹和约化范数不变。

迹型的双线性来自 \(\sigma\) 和约化迹的线性。再把刚证的迹不变性用于 \(\sigma(x)y\)，得到
\[
\operatorname{Trd}\bigl(\sigma(x)y\bigr)
=\operatorname{Trd}\Bigl(\sigma\bigl(\sigma(x)y\bigr)\Bigr)
=\operatorname{Trd}\bigl(\sigma(y)x\bigr),
\]
所以它对称。又因为 \(t_\sigma(x,y)=\beta(\sigma(x),y)\)，而 \(\sigma^2=\operatorname{id}\) 保证 \(\sigma\) 是可逆线性映射，将 \(\beta\) 的第一变量与它复合仍然非退化。
\end{proof}

对实矩阵代数的转置对合，\(t_\sigma(X,X)=\operatorname{tr}(X^{\mathsf T}X)\) 为所有矩阵条目的平方和；对四元数的标准共轭 \(\gamma\)，次数为二的标量迹公式则给出
\[
t_\gamma(x,x)=\operatorname{Trd}(\bar x x)=2N(x).
\]
这将乘法范数与一个双线性型联系起来。迹型依赖所选对合，不只是底层代数的属性；\cref{b2:ex:18-two-involution-trace-forms}会在同一个矩阵代数上比较两种不同的情形。

% !TeX root = ../../Book2.tex
% 纲要标题：### 18.4 交叉积与相对 Brauer 群
% 纲要位置：计划纲要.md，第 1652 行。

\section{交叉积与相对 Brauer 群}
\label{b2:sec:1804}

循环代数只需一个生成自同构和一个参数。一般有限 Galois 扩张有多个自同构；记录它们对应元素相乘时的标量误差，就得到\(2\)-余圈。本节用显式方程定义这些数据，证明它们与相对 Brauer 群的对应；第四卷再从一般上链复形研究群上同调。

% 纲要要点（供目录核对）：
%
% * 交叉积与2-余圈；结合律产生余圈方程，重选基元产生余边；正式证明因子集不依赖标量提升及矩阵模型
% * 有限 Galois 扩张的相对 Brauer 群及循环情形的范数商；对角幂等元对齐两份子域作用，取角证明群运算，同类同次数及 Skolem–Noether 证明单射，明确下降因子集的逆类约定
% * 与第四卷群上同调的衔接；循环余圈化为进位形式，正文证明交叉积约化迹只由单位指标分量贡献
% * Brauer 群不能与第五卷 Brauer 诱导定理混同
%

\subsection{\texorpdfstring{交叉积与 \(2\)-余圈}{交叉积与 2-余圈}}
\label{b2:subsec:1804:01}

固定有限 Galois 扩张 \(K/k\)，记 \(G=\operatorname{Gal}(K/k)\)、\(n=|G|=[K:k]\)。希望为每个 \(g\in G\) 配一个基元 \(u_g\)，使它对标量的作用由 \(g\) 给出，而两个基元相乘只与 \(u_{gh}\) 相差一个非零标量，即
\[
u_g a=g(a)u_g,\qquad u_g u_h=c(g,h)u_{gh}.
\]
这两个规则把结合律转成一个可计算的条件：
\[
(u_g u_h)u_t=c(g,h)c(gh,t)u_{ght},\qquad
u_g(u_h u_t)=g\bigl(c(h,t)\bigr)c(g,ht)u_{ght}.
\]
若再要求 \(u_1=1\)，就得到下面的余圈方程及归一化条件。

\begin{definition}\label{b2:def:normalized-two-cocycle}
设 \(K/k\) 为有限 Galois 扩张，\(G=\operatorname{Gal}(K/k)\)，\(G\) 按域自同构作用于 \(K^\times\)。若函数 \(c:G\times G\to K^\times\) 对任意 \(g,h,t\in G\) 满足
\begin{equation}\label{b2:eq:two-cocycle}
c(g,h)c(gh,t)=g\bigl(c(h,t)\bigr)c(g,ht),\qquad
c(1,g)=c(g,1)=1,
\end{equation}
则称 \(c\) 为\term[guiyihua-er-yuquan]{归一化 \(2\)-余圈}[normalized 2-cocycle]。若取元素族 \((b_g)_{g\in G}\subseteq K^\times\) 且 \(b_1=1\)，定义函数
\[
(\delta b)(g,h)=b_g\cdot g(b_h)b_{gh}^{-1}
\]
并称 \(\delta b\) 为\term[er-yubian]{\(2\)-余边}[2-coboundary]。
\end{definition}

直接将 \(\delta b\) 代入\eqref{b2:eq:two-cocycle}，两边都化为 \(b_g\cdot g(b_h)gh(b_t)b_{ght}^{-1}\)，所以余边确为余圈。

\ProseParagraphBreak
余边记录的是同一组基元的缩放。若将 \(u_g\) 换成 \(u'_g=b_g u_g\)，其中 \(b_1=1\)，则
\[
u'_g u'_h
=b_g g(b_h)c(g,h)b_{gh}^{-1}u'_{gh}
=\bigl(c\delta b\bigr)(g,h)u'_{gh}.
\]
因此改变这些基元会改变乘法系数，但只使余圈乘上一个余边。

\begin{definition}\label{b2:def:galois-crossed-product}
设 \(K/k\) 为有限 Galois 扩张，\(G=\operatorname{Gal}(K/k)\)，\(c:G\times G\rightarrow K^\times\) 为归一化 \(2\)-余圈。在左 \(K\) 空间 \(B_c=\bigoplus_{g\in G}Ku_g\) 上规定
\begin{equation}\label{b2:eq:crossed-product-multiplication}
(a u_g)(b u_h)=a\cdot g(b)c(g,h)u_{gh}.
\end{equation}
按双可加性延拓后，所得含幺 \(k\) 代数称为由 \(K/k\) 和 \(c\) 确定的\term[jiaocha-ji-daishu]{交叉积代数}[crossed product algebra]。
\end{definition}
\symidx{\(B_c=(K/k,G,c)\)：Galois 扩张的交叉积代数}

其中 \(K\) 通过 \(a\mapsto au_1\) 嵌入；乘法公式同时规定了 \(u_g b=g(b)u_g\) 与 \(u_g u_h=c(g,h)u_{gh}\)。前一条控制基元经过标量时发生的自同构，后一条控制基元之间的乘法误差。前述换基计算也给出一个同构 \(B_{c\delta b}\to B_c\)：把新代数中指标 \(g\) 的基元送到 \(b_g u_g\) 即可。

\begin{proposition}\label{b2:prop:crossed-product-central-simple}
设 \(K/k\) 为有限 Galois 扩张，\(G=\operatorname{Gal}(K/k)\)，\(n=[K:k]\)，\(c:G\times G\rightarrow K^\times\) 为归一化 \(2\)-余圈，则相应交叉积 \(B_c=\bigoplus_{g\in G}Ku_g\) 是次数 \(n\) 的中心单 \(k\) 代数，单位为 \(u_1\)。子域 \(K\) 分裂它；每个 \(u_g\) 都可逆。
\end{proposition}
\begin{proof}
比较三项乘积 \((a u_g)(b u_h)(d u_t)\)：除共同因子 \(a\cdot g(b)gh(d)\) 外，两种结合方式的系数分别是\eqref{b2:eq:two-cocycle}的两边。因此乘法结合。归一化条件使 \(u_1\) 为单位，底域标量在中心。向量空间维数为 \(n[K:k]=n^2\)。由余圈条件还有 \(c(g,g^{-1})=g\bigl(c(g^{-1},g)\bigr)\)，故
\[
u_g^{-1}=c(g^{-1},g)^{-1}u_{g^{-1}}.
\]
左右相乘均为一。

中心的计算与循环代数同样，先检查与 \(K\) 交换的条件，再检查与全部基元交换的条件。若 \(x=\sum_g a_g u_g\) 与全部 \(b\in K\) 交换，则比较系数得 \(a_g\bigl(g(b)-b\bigr)=0\)。对 \(g\ne1\) 可选使差非零的 \(b\)，所以只有单位指标的系数可能非零。再与各 \(u_g\) 交换，就使该系数被全部 \(G\) 固定，因而属于 \(k\)。故中心是 \(k\)。

单性的证明沿用循环代数中减少非零项数的方法。对一个非零双边理想，选非零项数最少的元素，并取其中一个非零项的指标 \(g\)。右乘 \(u_g^{-1}\) 使每个指标 \(h\) 变为 \(hg^{-1}\)，各系数只乘上非零标量，所以非零项数不变，而单位指标的系数变为非零。再左乘一个 \(K^\times\) 中的标量，使该系数为一。所得元素与任意 \(b\in K\) 的交换子在单位指标处为零，其他非零项只可能来自原有的指标；若交换子非零，就有更少的非零项，违反最少性。于是该元素与 \(K\) 交换，由刚才的系数比较只能为一，理想即为全代数。这证明单性。最后 \([K:k]=n=\deg B_c\)，由\cref{b2:prop:maximal-field-splits-csa}，\(K\) 分裂 \(B_c\)。
\end{proof}

若交叉积分裂为 \(n\) 阶矩阵代数，它在一个 \(n\) 维 \(k\) 空间上的作用限制到 \(K\)，就成为一维 \(K\) 空间上的作用。在这个空间上，每个 \(u_g\) 的作用满足 \(T_g(az)=g(a)T_g(z)\)；选择一个 \(K\) 基以后，它只能是域自同构 \(g\) 再乘上一个非零标量。这些标量正是余边公式中的 \(b_g\)。

\begin{proposition}\label{b2:prop:crossed-product-split-coboundary}
设 \(K/k\) 为有限 Galois 扩张，\(G=\operatorname{Gal}(K/k)\)，\(c:G\times G\rightarrow K^\times\) 为归一化 \(2\)-余圈，\(B_c\) 为相应交叉积代数，则 \(B_c\) 分裂，当且仅当 \(c\) 是 \(2\)-余边。
\end{proposition}
\begin{proof}
若 \(c(g,h)=b_g\cdot g(b_h)b_{gh}^{-1}\)，其中 \(b_1=1\)，在 \(k\) 空间 \(K\) 上让 \(K\) 左乘，让 \(u_g\) 按 \(T_g(z)=b_g\cdot g(z)\) 作用。计算得到
\[
T_gT_h(z)=b_g g(b_h)gh(z)=c(g,h)T_{gh}(z),\qquad
T_g(az)=g(a)T_g(z).
\]
因此这些算子满足交叉积的两个乘法规则，且 \(T_1=\operatorname{id}_K\)，给出含幺同态 \(B_c\to\operatorname{End}_k(K)\)。由单性及两端同为 \(n^2\) 维，它是同构，其中 \(n=[K:k]\)。

若 \(B_c\cong\operatorname{End}_k(W)\)，由 \(\deg B_c=n\) 有 \(\dim_k W=n\)。限制到 \(K\) 后，维数公式给出 \(\dim_KW=1\)。选一个 \(K\) 基，将 \(W\) 识别为 \(K\)。由于 \(u_g z=g(z)u_g\)，其作用满足 \(T_g(z)=g(z)T_g(1)=b_g g(z)\)，其中 \(b_g=T_g(1)\ne0\)，且 \(b_1=1\)。代入 \(u_g u_h=c(g,h)u_{gh}\)，得 \(b_g\cdot g(b_h)=c(g,h)b_{gh}\)，正是余边条件。此时把 \(u_g\) 换成 \(b_g^{-1}u_g\)，它们在 \(K\) 上就按 \(g\) 作用，乘法误差也同时消失。
\end{proof}

\subsection{Brauer 群与群上同调}
\label{b2:subsec:1804:02}

\(2\)-余圈按逐点相乘组成交换群，余边构成其子群。定义商群
\[
H^2(G,K^\times)=\{\text{归一化 }2\text{-余圈}\}/\{2\text{-余边}\}.
\]
这个商群是当前 \(G\) 作用下二阶\term[qun-shangtongdiao]{群上同调}[group cohomology]的显式模型。第四卷第 12.1 节从 bar 消解建立余链复形，第 12.2 节再用低阶余圈解释群扩张；这里把交换群 \(K^\times\) 的运算写成乘法，得到同一套二阶方程，下面的论证只用这些显式方程。
\symidx{\(H^2(G,K^\times)\)：乘法系数的二阶群上同调}

\ProseParagraphBreak
对于 \(g\in G\)，可加映射 \(T_g\) 若满足 \(T_g(av)=g(a)T_g(v)\)，称其为 \(g\)-\term[banxianxing-yingshe]{半线性映射}[semilinear map]。在矩阵代数上，逐条目施加 \(g\) 就给出这种映射。下面把一个已分裂代数的下降数据写成具体矩阵。

\begin{proposition}\label{b2:prop:split-algebra-factor-set}
设 \(K/k\) 为有限 Galois 扩张，\(G=\operatorname{Gal}(K/k)\)，\(A\) 为由 \(K\) 分裂的中心单 \(k\) 代数，\(m=\deg A\)，取 \(K\) 代数同构 \(\rho:A_K=A\otimes_kK\rightarrow M_m(K)\)。令 \(\alpha_g=\rho(\operatorname{id}_A\otimes g)\rho^{-1}\) 为 \(A_K\) 的自然 Galois 作用在矩阵代数上的表示，\(g(X)\) 表示对矩阵 \(X\) 逐条目施加 \(g\)，则存在 \(C_g\in\operatorname{GL}_m(K)\)、\(C_1=I_m\)，使 \(\alpha_g(X)=C_g\cdot g(X)C_g^{-1}\)。由
\[
C_g\cdot g(C_h)=c(g,h)C_{gh}
\]
定义的标量 \(c(g,h)\in K^\times\) 构成归一化 \(2\)-余圈，而且 \([A]=-[B_c]\)，其中 \(B_c\) 是相应交叉积代数。若将 \(C_g\) 换为 \(b_gC_g\)，其中 \(b_g\in K^\times\)、\(b_1=1\)，则余圈换为 \(c\delta b\)。若将分裂模型换为 \(\rho'=\operatorname{Ad}(P)\rho\)，其中 \(P\in\operatorname{GL}_m(K)\)、\(\operatorname{Ad}(P)(X)=PXP^{-1}\)，并取
\[
C'_g=PC_g g(P)^{-1},
\]
则得到原余圈 \(c\)。特别地，所得类 \([c]\in H^2(G,K^\times)\) 与 \(C_g\) 及分裂同构 \(\rho\) 的选择无关。
\end{proposition}
\begin{proof}
\(\alpha_g\) 在标量上作用为 \(g\)，即为 \(g\)-半线性代数自同构。将它与逐条目作用 \(g^{-1}\) 依次复合，所得 \(\alpha_g\circ g^{-1}\) 是一个 \(K\) 代数自同构；由\cref{b2:thm:skolem-noether}，存在 \(C_g\in\operatorname{GL}_m(K)\)，使
\begin{equation}\label{b2:eq:brauer-semilinear-matrices}
\alpha_g(X)=C_g\cdot g(X)C_g^{-1},\qquad C_1=I_m.
\end{equation}
\(\alpha_g\alpha_h=\alpha_{gh}\) 说明 \(C_g\cdot g(C_h)C_{gh}^{-1}\) 与所有矩阵交换，所以为唯一的标量 \(c(g,h)\in K^\times\)。比较同一个矩阵积的两种结合方式，得到
\[
\begin{aligned}
\bigl(C_g g(C_h)\bigr)gh(C_t)
&=c(g,h)c(gh,t)C_{ght},\\
C_g g\bigl(C_h h(C_t)\bigr)
&=g\bigl(c(h,t)\bigr)c(g,ht)C_{ght}.
\end{aligned}
\]
消去可逆矩阵 \(C_{ght}\)，即为\eqref{b2:eq:two-cocycle}；\(C_1=I_m\) 又给出归一化条件。

对同一个 \(\alpha_g\)，两个实现矩阵只相差非零标量，因为它们的商与全部矩阵交换。因此任意另一组归一化的实现矩阵都形如 \(b_gC_g\)、\(b_1=1\)，而
\[
(b_gC_g)g(b_hC_h)
=b_g g(b_h)c(g,h)C_{gh}
=\bigl(c\delta b\bigr)(g,h)(b_{gh}C_{gh}).
\]
再改变分裂模型。由 \(\rho'=\operatorname{Ad}(P)\rho\) 得
\[
\alpha'_g(X)
=P C_g g(P)^{-1}g(X)g(P)C_g^{-1}P^{-1}.
\]
所述 \(C'_g\) 因而实现 \(\alpha'_g\)，且 \(C'_1=I_m\)。它们满足
\[
C'_g g(C'_h)
=PC_g g(C_h)gh(P)^{-1}
=c(g,h)C'_{gh},
\]
所以余圈不变。任意两个 \(K\) 代数分裂同构之差都是 \(M_m(K)\) 的 \(K\) 自同构，由 Skolem--Noether 为这种内自同构；这就涵盖了全部分裂模型的选择。

在 \(k\) 空间 \(V=K^m\) 上，让 \(K\) 按标量乘法作用，让 \(u_g\) 按 \(v\mapsto C_g\cdot g(v)\) 作用，得到 \(B_c\) 的表示。它含幺且 \(B_c\) 单，因而忠实。与 \(K\) 作用交换的算子恰好是 \(K\) 线性算子，可写为 \(X\in M_m(K)\)；这样的 \(X\) 与所有 \(u_g\) 作用交换，当且仅当
\[
XC_g=C_g g(X),\qquad\text{即 }X=\alpha_g(X)\quad(g\in G).
\]
这些不动矩阵恰为 \(\rho(A\otimes1)\)：在 \(A_K\) 中按 \(A\) 的一组 \(k\) 基展开，自然 Galois 作用的不动性恰好使全部系数属于 \(K^G=k\)。因此 \(B_c\) 在 \(\operatorname{End}_k(V)\) 中的中心化子同构于 \(A\)。由\cref{b2:cor:csa-centralizer-factorization}的中心为 \(k\) 的情形，
\[
B_c\otimes_k A\cong\operatorname{End}_k(V).
\]
所以 \([A]=-[B_c]\)，得到所述关系。
\end{proof}

\begin{theorem}\label{b2:thm:relative-brauer-cocycles}
设 \(K/k\) 为有限 Galois 扩张，\(G=\operatorname{Gal}(K/k)\)。记 \(H^2(G,K^\times)\) 为归一化 \(2\)-余圈按\(2\)-余边取的商群，\(\operatorname{Br}(K/k)=\ker\bigl(\operatorname{Br}(k)\rightarrow\operatorname{Br}(K)\bigr)\)，则交叉积给出群同构
\[
H^2(G,K^\times)\longrightarrow\operatorname{Br}(K/k),
\qquad[c]\longmapsto[B_c].
\]
\end{theorem}
\begin{proof}
首先，\cref{b2:prop:crossed-product-central-simple}使像落在相对 Brauer 群；换基公式保证它不依赖余圈代表。

为核对乘法，取两个余圈 \(c,d\)，并记 \(n=[K:k]\)。\(B_c\otimes_kB_d\) 的 \(k\) 维数为 \(n^4\)，而 \(B_{cd}\) 的维数为 \(n^2\)；当 \(n>1\) 时两者尺寸不同，一般不能期待它们直接同构。要比较的是 Brauer 类，需要去掉不影响除代数代表的矩阵大小；下面通过非零幂等元取角来实现这一点。有限 Galois 扩张具有代数同构
\begin{equation}\label{b2:eq:galois-tensor-diagonal}
K\otimes_kK\longrightarrow\prod_{g\in G}K,
\qquad a\otimes b\longmapsto\bigl(a\cdot g(b)\bigr)_{g\in G}.
\end{equation}
具体地，第一卷定理16.4.2给出本原元素 \(K=k(\theta)\)；其可分最小多项式在 \(K\) 中分成互异因子 \(T-g(\theta)\)。对 \(K[T]/(f)\) 应用第一卷定理10.4.3的中国剩余定理，便得到上述同构。

选择单位指标，是因为该坐标把 \(a\otimes b\) 送到 \(ab\)：两份 \(K\) 在这一分量上的作用被对齐。记 \(e_t\in K\otimes_kK\) 为在指标 \(t\) 处坐标为一、其他坐标为零的幂等元，并取 \(e=e_1\)，就有 \(e(a\otimes1)e=e(1\otimes a)e\)。因此希望取角以后只保留对两份 \(K\) 同步施加同一自同构的项。在 \(E=B_c\otimes_kB_d\) 中，用 \(u_g,v_h\) 表示两侧的基元，记 \(z_{g,h}=u_g\otimes v_h\)。共轭 \(z_{g,h}\) 对 \(K\otimes_kK\) 作用为 \(g\otimes h\)。若 \(\lambda_t(a\otimes b)=a\,t(b)\)，则
\[
\lambda_t\bigl((g\otimes h)(a\otimes b)\bigr)
=g(a)\,th(b)
=g\bigl(\lambda_{g^{-1}th}(a\otimes b)\bigr).
\]
按可加性，这个等式对 \(K\otimes_kK\) 的全部元素成立。将它用于 \(e\)，可知 \(z_{g,h}e z_{g,h}^{-1}=e_{gh^{-1}}\)。不同坐标幂等元的乘积为零，因而
\[
e z_{g,h}e=e e_{gh^{-1}}z_{g,h}=0\quad\text{若 }g\ne h.
\]
当 \(g=h\) 时，\(z_{g,g}\) 与 \(e\) 交换，令 \(w_g=e z_{g,g}e=e z_{g,g}\)。取角后，系数环 \(e(K\otimes_kK)e\) 经 \(\lambda_1\) 识别为 \(K\)。原代数是自由左 \(K\otimes_kK\) 模，以全部 \(z_{g,h}\) 为基；取角消去了 \(g\ne h\) 的项，留下
\[
eEe=\bigoplus_{g\in G}e(K\otimes_kK)w_g
\cong\bigoplus_{g\in G}K w_g
\]
的左 \(K\) 模分解。这里各项仍然独立，因为它们在原自由模中有不同的指标；而 \(w_g\ne0\)，否则右乘 \(z_{g,g}^{-1}\) 就会使 \(e=0\)。计算乘法时，\(g\otimes g\) 在单位坐标上正是 \(g\)，并且
\[
w_g w_h=e\bigl(c(g,h)\otimes d(g,h)\bigr)z_{gh,gh}
=c(g,h)d(g,h)w_{gh}.
\]
同样，若用 \(e(b\otimes1)e\) 表示角代数中的标量 \(b\in K\)，便有
\[
w_g b=e(g(b)\otimes1)z_{g,g}=g(b)w_g.
\]
于是把 \(B_{cd}\) 的基元送到 \(w_g\)，把标量送到 \(e(K\otimes_kK)e\) 中的对应元素，得到含幺 \(k\) 代数同构 \(B_{cd}\to eEe\)，角代数的单位是 \(e=w_1\)。由\cref{b2:prop:brauer-division-representative}，中心单代数的非零幂等角与原代数有同一 Brauer 类。因此 \([B_c]+[B_d]=[B_{cd}]\)，映射保持群运算。\(c=1\) 时的交叉积由\cref{b2:prop:crossed-product-split-coboundary}分裂，确实对应零元。

若 \([B_c]=[B_d]\)，由\cref{b2:prop:brauer-division-representative}，可写 \(B_c\cong M_r(D)\)、\(B_d\cong M_s(D)\)，其中 \(D\) 为同一个中心除代数。两代数的次数同为 \(n\)，所以 \(r\deg D=s\deg D=n\)，从而 \(r=s\)，可取 \(k\) 代数同构 \(\phi:B_c\to B_d\)。记 \(\iota_c,\iota_d\) 为两代数中 \(K\) 的标准嵌入。由\cref{b2:thm:skolem-noether}，有 \(z\in B_d^\times\)，使
\[
\iota_d(a)=z\phi\bigl(\iota_c(a)\bigr)z^{-1}\qquad(a\in K).
\]
于是 \(\psi=\operatorname{Ad}(z)\phi\) 在两份标准 \(K\) 之间为恒等。将 \(x=\psi(u_g)\) 展开为 \(\sum_h b_{g,h}v_h\)，由 \(xa=g(a)x\) 得
\[
b_{g,h}h(a)=g(a)b_{g,h}\qquad(a\in K).
\]
若 \(b_{g,h}\ne0\)，可消去这个 \(K\) 中的系数，得到 \(h(a)=g(a)\) 对全部 \(a\in K\) 成立，故 \(h=g\)。因此 \(\psi(u_g)=b_gv_g\)，且因 \(u_g\) 可逆，有 \(b_g\in K^\times\)；单位条件给出 \(b_1=1\)。比较 \(\psi(u_g u_h)\) 与 \(\psi(u_g)\psi(u_h)\) 的系数，得到
\[
c(g,h)b_{gh}=b_g\cdot g(b_h)d(g,h).
\]
因此 \(c/d\) 为余边，证明单射性。

最后，对任意 \([A]\in\operatorname{Br}(K/k)\)，\(A_K\) 的 Brauer 类为零，由\cref{b2:prop:brauer-division-representative}在底域 \(K\) 上的结论，它同构于一个 \(K\) 上的全矩阵代数。因此可应用\cref{b2:prop:split-algebra-factor-set}，得到余圈 \(c\)，使 \([A]=-[B_c]\)。由已经证明的乘法相容性，\([B_{c^{-1}}]=-[B_c]\)，所以原类是 \([c^{-1}]\) 的像。这证明满射性，所需群同构成立。
\end{proof}

半线性映射作为 \(k\) 线性算子参与交叉积的表示，却同时改变 \(K\) 标量。乘法公式中的 \(g\bigl(c(h,t)\bigr)\) 正记录了这种改变。从余圈构造交叉积，与从分裂代数的半线性作用提取因子集，这两种构造的关系为
\[
[c]\longmapsto[B_c],\qquad
A\longmapsto c_A,\qquad
[A]=-[B_{c_A}]=[B_{c_A^{-1}}].
\]
前一个箭头是\cref{b2:thm:relative-brauer-cocycles}的同构；从一个已分裂的代数 \(A\) 出发按\cref{b2:prop:split-algebra-factor-set}构造因子集时，原类对应的是 \([c_A^{-1}]\)。

\begin{corollary}[循环情形的范数商]\label{b2:cor:cyclic-relative-brauer}
设 \(K/k\) 为次数 \(n\ge2\) 的循环 Galois 扩张，\(G=\operatorname{Gal}(K/k)=\langle\sigma\rangle\)，则有群同构
\[
H^2(G,K^\times)\cong\operatorname{Br}(K/k)
\cong k^\times/\operatorname{N}_{K/k}(K^\times).
\]
具体地，映射 \(a\operatorname{N}_{K/k}(K^\times)\mapsto[(K/k,\sigma,a)]\) 给出范数商到相对 Brauer 群的同构。
\end{corollary}
\begin{proof}
由\cref{b2:thm:relative-brauer-cocycles}，每个相对 Brauer 类都由某个交叉积 \(B_c\) 表示。在其中取 \(u=u_\sigma\)。反复使用乘法规则，可知 \(u^i\) 是 \(u_{\sigma^i}\) 的非零 \(K\) 倍数，特别地 \(u^n=a\in K^\times\)。因为 \(u\) 与其幂 \(a\) 交换，而 \(ua=\sigma(a)u\)，可消去 \(u\) 得 \(\sigma(a)=a\)，所以 \(a\in k^\times\)。将原基换成 \(1,u,\ldots,u^{n-1}\)，就是一次前述缩放基元的操作；它不改变余圈的上同调类，却将乘法误差化为
\[
u^i u^j=
\begin{cases}
u^{i+j},&i+j<n,\\
a\,u^{i+j-n},&i+j\ge n,
\end{cases}
\qquad(0\le i,j<n).
\]
因此相应余圈为
\[
c_a(\sigma^i,\sigma^j)=a^{\lfloor(i+j)/n\rfloor},
\]
且 \(B_c\cong(K/k,\sigma,a)\)。这证明参数给出从 \(k^\times\) 到 \(\operatorname{Br}(K/k)\) 的满射。进位余圈满足 \(c_a c_b=c_{ab}\)，由\cref{b2:thm:relative-brauer-cocycles}，该参数映射是群同态。

若再用 \(u'=bu\)、\(b\in K^\times\)，则
\[
(bu)^n=b\,\sigma(b)\cdots\sigma^{n-1}(b)\,u^n
=\operatorname{N}_{K/k}(b)a.
\]
因而乘以一个范数正对应再次改变这组基元。由\cref{b2:thm:cyclic-algebra-norm-splitting}，参数映射的核恰为 \(\operatorname{N}_{K/k}(K^\times)\)。商群同构定理遂给出所述范数商，连同\cref{b2:thm:relative-brauer-cocycles}便得到全部同构。
\end{proof}

交叉积的基元除了记录乘法，也使约化迹可以直接计算。左乘 \(u_h\) 将指标 \(g\) 移到 \(hg\)，因此只有单位指标的项会贡献矩阵的对角元；计算时仍须把左 \(K\) 系数转换成右 \(K\) 坐标。

\begin{proposition}\label{b2:prop:crossed-product-reduced-trace}
设 \(K/k\) 为有限 Galois 扩张，\(G=\operatorname{Gal}(K/k)\)，\(c:G\times G\to K^\times\) 为归一化 \(2\)-余圈，\(B_c=\bigoplus_{g\in G}Ku_g\) 为相应交叉积代数。对任意 \(x=\sum_{g\in G}a_g u_g\)、\(a_g\in K\)，有
\[
\operatorname{Trd}_{B_c/k}(x)=\operatorname{Tr}_{K/k}(a_1)
=\sum_{g\in G}g(a_1).
\]
\end{proposition}
\begin{proof}
记 \(n=[K:k]\)，将 \(B_c\) 的右 \(K\) 空间记为 \(W\)。由 \(a u_g=u_g g^{-1}(a)\)，\((u_g)_{g\in G}\) 也是 \(W\) 的一组基。根据\cref{b2:prop:maximal-field-splits-csa}中的分裂同构，
\[
B_c\otimes_k K\longrightarrow
\operatorname{End}_K(W),\qquad
x\otimes b\longmapsto(y\mapsto xyb),
\]
其中目标作用在原代数的右 \(K\) 空间上，而来源是标量扩张代数。由于 \(\dim_KW=n\)，右端是 \(M_n(K)\) 的一个模型。在这个模型中，\(x\otimes1\) 是左乘 \(x\) 的算子。对 \(h,g\in G\)，
\[
(a_hu_h)u_g
=a_h c(h,g)u_{hg}
=u_{hg}(hg)^{-1}\bigl(a_h c(h,g)\bigr).
\]
因此在列 \(g\) 中，这一项的矩阵系数只出现在行 \(hg\)。若 \(h\ne1\)，行指标不等于列指标，故它不贡献对角元；若 \(h=1\)，由 \(c(1,g)=1\)，对角元恰为 \(g^{-1}(a_1)\)。左乘 \(x\) 的矩阵迹于是为
\[
\sum_{g\in G}g^{-1}(a_1)=\sum_{g\in G}g(a_1)
=\operatorname{Tr}_{K/k}(a_1).
\]
有限 Galois 扩张的域迹等于全部共轭之和，而\cref{b2:def:reduced-trace-norm}中的约化迹就是任意分裂模型里的矩阵迹，故得所述公式。
\end{proof}

第五卷第 5.3 节的 Brauer 诱导定理研究有限群特征标，与本章按张量积运算的 Brauer 群是不同对象。

% !TeX root = ../../Book2.tex
% 纲要标题：### 18.5 周期与指数
% 纲要要点（供目录核对）：
%
% * 周期整除指数，指数整除次数
% * 十六维中心除代数的周期二、指数四例子
%
% <!-- 短节：不设小节 -->
%
% ---

\section{周期与指数}
\label{b2:sec:1805}

指数是除代数代表的次数，也等于有限分裂域扩张次数的最小值；周期则是 Brauer 类的加法阶。前节的半线性矩阵把两种量放在同一计算中：取矩阵行列式，会使因子集的一个幂成为余边。

\begin{theorem}\label{b2:thm:period-divides-index}
设 \(k\) 为域，\(A\) 为中心单 \(k\) 代数，则其 Brauer 类 \([A]\in\operatorname{Br}(k)\) 具有有限阶，且
\[
\operatorname{per}(A)\mid\operatorname{ind}(A)\mid\deg A.
\]
\end{theorem}
\begin{proof}
取同一 Brauer 类中的除代数 \(D\)，令 \(d=\deg D=\operatorname{ind}(A)\)，再取一个有限 Galois 分裂域 \(K/k\)。由\cref{b2:prop:split-algebra-factor-set}，对 \(D_K\cong M_d(K)\) 选矩阵 \(C_g\)，得到
\[
C_g\cdot g(C_h)=c(g,h)C_{gh},\qquad [D]=-[B_c].
\]
取矩阵行列式，令 \(b_g=\det C_g\)，则
\[
c(g,h)^d=b_g\cdot g(b_h)b_{gh}^{-1}.
\]
因此 \(c^d\) 为余边，在 \(H^2(G,K^\times)\) 中为零类。\cref{b2:thm:relative-brauer-cocycles}给出 \(d[B_c]=0\)，所以 \(d[A]=d[D]=0\)。元素的最小正阶于是存在，且整除任何零化它的正整数，特别整除 \(d\)。若 \(A\cong M_r(D)\)，则 \(\deg A=rd\)，还得到第二个整除关系。证明中没有除以 \(d\)，故包括正特征。
\end{proof}

\begin{example}[周期为二、指数为四]\label{b2:exm:period-two-index-four}
在域 \(k=\mathbb R((x))((y))\) 上，代数
\[
D=\mathbb H\otimes_{\mathbb R}k\ \otimes_k\ (x,y)_k
\]
为除代数，且 \(\operatorname{ind}(D)=4\)、\(\operatorname{per}(D)=2\)。这里 \((x,y)_k\) 为四元数代数。
\end{example}
\begin{proof}
给出一个直接的除环模型。先取中心不定元 \(u\)，组成 Laurent 级数除环 \(F=\mathbb H((u))\)。非零级数可以提出其首项，再将剩余的 \(1+h\)、\(h\) 只有正次数，按 \(1-h+h^2-\cdots\) 求逆；每个固定次数的系数只涉及有限多项，所以这个形式逆有意义。

令 \(\tau\) 固定 \(\mathbb H\)，将 \(u\) 送到 \(-u\)。在所有下端有界的形式级数 \(\sum_m a_m v^m\)、\(a_m\in F\) 上规定
\[
(a v^r)(b v^s)=a\tau^r(b)v^{r+s}.
\]
固定次数的系数仍只涉及有限和；\(\tau\) 是自同构，所以三项乘积的两种括号给出相同系数。这便定义结合环 \(E=F((v;\tau))\)。非零首项 \(a_m v^m\) 可逆，其逆为 \(\tau^{-m}(a_m^{-1})v^{-m}\)；提出首项后，再用同样的正次数几何级数构造逆。因此 \(E\) 是除环。

在 \(E\) 中令 \(x=u^2,y=v^2\)。二者都在中心，且 \(F\) 中按 \(u\) 的奇偶次数及四元数基逐项分组，得到左 \(\mathbb R((x))\) 空间的直和
\[
F=\bigoplus_{\epsilon=0}^{1}
\ \bigoplus_{h\in\{1,\mathbf i,\mathbf j,\mathbf k\}}
\mathbb R((x))\,h u^\epsilon.
\]
由于 \(\tau^2=\operatorname{id}\)，\(y=v^2\) 与 \(F\) 交换。再按 \(v\) 的奇偶次数分组，就有
\[
E=F((y))\oplus F((y))v
=\bigoplus_{\delta=0}^{1}
\ \bigoplus_{\epsilon=0}^{1}
\ \bigoplus_{h\in\{1,\mathbf i,\mathbf j,\mathbf k\}}
k\,h u^\epsilon v^\delta,
\qquad k=\mathbb R((x))((y)).
\]
两次形式 Laurent 展开的系数唯一，且 \(1,\mathbf i,\mathbf j,\mathbf k\) 在 \(\mathbb H\) 中线性无关，因此所列十六个元素确为 \(E\) 的 \(k\) 基。

中心也可逐项确定。若 \(z=\sum_m a_m v^m\) 与全部四元数常数交换，因 \(v\) 与这些常数交换，Laurent 展开的唯一性迫使每个 \(a_m\) 与 \(\mathbb H\) 交换。再逐项展开 \(a_m\) 的 \(u\) 系数，得到 \(a_m\in\mathbb R((u))\)。若 \(z\) 还与 \(u\) 交换，由 \(v^m u=(-1)^muv^m\) 及特征为零，所有奇数 \(m\) 的系数都为零。最后与 \(v\) 交换，使每个余下的 \(a_{2r}\) 被 \(\tau\) 固定；其 \(u\) 展开只能含偶数次幂，故 \(a_{2r}\in\mathbb R((x))\)。因此 \(Z(E)=\mathbb R((x))((y))=k\)。

四元数常数与 \(u,v\) 交换，而 \(u^2=x,v^2=y,vu=-uv\)。这些关系及所列基给出 \(E\cong\mathbb H_k\otimes_k(x,y)_k=D\)。因此 \(D\) 自身为十六维中心除代数，指数为四。由\eqref{b2:eq:quaternion-conjugation}，每个四元数代数都通过共轭同构于自身的反代数，其 Brauer 类的二倍为零，所以 \(2[D]=0\)。\(D\) 是不等于 \(k\) 的除代数，Brauer 类非零，故周期恰为二。
\end{proof}


\begin{chaptersummary}
Brauer 类由唯一的中心除代数代表确定，张量积给出加法，反代数给出负元。次数是代数在分裂域上成为全矩阵代数后的矩阵大小；指数记录除代数代表的次数，也是所有有限分裂域扩张次数中的最小值。周期是类的群阶，整除指数，但两者不必相等。

\ProseParagraphBreak
循环代数把自同构与参数转成扭曲乘法，参数是范数恰好等价于分裂。约化迹和约化范数在任意分裂模型中都给出相同的底域元素，与标量扩张相容；它们仍是具体代数中的元素数据，稳定加矩阵块会使 \(a\otimes I_r\) 的约化迹成为 \(r\operatorname{Trd}(a)\)，范数成为 \(\operatorname{Nrd}(a)^r\)。正则作用的迹和行列式还带有代数次数的倍数或幂次，不能不加区别地替代约化运算。第一类对合保持约化特征多项式，并与约化迹合成一个非退化对称型；这个迹型还会反映所选对合的信息。

\ProseParagraphBreak
交叉积的结合律正对应\(2\)-余圈方程，改变基元对应乘一个余边。对有限 Galois 扩张 \(K/k\)，相对 Brauer 群 \(\operatorname{Br}(K/k)\) 由这些余圈类描述；半线性矩阵的行列式进一步给出周期整除指数。周期二、指数四的例子说明，即使张量平方已经分裂，一个二次扩张仍未必足以分裂原代数。
\end{chaptersummary}

\BookExercises

\begin{exercise}\label{b2:ex:18-explicit-stabilization}
令 \(A=M_2(\mathbb H)\)、\(B=M_3(\mathbb H)\)。求使 \(M_r(A)\cong M_s(B)\) 的全部正整数对 \((r,s)\)，并分别列出 \(A,B\) 的次数、指数、周期。它们是否同构？
\end{exercise}
\begin{solution}
两边分别为 \(M_{2r}(\mathbb H)\)、\(M_{3s}(\mathbb H)\)。矩阵块大小相等时同构，反向维数相等迫使 \(2r=3s\)。所以全部解为 \((r,s)=(3t,2t)\)、\(t\ge1\)。两代数的次数分别为四、六，指数均为二，周期也均为二，因为除代数代表都是 \(\mathbb H\)，而四元数共轭表明其非零 Brauer 类的二倍为零。原代数实维数分别为十六、三十六，所以并不同构。
\end{solution}

\begin{exercise}\label{b2:ex:18-tensor-index-period-bounds}
设 \(k\) 为域，\(A,B\) 为中心单 \(k\) 代数。证明
\[
\operatorname{ind}(A\otimes_kB)\mid\operatorname{ind}(A)\operatorname{ind}(B),
\]
\[
\operatorname{per}(A\otimes_kB)\mid
\operatorname{lcm}\bigl(\operatorname{per}(A),\operatorname{per}(B)\bigr).
\]
用一个例子说明两项整除都可以严格。
\end{exercise}
\begin{solution}
取除代数代表 \(D,E\)，次数分别为 \(d,e\)。\(A\otimes B\) 与 \(D\otimes E\) 同类，后者次数为 \(de\)，其指数整除次数，得到第一式。记两个周期为 \(m,n\)，则 \(\operatorname{lcm}(m,n)\) 在 Brauer 群中同时零化 \([A]\) 和 \([B]\)，也零化它们的和，所以该和的周期整除此数。

取 \(A=B=\mathbb H\)。由四元数共轭及\cref{b2:thm:csa-enveloping-isomorphism}，\(\mathbb H\otimes_{\mathbb R}\mathbb H\cong M_4(\mathbb R)\)，其指数、周期均为一，而两式右侧分别为四和二。
\end{solution}

\begin{exercise}\label{b2:ex:18-index-after-extension}
设 \(k\) 为域，\(A\) 为中心单 \(k\) 代数，\(\operatorname{ind}(A)=12\)。若 \(L/k\) 为次数十八的有限扩张，列出\cref{b2:prop:index-after-base-change}所允许的 \(\operatorname{ind}(A_L)\) 值；这些约束是否已经保证每个值都能实现？若另一个有限扩张 \(F/k\) 使指数降为三，其次数须满足什么整除条件？再说明次数五的扩张能否改变原指数。
\end{exercise}
\begin{solution}
记 \(e=\operatorname{ind}(A_L)\)、\(s=12/e\)。正文命题要求 \(s\) 同时整除十二与十八，所以 \(s\in\{1,2,3,6\}\)，相应 \(e\in\{12,6,4,2\}\)。这是必要条件给出的候选值；它没有构造具有指定指数的扩域，不能单凭这些整数关系断言各候选值都能实现。

若指数降为三，则 \(12/3=4\) 必整除 \([F:k]\)，因此例如次数六的扩张不能使指数降为三。次数五与十二互素，正文命题使扩域后的指数仍为十二。
\end{solution}

\begin{exercise}\label{b2:ex:18-cyclic-splitting-matrices}
设 \(C=(K/k,\sigma,a)\)，\([K:k]=n\ge2\)。在原代数 \(C\) 的右 \(K\) 基 \(1,u,\ldots,u^{n-1}\) 中写出左乘 \(b\in K\) 与左乘 \(u\) 的矩阵。由此求 \(u\) 的约化特征多项式、约化迹和约化范数。
\end{exercise}
\begin{solution}
由 \(bu^j=u^j\sigma^{-j}(b)\)，左乘 \(b\) 的矩阵为
\[
D_b=\operatorname{diag}\bigl(b,\sigma^{-1}(b),\ldots,\sigma^{-(n-1)}(b)\bigr).
\]
左乘 \(u\) 把前面的基向量依次移到后一个，把最后一个送到 \(a\)，所以矩阵为 \(T^n-a\) 的友矩阵：其下次对角元为一，右上角为 \(a\)，其余为零。这些矩阵正来自\cref{b2:thm:cyclic-algebra-central-simple}证明中的分裂模型。

由\cref{b2:lem:companion-characteristic}，约化特征多项式是 \(T^n-a\)。因此 \(\operatorname{Trd}(u)=0\)、\(\operatorname{Nrd}(u)=(-1)^{n-1}a\)。这个符号来自特征多项式常数项的约定；尤其次数二时，约化范数并不等于参数 \(a\) 本身。
\end{solution}

\begin{exercise}\label{b2:ex:18-f9-cyclic-matrices}
取 \(K=\mathbb F_9=\mathbb F_3[i]\)、\(i^2=2\)，\(\sigma(z)=z^3\)。构造 \((K/\mathbb F_3,\sigma,2)\cong M_2(\mathbb F_3)\) 的显式同构，写出 \(i,u\) 的矩阵，并核对 \(u\) 的约化迹、范数。
\end{exercise}
\begin{solution}
\(\sigma(i)=2i=-i\)，而 \(b=1+i\) 满足
\[
\operatorname{N}_{K/\mathbb F_3}(b)=(1+i)(1-i)=1-2=2.
\]
在基 \(1,i\) 下，让 \(i\) 左乘，让 \(u\) 按 \(z\mapsto b\cdot\sigma(z)\) 作用，得到
\[
I=\begin{pmatrix}0&2\\1&0\end{pmatrix},\qquad
U=\begin{pmatrix}1&1\\1&2\end{pmatrix}.
\]
直接相乘有 \(I^2=U^2=2I_2\)、\(UI=-IU\)，所以它们满足循环代数的全部关系。所得含幺同态由来源单性为单射，两边都是四维，故为同构。\(\operatorname{tr}U=0\)、\(\det U=1=-2\)，与上一题的迹范数公式相符。
\end{solution}

\begin{exercise}\label{b2:ex:18-cyclic-generator-opposite}
在 \(C=(K/k,\sigma,a)\) 中，若整数 \(r\) 与 \(n=[K:k]\) 互素，证明改用 \(v=u^r\) 给出
\[
C\cong(K/k,\sigma^r,a^r).
\]
再证明 \(C^{\mathrm{op}}\cong(K/k,\sigma^{-1},a)\cong(K/k,\sigma,a^{-1})\)。
\end{exercise}
\begin{solution}
\(u\) 可逆，所以允许负整数幂。\(v\) 满足 \(vb=\sigma^r(b)v\)、\(v^n=a^r\)。取整数 \(s,t\) 使 \(rs+nt=1\)，则 \(u=v^s a^t\)，故 \(K,v\) 生成整个 \(C\)。两边维数相同，关系给出的满同态为同构。

在反代数中，使用原来的 \(u\) 有 \(u\mathbin{\cdot_{\mathrm{op}}}b=bu=u\sigma^{-1}(b)=\sigma^{-1}(b)\mathbin{\cdot_{\mathrm{op}}}u\)，且其 \(n\) 次幂仍为 \(a\)，得到第一个同构。再在该循环代数中把生成元取逆，按前一部分的 \(r=-1\) 得到第二个同构。
\end{solution}

\begin{exercise}\label{b2:ex:18-rational-norm-brauer-classes}
设 \(K=\mathbb Q(i)\)、\(\sigma(i)=-i\)。证明 \(C=(K/\mathbb Q,\sigma,3)\) 是除代数，确定其指数与周期，并证明参数改为 \(27\) 后代数同构。再把它与 \(H=(-1,-1)_{\mathbb Q}\) 比较：两者指数、周期相同，是否 Brauer 等价？
\end{exercise}
\begin{solution}
若三是范数，就有有理数 \(x,y\) 满足 \(x^2+y^2=3\)。清分母并约去公因子，得到互素整数 \(r,s,t\)、\(t\ne0\)，满足 \(r^2+s^2=3t^2\)。模三的平方只有零和一，所以 \(3\mid r,s\)。代回后又得 \(3\mid t\)，矛盾。因此由\cref{b2:thm:cyclic-algebra-norm-splitting}，\(C\) 不分裂；它是次数二的四元数代数，故由\cref{b2:prop:quaternion-division-or-split}为除代数，指数为二。共轭给出与反代数的同构，非零 Brauer 类的周期因而恰为二。

把生成元 \(u\) 换成 \(3u\)，其平方为 \(27\)，且与 \(K\) 的交换关系不变，所以得到所述代数同构。

\(H\) 的范数是四个有理数平方之和，非零向量的范数不为零，所以它也为除代数，指数、周期均为二。但是扩张到 \(\mathbb R\) 后，\(C\) 因正参数三是复数范数而分裂，\(H\) 则成为不分裂的 \(\mathbb H\)。若原来的 Brauer 类相同，扩域后的类也应相同，矛盾。因此相同的指数与周期不足以确定 Brauer 类。
\end{solution}

\begin{exercise}\label{b2:ex:18-reduced-trace-pairing}
设 \(k\) 为特征 \(p>0\) 的域，\(A=M_p(k)\)，\(\beta(X,Y)=\operatorname{tr}(XY)\)。求矩阵单位基的 \(\beta\) 对偶基，并求迹零子空间 \(H=\{X\in A:\operatorname{tr}X=0\}\) 上限制型的根空间。解释原配对非退化、\(\operatorname{tr}(I_p)=0\) 与限制型退化为何可以同时成立。
\end{exercise}
\begin{solution}
矩阵单位满足 \(\beta(E_{ij},E_{ab})=\delta_{ja}\delta_{ib}\)，所以 \(E_{ij}\) 的对偶向量为 \(E_{ji}\)，原配对非退化。

若 \(X\in A\) 与全部 \(H\) 配对为零，先取 \(E_{ij}\in H\)、\(i\ne j\)，得到 \(X_{ji}=0\)，故 \(X\) 对角。再取 \(E_{ii}-E_{11}\in H\)，得到 \(X_{ii}=X_{11}\)。因此 \(H^\perp=kI_p\)；反向包含由 \(\beta(cI_p,Y)=c\operatorname{tr}Y\) 成立。

特征 \(p\) 使 \(\operatorname{tr}(I_p)=p=0\)，所以 \(kI_p\subseteq H\)。限制型的根为 \(H\cap H^\perp=kI_p\)，维数一。原迹映射仍满射，因为 \(\operatorname{tr}(E_{11})=1\)；零迹单位元只是与这个迹零子空间正交，不是与全部 \(A\) 正交。
\end{solution}

\begin{exercise}\label{b2:ex:18-quaternion-block-reduced-polynomial}
令 \(A=M_2(\mathbb H)\)、\(a=\operatorname{diag}(\mathbf i,1+\mathbf j)\)。求 \(a\) 的约化特征多项式、约化迹与范数，并求左乘 \(a\) 在底层实空间上的迹和行列式。
\end{exercise}
\begin{solution}
扩张到 \(\mathbb C\) 后，两对角块各变为二阶矩阵。\(\mathbf i\) 的约化特征多项式为 \(T^2+1\)；\(1+\mathbf j\) 的约化迹为二、范数为二，故其多项式为 \(T^2-2T+2\)。因此
\[
\chi^{\mathrm{red}}_{A,a}(T)
=(T^2+1)(T^2-2T+2)
=T^4-2T^3+3T^2-2T+2.
\]
得到 \(\operatorname{Trd}(a)=2\)、\(\operatorname{Nrd}(a)=2\)。\(\deg A=4\)，底层实维数为十六，由\cref{b2:prop:regular-reduced-characteristic}，左乘算子的迹为八、行列式为 \(2^4=16\)。
\end{solution}

\begin{exercise}\label{b2:ex:18-reduced-tensor-formulas}
设 \(k\) 为域，\(A,B\) 为中心单 \(k\) 代数，\(\deg A=m\)、\(\deg B=n\)。对 \(a\in A,b\in B\) 证明
\[
\operatorname{Trd}_{A\otimes B}(a\otimes b)
=\operatorname{Trd}_A(a)\operatorname{Trd}_B(b),
\]
\[
\operatorname{Nrd}_{A\otimes B}(a\otimes b)
=\operatorname{Nrd}_A(a)^n\operatorname{Nrd}_B(b)^m.
\]
并说明稳定扩张 \(a\mapsto a\otimes I_r\) 如何改变这两个值。
\end{exercise}
\begin{solution}
在共同的代数闭分裂域中，两边成为矩阵的 Kronecker 积\footnote{克罗内克（Leopold Kronecker，1823-1891），德国数学家，研究数论、代数与矩阵。}。其对角项逐对相乘，迹为两个迹之积。分别把两个矩阵化为上三角形；可以使用\cref{b2:thm:jordan-canonical-form}，不要求可对角化。若对角元分别为 \(\lambda_i,\mu_j\)，张量矩阵的对角元为全部 \(\lambda_i\mu_j\)，其积为 \((\prod_i\lambda_i)^n(\prod_j\mu_j)^m\)，得到范数公式。由\cref{b2:prop:reduced-base-change}，这些等式下降到 \(k\)。

取 \(B=M_r(k),b=I_r\)，得到新约化迹为 \(r\operatorname{Trd}_A(a)\)，新约化范数为 \(\operatorname{Nrd}_A(a)^r\)。因此 Brauer 类忽略矩阵大小，但具体元素的约化运算仍记录代数次数。
\end{solution}

\begin{exercise}\label{b2:ex:18-opposite-reduced-polynomial}
设 \(k\) 为域，\(A\) 为中心单 \(k\) 代数，将 \(a\in A\) 对应到反代数中的元素 \(a^{\mathrm{op}}\)。证明其约化特征多项式、迹和范数均不变，并解释这与 \([A^{\mathrm{op}}]=-[A]\) 不矛盾。
\end{exercise}
\begin{solution}
若 \(\rho:A_L\to M_n(L)\) 是分裂同构，则 \(a^{\mathrm{op}}\mapsto\rho(a)^{\mathsf T}\) 是反代数的分裂同构，因为转置也反转乘法。矩阵与其转置有相同的特征多项式，所以约化特征多项式及两个指定系数不变。Brauer 群的负元通过整个代数的张量积刻画，并不是把每个元素的约化迹或范数取相反数；这两种运算的对象不同。
\end{solution}

\begin{exercise}\label{b2:ex:18-inseparable-field-trace}
设 \(k\) 为特征 \(p>0\) 的域，\(E=k(\alpha)\)、\(\alpha^p\in k\)、\([E:k]=p\)。通过正则作用 \(a\mapsto m_a\) 将 \(E\) 嵌入 \(A=\operatorname{End}_k(E)\)。在基 \(1,\alpha,\ldots,\alpha^{p-1}\) 下，令 \(T_j\) 将 \(\alpha^j\) 送到一、其余基向量送到零，\(0\le j<p\)。计算 \(\operatorname{Trd}(m_aT_j)\)，并求 \(E\) 对约化迹配对的正交空间维数。比较 \(p=2\) 与 \(p>2\) 时 \(E\subseteq E^\perp\) 是否相等。
\end{exercise}
\begin{solution}
写 \(a=\sum_{r=0}^{p-1}c_r\alpha^r\)。算子 \(m_aT_j\) 只有第 \(j\) 列可能非零，该列正是 \(a\) 的坐标列，所以其迹等于 \(c_j\)。\(A\) 已为分裂矩阵代数，约化迹就是这个矩阵迹，故
\[
\operatorname{Trd}(m_aT_j)=c_j.
\]
这些算子分别读取 \(E\) 的全部坐标，说明 \(A\to E^*\)、\(T\mapsto(a\mapsto\operatorname{Trd}(m_aT))\) 满射。其核为 \(E^\perp\)，因此 \(\dim_kE^\perp=p^2-p\)。

正文例子的幂基计算在这里仍适用：乘 \(\alpha^i\)、\(1\le i<p\)，将每个基向量的指数移到不同的模 \(p\) 余数类，没有对角项；乘标量 \(c\in k\) 的迹则为 \(pc=0\)。因此 \(E\) 上全部乘法算子的迹均为零，\(E\) 上的限制配对也恒为零，故 \(E\subseteq E^\perp\)。当 \(p=2\) 时两者维数都是二，所以相等；当 \(p>2\) 时 \(p^2-p>p\)，所以包含严格。尽管 \(E\) 内部无法用迹检测元素，\(T_j\) 这些外部算子却能恢复每一个坐标。
\end{solution}

\begin{exercise}\label{b2:ex:18-norm-nonlinearity-nilpotents}
设 \(k\) 为域，在 \(M_2(k)\) 中令 \(X_t=\begin{psmallmatrix}t&1\\0&0\end{psmallmatrix}\)、\(t\in k\)。求其约化迹、范数，计算 \(\operatorname{Nrd}(X_t+E_{22})-\operatorname{Nrd}(X_t)-\operatorname{Nrd}(E_{22})\)，并判断哪些 \(X_t\) 幂零。再证明任意中心单代数中的幂零元素约化迹、范数均为零，并给出两者同时为零却不幂零的矩阵。
\end{exercise}
\begin{solution}
矩阵代数的约化运算为通常运算，所以 \(\operatorname{Trd}(X_t)=t\)、\(\operatorname{Nrd}(X_t)=0\)。加上 \(E_{22}\) 后行列式为 \(t\)，而 \(\operatorname{Nrd}(E_{22})=0\)，所求差正是 \(t\)。例如 \(t=1\) 就显示范数不保持加法。

\(X_t^2=tX_t\)，归纳得 \(X_t^r=t^{r-1}X_t\)、\(r\ge1\)。因此 \(t=0\) 时平方为零，\(t\ne0\) 时每个正次幂均非零，恰好只有 \(X_0\) 幂零。

若中心单代数中 \(a^r=0\)，扩张到代数闭分裂域后所得矩阵的全部特征值都为零，所以约化特征多项式为 \(T^n\)，约化迹与范数均为零。反向失败：任意域上的 \(\operatorname{diag}(1,-1,0)\) 迹与行列式均为零，但每个正次幂仍有首个对角元一，故不幂零。
\end{solution}

\begin{exercise}\label{b2:ex:18-factor-set-changes}
设 \(K/k\) 为次数 \(n\ge2\) 的循环 Galois 扩张，生成元为 \(\sigma\)，\(a\in k^\times\)，\(A=(K/k,\sigma,a)\)。沿用\cref{b2:ex:18-cyclic-splitting-matrices}的右 \(K\) 基分裂模型，记左乘 \(b\in K\) 的矩阵为 \(D_b\)，左乘 \(u\) 的矩阵为 \(U\)。证明自然 Galois 作用可取实现矩阵 \(C_{\sigma^i}=U^{-i}\)、\(0\le i<n\)，计算其因子集，并说明它如何检验\cref{b2:prop:split-algebra-factor-set}中的负号。
\end{exercise}
\begin{solution}
这里 \(D_b\) 的第 \(j\) 个对角元为 \(\sigma^{-j}(b)\)、\(0\le j<n\)，而 \(U\) 将基向量循环移到下一位，越过末位时乘以 \(a\)。所以
\[
U D_bU^{-1}=D_{\sigma(b)},\qquad
\sigma(D_b)=D_{\sigma(b)},\qquad \sigma(U)=U.
\]
因此 \(X\mapsto U^{-1}\sigma(X)U\) 固定原代数的全部 \(D_b,U\)，并在标量上施加 \(\sigma\)，正是自然半线性作用。其第 \(i\) 次迭代由 \(U^{-i}\) 实现。

所有 \(C_{\sigma^i}\) 的条目都在 \(k\) 中。若 \(i+j=qn+r\)、\(0\le r<n\)，则 \(U^n=aI_n\) 给出
\[
C_{\sigma^i}\sigma^i(C_{\sigma^j})
=U^{-(i+j)}=a^{-q}U^{-r}.
\]
故因子集为 \(c(\sigma^i,\sigma^j)=a^{-\lfloor(i+j)/n\rfloor}\)，它是参数 \(a^{-1}\) 的进位余圈，而不是参数 \(a\) 的进位余圈。相应交叉积为 \((K/k,\sigma,a^{-1})\)，所以 \([A]=-[B_c]\) 恰好把该逆参数还原为原来的 Brauer 类。
\end{solution}

\begin{exercise}\label{b2:ex:18-cyclic-cohomology}
设 \(K/k\) 为次数 \(n=rm\) 的循环 Galois 扩张，\(r,m\ge2\)，生成元为 \(\sigma\)。令 \(H=\langle\sigma^r\rangle\)、\(F=K^H\)。将参数 \(a\in k^\times\) 的进位余圈限制到 \(H\)，求所得参数。证明自然包含 \(k^\times\hookrightarrow F^\times\) 诱导范数商群之间的同态，并说明它对应余圈的限制。
\end{exercise}
\begin{solution}
\(H\) 的阶为 \(m\)，且 \(K/F\) 的生成自同构为 \(\sigma^r\)。对 \(0\le i,j<m\)，限制后的余圈满足
\[
c_a(\sigma^{ri},\sigma^{rj})
=a^{\lfloor r(i+j)/(rm)\rfloor}
=a^{\lfloor(i+j)/m\rfloor},
\]
仍是参数 \(a\) 的进位余圈，只是循环长度由 \(n\) 变成 \(m\)。

若 \(a\) 在原范数陪集中乘上 \(N_{K/k}(b)\)，令 \(z=\prod_{j=0}^{r-1}\sigma^j(b)\)，则
\[
N_{K/F}(z)
=\prod_{i=0}^{m-1}\prod_{j=0}^{r-1}\sigma^{ri+j}(b)
=N_{K/k}(b).
\]
因此原范数子群包含于 \(N_{K/F}(K^\times)\)，得到良定义群同态
\[
k^\times/N_{K/k}(K^\times)
\longrightarrow F^\times/N_{K/F}(K^\times),
\qquad [a]\longmapsto[a].
\]
余圈和余边限制到子群后仍满足同一方程；所算的进位代表说明，在正文的循环范数商识别下，这个同态正对应限制 \(H^2(\langle\sigma\rangle,K^\times)\to H^2(H,K^\times)\)。
\end{solution}

\begin{exercise}\label{b2:ex:18-real-brauer-group}
利用\cref{b2:cor:cyclic-relative-brauer}及复数范数，计算 \(\operatorname{Br}(\mathbb R)\)，并据此确定全部有限维实中心除代数的同构类型。不使用实除代数分类定理。
\end{exercise}
\begin{solution}
任意实中心单代数扩域到 \(\mathbb C\) 后仍中心单，而代数闭域上的中心单代数为全矩阵代数，见\cref{b2:thm:csa-base-change}及\cref{b2:prop:csa-division-description}。因此所有实 Brauer 类都由 \(\mathbb C\) 分裂，即 \(\operatorname{Br}(\mathbb R)=\operatorname{Br}(\mathbb C/\mathbb R)\)。由\cref{b2:thm:relative-brauer-cocycles}和\cref{b2:cor:cyclic-relative-brauer}，
\[
\operatorname{Br}(\mathbb R)
\cong\mathbb R^\times/\operatorname{N}_{\mathbb C/\mathbb R}(\mathbb C^\times)
=\mathbb R^\times/\mathbb R_{>0}
\cong\mathbb Z/2\mathbb Z.
\]
正参数给零类，负参数给 \(\mathbb H\) 的类。由\cref{b2:prop:brauer-division-representative}的除代数代表唯一性，两个类的中心除代数代表分别为 \(\mathbb R\) 和 \(\mathbb H\)。因此全部实中心单代数是某个 \(M_r(\mathbb R)\) 或 \(M_r(\mathbb H)\)，结论来自 Brauer 对应及已知分裂性。
\end{solution}

\begin{exercise}\label{b2:ex:18-quaternion-factor-set}
在 \(\mathbb H\otimes_{\mathbb R}\mathbb C\cong M_2(\mathbb C)\) 中取
\[
\mathbf i\longmapsto I=\begin{pmatrix}i&0\\0&-i\end{pmatrix},\qquad
\mathbf j\longmapsto J=\begin{pmatrix}0&-1\\1&0\end{pmatrix}.
\]
令 \(s\) 为复共轭。证明自然半线性作用为 \(\alpha_s(X)=J\overline XJ^{-1}\)，并计算因子集的 \(c(s,s)\)。
\end{exercise}
\begin{solution}
逐条目共轭将 \(I\) 送到 \(-I\)，将 \(J\) 固定；而 \(JI=-IJ\)，故再与 \(J\) 共轭以后，\(I,J\) 都被固定。该半线性作用在标量上是复共轭，因而正是固定原四元数代数的自然 Galois 作用。

取 \(C_s=J,C_1=I_2\)，则 \(C_s\overline{C_s}=J^2=-I_2\)，所以 \(c(s,s)=-1\)。相应交叉积是参数负一的循环代数，即 \(\mathbb H\)。在这个二阶 Brauer 类上，\([\mathbb H]=-[\mathbb H]\)，因此一般公式中的负号在该例上不改变结果。
\end{solution}

\begin{exercise}\label{b2:ex:18-crossed-product-reduced-trace}
设 \(K/k\) 为有限 Galois 扩张，\(G=\operatorname{Gal}(K/k)\)，\(c\) 为归一化 \(2\)-余圈，\(B_c=(K/k,G,c)\)。计算约化迹配对在齐次分量 \(Ku_g\times Ku_h\) 上的值。给定 \(K/k\) 的基 \((a_i)\) 及其域迹对偶基 \((a_i^*)\)，为 \(B_c\) 的基 \((a_i u_g)_{i,g}\) 写出约化迹对偶基。
\end{exercise}
\begin{solution}
由交叉积乘法与\cref{b2:prop:crossed-product-reduced-trace}，
\[
\operatorname{Trd}\bigl((a u_g)(b u_h)\bigr)=
\begin{cases}
0,&gh\ne1,\\
\operatorname{Tr}_{K/k}\bigl(a\,g(b)c(g,g^{-1})\bigr),&h=g^{-1}.
\end{cases}
\]
因此分量 \(Ku_g\) 只与逆指标分量 \(Ku_{g^{-1}}\) 配对。将
\[
w_{i,g}=g^{-1}\bigl(a_i^*c(g,g^{-1})^{-1}\bigr)u_{g^{-1}}
\]
代入，得到 \(\operatorname{Trd}((a_j u_h)w_{i,g})=\delta_{h,g}\operatorname{Tr}_{K/k}(a_j a_i^*)=\delta_{h,g}\delta_{j,i}\)。故 \((w_{i,g})\) 是所求对偶基。域迹对偶基的存在来自 \(K/k\) 可分时域迹配对非退化，见第一卷定理18.3.3及推论18.3.4。
\end{solution}

\begin{exercise}\label{b2:ex:18-period-index-minimal-splitting}
沿用\cref{b2:cor:period-index-biquadratic-tower}中的 \(k,D,i,s\)。设 \(F/k\) 为奇数次数的有限扩张，置于同一个代数闭包中。求 \(D_F\) 的指数与周期，并求 \(D_{F(i)}\)、\(D_{F(i,s)}\) 的指数和两级扩张次数。证明 \(F(i,s)/F\) 仍是最小有限分裂扩张。
\end{exercise}
\begin{solution}
令 \(e=\operatorname{ind}(D_F)\)。\cref{b2:prop:index-after-base-change}使 \(4/e\) 同时整除四与奇数 \([F:k]\)，所以 \(e=4\)。原来的 \(2[D]=0\) 经标量扩张仍成立，而指数四排除零类，故周期仍为二。

\(-1\) 不可能已经在 \(F\) 中有平方根，否则 \(k(i)\subseteq F\) 会使二整除 \([F:k]\)。因此 \([F(i):F]=2\)。这个扩张使四元数因子 \(\mathbb H_F\) 分裂，余类由次数二的另一个因子代表，指数至多为二；但不能为一，否则指数四会整除二次分裂次数，所以恰为二。

再添入 \(s\) 使第二因子也分裂，\(D_{F(i,s)}\) 的指数为一。\([F(i,s):F(i)]\le2\)，且不能为一，否则前一步已经分裂，故等于二。合成域的次数为四，指数四又排除次数更小的有限分裂域，因此它仍为最小分裂扩张。
\end{solution}

\begin{exercise}\label{b2:ex:18-finite-field-cocycles}
设 \(q\) 为素数幂、\(n\ge1\)，\(K=\mathbb F_{q^n}\)、\(k=\mathbb F_q\)。用循环群的范数商证明 \(H^2\bigl(\operatorname{Gal}(K/k),K^\times\bigr)=0\)，并说明如何把一个进位余圈的参数改成一。
\end{exercise}
\begin{solution}
\(n=1\) 时群平凡，归一化余圈只有平凡类。设 \(n\ge2\)。第一卷推论17.2.1和定理17.2.2分别说明 \(K^\times\) 循环、Galois 群由 \(z\mapsto z^q\) 生成。范数为
\[
z\longmapsto z^{1+q+\cdots+q^{n-1}}.
\]
若 \(z\) 为 \(K^\times\) 的生成元，其范数的阶为 \(q-1\)，所以范数满射到 \(k^\times\)。由\cref{b2:cor:cyclic-relative-brauer}，所求上同调群为零。

具体地，进位参数为 \(a\) 时，选 \(b\in K^\times\) 满足 \(\operatorname{N}_{K/k}(b)=a^{-1}\)，将循环生成元换成 \(bu\)，其 \(n\) 次幂便为一。再用它的各次幂作全部基向量，所有进位系数都成为一，所以余圈化为平凡余圈。
\end{solution}

\begin{exercise}\label{b2:ex:18-two-involution-trace-forms}
在 \(A=M_2(\mathbb R)\) 上分别取 \(\sigma_+(X)=X^{\mathsf T}\) 与 \(\sigma_J(X)=J^{-1}X^{\mathsf T}J\)，其中 \(J=\begin{psmallmatrix}0&1\\-1&0\end{psmallmatrix}\)。求两种迹型 \(t_\sigma(X,Y)=\operatorname{tr}\bigl(\sigma(X)Y\bigr)\) 的惯性，说明它们不是等距型。
\end{exercise}
\begin{solution}
写 \(X=\begin{psmallmatrix}a&b\\c&d\end{psmallmatrix}\)。第一种有 \(t_{\sigma_+}(X,X)=a^2+b^2+c^2+d^2\)，惯性为 \((4,0,0)\)。第二种直接计算为
\[
\sigma_J(X)=\begin{pmatrix}d&-b\\-c&a\end{pmatrix},\qquad
t_{\sigma_J}(X,X)=2(ad-bc).
\]
令 \(a=s+t,d=s-t,b=u+v,c=u-v\)，这是可逆实线性换元，所得型为 \(2s^2-2t^2-2u^2+2v^2\)，惯性为 \((2,2,0)\)。由实惯性定律，两型不等距。底层代数始终是同一个 \(M_2(\mathbb R)\)，差别来自对合的选择。
\end{solution}

\begin{exercise}\label{b2:ex:18-prime-index-period}
设 \(k\) 为域，\(p\) 为素数。证明指数为 \(p\) 的中心单 \(k\) 代数，其周期必为 \(p\)。由此给出次数为素数的中心单代数的两种可能。反过来，周期为素数是否迫使指数等于周期？
\end{exercise}
\begin{solution}
由\cref{b2:thm:period-divides-index}，周期整除 \(p\)，只能为一或 \(p\)。周期为一表示 Brauer 类为零，除代数代表就是底域，指数应为一，与题设矛盾。所以周期为 \(p\)。

次数为素数时，写 \(A=M_r(D)\)，有 \(p=r\deg D\)。若 \(\deg D=1\)，则 \(A=M_p(k)\) 分裂；否则 \(r=1\)，\(A\) 本身是次数 \(p\) 的除代数，指数、周期均为 \(p\)。逆向断言则被\cref{b2:exm:period-two-index-four}否定：其中周期二为素数，指数却为四。
\end{solution}

% !TeX root = ../../Book2.tex
% 纲要标题：## 第19章*　非交换局部化、多项式恒等式与增长
% 纲要位置：计划纲要.md，第 1668 行。

\OptionalChapter{非交换局部化、多项式恒等式与增长}{b2:ch:19}
\BookChapterNumber{b2:ch:19}

\begin{chapterabstract}
本章研究非交换环中允许分母换写的 Ore 条件，构造正则分母的局部化，并证明半素右 Goldie 环取得半单经典分式环的条件。矩阵的 Amitsur–Levitzki 多项式恒等式与有限生成代数的 Gelfand–Kirillov 增长维数，分别从恒等式和增长两方面考察非交换代数。
\end{chapterabstract}

\begin{learninggoals}
\begin{itemize}
\item 区分左右 Ore 方程，核对共同分母、分式等价关系、乘法和泛性质中的次序。
\item 掌握本质子模与一致维数，完整证明 Goldie 定理的充分方向，说明半素性及两项右 Goldie 条件的用途。
\item 证明矩阵的标准恒等式，说明有理系数证明怎样通过整数泛型矩阵推广到任意特征。
\item 计算有限生成代数的增长，证明 GK 维数不依赖生成空间，并区分增长、PI 性和分式存在性。
\end{itemize}
\end{learninggoals}

设 \(R=k\langle x,y\rangle\) 为域 \(k\) 上的自由结合代数。非零元素 \(x\) 与 \(y\) 没有非零的共同右倍数：\(xR\) 中各个字都以 \(x\) 开头，\(yR\) 中各个字都以 \(y\) 开头。若想把所有非零元素当成右分母，就必须能为这两个元素找到非零的共同右倍数，因而在此处已经失败。非交换分式不能仅靠“分母非零”来定义；分母能否越过分子，是必须先解决的乘法问题。

\ProseParagraphBreak

Ore 条件把分母换写要求表达为环内的方程；Goldie 理论则寻找使这些方程普遍可解的右理想条件，并判断分式环何时半单。本章另有两类相对独立的问题：一个非交换多项式能否在某个代数中对所有代入都取零，以及长度有限的乘积能张成多大的空间。前者检查所有代入共同满足的乘法关系，后者计算短字张成的线性空间。分母方程是否可解、是否存在这样的共同恒等式、短字怎样增长，是三个需要分别研究的问题；自由代数在分母方程处失败，并不把后两类问题变成前一类问题的推论。

\ProseParagraphBreak
非交换分式环与 Goldie 理论的进一步阅读，可参阅 T. Y. Lam 的《Lectures on Modules and Rings》\cite[\S 10A--10B, ``Ore Localizations'' and ``Right Ore Rings and Domains''; \S 11, ``Right Goldie Rings and Goldie's Theorems'']{Lam1999ModulesRings}。不同局部化理论对分母是否允许零因子、原环能否单射地映入局部化有不同要求。

% !TeX root = ../../Book2.tex
% 纲要标题：### 19.1 Ore 局部化
% 纲要位置：计划纲要.md，第 1670 行。

\section{Ore 局部化}
\label{b2:sec:1901}

交换环中把分母移到乘积的一侧，依靠的是元素可以交换。非交换环里，写下 \(as^{-1}\) 之后，还须解释 \(s^{-1}b\) 如何重新写成分母在右侧的形式。Ore 条件正是为这一步提供可解的方程。本节只讨论由双侧非零因子组成的分母集，从而原环能够单射到分式环。\cref{b2:prop:zero-divisor-inversion}将说明，若把原环中的零因子变成可逆元，单射性就必然丢失。

% 纲要要点（供目录核对）：
%
% * 分母在非交换情形的障碍；由实际右分式乘积产生 Ore 方程，正文证明只逆一个正则元素也会失败、逆零因子必破坏单射性
% * 左、右 Ore 条件；按最终分母位置区分，共同分母属于分母集而正则乘子不必属于其中
% * 典型局部化及其泛性质；先构造加法分式类，再由左乘算子复合给出结合乘法，解释被迫可逆的乘子及非中心矩阵计算
%

\subsection{分母在非交换情形的障碍}
\label{b2:subsec:1901:01}

本节设 \(R\ne0\) 为含单位元的环。若 \(s\in R\) 满足
\[
sx=0\Longrightarrow x=0,\qquad xs=0\Longrightarrow x=0,
\]
称 \(s\) 为\term[zhengze-yuansu]{正则元素}[regular element]，即它既不是左零因子，也不是右零因子。正则元素的乘积仍正则，一也是正则元素。

\ProseParagraphBreak
假设 \(R\) 已嵌入某个环 \(Q\)，分母集 \(S\) 中的元素在 \(Q\) 中可逆，并且每个 \(Q\) 元素都能写为右分式。对 \(p,a\in R\) 及 \(s,r\in S\)，乘积为
\[
(ps^{-1})(ar^{-1})=p(s^{-1}a)r^{-1}.
\]
要使中间的 \(s^{-1}a\) 也成为右分式，必须将它写成 \(bt^{-1}\)，其中 \(b\in R,t\in S\)。这个等式两侧分别左乘 \(s\)、右乘 \(t\)，就得到
\[
at=sb.
\]
因此共同倍数方程正是乘法中移换分母的要求。取得这样的 \(b,t\) 后，乘积便能整理为 \((pb)(rt)^{-1}\)。

\ProseParagraphBreak
以自由代数 \(R=k\langle x,y\rangle\) 为例。由\cref{b2:prop:free-associative-algebra}，有限字构成基。对非零多项式按字长、同长度内再按字典序选最大字，乘积中字长达到最大时，两因子都须各自达到最大字长；在这些固定长度的字中，任一因子按字典序变小，连接后的字也变小。因此乘积的最大字只能由两个最大字连接而成，系数为两个非零首系数的乘积，故乘积非零。因此所有非零元素正则。

\ProseParagraphBreak
但 \(xR\cap yR=0\)，因为第一项中的非零字都以 \(x\) 开头，第二项中的非零字都以 \(y\) 开头。若 \(a=x,s=y\) 能满足上述关系且 \(t\ne0\)，就有 \(0\ne xt=yb\in xR\cap yR\)，矛盾。所以这个非交换整环不能通过“所有非零元素作右分母，且每个元素是一个右分式”的方式得到分式除环。这不排除其他类型的扩张环，只说明通常的单分式表达需要额外条件。即使把分母限制为一个正则元素 \(x\) 的各次幂，所需方程也未必可解，见\cref{b2:prop:single-element-ore-obstruction}。

\subsection{左、右 Ore 条件}
\label{b2:subsec:1901:02}

\begin{definition}\label{b2:def:ore-set}
设 \(R\ne0\) 为含幺环，\(S\subseteq R\) 含一、乘法封闭，且全部元素均双侧正则。对任意 \(a\in R,s\in S\)，若能取 \(t\in S,b\in R\) 使 \(at=sb\)，则称 \(S\) 满足\term[you-Ore-tiaojian]{右 Ore 条件}[right Ore condition]，并称这样的 \(S\) 为\term[zhengze-you-Ore-ji]{正则右 Ore 集}[regular right Ore set]；在本节的正则情形，也称右 Ore 分母集。若对任意 \(a\in R,s\in S\) 都能取 \(t\in S,b\in R\) 使 \(ta=bs\)，则称 \(S\) 为正则左 Ore 集。两个方程及其分式方向为
\[
\begin{aligned}
\text{右 Ore：}&\quad at=sb,\qquad as^{-1};\\
\text{左 Ore：}&\quad ta=bs,\qquad s^{-1}a.
\end{aligned}
\]
这里的左右按最终分式中分母的位置命名。左条件就是反环中的右条件；即使分母都正则，两条件也不自动相同。
\end{definition}

\begin{lemma}\label{b2:lem:ore-common-denominators}
设 \(R\ne0\) 为含幺环，\(S\subseteq R\) 为正则右 Ore 集。
\begin{enumerate}
\item 若 \(s,su\in S\)，则 \(u\) 正则，但不一定属于 \(S\)。
\item 任意 \(s,t\in S\) 有共同右倍数 \(su=tv\in S\)，其中乘子 \(u,v\) 均正则。任意有限多个分母也有这样的共同右倍数。
\end{enumerate}
\end{lemma}
\begin{proof}
若 \(ux=0\)，则 \((su)x=0\)，由 \(su\) 正则得 \(x=0\)。对分子 \(s\) 和分母 \(su\) 使用 Ore 条件，取 \(v\in S,b\in R\) 使 \(sv=(su)b\)。消去左侧正则因子 \(s\)，得 \(v=ub\)。若 \(xu=0\)，则 \(xv=xub=0\)，由 \(v\) 正则得 \(x=0\)。故 \(u\) 双侧正则。

对分子 \(t\)、分母 \(s\) 使用 Ore 条件，得 \(tv=su\)，其中 \(v\in S\)。共同值属于 \(S\)，第一部分保证 \(u\) 也正则。有限多个分母逐次与已经取得的共同分母再作此操作即可；每个旧分母都仍右整除新的共同值。
\end{proof}

共同分母必须属于 \(S\)，放大分母所用的乘子则只需正则，以便在比较分子时消去。例如在 \(\mathbb Z\) 中取 \(S=\{2^i6^j:i,j\ge0\}\)，有 \(2,6\in S\) 及 \(2\cdot3=6\)，但乘子三不属于 \(S\)。

\ProseParagraphBreak
在交换环中，或更一般地在分母都位于中心时，取 \(t=s\)、\(b=a\) 即满足右 Ore 条件，左条件也成立。一般情形的共同倍数则可能需要非平凡的乘子；不能直接把共同分母写成 \(st\)，因为 \(t\) 未必右整除 \(st\)。

\subsection{典型局部化及其泛性质}
\label{b2:subsec:1901:03}

共同右分母足以定义分式的加法，而非交换乘法还要处理分母与分子相遇时的次序。先构造加法群，再将分式实现成这个群上的左乘算子，就能用算子复合取得结合乘法。若 \(su=tv=w\in S\)，预期的换分母等式是 \(as^{-1}=(au)w^{-1}\)、\(bt^{-1}=(bv)w^{-1}\)；因此把两分式视为相同，正是要求它们通分后有相同的分子。

\begin{lemma}[右分式类的加法]\label{b2:lem:ore-additive-classes}
设 \(R\ne0\) 为含幺环，\(S\subseteq R\) 为正则右 Ore 集。在 \(R\times S\) 上规定
\begin{equation}\label{b2:eq:ore-equivalence}
(a,s)\sim(b,t)
\quad\Longleftrightarrow\quad
\exists u,v\in R:\ su=tv\in S,\quad au=bv.
\end{equation}
则 \(\sim\) 是等价关系。其商集 \(F=(R\times S)/{\sim}\) 上的公式
\[
[a,s]+[b,t]=[au+bv,w]\quad(su=tv=w\in S),
\qquad -[a,s]=[-a,s]
\]
给出交换群结构。其零元满足
\begin{equation}\label{b2:eq:ore-zero-numerator}
[a,s]=0\quad\Longleftrightarrow\quad a=0.
\end{equation}
\end{lemma}
\begin{proof}
自反性与对称性直接成立。证明传递性时，设 \(su=tv=w\)、\(au=bv\)，又有 \(tu'=rv'=w'\)、\(bu'=cv'\)。取共同右倍数 \(w\alpha=w'\beta\in S\)。从 \(tv\alpha=tu'\beta\) 消去左侧正则元 \(t\)，得到 \(v\alpha=u'\beta\)，于是
\[
s u\alpha=r v'\beta,\qquad
a u\alpha=bv\alpha=bu'\beta=cv'\beta.
\]
故 \((a,s)\sim(c,r)\)。

先核对一个换分母事实。若 \((a,s)\sim(b,t)\)，则对任意共同右分母 \(w=su=tv\in S\)，都有 \(au=bv\)。事实上，取等价关系中原有的共同分母 \(w_0=su_0=tv_0\)，再取 \(w_0\alpha=w\beta\in S\)。消去 \(s,t\) 得 \(u_0\alpha=u\beta\)、\(v_0\alpha=v\beta\)，故
\[
(au-bv)\beta=(au_0-bv_0)\alpha=0.
\]
由\cref{b2:lem:ore-common-denominators}，乘子 \(\beta\) 正则，因而 \(au=bv\)。

固定两对代表元，若分别选 \(w=su=tv\) 与 \(w'=su'=tv'\)，取 \(W=w\alpha=w'\beta\in S\)。消去 \(s,t\) 得 \(u\alpha=u'\beta\) 及 \(v\alpha=v'\beta\)，所以两个候选和扩张到 \(W\) 后分子相同。这证明共同分母的选择无关。

若把第一对代表元 \((a,s)\) 换为等价的 \((a',s')\)，而 \((b,t)\) 保持不变，分别取 \(w=su=tv\)、\(w'=s'u'=tv'\)，再取 \(W=w\alpha=w'\beta\in S\)。由刚证的换分母事实，\(a u\alpha=a'u'\beta\)；由 \(tv\alpha=tv'\beta\) 消去 \(t\)，得 \(v\alpha=v'\beta\)。因此
\[
(au+bv)\alpha=(a'u'+bv')\beta,
\]
两个候选和等价。若改换第二对代表元 \((b,t)\sim(b',t')\)，固定 \((a,s)\)，分别取 \(w=su=tv\)、\(w'=su'=t'v'\)，再取 \(W=w\alpha=w'\beta\in S\)。消去 \(s\) 得 \(u\alpha=u'\beta\)；换分母事实给出 \(bv\alpha=b'v'\beta\)。于是
\[
(au+bv)\alpha=(au'+b'v')\beta,
\]
两个候选和仍等价。依次换两对代表元即可处理任意代表元选择。

由\cref{b2:lem:ore-common-denominators}，三个类可取同一共同右分母，结合律与交换律于是归结为 \(R\) 的加法；\([0,1]\) 是零元，\([-a,s]\) 是 \([a,s]\) 的负元。最后，\([a,s]=[0,1]\) 时存在 \(u,v\in R\) 使 \(su=v\in S\)、\(au=0\)；由 \(u\) 正则得 \(a=0\)。反向由共同分母 \(s\) 直接得到。
\end{proof}

\begin{theorem}[正则分母的右 Ore 局部化]\label{b2:thm:right-ore-localization}
设 \(R\ne0\) 为含幺环，\(S\subseteq R\) 为正则右 Ore 集，则存在含幺环 \(Q=RS^{-1}\) 及保幺单射 \(R\rightarrow Q\)，使 \(S\) 中每个元素在 \(Q\) 中可逆，且每个 \(Q\) 元素都可写为 \(as^{-1}\)，其中 \(a\in R,s\in S\)。对任意保幺环同态 \(f:R\rightarrow B\)，若全部 \(f(s)\) 在 \(B\) 中可逆，则存在唯一保幺环同态
\[
\widetilde f:Q\longrightarrow B,\qquad
\widetilde f(as^{-1})=f(a)f(s)^{-1}.
\]
\end{theorem}
\begin{proof}
由\cref{b2:lem:ore-additive-classes}得到右分式类的加法群 \(F\)，在它的自同态环 \(\operatorname{End}_{\mathbb Z}(F)\) 中构造分式环。

对 \(a\in R\) 定义加法群自同态 \(L_a[b,t]=[ab,t]\)。等价关系和加法公式说明它良定义且可加，且 \(L_{a+b}=L_a+L_b\)、\(L_aL_b=L_{ab}\)、\(L_1=\operatorname{id}\)。对分母 \(s\in S\)，还须证明 \(L_s\) 是自同构。它单射，因为 \([sb,t]=0\) 迫使 \(sb=0\)，进而 \(b=0\)。它也满射：对 \([b,t]\)，由 Ore 条件取 \(v\in S,u\in R\) 使 \(bv=su\)。因 \(tv\in S\)，可以通分为
\[
[b,t]=[bv,tv]=L_s[u,tv].
\]
故 \(L_s\) 可逆。

在 \(\operatorname{End}_{\mathbb Z}(F)\) 中考虑所有算子 \(L_aL_s^{-1}\)。若 \(su=tv=w\in S\)，由 \(L_w=L_sL_u=L_tL_v\) 可知 \(L_u=L_s^{-1}L_w\)、\(L_v=L_t^{-1}L_w\) 也可逆。这一步用的是 \(s,t,w\) 都已被逆转，而不只是乘子 \(u,v\) 正则。于是
\[
L_aL_s^{-1}=L_{au}L_w^{-1},\qquad
L_bL_t^{-1}=L_{bv}L_w^{-1}.
\]
因此它们对加法与负元封闭，且等价分式给出相同算子。若 \(bv=su\)、\(v\in S\)，则
\[
L_s^{-1}L_b=L_uL_v^{-1},
\]
从而
\begin{equation}\label{b2:eq:right-ore-product}
\begin{aligned}
(L_aL_s^{-1})(L_bL_t^{-1})
&=L_{au}L_v^{-1}L_t^{-1}\\
&=L_{au}L_{tv}^{-1}.
\end{aligned}
\end{equation}
所以这些算子组成含单位元的子环。它们与 \(F\) 的元素一一对应：\(L_aL_s^{-1}\) 在 \([1,1]\) 处的值恰为 \([a,s]\)，因为 \(L_s[1,s]=[1,1]\)。故相同算子必来自相同分式类，反向已由共同分母证明。将这个子环的乘法转移到 \(F\)，得到结合环 \(Q\)。乘法结合律由算子复合继承，分配律也不需另作分母猜测。

映射 \(a\mapsto L_a\) 为保单位元环同态，且由在 \([1,1]\) 处取值及式\eqref{b2:eq:ore-zero-numerator}可知它单射。每个 \(L_s^{-1}=L_1L_s^{-1}\) 也在 \(Q\) 中，所以分母可逆。于是将 \([a,s]\) 记作 \(as^{-1}\) 后，式\eqref{b2:eq:right-ore-product}给出具体乘法
\[
(as^{-1})(bt^{-1})=(au)(tv)^{-1},
\qquad bv=su,\quad v\in S.
\]

最后核对泛性质。若 \(su=tv=w\in S\)，则 \(f(u)=f(s)^{-1}f(w)\) 可逆：即使 \(u\notin S\)，任何把 \(s,w\) 变成单位的同态也会迫使 \(u\) 变成单位。并有
\[
f(a)f(s)^{-1}=f(au)f(w)^{-1}.
\]
因此式\eqref{b2:eq:ore-equivalence}的相同分子保证 \(\widetilde f\) 良定义，共同分母公式保证可加。由 \(bv=su\)，得到 \(f(s)^{-1}f(b)=f(u)f(v)^{-1}\)，恰好验证乘法公式被保持。单位也被保持。每个元素都是 \(as^{-1}\)，其像已由 \(f\) 及分母求逆决定，故延拓唯一。
\end{proof}
\symidx{\(RS^{-1}\)：右 Ore 分式环}

这称为\term[Ore-jubuhua]{Ore 局部化}[Ore localization]。对反环应用同一定理，得到左 Ore 分式环 \(S^{-1}R\)。若同一个 \(S\) 同时满足左右条件，两种环都满足使 \(S\) 中元素可逆的同一个泛性质。分别延拓原环到另一分式环的映射，得到两个方向的同态；其复合都在 \(R\) 上为恒等，泛性质中的唯一性便使复合在整个分式环上也为恒等。因此两种构造具有保持 \(R\) 的互逆同构。这里所用的唯一性论证与\cref{b2:prop:universal-object-uniqueness}中积、余积的情形相同。

\begin{proposition}\label{b2:prop:ore-clearing-denominators}
设 \(R\ne0\) 为含幺环，\(S\subseteq R\) 为正则右 Ore 集，\(Q=RS^{-1}\) 为其右 Ore 局部化。对任意有限组 \(q_1,\ldots,q_m\in Q\)，存在共同右分母 \(s\in S\) 使全部 \(q_i s\in R\)。若 \(I\) 为 \(R\) 的右理想，则
\[
IQ=\{\,as^{-1}:a\in I,\ s\in S\,\}.
\]
每个非零右 \(R\) 子模 \(X\subseteq Q\) 都与 \(R\) 有非零交。
\end{proposition}
\begin{proof}
先把各分母取共同右倍数，分子随之扩张，就得到第一项。\(IQ\) 中元素为有限和 \(\sum_i a_iq_i\)、\(a_i\in I\)，通分后分子仍属于右理想 \(I\)，得到第二式的一向；另一向由 \(as^{-1}\in IQ\) 直接成立。最后，若 \(0\ne as^{-1}\in X\)，右乘 \(s\in R\) 得到 \(0\ne a\in X\cap R\)，其中非零性来自分母在 \(Q\) 中可逆。
\end{proof}

例如令 \(C\) 为交换整环，\(R=M_n(C)\)，分母取非零标量矩阵 \(sI_n\)。它们在中心且正则，局部化得到
\[
RS^{-1}\cong M_n(\operatorname{Frac}C).
\]
\symidx{\(\operatorname{Frac}C\)：交换整环的分式域}
每个右分式逐条目取标量分母，反向对一个分式域矩阵的有限多个条目取共同分母，便写成一个矩阵右分式；两种操作互逆且保持加乘。原环可能有许多矩阵零因子，但这些元素没有被放入分母集。

\ProseParagraphBreak
也可以让分母不在中心，此时 \(s^{-1}a\) 与 \(as^{-1}\) 可能不同，换分母还须改变分子。取 \(R=M_2(\mathbb Z)\)，令 \(S\) 为其中行列式非零的全部矩阵。它们在 \(M_2(\mathbb Q)\) 中可逆，因此在 \(R\) 中双侧正则，且 \(S\) 含一、乘法封闭。给定 \(a\in R,s\in S\)，选正整数 \(q\) 清除 \(s^{-1}a\) 的全部条目分母，令 \(t=qI_2\in S\)、\(b=s^{-1}at\in R\)，便有 \(at=sb\)。所以 \(S\) 满足右 Ore 条件。所得分式环仍为 \(M_2(\mathbb Q)\)：每个右分式可在该矩阵环中解释，而每个有理矩阵都能用非零整数标量矩阵作共同分母；若右分式在其中为零，其分子也必为零。

\ProseParagraphBreak
具体取
\[
s=\begin{pmatrix}1&1\\0&2\end{pmatrix},\qquad
a=E_{21},\qquad t=2I_2,\qquad
b=\begin{pmatrix}-1&0\\1&0\end{pmatrix}.
\]
直接相乘得到
\[
at=\begin{pmatrix}0&0\\2&0\end{pmatrix}=sb.
\]
于是非中心分母可以按 Ore 方程移到右侧，但分子也随之改变：
\[
s^{-1}a=bt^{-1}
=\begin{pmatrix}-\frac12&0\\\frac12&0\end{pmatrix},
\qquad
as^{-1}=\begin{pmatrix}0&0\\1&-\frac12\end{pmatrix}.
\]
例如两个右分式的乘积为
\[
\begin{aligned}
(E_{11}s^{-1})(as^{-1})
&=E_{11}b\,t^{-1}s^{-1}\\
&=E_{11}b(st)^{-1}\\
&=-E_{11}(2s)^{-1}
=\begin{pmatrix}-\frac12&\frac14\\0&0\end{pmatrix}.
\end{aligned}
\]
这里先用 \(s^{-1}a=bt^{-1}\) 处理两个分式之间相遇的分母与分子，再用 \(t^{-1}s^{-1}=(st)^{-1}\) 合并分母；次序来自实际乘法。

\begin{proposition}\label{b2:prop:single-element-ore-obstruction}
设 \(k\) 为域，\(R=k\langle x,y\rangle\) 为两个生成元上的自由结合含幺 \(k\) 代数，\(S=\{x^n:n\ge0\}\)。则 \(S\) 含一、乘法封闭，全部元素双侧正则，但不满足右 Ore 条件。
\end{proposition}
\begin{proof}
由\cref{b2:prop:free-associative-algebra}，\(x,y\) 上的有限字是 \(R\) 的一组 \(k\) 基，空字表示一。对每个 \(n\ge0\)，在字的左侧或右侧连接 \(x^n\)，都把不同的字送到不同的字。因此左乘与右乘 \(x^n\) 都是单射，\(x^n\) 双侧正则。\(x^0=1\) 及 \(x^i x^j=x^{i+j}\) 还说明 \(S\) 含一、乘法封闭。

若满足右 Ore 条件，对 \(a=y,s=x\) 就能取 \(t=x^n\in S\)、\(b\in R\)，使 \(yx^n=xb\)。左侧是以 \(y\) 开头的非零基字，右侧的每个非零基字都以 \(x\) 开头。字基的线性无关性排除该等式，所以右 Ore 条件失败。
\end{proof}

\begin{proposition}\label{b2:prop:zero-divisor-inversion}
设 \(R\ne0\)、\(B\) 为含幺环，\(f:R\to B\) 为保幺单射。若 \(s\in R\) 的像 \(f(s)\) 可逆，则 \(s\) 双侧正则。

具体地，设 \(k\) 为域，\(R=k\times k\)、\(e=(1,0)\)。第一坐标投影 \(\pi:R\to k\) 把 \(e\) 送到单位元，且对每个含幺环 \(B\) 及每个使 \(f(e)\) 可逆的保幺环同态 \(f:R\to B\)，都存在唯一保幺环同态 \(\bar f:k\to B\) 使 \(f=\bar f\pi\)。但是 \(\ker\pi=0\times k\ne0\)，任何这样的 \(f\) 都不是单射。
\end{proposition}
\begin{proof}
若 \(sx=0\)，则 \(f(s)f(x)=0\)，乘以 \(f(s)^{-1}\) 得 \(f(x)=0\)，再由单射性得 \(x=0\)。若 \(xs=0\)，在右侧乘以 \(f(s)^{-1}\) 同样得到 \(x=0\)。所以可嵌入后变成单位的元素必须双侧正则。

在第二个例子中，\(e^2=e\)。可逆的 \(f(e)\) 满足 \(f(e)^2=f(e)\)，消去一个因子得到 \(f(e)=1_B\)，于是 \(f(1-e)=0\)。每个第二坐标元素都为 \((0,b)=(1-e)(0,b)\)，故被 \(f\) 送到零。定义 \(\bar f(a)=f(a,0)\)，它保持加法与乘法，并把一送到 \(f(e)=1_B\)，因而为保幺环同态。又有 \(f(a,b)=f(a,0)=\bar f(\pi(a,b))\)，所以 \(f=\bar f\pi\)。投影 \(\pi\) 满射，保证这种分解唯一；它本身满足 \(\pi(e)=1\)，且核恰为非零的第二坐标因子。因此使 \(e\) 可逆必然零化原环中的非零元素。
\end{proof}

% !TeX root = ../../Book2.tex
% 纲要标题：### 19.2 Goldie 理论
% 纲要位置：计划纲要.md，第 1676 行。

\section{Goldie 理论}
\label{b2:sec:1902}

Ore 条件逐个要求方程 \(at=sb\) 可解，直接核验往往困难。Goldie 理论用环的右理想结构保证这些方程有解，并进一步说明所得分式环为何半单。本节明确改用右模：一致维数、子模和零化子链均按所写的作用方向理解。证明先建立本质右理想中正则元素的存在性，再进行局部化。

% 纲要要点（供目录核对）：
%
% * 一致维数及零化子条件；以整数、向量空间和截断多项式模区别本质、一致、长度，解释有限本质一致直和的构造与逆向证明
% * 半素右 Goldie 环的分式环；分别追踪半素性、右零化子升链条件及有限一致维数，补清本质理想含正则元、Ore 乘子的原像和右半单转为左半单的步骤
% * 定理假设与交换整环的比较；正文证明整环右 Ore 等价于右正则模一致，正式证明上三角环中的本质理想无正则元，区分矩阵的行模数与基域维数
%

\subsection{一致维数与零化子条件}
\label{b2:subsec:1902:01}

\begin{definition}\label{b2:def:essential-uniform-dimension}
设 \(R\) 为含幺环，\(M\) 为幺性右 \(R\) 模，\(E\subseteq M\) 为子模。若每个非零子模 \(N\subseteq M\) 都满足 \(N\cap E\ne0\)，则称 \(E\) 为 \(M\) 的\term[benzhi-zimo]{本质子模}[essential submodule]。若非零幺性右 \(R\) 模 \(U\) 的任意两个非零子模交非零，则称 \(U\) 为\term[yizhi-mo]{一致模}[uniform module]；等价地，\(U\) 的每个非零子模都在 \(U\) 中本质。

若 \(M\) 的非零子模 \(N_1,\ldots,N_r\) 的和为内部直和，则称它们独立。定义 \(M\) 的\term[yizhi-weishu]{一致维数}[uniform dimension] \(\operatorname{udim}M\) 为这类有限独立族大小的上确界，允许为无穷；零模的值为零。一致维数有限时也称为 \term[Goldie-zhi]{Goldie 秩}[Goldie rank]。
\end{definition}
\symidx{\(\operatorname{udim}M\)：右模的一致维数}

本质包含要求每个非零子模都与给定子模相交，并不要求给定子模等于全模。例如 \(2\mathbb Z\) 在 \(\mathbb Z\) 中本质，因为任意非零子模 \(d\mathbb Z\) 都与它有非零公共元素 \(2d\)。但域 \(k\) 上向量空间 \(V\) 的真子空间 \(W\) 不本质：取 \(v\notin W\)，则 \(kv\cap W=0\)。

\ProseParagraphBreak
一致模不能同时容纳两个交零的非零子模，所以对非零模 \(U\)，一致性等价于 \(\operatorname{udim}U=1\)。对除环上的右向量空间，这恰好是维数为一的条件：一维空间只有零子空间与自身，而两个线性无关向量张成的直线交于零。简单模因此一致，但一致模未必简单。例如交换整环 \(R\) 的右正则模一致：两个非零理想分别取非零元素 \(a,b\)，非零乘积 \(ab\) 落在两者的交中。整数模因而一致，虽然有无限严格降链 \(\mathbb Z\supsetneq2\mathbb Z\supsetneq4\mathbb Z\supsetneq\cdots\)，并不具有有限长度。

\ProseParagraphBreak
一致维数数的是同时独立的直和项，合成长度则还计算这些项内部的简单层。取域 \(k\)、整数 \(m\ge1\) 及 \(R=k[t]/(t^m)\)，把 \(t\) 的剩余类仍记为 \(t\)。右理想与 \(k[t]\) 中包含 \((t^m)\) 的理想对应；\(k[t]\) 为主理想整环，这些理想的生成元整除 \(t^m\)，故可取为 \(t^i\)。所以 \(R\) 的右理想恰为 \(t^iR\)、\(0\le i\le m\)。它们彼此可比，任意两个非零右理想的交非零，因此 \(\operatorname{udim}R_R=1\)。另一方面，链
\[
0=t^mR\subsetneq t^{m-1}R\subsetneq\cdots\subsetneq tR\subsetneq R
\]
的每个相邻商 \(t^iR/t^{i+1}R\) 都由 \(t^i\) 的类张成，且 \(t\) 在该商上作用为零，所以同构于简单右模 \(R/tR\cong k\)。这给出长度为 \(m\) 的合成列，故 \(\ell(R_R)=m\)。

\begin{lemma}\label{b2:lem:essential-tools}
设 \(R\) 为含幺环，所涉右 \(R\) 模均为幺性模。若 \(M\) 的子模 \(E\subseteq N\subseteq M\) 满足 \(E\) 在 \(N\) 中本质、\(N\) 在 \(M\) 中本质，则 \(E\) 在 \(M\) 中本质。若 \(E\subseteq N\) 是右 \(R\) 模中的本质包含，\(f:M\rightarrow N\) 是右 \(R\) 模同态，则 \(f^{-1}(E)\) 在 \(M\) 中本质。对有限族右 \(R\) 模中的本质包含 \(E_i\subseteq M_i\)，还有本质包含 \(\bigoplus_iE_i\subseteq\bigoplus_iM_i\)。
\end{lemma}
\begin{proof}
若 \(E\) 本质于 \(N\)，而 \(N\) 本质于 \(M\)，任意非零 \(L\subseteq M\) 先与 \(N\) 非零相交，再与 \(E\) 非零相交，故传递性成立。

对非零子模 \(L\subseteq M\)，若 \(f(L)=0\)，则 \(L\subseteq f^{-1}(E)\)；否则 \(f(L)\cap E\ne0\)，取一个非零像的原像 \(x\in L\)，便有 \(0\ne x\in L\cap f^{-1}(E)\)。这证明原像的本质性。

在有限直和中，第 \(i\) 个投影的原像 \(\pi_i^{-1}(E_i)\) 由上段本质。有限个本质子模的交仍本质：先与第一个相交，再逐个与其余相交，每一步保持非零。它们的交恰为 \(\bigoplus_iE_i\)，得到最后结论。
\end{proof}

\begin{lemma}\label{b2:lem:uniform-submodule-existence}
设 \(R\) 为含幺环，\(M\) 为非零幺性右 \(R\) 模。若 \(M\) 的一致维数有限，或 \(M\) 为 Noether 模，则 \(M\) 的每个非零子模都含有一致子模。
\end{lemma}
\begin{proof}
假设某个非零子模 \(N\subseteq M\) 不含一致子模。令 \(N_0=N\)。由于 \(N_0\) 不一致，可取非零子模 \(N_1,W_0\subseteq N_0\)，使 \(N_1\cap W_0=0\)。\(N_1\) 也不含一致子模，故可继续在 \(N_1\) 中取非零子模 \(N_2,W_1\)，使 \(N_2\cap W_1=0\)，依此类推。每一步都有 \(N_{i+1}\oplus W_i\subseteq N_i\)，归纳可知 \(W_0,W_1,\ldots\) 独立。

若 \(\operatorname{udim}M<\infty\)，任意多个 \(W_i\) 构成的有限独立族便与一致维数的上界矛盾。若 \(M\) 为 Noether 模，则 \(W_0\subsetneq W_0\oplus W_1\subsetneq\cdots\) 是子模的无限严格升链，同样矛盾。
\end{proof}

\begin{proposition}\label{b2:prop:uniform-dimension-structure}
设 \(R\) 为含幺环，\(M\) 为非零幺性右 \(R\) 模。若 \(\operatorname{udim}M=d<\infty\)，则 \(M\) 含有本质子模 \(E=U_1\oplus\cdots\oplus U_d\)，其中各 \(U_i\) 是一致右 \(R\) 模。反过来，只要 \(M\) 含有这样的有限本质直和，就有 \(\operatorname{udim}M=d\)。若子模 \(N\subseteq M\) 且 \(M\) 一致维数有限，则
\[
\operatorname{udim}N=\operatorname{udim}M
\quad\Longleftrightarrow\quad N\text{ 在 }M\text{ 中本质}.
\]
每个 Noether 右 \(R\) 模都有有限一致维数。
\end{proposition}
\begin{proof}
先设 \(\operatorname{udim}M=d<\infty\)。由\cref{b2:lem:uniform-submodule-existence}取一致子模 \(U_1\)。若已取的有限直和 \(U_1\oplus\cdots\oplus U_r\) 不本质，就有非零子模与它交于零，在其中再取一致子模 \(U_{r+1}\)。一致维数限制新增项数，所以过程在有限步后停止，得到本质直和 \(E\)。

现在证明：本质直和 \(E=\bigoplus_{i=1}^dU_i\) 存在时，任意独立非零族 \(N_1,\ldots,N_s\subseteq M\) 都满足 \(s\le d\)。先以 \(N_j\cap E\) 代替 \(N_j\)，这些交仍非零且独立，故可假设它们在 \(E\) 中。对 \(d\) 归纳，\(d=0\) 时没有非零子模。若 \(s>0\)，令 \(W=N_2\oplus\cdots\oplus N_s\)。它与非零 \(N_1\) 交于零，所以不在 \(E\) 中本质。要减少一致直和的项数，可以找一个与 \(W\) 交零的 \(U_i\)，再将它商去；关键是证明这样的 \(U_i\) 必存在。如果每个 \(U_i\cap W\) 都非零，一致性及\cref{b2:lem:essential-tools}将使这些交的直和在 \(E\) 中本质，而该直和包含于 \(W\)，迫使 \(W\) 与 \(N_1\) 非零相交，矛盾。因此某个 \(U_i\cap W=0\)。商映射 \(E\rightarrow E/U_i\) 在 \(W\) 上单射，故其余 \(s-1\) 项在 \(E/U_i\cong\bigoplus_{j\ne i}U_j\) 中仍独立。归纳给出 \(s-1\le d-1\)。原来的 \(U_i\) 又给出大小 \(d\) 的独立族，故一致维数恰为 \(d\)。因此，前段构造的本质直和恰有 \(d\) 项。

若 \(M\) 为 Noether 模，\cref{b2:lem:uniform-submodule-existence}同样保证每次都能取出一致子模。若抽取过程不停，所得部分直和就构成严格升链 \(U_1\subsetneq U_1\oplus U_2\subsetneq\cdots\)，违反第十二章的子模升链条件。因此过程在有限步后停止，所得本质直和有有限项数；由刚证的逆向结论，\(M\) 的一致维数有限。

最后，若 \(N\) 本质，把 \(M\) 中任意最大独立族的每项与 \(N\) 相交，得到 \(N\) 中同样大的独立族，所以两维数相等。若 \(N\) 不本质，选 \(0\ne W\subseteq M\) 与 \(N\) 交零。将它加入 \(N\) 的一个最大独立族，得到 \(\operatorname{udim}M\ge\operatorname{udim}N+1\)，故两者不等。
\end{proof}

对 \(X\subseteq R\)，分别记
\[
\begin{aligned}
\operatorname{ann}_r(X)&=\{\,a\in R:xa=0\text{ 对所有 }x\in X\,\},\\
\operatorname{ann}_{\ell}(X)&=\{\,a\in R:ax=0\text{ 对所有 }x\in X\,\}.
\end{aligned}
\]
它们分别为右理想与左理想，称为右零化子与左零化子。
\symidx{\(\operatorname{ann}_r(X)\)：子集的右零化子}
\symidx{\(\operatorname{ann}_{\ell}(X)\)：子集的左零化子}

\begin{definition}\label{b2:def:semiprime-right-goldie}
设 \(R\) 为含幺环。若 \(R\) 没有非零的平方为零的双边理想，则称 \(R\) 为\term[bansu-huan]{半素环}[semiprime ring]。对 \(X\subseteq R\)，记 \(\operatorname{ann}_r(X)=\{\,r\in R\mid xr=0\text{ 对所有 }x\in X\,\}\)。若右正则模 \(R_R\) 一致维数有限，且这些右零化子组成的每条升链都最终稳定，则称 \(R\) 为\term[you-Goldie-huan]{右 Goldie 环}[right Goldie ring]。
\end{definition}

半素也等价于没有非零幂零双边理想：若 \(I^m=0\) 且 \(m\ge2\) 最小，则 \(I^{m-1}\ne0\) 而它的平方为零。这里排除的是双边理想，不是幂零元素；例如对域 \(k\)，单环 \(M_2(k)\) 半素，却有 \(E_{12}\ne0\)、\(E_{12}^2=0\)。半素环中的非零左理想或右理想，其平方也不能为零。比如左理想 \(L\) 若满足 \(L^2=0\)，双边理想 \(LR\) 满足
\((LR)^2\subseteq L^2R=0\)，故 \(LR=0\)，进而 \(L=0\)；右理想采用 \(RL\) 的对应论证。这里零化子升链条件只针对这些特殊右理想，比全部右理想都满足升链条件弱。

\subsection{半素右 Goldie 环的分式环}
\label{b2:subsec:1902:02}

下面要从半素右 Goldie 条件推出每个本质右理想都包含正则元素。半素性使非零一侧理想中能找到乘积非零的元素，右零化子的升链条件允许选择不能继续扩大的零化子；有限一致维数则限制可以加入的独立子模数目，使寻找单射左乘作用的过程在有限步后完成。三个条件将分别用于这些具体论证。

\begin{lemma}\label{b2:lem:essential-annihilator-zero}
设 \(R\) 为非零含幺半素右 Goldie 环，\(E\subseteq R_R\) 为本质右理想，则其左零化子 \(\operatorname{ann}_{\ell}(E)=\{\,r\in R\mid re=0\text{ 对所有 }e\in E\,\}\) 为零。
\end{lemma}
\begin{proof}
若存在本质右理想 \(E\) 满足 \(\operatorname{ann}_{\ell}(E)\ne0\)，取非零左理想 \(L=\operatorname{ann}_{\ell}(E)\)，则 \(\operatorname{ann}_r(L)\supseteq E\)，因而也是本质右理想。故下列右零化子的集合非空：
\[
\mathcal A=
\bigl\{\operatorname{ann}_r(L):
0\ne L\subseteq R\text{ 为左理想，且 }\operatorname{ann}_r(L)\text{ 在 }R_R\text{ 中本质}\bigr\}.
\]
沿用\cref{b2:prop:chain-maximal-minimal}证明中的升链论证，若 \(\mathcal A\) 没有极大元，就可在其中逐次严格增大，得到一条永不稳定的右零化子升链，违反右 Goldie 条件。因此可取 \(A=\operatorname{ann}_r(L)\in\mathcal A\) 为按包含关系极大的元素。

半素性给出 \(L^2\ne0\)，取 \(x,y\in L\) 使 \(xy\ne0\)。\(yR\) 是非零右理想，本质性使其中有非零元素 \(ya\in A\)。因为 \(x\in L\)、\(ya\in\operatorname{ann}_r(L)\)，有 \(xya=0\)，而 \(ya\ne0\)；也就是说，右乘子 \(a\) 零化了 \(xy\)，却没有零化 \(y\)。所以
\[
A\subseteq\operatorname{ann}_r(y)\subsetneq\operatorname{ann}_r(xy).
\]
这里 \(y,xy\in L\)，所以两者都包含 \(A\)；并且
\[
\operatorname{ann}_r(y)=\operatorname{ann}_r(Ry),\qquad
\operatorname{ann}_r(xy)=\operatorname{ann}_r(Rxy).
\]
由于 \(Ry,Rxy\) 均为非零左理想，而包含 \(A\) 的右理想仍本质，这两个零化子都属于 \(\mathcal A\)。\(A\) 的极大性要求它们都等于 \(A\)，与严格包含矛盾。
\end{proof}

\begin{lemma}\label{b2:lem:one-sided-regular-goldie}
设 \(R\) 为非零含幺半素右 Goldie 环，\(s\in R\)。若 \(\operatorname{ann}_r(s)=\{\,r\in R\mid sr=0\,\}=0\)，则 \(s\) 双侧正则，且 \(sR\) 为 \(R_R\) 的本质右理想。
\end{lemma}
\begin{proof}
左乘 \(s\) 给出右模同构 \(R_R\cong sR\)，所以两者一致维数相同。由\cref{b2:prop:uniform-dimension-structure}，\(sR\) 在 \(R\) 中本质。上一引理给出 \(\operatorname{ann}_{\ell}(sR)=0\)，而它正是 \(\operatorname{ann}_{\ell}(s)\)。两侧零化子均为零，所以 \(s\) 正则。
\end{proof}

\begin{lemma}\label{b2:lem:goldie-ideal-regular-element}
设 \(R\) 为非零含幺半素右 Goldie 环，则每个右理想 \(N\subseteq R\) 都含有某个 \(v\in N\)，使 \(\operatorname{ann}_r(v)\cap N=0\)。因此每个本质右理想都含有双侧正则元素。
\end{lemma}
\begin{proof}
先设 \(U\) 为非零一致右理想。由半素性，取 \(x,y\in U\) 使 \(xy\ne0\)。若 \(W=\operatorname{ann}_r(x)\cap U\ne0\)，则 \(W\) 在 \(U\) 中本质。对所有满足 \(ya\in W\) 的右乘子 \(a\)，都有 \(xya=0\)；为了把这些乘子组成的右理想与本质性联系起来，考虑右线性映射 \(\lambda_y:R\rightarrow U\)、\(a\mapsto ya\)。这些乘子恰好组成 \(\lambda_y^{-1}(W)\)，由\cref{b2:lem:essential-tools}它是本质右理想，却被非零元素 \(xy\) 从左零化，违反\cref{b2:lem:essential-annihilator-zero}。故 \(W=0\)，一致右理想中的结论成立。

若 \(N\ne0\)，它含有一致右理想，上一段因而保证存在满足
\[
v\in V\subseteq N,\qquad \operatorname{ann}_r(v)\cap V=0
\]
的非零子模 \(V\) 及其中的元素 \(v\)。选择一对 \((V,v)\)，使 \(\operatorname{udim}V\) 最大；这些维数是被 \(\operatorname{udim}R_R\) 限制的正整数，所以最大值存在。假设仍有 \(\operatorname{ann}_r(v)\cap N\ne0\)，在其中取一致右理想 \(U\)，再由上段取 \(u\in U\) 使 \(\operatorname{ann}_r(u)\cap U=0\)。因为 \(U\subseteq\operatorname{ann}_r(v)\)，有 \(U\cap V=0\)。

令 \(x=u+v\)，是为了在保留 \(v\) 对 \(V\) 的单射作用时，用 \(u\) 补上它在 \(U\) 上的零作用。这里 \(U,V\) 都是右理想，所以 \(u(U\oplus V)\subseteq U\)、\(vV\subseteq V\)，又有 \(vU=0\)。因此左乘 \(x\) 确实给出右模自同态 \(\lambda_x:U\oplus V\rightarrow U\oplus V\)，相对于这个分解写为
\[
\lambda_x=
\begin{pmatrix}
\lambda_u|_U&\lambda_u|_V\\
0&\lambda_v|_V
\end{pmatrix}.
\]
两条对角映射按 \(u,v\) 的选择都是单射，左下角则为零，所以相加不会产生新的核。具体地，若 \(z=a+b\)、\(a\in U,b\in V\)，且 \(xz=0\)，先看 \(V\) 分量得到 \(vb=0\)，从而 \(b=0\)；再看 \(U\) 分量得到 \(ua=0\)，从而 \(a=0\)。故 \(\operatorname{ann}_r(x)\cap(U\oplus V)=0\)。但 \(x\in U\oplus V\subseteq N\)，而把 \(V\) 中大小为 \(\operatorname{udim}V\) 的独立族与非零项 \(U\) 并在一起，就使一致维数至少增加一，与选择的最大值矛盾。因此 \(\operatorname{ann}_r(v)\cap N=0\)。\(N=0\) 时取 \(v=0\) 即可。

若 \(N\) 本质，这个零交迫使右理想 \(\operatorname{ann}_r(v)\) 为零。由\cref{b2:lem:one-sided-regular-goldie}，\(v\) 双侧正则。
\end{proof}

\begin{theorem}[Goldie 定理的充分方向]\label{b2:thm:goldie-sufficiency}
设 \(R\) 为非零含幺半素右 Goldie 环，\(S\subseteq R\) 为全部双侧正则元素，则 \(S\) 满足右 Ore 条件，经典右分式环 \(Q=RS^{-1}\) 为半单 Artin 环。
\end{theorem}
\begin{proof}
一属于 \(S\)，正则元素乘法封闭。给定 \(a\in R,s\in S\)，要解 \(at=sb\)，就是要找到正则的右乘子 \(t\)，使 \(at\in sR\)。这些候选乘子组成右模同态 \(\lambda_a:R_R\rightarrow R_R\)、\(x\mapsto ax\) 的原像 \(\lambda_a^{-1}(sR)\)。由\cref{b2:lem:one-sided-regular-goldie}，\(sR\) 本质，所以\cref{b2:lem:essential-tools}使这个原像也本质。再由\cref{b2:lem:goldie-ideal-regular-element}取其中的正则元素 \(t\)，便有 \(at=sb\) 对某个 \(b\in R\) 成立，得到右 Ore 条件。\cref{b2:thm:right-ore-localization}给出 \(Q\)。

局部化把正则元素变成单位，这将排除 \(Q\) 中的真本质右理想。设 \(E\) 为 \(Q\) 的本质右理想，\(I\) 为 \(R\) 的任意非零右理想，则 \(IQ\) 是非零右 \(Q\) 理想，所以 \(E\cap IQ\) 含有非零元素 \(q\)。由\cref{b2:prop:ore-clearing-denominators}，可取 \(s\in S\) 使 \(qs\in I\)，且 \(s\) 在 \(Q\) 中可逆保证 \(qs\ne0\)。因为 \(E\) 是右理想，\(qs\) 也在 \(E\)，因此 \(E\cap R\) 是 \(R\) 的本质右理想。它含有正则元素 \(t\)，而 \(t^{-1}\in Q\)；右理想 \(E\) 因而含有 \(tt^{-1}=1\)，只能等于 \(Q\)。

要将本质理想的结论变成直和分解，还需为任意右理想找到与它交零、且二者之和本质的子模。任意独立的非零右 \(Q\) 理想与 \(R\) 的交均非零，且作为右 \(R\) 子模仍独立，所以 \(\operatorname{udim}Q_Q\le\operatorname{udim}R_R<\infty\)。对给定右 \(Q\) 理想 \(J\)，在与 \(J\) 交零的右理想中选一致维数最大的 \(K\)：零理想保证候选非空，而这些维数是有上界的非负整数，保证最大值存在。若 \(J\oplus K\) 不本质，取非零右理想 \(L\) 与之交零；则 \(K\oplus L\) 仍与 \(J\) 交零，且将 \(L\) 加入 \(K\) 的最大独立族使一致维数至少增加一，违反所选最大值。故 \(J\oplus K\) 本质。刚证 \(Q\) 没有真本质右理想，所以 \(J\oplus K=Q\)。

因此右正则模 \(Q_Q\) 的每个子模都有补模。此前的\cref{b2:thm:semisimple-module-characterizations}以左模陈述，而半单环也按左正则模定义；把右 \(Q\) 模视为左 \(Q^{\mathrm{op}}\) 模，便可对同一个空间及其子模应用这一刻画，得到
\[
Q_Q\text{ 半单}
\quad\Longleftrightarrow\quad
{}_{Q^{\mathrm{op}}}Q\text{ 半单}
\quad\Longleftrightarrow\quad
Q^{\mathrm{op}}\text{ 为半单环}.
\]
对 \(Q^{\mathrm{op}}\) 应用\cref{b2:thm:artin-wedderburn}，存在有限多个除环 \(D_i\) 和正整数 \(n_i\)，使
\[
Q^{\mathrm{op}}\cong\prod_i M_{n_i}(D_i).
\]
取反环，并以矩阵转置识别 \(M_{n_i}(D_i)^{\mathrm{op}}\cong M_{n_i}(D_i^{\mathrm{op}})\)，便得
\[
Q\cong\prod_i M_{n_i}(D_i^{\mathrm{op}}).
\]
故 \(Q\) 是半单 Artin 环。
\end{proof}

这给出本章所用的\term[Goldie-dingli]{Goldie 定理}[Goldie's theorem]：由半素右 Goldie 条件得到半单经典分式环。本章证明这一充分方向，反向刻画留待进一步阅读。

\begin{corollary}\label{b2:cor:goldie-rank}
设 \(R\) 为非零含幺半素右 Goldie 环，\(S\subseteq R\) 为全部双侧正则元素，\(Q=RS^{-1}\) 为经典右分式环，则 \(\operatorname{udim}R_R=\ell(Q_Q)\)。
\end{corollary}
\begin{proof}
若 \(I_1,\ldots,I_d\) 为 \(R\) 中独立非零右理想，则 \(I_iQ\) 仍独立：一条和为零的关系中的各项，通到共同右分母后成为各 \(I_i\) 中和为零的分子，原独立性使每个分子为零。反过来，\(Q\) 中独立非零右理想与 \(R\) 的交都非零，因为单个非零分式可清分母；这些交仍独立。因此 \(R_R,Q_Q\) 一致维数相同。\(Q_Q\) 是有限个简单右模的直和，简单模一致，\cref{b2:prop:uniform-dimension-structure}说明其一致维数正好等于简单分量个数，也就是长度。
\end{proof}

\subsection{Goldie 假设与交换整环}
\label{b2:subsec:1902:03}

若 \(R\) 右 Noether，则全部右理想满足升链条件，特别是右零化子链稳定；\cref{b2:prop:uniform-dimension-structure}又保证右正则模一致维数有限。因此半素右 Noether 环满足 Goldie 定理。对允许非交换的整环，非零双边理想中的 \(a\ne0\) 满足 \(a^2\ne0\)，所以半素性自动成立，右 Noether 整环因而有经典右分式环；分母包含全部非零元素，使它成为除环，具体见\cref{b2:cor:domain-uniform-ore}。

\ProseParagraphBreak
交换整环甚至不必 Noether：其右正则模一致，而含非零元素的子集零化子为零，故零化子链自动稳定。Goldie 定理便恢复通常的分式域。

\ProseParagraphBreak
在 \(R=M_n(C)\)、\(C\) 为交换整环、\(n\ge1\) 的例子中，正则矩阵恰为行列式非零的矩阵。行列式非零时，伴随矩阵公式及 \(C\) 中的消去律排除两侧零化元；行列式为零时，在分式域中取非零核列，清分母得非零的 \(C\) 系数核列，将它放入一个矩阵的第一列，其余列为零，就得到非零右零化元。因此所有正则矩阵都在 \(Q=M_n(\operatorname{Frac}C)\) 中可逆，而该环的每个元素已可用非零标量矩阵作一个右分母，故它正是经典分式环。由\cref{b2:thm:artin-wedderburn}证明中的行分解，\(Q_Q=E_{11}Q\oplus\cdots\oplus E_{nn}Q\)，每个行理想都是简单右模，故 \(\ell(Q_Q)=n\)。清分母保持独立理想族的同一论证给出 \(\operatorname{udim}R_R=n\)。这里数的是 \(n\) 个简单行模，不是作为底域向量空间的 \(n^2\) 个矩阵坐标。

\ProseParagraphBreak
半素性不能删去。\cref{b2:prop:essential-without-semiprime}将对二阶上三角矩阵环给出一个本质右理想，其中没有正则元素；该环仍满足右 Goldie 条件，其经典分式环也仍是原来的非半单环。

\ProseParagraphBreak
有限一致维数也不能由另外两个条件替代。取域 \(k\)，自由代数 \(R=k\langle x,y\rangle\) 是整环，故半素；其含有非零元素的子集右零化子均为零，所以零化子升链条件成立。但右理想
\[
yR,\ xyR,\ x^2yR,\ldots
\]
独立，因为各自的基字以前缀 \(x^iy\) 开始，这些前缀彼此不会延长成另一个。故一致维数无穷，不能应用 Goldie 定理。这与上一节发现它不满足右 Ore 条件相吻合。

\begin{corollary}\label{b2:cor:domain-uniform-ore}
设 \(R\) 为非零含幺整环，不要求交换，并令 \(S=R\setminus\{0\}\)。则 \(S\) 满足右 Ore 条件，当且仅当幺性右正则模 \(R_R\) 一致。条件成立时，经典右分式环 \(RS^{-1}\) 为除环。
\end{corollary}
\begin{proof}
整环中每个非零元素都双侧正则，\(S\) 含一且乘法封闭。若 \(S\) 满足右 Ore 条件，在任意两个非零右理想 \(I,J\) 中分别取 \(a,s\ne0\)，再取 \(t\in S,b\in R\) 使 \(at=sb\)。整环性保证 \(at\ne0\)，所以 \(0\ne at\in I\cap J\)，右正则模一致。

反过来，若 \(R_R\) 一致，对 \(a,s\ne0\)，两个非零右理想 \(aR,sR\) 的交中有非零元素 \(at=sb\)。这一元素非零，保证 \(t\ne0\)，即 \(t\in S\)。若 \(a=0\)，取 \(t=1,b=0\) 即可。因此对全部 \(a\in R,s\in S\) 都满足右 Ore 条件。

由\cref{b2:thm:right-ore-localization}得到单射 \(R\hookrightarrow Q=RS^{-1}\)。每个非零分式 \(as^{-1}\) 的分子也非零，故 \(a\) 与 \(s\) 都在 \(Q\) 中可逆，其逆为 \(sa^{-1}\)。所以 \(Q\) 是除环。
\end{proof}

\begin{proposition}\label{b2:prop:essential-without-semiprime}
设 \(k\) 为域，\(R=T_2(k)\) 为含幺二阶上三角矩阵环，并令
\[
E=\left\{\begin{pmatrix}0&b\\0&c\end{pmatrix}:b,c\in k\right\}.
\]
则 \(E\) 是 \(R_R\) 的本质右理想，却不含双侧正则元素。环 \(R\) 满足右 Goldie 条件，但不半素；它的正则元素恰为单位，经典右分式环仍为 \(R\)，因而不半单。
\end{proposition}
\begin{proof}
右乘上三角矩阵保持第一列为零，所以 \(E\) 是右理想。给定非零右理想 \(I\)，取
\[
0\ne A=\begin{pmatrix}a&b\\0&c\end{pmatrix}\in I.
\]
若 \(a=0\)，则 \(A\in I\cap E\) 已非零；若 \(a\ne0\)，则 \(AE_{12}=aE_{12}\ne0\) 也属于 \(I\cap E\)。因此 \(E\) 本质。但对每个 \(A\in E\) 都有 \(AE_{11}=0\)，而 \(E_{11}\ne0\)，所以 \(E\) 没有正则元素。

环 \(R\) 作为 \(k\) 空间只有三维，故右理想升链必稳定，即 \(R\) 右 Noether；这保证右零化子升链稳定，并由\cref{b2:prop:uniform-dimension-structure}保证一致维数有限，所以 \(R\) 右 Goldie。严格上三角双边理想 \(kE_{12}\) 非零且平方为零，故 \(R\) 不半素。半单环的 Artin--Wedderburn 分解不允许非零平方为零的双边理想，因此 \(R\) 也不半单。

最后，对上面的 \(A\)，若 \(a,c\ne0\)，就有上三角逆，所以 \(A\) 正则。若 \(a=0\)，则 \(AE_{11}=0\)；若 \(c=0\)，则 \(E_{22}A=0\)。因此正则元素恰为单位。对任意 \(r\in R,s\in R^\times\)，取 \(t=1\)、\(b=s^{-1}r\)，就有 \(rt=sb\)，即满足右 Ore 条件。将已有单位求逆不增加元素，故经典右分式环就是 \(R\) 本身。
\end{proof}

% !TeX root = ../../Book2.tex
% 纲要标题：### 19.3 多项式恒等式
% 纲要位置：计划纲要.md，第 1682 行。

\section{多项式恒等式}
\label{b2:sec:1903}

上一节通过右理想和正则元寻找可以使用的分母，本节考察所有代入共同满足的乘法关系。矩阵乘法一般不交换，但固定阶数的矩阵仍满足某些共同的非交换多项式关系。多项式恒等式研究的正是这种“对任意代入都成立”的限制。把待检验的多项式写在自由代数中，就能区分一次偶然代入为零与整个代数的恒等式；外代数和 Cayley–Hamilton 定理将给出矩阵的标准恒等式。

% 纲要要点（供目录核对）：
%
% * PI 代数与标准多项式；求值同态核的交表达所有代入的共同关系，任意特征下的交替性联系外幂消失
% * 矩阵代数的 Amitsur–Levitzki 定理；Newton 生成函数明确整数可逆条件，正特征反例并入原引理；偶次外代数恢复普通行列式，整数泛型恒等式再特化至任意系数环
% * 有限维性与 PI 条件的区别；分别排除 PI 蕴含有限维或交换的误解，自由代数通过可唯一解码的字代入证明非 PI
%

\subsection{PI 代数与标准多项式}
\label{b2:subsec:1903:01}

\begin{definition}\label{b2:def:pi-algebra}
设 \(k\) 为域，\(A\) 为结合含幺 \(k\) 代数，\(m\geq1\)。对每组 \(\mathbf a=(a_1,\ldots,a_m)\in A^m\)，令
\[
\operatorname{ev}_{\mathbf a}:k\langle X_1,\ldots,X_m\rangle\longrightarrow A,
\qquad X_i\longmapsto a_i
\]
为固定 \(k\) 标量的求值同态。若非零多项式
\[
f\in k\langle X_1,\ldots,X_m\rangle
\]
对任意 \(\mathbf a\in A^m\) 都满足 \(\operatorname{ev}_{\mathbf a}(f)=f(a_1,\ldots,a_m)=0\)，即
\[
f\in\bigcap_{\mathbf a\in A^m}\ker\operatorname{ev}_{\mathbf a},
\]
则称 \(f\) 为 \(A\) 的一个\term[duoxiangshi-hengdengshi]{多项式恒等式}[polynomial identity]。若对某个 \(m\geq1\) 存在这样的非零 \(f\)，则称 \(A\) 为\term[PI-daishu]{PI 代数}[polynomial identity algebra]。
\end{definition}

这里讨论的是自由结合代数中的普通多项式恒等式，不把迹、约化迹等运算加入多项式。Cayley–Hamilton 等式中的系数随代入的矩阵而变，也不是上述意义下的一个固定多项式恒等式。

\ProseParagraphBreak
例如交换代数满足 \(X_1X_2-X_2X_1=0\)，而这个多项式在自由代数中非零。相反，“某个元素 \(a\) 满足 \(a^2=0\)”只涉及一次代入，不说明 \(X^2\) 是整个代数的恒等式。非零单位代数甚至不可能让 \(X^2\) 在所有元素上为零，因为可取 \(X=1\)。

\begin{definition}\label{b2:def:standard-polynomial}
设 \(d\geq1\) 为整数。在整数系数自由结合代数 \(\mathbb Z\langle X_1,\ldots,X_d\rangle\) 中令
\[
s_d(X_1,\ldots,X_d)
=\sum_{\sigma\in S_d}\operatorname{sgn}(\sigma)
X_{\sigma(1)}\cdots X_{\sigma(d)}.
\]
称 \(s_d\) 为 \(d\) 次\term[biaozhun-duoxiangshi]{标准多项式}[standard polynomial]。通过 \(\mathbb Z\) 到任意交换系数环的结构映射，也得到该环上的标准多项式。
\end{definition}
\symidx{\(s_d(X_1,\ldots,X_d)\)：标准多项式}

各个排列给出不同的基字，所以 \(s_d\) 在任意域上都是非零多项式。它对每个变量线性，而且交替：若两个输入相同，将每个排列与交换这两个标签后的排列配对，两个乘积相同、符号相反，因而相消。这一论证在特征二也成立，不通过除以二来判断交替性。

\begin{proposition}\label{b2:prop:finite-dimensional-pi}
设 \(k\) 为域，\(A\) 为非零结合含幺 \(k\) 代数。若 \(\dim_kA=d<\infty\)，则 \(A\) 满足标准恒等式 \(s_{d+1}=0\)。
\end{proposition}
\begin{proof}
把 \(s_{d+1}\) 的求值看成底向量空间上的交替 \((d+1)\) 线性映射 \(A^{d+1}\to A\)。由\cref{b2:prop:exterior-alternating-universal}，它通过 \(\bigwedge^{d+1}_kA\) 分解，而\cref{b2:thm:exterior-power-basis}给出 \(\bigwedge^{d+1}_kA=0\)，所以全部求值为零。交替性使重复输入的项消失；高于维数的外幂为零，正把这一计算推广到任意输入。标准多项式本身非零，故它是一个恒等式。
\end{proof}

这个维数界不是矩阵代数的最佳界。\(M_n(k)\) 的向量空间维数为 \(n^2\)，下一小节将把标准恒等式的次数从 \(n^2+1\) 降到 \(2n\)。

\subsection{Amitsur–Levitzki 定理}
\label{b2:subsec:1903:02}

标准多项式把所有排列的矩阵乘积按排列符号相加。外代数恰能把这些符号编码进乘法：在交换含幺 \(\mathbb Q\) 代数 \(B\) 上取矩阵 \(A_1,A_2\in M_n(B)\)，并在自由 \(B\) 模的外代数中取基元 \(e_1,e_2\)。矩阵 \(A_i\) 的条目属于零次部分 \(B\)，与外代数的基元交换，因而
\[
(e_1A_1+e_2A_2)^2
=e_1e_2(A_1A_2-A_2A_1).
\]
因为 \(e_i^2=0\)，重复输入的项消失；因为 \(e_2e_1=-e_1e_2\)，交换输入次序自动产生负号。对 \(2n\) 个输入组成 \(X=\sum_i e_iA_i\) 时，同样的展开使 \(X^{2n}\) 中 \(e_1\cdots e_{2n}\) 的系数恰为 \(s_{2n}(A_1,\ldots,A_{2n})\)。因此证明将先建立 \(X^{2n}=0\)，再读取这个最高外积系数。

\ProseParagraphBreak
要使 \(X^{2n}\) 为零，需要把外代数中的分次符号转为普通矩阵的特征多项式信息。矩阵 \(X^2\) 的条目属于交换的偶次部分，因而可以对它使用行列式与 Cayley–Hamilton 定理；它的各次幂迹将使特征多项式只剩最高次项。下面先在含有有理数的系数环上完成这一计算，最后用整数系数的泛型矩阵去除特征限制。这一证明方法见 Rosset 的《A new proof of the Amitsur-Levitski identity》\cite[pp.~187--188]{Rosset1976AmitsurLevitski}。

\ProseParagraphBreak
幂迹与特征多项式系数之间的关系由 Newton 恒等式表达。这里所需的只是幂迹全部为零的情形，可以用行列式生成函数 \(d(t)=\det(I-tY)\) 直接证明：它的系数给出特征多项式，而它的导数由各次幂迹决定。

\begin{lemma}\label{b2:lem:trace-zero-nilpotence}
设 \(C\) 为交换含幺 \(\mathbb Q\) 代数，\(n\geq1\) 为整数，\(Y\in M_n(C)\)。若 \(\operatorname{tr}(Y^r)=0\) 对每个整数 \(r\geq1\) 成立，则 \(\det(TI_n-Y)=T^n\)，从而 \(Y^n=0\)。另一方面，若 \(k\) 为特征 \(p>0\) 的域，则 \(I_p\in M_p(k)\) 满足全部正次幂的迹为零，却不幂零。
\end{lemma}
\begin{proof}
在 \(C[\![t]\!]\) 中，\(A(t)=I-tY\) 的逆为 \(\sum_{j\ge0}t^jY^j\)。令 \(d(t)=\det A(t)\)。对行列式的置换展开逐项求导，每次恰好微分一个条目，整理该条目的余子式，得到
\[
d'(t)=\operatorname{tr}\Bigl(\operatorname{adj}\bigl(A(t)\bigr)A'(t)\Bigr).
\]
这是多项式求导恒等式，不要求条目是数域中的数。由\cref{b2:prop:adjugate-identity}以及 \(A(t)\) 可逆，
\[
\operatorname{adj}\bigl(A(t)\bigr)=d(t)A(t)^{-1}.
\]
因此
\[
d'(t)=-d(t)\sum_{j\ge0}t^j\operatorname{tr}(Y^{j+1})=0.
\]
写
\[
d(t)=1+c_1t+\cdots+c_nt^n,
\qquad
d'(t)=c_1+2c_2t+\cdots+nc_nt^{n-1}.
\]
因此 \(d'(t)=0\) 给出 \(i c_i=0\)、\(1\le i\le n\)。这里用到 \(C\) 是 \(\mathbb Q\) 代数：每个 \(i\cdot1_C\) 可逆，所以 \(c_i=0\)，得到 \(d(t)=1\)。将 \(\det(TI-Y)\) 的系数与 \(T^n\cdot d(T^{-1})\) 比较，得到特征多项式为 \(T^n\)。对交换环 \(C\) 上的矩阵应用\cref{b2:thm:exterior-cayley-hamilton}，得 \(Y^n=0\)。

在特征 \(p>0\) 的域 \(k\) 上，取 \(Y=I_p\)，则对全部 \(r\ge1\) 都有 \(Y^r=I_p\)、\(\operatorname{tr}(Y^r)=p\cdot1_k=0\)，而 \(Y^r\ne0\)，所以它不幂零。此时
\[
d(t)=\det(I_p-tI_p)=(1-t)^p=1-t^p
\]
不是常数，导数却为零。失败的正是从 \(p c_p=0\) 消去 \(p\) 的一步。
\end{proof}

上述系数比较实际只需要 \(1,\ldots,n\) 在交换系数环中可逆；\(\mathbb Q\) 代数的假设保证了这一条件。

\begin{theorem}[Amitsur–Levitzki]\label{b2:thm:amitsur-levitzki}
设 \(R\) 为交换含幺环，\(n\geq1\) 为整数，则对任意 \(A_1,\ldots,A_{2n}\in M_n(R)\)，有
\[
s_{2n}(A_1,\ldots,A_{2n})=0.
\]
特别地，任意域上的 \(n\) 阶矩阵代数都是 PI 代数。
\end{theorem}
\begin{proof}
先令 \(B\) 为交换含幺 \(\mathbb Q\) 代数，取 \(A_1,\ldots,A_{2n}\in M_n(B)\)。在自由 \(B\) 模的外代数 \(\Lambda_B(e_1,\ldots,e_{2n})\) 上组成矩阵
\[
X=\sum_{i=1}^{2n}e_iA_i.
\]
由\cref{b2:prop:exterior-graded-commutativity}，次数 \(p,q\) 的齐次元素交换时产生符号 \((-1)^{pq}\)。因此对条目分别齐次为次数 \(p,q\) 的矩阵 \(U,V\)，有
\[
\operatorname{tr}(UV)
=\sum_{i,j}u_{ij}v_{ji}
=(-1)^{pq}\sum_{i,j}v_{ji}u_{ij}
=(-1)^{pq}\operatorname{tr}(VU).
\]
取 \(U=X\)、\(V=X^{2r-1}\)，两者次数均为奇数，得到
\(\operatorname{tr}(X^{2r})=-\operatorname{tr}(X^{2r})\)。因为二可逆，全部这些迹为零。

记 \(C=\Lambda_B^{\mathrm{even}}(e_1,\ldots,e_{2n})\)，即全部偶次分量的直和。偶数次数相加仍为偶数，零次单位也属于 \(C\)；两项偶次齐次元素交换时的符号为一，所以 \(C\) 是交换含幺 \(\mathbb Q\) 代数。\(X\) 的条目是一次元素，故 \(Y=X^2\) 的条目是二次元素，\(Y\in M_n(C)\)。这使普通行列式、特征多项式及 Cayley–Hamilton 定理可以用于 \(Y\)，而不是直接用于条目位于整个外代数中的 \(X\)。又因 \(\operatorname{tr}(Y^r)=0\) 对每个 \(r\ge1\) 成立，故可在 \(M_n(C)\) 中应用\cref{b2:lem:trace-zero-nilpotence}，得到 \(Y^n=0\)，即 \(X^{2n}=0\)。展开 \(X^{2n}\) 后，重复的 \(e_i\) 使相应项为零，其余项按排列整理；\(e_1\cdots e_{2n}\) 的矩阵系数恰为
\[
\sum_{\sigma\in S_{2n}}\operatorname{sgn}(\sigma)
A_{\sigma(1)}\cdots A_{\sigma(2n)}.
\]
外幂基定理\cref{b2:thm:exterior-power-basis}保证这个基系数必须为零。因此结论对任意交换 \(\mathbb Q\) 代数成立。

最后取 \(2n\) 个泛型整数矩阵 \(\mathcal A_i=(z_{i,rs})_{r,s=1}^n\)，其中全部条目为独立不定元，系数环记为 \(\mathbb Z[\mathbf z]\)。所得 \(s_{2n}(\mathcal A_1,\ldots,\mathcal A_{2n})\) 的每个矩阵条目都是整数系数多项式 \(F_{rs}(\mathbf z)\)。在单射
\[
\mathbb Z[\mathbf z]\hookrightarrow\mathbb Q[\mathbf z]
\]
下，上一部分证明每个 \(F_{rs}\) 的像是零多项式，因而它原来的全部整数系数就已经为零。此时对给定 \(A_i\in M_n(R)\)，使用环同态
\[
\mathbb Z[\mathbf z]\longrightarrow R,\qquad
z_{i,rs}\longmapsto (A_i)_{rs}
\]
进行特化，零整数多项式仍送到零，便得到任意特征、任意交换环上的结论。证明中除以二和正整数的步骤只发生在核验整数多项式为零的中间环 \(\mathbb Q[\mathbf z]\) 上；最终特化使用的是整数系数环到 \(R\) 的映射。
\end{proof}

这称为\term[Amitsur-Levitzki-dingli]{Amitsur–Levitzki 定理}[Amitsur–Levitzki theorem]。例如二阶矩阵满足四次标准恒等式，即二十四个有序乘积按符号求和为零；它们通常不满足二次交换恒等式。本节不讨论标准恒等式次数的最小性，只使用已经证明的 \(2n\) 次恒等式。

\subsection{有限维性与 PI 条件}
\label{b2:subsec:1903:03}

PI 性对取子代数和取商代数都保持：相同恒等式在子代数的代入只是原代入的一部分；在商代数中，先提升各输入到原代数再取商即可。因此 \(M_n(k)\) 的任意子代数、以及这些子代数的任意商，都继承至少一个非零多项式恒等式。

\ProseParagraphBreak
有限维蕴含 PI，反向不成立：多项式代数 \(k[t]\) 无限维，却满足 \(X_1X_2-X_2X_1=0\)。甚至对任意 \(d\)，\(k[t_1,\ldots,t_d]\) 都满足同一个交换恒等式，而它们的增长可以不同。PI 也不蕴含交换性：当 \(n\ge2\) 时，\(M_n(k)\) 满足刚证的标准恒等式，却有 \(E_{12}E_{21}=E_{11}\ne E_{22}=E_{21}E_{12}\)。PI 条件限制共同的乘法关系，既不强制底向量空间有限维，也不强制任意两个元素交换。

\begin{proposition}\label{b2:prop:free-algebra-not-pi}
设 \(k\) 为域，则自由结合 \(k\) 代数 \(k\langle x,y\rangle\) 不满足任何非零多项式恒等式。
\end{proposition}
\begin{proof}
为了在两个生成元中保存任意候选恒等式的全部字，需要把有限个字母 \(X_i\) 编码成 \(x,y\) 字，并使连接后的字仍能唯一解码。给定任意非零 \(f\in k\langle X_1,\ldots,X_m\rangle\)，取代入
\[
X_i\longmapsto xy^ix,\qquad 1\le i\le m.
\]
每个块 \(xy^ix\) 内部只有两个端点为 \(x\)，相邻块之间出现的 \(xx\) 正好标记分界；\(i\ge1\) 保证块内没有这种分界。因此可在每个 \(xx\) 的两个 \(x\) 之间切开，读出各块中的 \(y\) 个数，恢复原下标串。不同的 \(X_i\) 字因而被送到不同的 \(x,y\) 字，空字也只映到空字。

所以原多项式中不同基字不会在代入后合并，更不会互相抵消；至少一个非零系数仍保留，故 \(f(xyx,xy^2x,\ldots,xy^mx)\ne0\)。这为每个非零候选恒等式给出一次不为零的代入，证明结论。
\end{proof}

同一个自由代数在前面已经给出 Ore 条件的失败例子，这里又显示其不存在多项式恒等式。但两个失败的证明不同：一个比较右理想的共同倍数，另一个利用字编码保存所有多项式信息，不能从其中一项结论直接代替另一项论证。

% !TeX root = ../../Book2.tex
% 纲要标题：### 19.4 代数的增长与 Gelfand–Kirillov 维数
% 纲要位置：计划纲要.md，第 1688 行。

\section{代数的增长与 Gelfand–Kirillov 维数}
\label{b2:sec:1904}

有限生成代数也可能无限维，此时可以逐步数出短乘积产生多少个线性无关方向。生成元改变会改变每一步的具体数值，却只把“允许的乘积长度”扩大一个固定倍数。Gelfand–Kirillov 维数利用这种比较，得到不依赖有限生成组的增长指标。本节固定域 \(k\)，代数均为非零含单位元的有限生成 \(k\) 代数。

关于代数增长与 GK 维数的系统论述，参见 Krause–Lenagan 的《Growth of Algebras and Gelfand-Kirillov Dimension》\cite[Chapters~1--2]{KrauseLenagan2000GrowthAlgebras}。

% 纲要要点（供目录核对）：
%
% * 有限生成代数的增长；一允许补齐较短乘积，生成元换写只使长度扩大固定倍数，区别增长指数与向量空间维数
% * Gelfand–Kirillov 维数的基本例子；常数倍可比的幂次增长、单项式格点计数、矩阵单位比较与基本 Morita 模型，综合题区分分式环、一致维数与增长
% * 不把一般非交换环都视为可按交换方式局部化；局部化先检查有限生成性，完整说明有限分母不能生成有理函数域
% ---
%

\subsection{有限生成代数的增长}
\label{b2:subsec:1904:01}

选取有限维子空间 \(V\subseteq A\)，要求 \(1\in V\)，且 \(V\) 作为代数生成 \(A\)。令
\[
V^0=k1,\qquad
V^n=\operatorname{span}_k\{\,v_1\cdots v_n:v_i\in V\,\},
\qquad
\gamma_V(n)=\dim_kV^n.
\]
因为 \(1\in V\)，较短的乘积可以补上若干个一而写成恰有 \(n\) 项的乘积。因此 \(V^n\) 记录的是长度至多 \(n\) 的全部乘积，且
\[
V^0\subseteq V^1\subseteq\cdots,\qquad
A=\bigcup_{n\ge0}V^n,\qquad V^iV^j\subseteq V^{i+j}.
\]
每一步有限维；这组空间按乘积长度逐步覆盖整个代数。
\symidx{\(\gamma_V(n)=\dim_kV^n\)：给定生成空间的增长函数}

若改用另一组有限生成元，每个旧生成元都能写成新生成元的有限个字的线性组合。旧生成元只有有限个，因而这些字的长度有共同上界 \(c\)；一个长为 \(n\) 的旧字换写后，新字的长度至多为 \(cn\)。下面的命题把这个换写过程表述为生成空间之间的包含关系。

\begin{proposition}\label{b2:prop:growth-comparison}
设 \(k\) 为域，\(A\) 为非零有限生成含幺 \(k\) 代数，\(V,W\subseteq A\) 均为含 \(1\)、有限维且生成 \(A\) 的 \(k\) 子空间。记 \(\gamma_V(n)=\dim_kV^n\)、\(\gamma_W(n)=\dim_kW^n\)，其中 \(V^n,W^n\) 由各自空间内的 \(n\) 项乘积张成，则存在正整数 \(c,d\)，使对所有 \(n\ge1\)，
\[
\gamma_V(n)\le\gamma_W(cn),\qquad
\gamma_W(n)\le\gamma_V(dn).
\]
\end{proposition}
\begin{proof}
取 \(V\) 的有限基，每个基向量都在某个 \(W^{j_i}\) 中。因为 \(W^j\) 递增，取正整数 \(c=\max\{1,j_1,\ldots,j_r\}\)，即可使 \(V\subseteq W^c\)。于是 \(V^n\subseteq W^{cn}\)，得到第一个维数不等式。交换两者得到第二个。
\end{proof}

例如 \(A=k[t]\)，取 \(V=\operatorname{span}_k\{1,t\}\) 时，\(\gamma_V(n)=n+1\)；取 \(W=\operatorname{span}_k\{1,t,t^2\}\) 时，\(\gamma_W(n)=2n+1\)。具体函数不同，但二者都随乘积长度呈一次多项式增长，后者只相当于在前者中允许两倍长度。

\subsection{Gelfand–Kirillov 维数}
\label{b2:subsec:1904:02}

\begin{definition}\label{b2:def:gk-dimension}
设 \(k\) 为域，\(A\) 为非零有限生成含幺 \(k\) 代数。取含 \(1\)、有限维且生成 \(A\) 的 \(k\) 子空间 \(V\)，令 \(V^n\) 为 \(V\) 中 \(n\) 项乘积张成的空间，\(\gamma_V(n)=\dim_kV^n\)。定义 \(A\) 的\term[Gelfand-Kirillov-weishu]{Gelfand–Kirillov 维数}[Gelfand–Kirillov dimension]为
\[
\operatorname{GKdim}_kA
=\limsup_{n\to\infty}\frac{\log\gamma_V(n)}{\log n},
\]
该数值与 \(V\) 的选择无关；取值允许为 \(+\infty\)，简称 GK 维数。
\end{definition}
\symidx{\(\operatorname{GKdim}_kA\)：代数的 Gelfand–Kirillov 维数}

对数的底数可以任取大于一的实数，因为分子、分母在换底时乘上同一个常数。若对某个 \(d\ge0\) 和正数 \(c,C\)，充分大的 \(n\) 满足 \(cn^d\le\gamma_V(n)\le Cn^d\)，则取对数后两个界都趋于 \(d\)，故 GK 维数为 \(d\)。这里要求的是增长函数与 \(n^d\) 在常数倍范围内可比，并不要求增长函数本身是多项式。若增长有下界 \(c\lambda^n\)、\(\lambda>1\)，则同一计算给出无穷 GK 维数。

\begin{proposition}\label{b2:prop:gk-independent}
设 \(k\) 为域，\(A\) 为非零有限生成含幺 \(k\) 代数，\(V\subseteq A\) 为含 \(1\)、有限维且生成 \(A\) 的 \(k\) 子空间，\(\gamma_V(n)=\dim_kV^n\)，则定义 \(\operatorname{GKdim}_kA\) 的上极限与所选 \(V\) 无关。若存在常数 \(C>0\)、\(d\geq0\) 使 \(\gamma_V(n)\leq Cn^d\) 对所有充分大的 \(n\) 成立，则 \(\operatorname{GKdim}_kA\le d\)。
\end{proposition}
\begin{proof}
由\cref{b2:prop:growth-comparison}，对某个固定 \(c\ge1\)，
\[
\frac{\log\gamma_V(n)}{\log n}
\le
\frac{\log\gamma_W(cn)}{\log(cn)}
\cdot\frac{\log(cn)}{\log n}.
\]
第二因子趋于一，第一因子的上极限不超过在全部整数指标上取得的上极限；各项非负，故得到用 \(V\) 定义的数不超过用 \(W\) 定义的数，允许右侧为无穷。交换 \(V,W\) 得到反向不等式。若 \(\gamma_V(n)\le Cn^d\)，取对数后除以 \(\log n\)，常数项 \(\log C/\log n\) 趋于零，便得最后结论。
\end{proof}

\begin{proposition}\label{b2:prop:gk-basic-examples}
设 \(k\) 为域。每个非零有限维含幺 \(k\) 代数的 GK 维数为零；对整数 \(d\geq0\)，多项式代数 \(k[t_1,\ldots,t_d]\) 的 GK 维数为 \(d\)；对整数 \(r\geq2\)，自由结合代数 \(k\langle x_1,\ldots,x_r\rangle\) 的 GK 维数为无穷。
\end{proposition}
\begin{proof}
有限维时 \(1\le\gamma_V(n)\le\dim_kA\)，因此定义中的比值趋于零。多项式环取 \(V=\operatorname{span}_k\{1,t_1,\ldots,t_d\}\)。\(V^n\) 的基是全次数不超过 \(n\) 的单项式，它们与满足 \(a_1+\cdots+a_d\le n\) 的非负整数格点一一对应。添入剩余次数 \(a_0=n-a_1-\cdots-a_d\)，便转为计数 \(a_0+\cdots+a_d=n\) 的非负整数解，得到
\[
\gamma_V(n)=\binom{n+d}{d}.
\]
\(d>0\) 时，这关于 \(n\) 是首项为 \(n^d/d!\) 的实数值多项式，所以对数比值趋于 \(d\)；\(d=0\) 时代数为 \(k\)，回到第一种情形。

对 \(r\ge2\) 个生成元的自由代数，取由一与生成元线性生成的 \(V\)。不同字为基，故
\[
\gamma_V(n)=1+r+\cdots+r^n\ge r^n.
\]
于是 \(\log\gamma_V(n)/\log n\ge n\log r/\log n\rightarrow+\infty\)，得到无穷 GK 维数。
\end{proof}

\begin{proposition}\label{b2:prop:gk-matrix-quotient}
设 \(k\) 为域，\(A\) 为非零有限生成含幺 \(k\) 代数。对整数 \(m\geq1\)，有
\[
\operatorname{GKdim}_k M_m(A)=\operatorname{GKdim}_kA.
\]
任意含幺的有限生成 \(k\) 子代数 \(B\subseteq A\) 及任意非零商 \(k\) 代数 \(C=A/I\) 还满足
\[
\operatorname{GKdim}_kB\leq\operatorname{GKdim}_kA,\qquad
\operatorname{GKdim}_kC\leq\operatorname{GKdim}_kA.
\]
\end{proposition}
\begin{proof}
取 \(W=M_m(V)\)，它有限维、含单位矩阵，并生成矩阵代数。对 \(n\ge1\)，有 \(W^n=M_m(V^n)\)：一方向由矩阵乘法展开；反方向把每个 \(n\) 项乘积 \(v_1\cdots v_n\) 放到任意矩阵单位位置，用中间指标一将它写成 \(n\) 个 \(M_m(V)\) 矩阵的乘积。故 \(\dim_kW^n=m^2\cdot\gamma_V(n)\)。取对数时只多出常数 \(\log m^2\)，不改变上极限。

对商代数，\(V\) 的像生成它，且每次乘积空间的维数不增加，故结论成立。若子代数 \(B\) 由有限维含一空间 \(U\) 生成，取 \(c\) 使 \(U\subseteq V^c\)，则 \(\dim_kU^n\le\gamma_V(cn)\)。使用与\cref{b2:prop:gk-independent}相同的对数比较，得到子代数的不等式。
\end{proof}

\(A\) 与 \(M_m(A)\) 在第十六章给出了 Morita 等价的基本模型，这里计算出该模型中的增长指数相同；以上证明直接使用矩阵单位，并未讨论一般 Morita 等价怎样影响增长。也须区别 GK 维数与向量空间维数：\(k[t]\) 的向量空间维数无穷，GK 维数却为一；矩阵代数使每个乘积空间的维数增大固定倍数，也不改变这个增长指数。

\subsection{局部化可能改变有限生成性}
\label{b2:subsec:1904:03}

本节的 GK 维数以有限生成代数为对象，所以比较局部化前后的增长之前，须先检查局部化是否仍为有限生成代数。只加入有限个分母的逆与逆转所有非零元，在这一点上可能不同。

对 \(k[t,t^{-1}]\)，取 \(V=\operatorname{span}_k\{1,t,t^{-1}\}\)，则 \(V^n\) 以 \(t^{-n},\ldots,t^n\) 为基，增长函数为 \(2n+1\)，故 GK 维数仍为一。这个具体中心局部化保留了有限生成性，并使增长函数容易直接计算。

\ProseParagraphBreak
但把所有非零多项式都变成单位，得到的 \(k(t)\) 不是有限生成 \(k\) 代数。若有限个有理函数生成它，将其分母相乘记为非零多项式 \(d\)，所得子代数包含在 \(k[t,1/d]\) 内。多项式 \(td+1\) 非常数，且任何同时整除它与 \(d\) 的多项式也整除 \((td+1)-td=1\)，所以 \(\gcd(d,td+1)=1\)。取 \(td+1\) 的一个不可约因子 \(p\)，便有 \(p\nmid d\)。若 \(1/p\in k[t,1/d]\)，写成 \(a/d^n\) 后就有 \(d^n=pa\)，与 \(p\nmid d\) 矛盾。因此有限个分母不能产生整个分式域，不能将这里的 GK 维数定义直接用于 \(k(t)\)。

\ProseParagraphBreak
在非交换情形，还须先通过 Ore 条件才能得到本章的单分式环。自由代数的指数增长计算没有自动给出分式构造；半素右 Noether 环之所以能作经典右局部化，依据的是\cref{b2:thm:goldie-sufficiency}及其链条件论证。遇到新的代数时，生成关系、增长、PI 性及分母方程需要分别核对。下一章将给出 Weyl 代数和量子平面的具体基与增长计算，作为这些方法的进一步应用。


\begin{chaptersummary}
右 Ore 条件使分母可以通过方程 \(at=sb\) 整理到右侧。共同分母给出分式的加法，再由作用在这个加法群上的左乘算子建立结合乘法，避免未经验证地沿用交换分数法则。正则分母保证原环不被压缩；若允许零因子，单射性便可能丢失。

\ProseParagraphBreak
Goldie 理论以右理想结构保证足够多的正则元素。有限一致维数限制独立方向，零化子升链条件与半素性排除本质理想的非零左零化子；这些步骤共同推出 Ore 条件，并使分式环没有真本质右理想，最终得到半单 Artin 结构。一致维数对应分式环正则模的长度，而非原代数的基域维数。

\ProseParagraphBreak
PI 条件要求同一个非交换多项式对全部代入为零。Amitsur–Levitzki 恒等式把矩阵阶数变成这样的乘法限制，正特征结论通过整系数特化获得。GK 维数则记录短字产生的线性空间的增长，衡量另一类有限性。一致维数与增长也记录不同数据：例如 \(k[x,y]/(xy)\) 的经典分式环为 \(k(x)\times k(y)\)，一致维数为二，增长指数却为一，其具体计算见\cref{b2:ex:19-crossing-lines}。这些恒等式和增长计算都不能替代分母方程的核验；下一章将计算 Weyl 代数和量子平面的基与增长。
\end{chaptersummary}

\BookExercises

\begin{optionalexercise}\label{b2:ex:19-central-matrix-fractions}
设 \(k\) 为域。在 \(R=M_2\bigl(k[t]\bigr)\) 中，以非零标量矩阵为分母。计算
\[
E_{12}(tI)^{-1}+E_{21}\bigl((t+1)I\bigr)^{-1},\qquad
E_{12}(tI)^{-1}E_{21}\bigl((t+1)I\bigr)^{-1},
\]
并比较反序乘积。
\end{optionalexercise}
\begin{solution}
两个分母在中心，共同分母可取 \(t(t+1)I\)。和为
\[
\bigl((t+1)E_{12}+tE_{21}\bigr)\bigl(t(t+1)I\bigr)^{-1}.
\]
乘积因 \(E_{12}E_{21}=E_{11}\)，为 \(E_{11}\bigl(t(t+1)I\bigr)^{-1}\)；反序则为 \(E_{22}\bigl(t(t+1)I\bigr)^{-1}\)。分母可以交换不表示分子可以交换。由于原环嵌入分式环，两个不同矩阵单位仍不同。
\end{solution}

\begin{optionalexercise}\label{b2:ex:19-normal-element-ore}
设 \(R\) 为非零含幺环，\(s\in R\) 正则，\(\sigma:R\to R\) 为含幺环自同构，满足 \(as=s\cdot\sigma(a)\) 对每个 \(a\in R\) 成立。证明 \(S=\{s^n:n\ge0\}\) 同时满足左右 Ore 条件，并在局部化中计算 \((as^{-i})(bs^{-j})\)，其中 \(a,b\in R\)、\(i,j\ge0\) 为整数。
\end{optionalexercise}
\begin{solution}
归纳得 \(as^n=s^n\sigma^n(a)\)，因此对分子 \(a\)、分母 \(s^n\)，选择新的右分母 \(s^n\) 就满足右 Ore 方程。将 \(a\) 换成 \(\sigma^{-n}(a)\)，又得 \(s^na=\sigma^{-n}(a)s^n\)，给出左条件。正则性在取幂后保持，所以分母集满足构造定理的其余条件。

原关系在局部化中等价于 \(s^{-1}a=\sigma(a)s^{-1}\)。迭代得到 \(s^{-i}b=\sigma^i(b)s^{-i}\)，故
\[
(as^{-i})(bs^{-j})=a\sigma^i(b)s^{-(i+j)}.
\]
若忽略这里的 \(\sigma^i\)，一般会得到错误乘积。
\end{solution}

\begin{optionalexercise}\label{b2:ex:19-one-denominator-failure}
设 \(k\) 为域，\(R=k\langle x,y\rangle\)，并令
\[
B=k\langle x,y,z\rangle/(xz-1,zx-1).
\]
证明 \(B\) 满足把 \(x\) 变成单位的含幺 \(k\) 代数泛性质，且自然映射 \(R\to B\) 单射。再证明 \(B\) 中的 \(zy\) 不能写成 \(ax^{-n}\)、\(a\in R,n\ge0\)。这与\cref{b2:prop:single-element-ore-obstruction}有何关系？
\end{optionalexercise}
\begin{solution}
若含幺同态 \(f:R\to A\) 使 \(f(x)\) 可逆，则必须把 \(z\) 送到 \(f(x)^{-1}\)。这个指定满足两条关系，因而由自由代数取商的泛性质唯一给出 \(B\to A\)。

为核对单射性，把一个字按 \(y\) 的出现位置分成若干个只含 \(x,z\) 的区间。在每个区间内消去相邻的 \(xz\) 或 \(zx\)，最后只剩一个整数次幂 \(x^a\)，其中负次幂用 \(z\) 表示。得到的字唯一写成
\[
x^{a_0}yx^{a_1}y\cdots yx^{a_r},\qquad
r\ge0,\quad a_0,\ldots,a_r\in\mathbb Z.
\]
这些字确实组成商代数的基：以这样的整数列为形式基，乘法由连接两列并将接点的两个指数相加定义，再双线性延拓。三个列连接时，若中间列没有 \(y\)，两种结合次序都只把同一批整数相加；否则两端接点分别相加，因而乘法结合。该代数由 \(x,y,x^{-1}\) 生成并满足给定关系，反向也由消去过程映回 \(B\)。两映射在生成元上互逆，证明所列字线性无关。

原自由代数中的字恰好对应所有指数都非负的上述字，所以 \(R\to B\) 单射。若把 \(a\in R\) 的每个字右乘 \(x^{-n}\)，只会改变最后一个指数；每个 \(y\) 之前的指数仍非负。\(zy=x^{-1}y\) 却在第一个 \(y\) 之前有指数 \(-1\)，故由字基独立性不能得到这种表示。逆转 \(x\) 的泛性质仍可通过加生成元与关系实现，但它不能给出全部元素都是单个右分式的 Ore 构造。
\end{solution}

\begin{optionalexercise}\label{b2:ex:19-zero-divisor-localization}
设 \(R\) 为非零含幺环，\(0\ne e\in R\) 为中心幂等元。证明映射 \(R\to eR\)、\(a\mapsto ea\) 满足使 \(e\) 可逆的含幺环泛性质，并求其核，其中 \(eR\) 的单位元取为 \(e\)。再设 \(k\) 为域，证明不存在从 \(M_2(k)\) 到非零含幺环的含幺同态把 \(E_{11}\) 送成单位。
\end{optionalexercise}
\begin{solution}
中心性保证 \((ea)(eb)=eab\)，所以所给映射为满射含幺环同态。其核是 \((1-e)R\)：若 \(ea=0\)，则 \(a=(1-e)a\)，反向也显然。若含幺同态 \(f:R\to A\) 使 \(f(e)\) 为单位，幂等关系迫使 \(f(e)=1\)，故 \(f((1-e)R)=0\)。因此 \(f\) 唯一通过这个商映射分解；映射自身把 \(e\) 送到 \(eR\) 的单位元。

在矩阵情形，若 \(f(E_{11})\) 可逆，也有 \(f(E_{11})=1\)，于是 \(f(E_{22})=0\)。但
\[
E_{12}E_{22}E_{21}=E_{11},
\]
又迫使 \(f(E_{11})=0\)，与目标非零含幺矛盾。这里 \(E_{11}\) 非中心，不能把矩阵环像直积环那样仅保留一个坐标因子。
\end{solution}

\begin{optionalexercise}\label{b2:ex:19-essential-lattice}
设 \(C\) 为交换整环、\(K\) 为其分式域，\(r\ge1\) 为整数，\(N\subseteq C^r\) 为 \(C\) 子模。证明 \(N\) 本质于 \(C^r\)，当且仅当它在 \(K^r\) 中的 \(K\) 线性生成空间为全体，并求 \(\operatorname{udim}C^r\)。
\end{optionalexercise}
\begin{solution}
若 \(N\) 的 \(K\) 线性生成空间真包含于 \(K^r\)，取不在其中的非零向量，清分母后得到 \(v\in C^r\)。若 \(0\ne av\in N\)、\(a\in C\)，则 \(v=a^{-1}(av)\) 又在该线性生成空间中，矛盾。因此 \(Cv\cap N=0\)，\(N\) 不本质。

反过来，若生成空间为全体，把每个标准向量写成有限个 \(N\) 元素的 \(K\) 线性组合，再取全部系数的共同非零分母 \(d\)，得到 \(dC^r\subseteq N\)。任意非零子模 \(L\subseteq C^r\) 中取 \(0\ne x\)，则 \(dx\ne0\) 且属于 \(L\cap N\)，所以 \(N\) 本质。每个坐标模 \(C\) 一致，整个 \(C^r\) 为 \(r\) 个一致模的直和，故由\cref{b2:prop:uniform-dimension-structure}，一致维数为 \(r\)。
\end{solution}

\begin{optionalexercise}\label{b2:ex:19-uniform-versus-length}
设 \(k\) 为域，\(m,r\ge1\) 为整数，\(R=k[t]/(t^m)\)，\(M=R^r\)，并令 \(S=\bar t^{m-1}M\)。证明 \(S=\{v\in M:v\bar t=0\}\)，且右子模 \(N\subseteq M\) 本质于 \(M\) 当且仅当 \(S\subseteq N\)。据此求本质子模 \(N\) 的一致维数，并用 \(S,M\) 的长度比较说明一致维数没有记录什么信息。
\end{optionalexercise}
\begin{solution}
逐坐标展开截断多项式，乘 \(\bar t\) 为零恰好意味着每个坐标都是 \(\bar t^{m-1}\) 的标量倍，故所给两种描述相同。\(S\) 是 \(r\) 个简单右模的直和，每个 \(k\) 直线也是简单右子模。若 \(N\) 本质，它与每条这样的直线的交非零，因而包含该直线，于是包含 \(S\)。

反过来，若 \(S\subseteq N\)，任取非零子模 \(L\subseteq M\) 及 \(0\ne v\in L\)。在 \(v,v\bar t,\ldots,v\bar t^{m-1}\) 中取最后一个非零项，它被 \(\bar t\) 零化，故属于 \(L\cap S\subseteq L\cap N\)。这证明本质性。

由正文中 \(R_R\) 一致及\cref{b2:prop:uniform-dimension-structure}，\(M\) 的一致维数为 \(r\)。\(S\subseteq N\subseteq M\) 给出 \(\operatorname{udim}N=r\)。\(S\) 的长度为 \(r\)，而 \(M\) 的每一层 \(\bar t^iM/\bar t^{i+1}M\) 都是 \(r\) 个简单模的直和，共有 \(m\) 层，所以 \(\ell(M)=mr\)。当 \(m>1\) 时，两个本质子模有相同的一致维数而长度不同：独立方向内部还可以有不同数目的商层。
\end{solution}

\begin{optionalexercise}\label{b2:ex:19-matrix-annihilators}
设 \(k\) 为域，\(n\ge1\) 为整数。对 \(R=M_n(k)\) 和 \(X\subseteq R\)，用共同核 \(W=\bigcap_{A\in X}\ker A\) 描述 \(\operatorname{ann}_r(X)\)。据此检验右零化子升链条件，并计算 \(R_R\) 的一致维数和经典分式环。
\end{optionalexercise}
\begin{solution}
\(AB=0\) 对所有 \(A\in X\) 成立，当且仅当 \(B\) 的每个列都在 \(W\) 中。所以右零化子恰为这些列矩阵的集合，且可从其中允许的第一列恢复 \(W\)。严格增大的零化子链对应严格增大的子空间链，最多有 \(n\) 次严格增加，故升链条件成立。

正则右模是 \(E_{11}R,\ldots,E_{nn}R\) 的直和；每个非零行模在右矩阵作用下简单，因而一致，所以一致维数为 \(n\)。矩阵正则当且仅当可逆，奇异矩阵有非零核列并由此产生非零右零化矩阵。因此分式构造没有增加新元素，经典分式环就是 \(R\) 本身。
\end{solution}

\begin{optionalexercise}\label{b2:ex:19-product-domains}
设 \(r\ge1\) 为整数，\(C_1,\ldots,C_r\) 为交换整环，求 \(R=\prod_{i=1}^rC_i\) 的正则元素、经典分式环和一致维数。再设 \(k\) 为域，解释无限乘积 \(\prod_{i\ge1}k\) 为何不能由 Goldie 定理推出半单 Artin 分式环。
\end{optionalexercise}
\begin{solution}
一个元素正则当且仅当每个坐标都非零：若某坐标为零，相应坐标幂等元将它零化；若全部非零，逐坐标消去律排除零化元。于是分式环自然为 \(\prod_i\operatorname{Frac}(C_i)\)，有限个坐标的分式可将其分子、分母分别组成一个元素。各坐标理想一致且直和为全体，所以一致维数为 \(r\)。

无限乘积域中，任意多个非零坐标理想独立，所以一致维数无穷。正则元素恰好是所有坐标非零的元素，它们已经可逆，因此经典分式环仍为原无限乘积。令 \(I_n\) 为前 \(n\) 个坐标为零的理想，得到无限严格降链，故该环不是 Artin 环。有限一致维数条件在这里确实缺失。
\end{solution}

\begin{optionalexercise}\label{b2:ex:19-domain-uniform-ore}
设 \(R\) 为非零含幺整环，其非零元素满足右 Ore 条件，\(Q\) 为经典右分式除环。证明 \(R_R\) 本质于 \(Q_R\)，\(Q_R\) 一致，并证明左乘给出环同构 \(Q\cong\operatorname{End}_R(Q_R)\)。这里 \(Q_R\) 表示将右 \(Q\) 模 \(Q\) 限制为右 \(R\) 模。
\end{optionalexercise}
\begin{solution}
非零右 \(R\) 子模 \(L\subseteq Q\) 中取 \(0\ne q=as^{-1}\)。则 \(qs=a\ne0\) 属于 \(L\cap R\)，证明本质性。由\cref{b2:cor:domain-uniform-ore}，\(R_R\) 一致。若 \(L_1,L_2\subseteq Q_R\) 非零，则它们与 \(R\) 的交是非零右理想，故交中还有非零公共元素，所以 \(Q_R\) 也一致。

对任意右 \(R\) 线性映射 \(f:Q\to Q\)，令 \(c=f(1)\)。若 \(q=as^{-1}\)，则
\[
f(q)s=f(qs)=f(a)=ca,
\]
在 \(Q\) 中右乘 \(s^{-1}\) 得 \(f(q)=cq\)。因此每个这样的自同态都是唯一的左乘映射；反向左乘任意 \(c\in Q\) 都右 \(R\) 线性。复合左乘 \(c\) 与左乘 \(d\) 等于左乘 \(cd\)，故得到所称的环同构，无须改成反环。
\end{solution}

\begin{optionalexercise}\label{b2:ex:19-essential-without-regular}
设 \(k\) 为域，\(n\ge2\) 为整数，\(R=T_n(k)\)，\(E\subseteq R\) 为仅最后一列可能非零的矩阵所成的右理想。证明 \(E\) 本质于 \(R_R\) 且不含正则元素，计算 \(E,R_R,R/E\) 的一致维数，并说明一致维数为何不具有合成长度对短正合列的可加性。
\end{optionalexercise}
\begin{solution}
右乘上三角矩阵将 \(E\) 中的矩阵乘以其最后一个对角条目，故 \(E\) 为右理想。非零右理想 \(I\) 中取 \(A\ne0\)，选第 \(j\) 列非零，则 \(AE_{jn}\ne0\)，且仅最后一列非零。因此 \(I\cap E\ne0\)，\(E\) 本质。\(E\) 中每个矩阵都满足 \(AE_{11}=0\)，所以没有正则元素。

\(E=\bigoplus_{i=1}^n kE_{in}\) 是 \(n\) 个简单右模的直和，故一致维数为 \(n\)。它本质于 \(R_R\)，由\cref{b2:prop:uniform-dimension-structure}，\(R_R\) 的一致维数也为 \(n\)。取左上 \((n-1)\) 阶块的满射环同态 \(R\to T_{n-1}(k)\) 的核为 \(E\)，所以 \(R/E\) 的右子模就是 \(T_{n-1}(k)\) 的右理想。重复最后一列的本质子模计算，得到 \(\operatorname{udim}(R/E)=n-1\)，其中 \(n=2\) 时商为域，一致维数为一。因此在
\[
0\longrightarrow E\longrightarrow R\longrightarrow R/E\longrightarrow0
\]
中，中项一致维数为 \(n\)，两端之和却为 \(2n-1\)。本质子模已经占满原模的独立方向，商仍可能非零；这与有限长度的可加性不同。
\end{solution}

\begin{optionalexercise}\label{b2:ex:19-standard-three}
设 \(k\) 为域。在 \(M_2(k)\) 中计算 \(s_3(E_{11},E_{12},E_{21})\)，证明三次标准多项式在任意特征下都不是二阶矩阵的恒等式。四次标准多项式的情形如何？
\end{optionalexercise}
\begin{solution}
按六个排列展开，非零项分别是
\[
E_{11}E_{12}E_{21}=E_{11},\qquad
E_{12}E_{21}E_{11}=E_{11},\qquad
E_{21}E_{11}E_{12}=E_{22},
\]
这三项符号都为正，其余三项为零。因此结果为 \(2E_{11}+E_{22}\)，其右下条目为一，任意特征下都非零。由\cref{b2:thm:amitsur-levitzki}，\(s_4\) 却对所有二阶矩阵代入为零。这里具体检验的是标准多项式，不把结论扩大到任意形式的三次多项式。
\end{solution}

\begin{optionalexercise}\label{b2:ex:19-standard-recursion}
设 \(k\) 为域，\(d\ge2\) 为整数。按排列的首项分组，证明
\[
s_{d+1}(X_1,\ldots,X_{d+1})
=\sum_{i=1}^{d+1}(-1)^{i-1}X_i
s_d(X_1,\ldots,\widehat X_i,\ldots,X_{d+1}),
\]
其中帽号表示删去该变量。由此证明若 \(s_d\) 为恒等式，则所有更高次标准多项式也是恒等式。另计算 \(s_d(1,A_2,\ldots,A_d)\)。
\end{optionalexercise}
\begin{solution}
首项为 \(X_i\) 的排列，先把第 \(i\) 个变量移到最前产生符号 \((-1)^{i-1}\)，剩余变量的排列符号正好组成右端对应的 \(s_d\)。这证明公式，代入后每个求和项都含一次 \(s_d\) 的值，因此它为恒等式时，下一次也为恒等式，再归纳即可。

若把一放入第 \(j\) 个位置，乘积不变，而相对于将它放在首位的符号为 \((-1)^{j-1}\)。所以结果是 \(\sum_{j=1}^d(-1)^{j-1}\) 乘 \(s_{d-1}(A_2,\ldots,A_d)\)：\(d\) 偶数时为零，\(d\) 奇数时为 \(s_{d-1}\)。此处 \(d\ge2\)，所有配对关系在任意特征下都成立。
\end{solution}

\begin{optionalexercise}\label{b2:ex:19-no-one-variable-identity}
设 \(k\) 为无限域，\(A\ne0\) 为含单位元的 \(k\) 代数。证明 \(A\) 没有非零的一元多项式恒等式。解释这为什么不违背有限维矩阵满足 Cayley–Hamilton 定理。
\end{optionalexercise}
\begin{solution}
若 \(f(X)\in k[X]\) 在全部 \(A\) 元素上为零，特别在每个标量 \(\lambda1\) 上为零，则 \(f(\lambda)1=0\)。非零单位代数中的 \(k\) 作用忠实，所以 \(f(\lambda)=0\) 对所有 \(\lambda\in k\) 成立。无限域上的非零一元多项式只有有限多个根，因此 \(f=0\)，矛盾。Cayley–Hamilton 中使用的是每个矩阵自身的特征多项式，其系数随矩阵变化，并非同一个固定的一元多项式对全部矩阵成立。
\end{solution}

\begin{optionalexercise}\label{b2:ex:19-triangular-pi}
设 \(k\) 为域。证明 \(T_2(k)\) 满足恒等式
\[
(X_1X_2-X_2X_1)(X_3X_4-X_4X_3)=0,
\]
而 \(M_2(k)\) 不满足它。
\end{optionalexercise}
\begin{solution}
两个上三角矩阵的交换子在对角线上为零，因为对角条目在域中交换，所以它是严格上三角矩阵。两个严格上三角二阶矩阵的乘积为零，得到恒等式。对全矩阵环，令 \(X_1=X_3=E_{12}\)、\(X_2=X_4=E_{21}\)，两个交换子均为 \(E_{11}-E_{22}\)，其乘积为 \(I\ne0\)。因此子代数可以满足比全矩阵环更多的恒等式。
\end{solution}

\begin{optionalexercise}\label{b2:ex:19-trace-positive-characteristic}
设 \(k\) 为特征 \(p>0\) 的域，\(1\le n<p\)，\(Y\in M_n(k)\)。若 \(\operatorname{tr}(Y^r)=0\) 对 \(1\le r\le n\) 成立，证明 \(Y\) 幂零。解释为什么维数限制不能改为 \(n\le p\)，并在 \(p>2,n=2\) 时说明只检查 \(\operatorname{tr}Y=0\) 不够。
\end{optionalexercise}
\begin{solution}
写 \(d(t)=\det(I-tY)=1+\sum_{j=1}^n c_jt^j\)，\(c_0=1\)。迹零引理证明中的伴随矩阵与形式幂级数计算不需要除法，仍给出
\[
d'(t)=-d(t)\sum_{r\ge1}\operatorname{tr}(Y^r)t^{r-1}.
\]
比较 \(t^{j-1}\) 的系数，得到
\[
jc_j=-\sum_{r=1}^j c_{j-r}\operatorname{tr}(Y^r)=0,
\qquad 1\le j\le n.
\]
因 \(j<p\)，每个 \(j\) 在 \(k\) 中可逆，所以 \(c_j=0\)。于是 \(\chi_Y(T)=T^n\)，Cayley--Hamilton 定理给 \(Y^n=0\)。只需要前 \(n\) 个迹，因为决定 \(d(t)\) 的系数也只有这 \(n\) 个。

当 \(n=p\) 时，正文中的 \(I_p\) 反例说明最后一步不能成立：系数 \(c_p\) 不再由 \(pc_p=0\) 确定。若 \(p>2,n=2\)，矩阵 \(\operatorname{diag}(1,-1)\) 的迹为零但可逆，故只检查一次幂的迹不足以推出幂零。
\end{solution}

\begin{optionalexercise}\label{b2:ex:19-crossing-lines}
设 \(k\) 为域，令 \(R=k[x,y]/(xy)\)。计算它的经典分式环、一致维数和 GK 维数。可以先把它识别为 \(k[x]\times k[y]\) 中两坐标常数项相同的子环。
\end{optionalexercise}
\begin{solution}
商中的元素唯一写成 \(c+x\cdot p(x)+y\cdot q(y)\)，故取值于两坐标轴的映射把它识别为所述子环。一个元素正则当且仅当两个坐标都非零：二者非零时逐坐标用整环消去；若第一坐标为零，则被非零的 \((x,0)\) 零化，若第二坐标为零，则被 \((0,y)\) 零化。

所以分式环嵌入 \(k(x)\times k(y)\)。它也是满的：对 \(h(x)=f(x)/g(x)\)，其中 \(g\ne0\)，分式
\[
\frac{\bigl(x\cdot f(x),0\bigr)}{\bigl(x\cdot g(x),y\bigr)}
\]
的分子、分母都在原环中，分母正则，像为 \(\bigl(h(x),0\bigr)\)。另一坐标相同构造，再相加即可得到任意一对有理函数。因此经典分式环为 \(k(x)\times k(y)\)。

右理想 \((x)\oplus(y)\) 是本质子模：非零元素至少有一个非零坐标，乘 \(x\) 或 \(y\) 后就得到落在此和中的非零元素。两个分量分别同构于整环 \(k[x],k[y]\) 的相应一致模，所以一致维数为二。取生成空间 \(\operatorname{span}_k\{1,x,y\}\)，长度至多 \(n\) 的独立单项式为 \(1,x,\ldots,x^n,y,\ldots,y^n\)，增长为 \(2n+1\)，故 GK 维数为一。一致维数与增长指数在这个例子中不同。
\end{solution}

\begin{optionalexercise}\label{b2:ex:19-partial-laurent-growth}
设 \(k\) 为域。对 \(A=k[x,x^{-1},y]\)，取 \(V=\operatorname{span}_k\{1,x,x^{-1},y\}\)，精确计算 \(\dim_kV^n\)，并求 GK 维数。
\end{optionalexercise}
\begin{solution}
\(V^n\) 的基是 \(x^iy^j\)，其中 \(i\in\mathbb Z,j\ge0,|i|+j\le n\)。它们都可由至多 \(n\) 个指定生成元相乘得到；反向任意这样的乘积整理后也满足此界。不同 Laurent 单项式线性无关，清去一个共同的 \(x\) 幂分母即可归结为多项式单项式的无关性。因此
\[
\dim_kV^n=\sum_{j=0}^n\bigl(2(n-j)+1\bigr)=(n+1)^2,
\]
GK 维数为二。
\end{solution}

\begin{optionalexercise}\label{b2:ex:19-triangular-growth}
设 \(k\) 为域，\(d\ge0\) 为整数，\(C=k[t_1,\ldots,t_d]\)。求上三角矩阵代数 \(T_2(C)\) 的 GK 维数，说明它虽含非零幂零理想，仍可有与 \(C\) 相同的增长指数。
\end{optionalexercise}
\begin{solution}
它由有限个标量矩阵 \(t_iI\) 及矩阵单位 \(E_{11},E_{12},E_{22}\) 生成。作为 \(M_2(C)\) 的有限生成子代数，其 GK 维数不超过 \(d\)；而标量子代数 \(CI\cong C\) 又给出不小于 \(d\) 的下界，均由\cref{b2:prop:gk-matrix-quotient}得到。因此 GK 维数为 \(d\)。严格上三角矩阵理想非零且平方为零，这并不改变上述增长比较，也说明增长指数本身不排除幂零理想。
\end{solution}

\begin{optionalexercise}\label{b2:ex:19-monomial-quotient-growth}
设 \(k\) 为域。分别计算 \(A=k\langle x,y\rangle/(x^2,y^2)\) 与 \(B=k\langle x,y\rangle/(yx)\) 的 GK 维数。证明所使用的字确实组成基，而不只是一组生成元。
\end{optionalexercise}
\begin{solution}
由字关系生成的双边理想，恰是包含相应禁用连续子字的全部字的线性生成空间：左右乘字不会消除该子字，反向每个含子字的字都可分解成左字、关系、右字。因此剩余字在商中线性无关并生成。

在 \(A\) 中，剩余字必须交替出现 \(x,y\)，每个正长度恰有两个字，所以以一、两个生成元作生成空间时，增长为 \(1+2n\)，GK 维数为一。在 \(B\) 中，剩余字恰为 \(x^iy^j\)、\(i,j\ge0\)，因为出现 \(y\) 后不能再出现 \(x\)。长度至多 \(n\) 的数目为 \(\binom{n+2}{2}\)，GK 维数为二。两个代数都有两个生成元，但关系造成了不同增长。
\end{solution}

\begin{optionalexercise}\label{b2:ex:19-finite-denominator-growth}
设 \(k\) 为域，\(0\ne d(t)\in k[t]\)。证明 \(A=k[t,1/d]\) 的 GK 维数为一。须给出关于某个有限生成空间的上下界，而不是从它包含于分式域直接断言维数。
\end{optionalexercise}
\begin{solution}
取 \(V=\operatorname{span}_k\{1,t,1/d\}\)，记 \(m=\deg d\)。\(V^n\) 由 \(t^i/d^j\)、\(i+j\le n\) 线性生成。统一分母为 \(d^n\)，分子次数至多 \((m+1)n\)，所以
\[
\dim_kV^n\le(m+1)n+1.
\]
另一方面 \(1,t,\ldots,t^n\) 都在 \(V^n\) 中且线性无关，故维数至少为 \(n+1\)。取对数比值，两个界都给出极限一，因此 GK 维数为一。有限多个非零多项式作分母也可将它们相乘化为这一情形。
\end{solution}

% !TeX root = ../../Book2.tex
% 纲要标题：## 第20章*　分次代数、滤过与非交换实例
% 纲要位置：计划纲要.md，第 1696 行。

\OptionalChapter{分次代数、滤过与非交换实例}{b2:ch:20}
\BookChapterNumber{b2:ch:20}

\begin{chapterabstract}
本章建立滤过、伴随分次和 Rees 代数，证明若干有限性提升结果，并计算 Weyl 代数与量子平面的正规单项式基。交换子引出 Lie 代数及其包络代数；Ore 扩张与 Diamond 引理分别处理扭曲乘法和重写。最后的例子表明，代数有有限展示，也未必在指定字序下有有限 Gröbner 基。
\end{chapterabstract}

\begin{learninggoals}
\begin{itemize}
\item 构造伴随分次与 Rees 代数，说明穷尽性和滤过下界如何保证降次论证。
\item 在任意特征下证明 Weyl 代数的正规单项式基，比较特征零的单性与正特征的中心。
\item 使用扭曲乘法计算量子平面，辨认参数的阶如何影响中心。
\item 用交换子和矩阵核验 Lie 公理，构造伴随表示，并把 Lie 表示对应为泛包络代数上的模。
\item 从算子模型构造 Ore 扩张，按左模约定解释微分方程。
\item 检查字序、包含及重叠歧义，证明正规式唯一，并区分约化终止与 Gröbner 基的有限性。
\end{itemize}
\end{learninggoals}

在第一 Weyl 代数里，关系 \(\partial x=x\partial+1\) 允许把 \(\partial x^2\) 改写为 \(x^2\partial+2x\)。按字长取 \(\deg x=\deg\partial=1\)，两边主项的次数都是三；特征不为二时，修正项 \(2x\) 的次数是一。在相应的最高次商片中，低次项成为零，于是原来不交换的两个生成元的主符号交换。这里舍去的是低次乘法信息，原代数的哪些性质仍可从这些商片恢复，需要另作证明。

\ProseParagraphBreak

也可以直接规定所有 \(x\) 排在 \(\partial\) 之前，反复使用关系重排。字长和逆序数的下降能保证过程停止，并使每个元素成为 \(x^i\partial^j\) 的线性组合；这还没有证明表达式唯一。若不同重排顺序强迫出额外关系，预期的单项式仍可能线性相关。后面的坏 PBW 关系正会使一个原本写出的生成元变成零。

\ProseParagraphBreak

证明独立性需要比重排更多的信息。Weyl 代数把 \(\partial\) 实现为 \(\partial_X+Y\)，使 \(x^i\partial^j\) 作用于一得到独立单项式 \(X^iY^j\)；量子平面直接在预定基上构造扭曲乘法，Ore 扩张则在自由左模上构造满足关系的算子。一般泛包络代数由 PBW 定理确定伴随分次为对称代数，再把独立性提升回来。Diamond 引理使用另一种检验：给出下降字序，并证明所有包含和重叠歧义的两支最终相同。它们所提供的，都是不能仅从生成元重排中得到的独立性依据。

\ProseParagraphBreak

滤过和包络代数的构造使用第八章的张量代数、对称代数及按关系取商；有限性提升还用到第十二章的 Noether 模概念。

\ProseParagraphBreak
Weyl 代数及其表示可参阅 Etingof 等人的《Introduction to Representation Theory》\cite[\S 2.7, pp. 17--18]{EtingofEtAl2011RepresentationTheory}；该处也提醒正特征中的自然多项式表示不忠实。本章的量子平面不把生成元取逆，与该书同节的量子 Laurent 型例子范围不同。字序、交换单项式序与含分母的局部化适用的计算对象也不同。

% !TeX root = ../../Book2.tex
% 纲要标题：### 20.1 分次与滤过
% 纲要位置：计划纲要.md，第 1698 行。

\section{分次与滤过}
\label{b2:sec:2001}

给多项式记录全次数，可以在乘法中保留最高次部分。非交换关系若还含有较低次项，直接分次可能失效，但次数不超过给定界的子空间仍然有用。这就是滤过的出发点。以下固定域 \(k\)，代数均结合含幺；本节的滤过均递增、穷尽，并在负次数为零。

% 纲要要点（供目录核对）：
%
% * 分次模和分次代数；区分恰为一次与次数至多一，非齐次交换关系自然给出滤过
% * 伴随分次与 Rees 构造；主符号用 symb，明确乘积降次、内部理想与环境求值核、非零参数求值的次数权重和平坦性范围
% * 滤过上的有限性问题；匹配主符号后逐层降次，子模用诱导滤过；整环命题的逆向反例在正文证明，原题改为粗化滤过的 Rees 展示与纤维
%

\subsection{分次模和分次代数}
\label{b2:subsec:2001:01}

\cref{b2:def:associative-graded-algebra}已经定义非负分次代数 \(A=\bigoplus_{n\ge0}A_n\)，其乘法满足 \(A_iA_j\subseteq A_{i+j}\)。一个左 \(A\) 模的分次是直和分解 \(M=\bigoplus_{n\in\mathbb Z}M_n\)，使 \(A_iM_j\subseteq M_{i+j}\)；带有这种分解的模称为\term[fenci-mo]{分次模}[graded module]，保持次数的同态称为分次同态。每个元素都只有有限多个非零齐次分量。

\begin{definition}\label{b2:def:filtered-algebra-module}
设 \(k\) 为域，\(A\) 为含幺结合 \(k\) 代数。若一列 \(k\) 子空间
\[
0=F_{-1}A\subseteq F_0A\subseteq F_1A\subseteq\cdots,
\quad \bigcup_{n\ge0}F_nA=A,
\]
满足 \(1\in F_0A\) 及 \(F_iA\,F_jA\subseteq F_{i+j}A\)，则称 \((F_nA)_{n\geq0}\) 为 \(A\) 的非负\term[lvguo]{滤过}[filtration]；负指标的滤过项约定为零。若 \(M\) 为左 \(A\) 模，\((F_nM)_{n\in\mathbb Z}\) 为递增的 \(k\) 子空间列，满足 \(\bigcup_nF_nM=M\)、\(F_nM=0\) 对 \(n<0\)，以及 \(F_iA\,F_jM\subseteq F_{i+j}M\)，则称它为与 \(A\) 的滤过相容的模滤过。
\end{definition}

对有统一整数下界的模滤过，平移模的指标便可使用本节的论证。这里的穷尽性保证每个元素最终出现，下界则保证降次归纳能到达零；两项条件各有用途。

\ProseParagraphBreak
分次给出滤过 \(F_nA=\bigoplus_{i\le n}A_i\)。分次中的 \(A_n\) 记录恰为 \(n\) 次的齐次分量，滤过中的 \(F_nA\) 则记录次数不超过 \(n\) 的全部元素，因此滤过项彼此包含，而不是彼此直和。反过来，滤过不一定来自这样一个与乘法相容的直和。例如关系 \(\mathrm{d}x-x\mathrm{d}=1\) 混合二次与零次，但按字长不超过 \(n\) 所生成的子空间仍构成滤过。后面将严格构造这个例子。

\ProseParagraphBreak
若 \(M\) 由 \(m_1,\ldots,m_r\) 生成，并给每个生成元指定非负次数 \(d_i\)，可以取
\[
F_nM=\sum_i(F_{n-d_i}A)m_i.
\]
这确实相容且穷尽。这样的表达把有限生成元及其所需次数一并记录下来，而不仅记录一个抽象生成集合。

\subsection{伴随分次与 Rees 构造}
\label{b2:subsec:2001:02}

\begin{definition}\label{b2:def:associated-graded-rees}
设 \(k\) 为域，\(A\) 为带有非负穷尽滤过 \((F_nA)\) 的含幺结合 \(k\) 代数。定义 \(A\) 的\term[bansui-fenci-daishu]{伴随分次代数}[associated graded algebra]为
\[
\operatorname{gr}_F A=\bigoplus_{n\ge0}F_nA/F_{n-1}A,
\]
乘法由 \((a+F_{i-1}A)(b+F_{j-1}A)=ab+F_{i+j-1}A\) 给出。若左 \(A\) 模 \(M\) 带有相容滤过 \((F_nM)\)，则定义 \(\operatorname{gr}_F M=\bigoplus_{n\geq0}F_nM/F_{n-1}M\)。若非零 \(a\in A\) 首次属于 \(F_nA\)，则称 \(n\) 为其滤过次数，称 \(\operatorname{symb}(a)=a+F_{n-1}A\) 为其\term[zhufuhao]{主符号}[principal symbol]。
\end{definition}

改变代表元所产生的乘积误差落在 \(F_{i+j-1}A\)，故上述乘法良定义；结合律由原代数继承，模作用的验证相同。若非零 \(a,b\) 的滤过次数分别为 \(i,j\)，则
\[
\operatorname{symb}(a)\operatorname{symb}(b)
=ab+F_{i+j-1}A\in F_{i+j}A/F_{i+j-1}A.
\]
当 \(ab\ne0\) 且次数为 \(i+j\) 时，这才等于 \(\operatorname{symb}(ab)\)；若 \(ab=0\)，或其实际次数小于 \(i+j\)，上述乘积为零。因此原代数没有零因子，也不能保证主符号相乘不为零。

\ProseParagraphBreak
令 \(z\) 为与 \(A\) 中全部元素交换的不定元。在中心多项式环 \(A[z]\) 中定义\term[Rees-daishu]{Rees 代数}[Rees algebra]
\[
\mathcal R_F(A)=\bigoplus_{n\ge0}(F_nA)z^n.
\]
滤过的乘法条件保证它为子代数。由于 \(1\in F_0A\subseteq F_1A\)，\(z\) 也属于其中且居中。
\symidx{\(\operatorname{gr}_F A\)、\(\mathcal R_F(A)\)：伴随分次代数与 Rees 代数}

求零纤维时，\((z)\) 指 Rees 代数内部的理想 \(z\mathcal R_F(A)\)。它的第 \(n\) 次分量是 \((F_{n-1}A)z^n\)：要在这个子代数内部提出一个 \(z\)，剩下的系数必须已属于较低的滤过项。环境多项式环中的求值 \(z=0\) 限制到 Rees 代数后，核却是 \(\mathcal R_F(A)\cap zA[z]\)，包含全部正次项。总有
\[
z\mathcal R_F(A)\subseteq\mathcal R_F(A)\cap zA[z],
\]
但这个包含可以严格，因而不能用环境多项式环中的求值代替内部取商。

\ProseParagraphBreak
例如 \(A=k[x]\) 配通常次数滤过，则
\[
\mathcal R_F(A)=k[z,xz]\cong k[z,w],\qquad w\longmapsto xz.
\]
这里 \(z^iw^j\) 的像为 \(x^jz^{i+j}\)，不同的 \((i,j)\) 给出不同单项式，所以所写同态确为同构。元素 \(w=xz\) 不属于 Rees 代数内部的 \((z)\)：若 \(xz=zq\) 且 \(q\in\mathcal R_F(A)\)，在 \(k[x,z]\) 中消去 \(z\) 就有 \(q=x\)，但 \(x\notin\mathcal R_F(A)\)，因为其常数次系数不在 \(F_0A=k\) 中。因此模掉 \(z\) 后，\(w\) 仍保留下来，所得商为 \(k[w]\)，正好记录 \(x\) 的最高次信息。下面分别计算参数为零和非零时的纤维；平坦性另用到第七章的\cref{b2:thm:pid-flat-torsionfree}。

\begin{theorem}\label{b2:thm:rees-fibres}
设 \(k\) 为域，\(A\) 为带有非负穷尽滤过 \((F_nA)\) 的含幺结合 \(k\) 代数，令 \(z\) 为与 \(A\) 交换的不定元，\(\mathcal R_F(A)=\bigoplus_{n\geq0}(F_nA)z^n\subseteq A[z]\) 为 Rees 代数，则有自然 \(k\) 代数同构
\[
\mathcal R_F(A)/(z)\cong\operatorname{gr}_F A,
\qquad\mathcal R_F(A)/(z-\lambda)\cong A\quad(\lambda\in k^\times).
\]
此外，\(\mathcal R_F(A)\) 作为 \(k[z]\) 模平坦。
\end{theorem}
\begin{proof}
将 \(a_nz^n\) 送到 \(a_n+F_{n-1}A\)，得到到伴随分次的满同态。其核的第 \(n\) 次分量为 \((F_{n-1}A)z^n\)，正是 \(z\mathcal R_F(A)\)，所以得到第一个同构。

在 \(z=1\) 求值得到到 \(A\) 的同态；穷尽性保证它满射。若 \(p(z)=\sum_{n=0}^N a_nz^n\) 的像为零，则 \(\sum_n a_n=0\)。令
\[
q(z)=-\sum_{n=0}^{N-1}\left(\sum_{i=0}^n a_i\right)z^n.
\]
第 \(n\) 个系数属于 \(F_nA\)，所以 \(q\in\mathcal R_F(A)\)，而逐项相减给出 \(p=(z-1)q\)。反向显然，故求值同态的核恰为 \((z-1)\)。

对任意 \(\lambda\in k^\times\)，逐次乘以 \(\lambda^n\) 给出 \(\mathcal R_F(A)\) 的自同构
\[
\theta_\lambda\colon\sum_n a_nz^n\longmapsto\sum_n\lambda^na_nz^n.
\]
它与 \(z=1\) 处的求值复合，得到 \(z=\lambda\) 处的求值 \(\sum_n a_nz^n\mapsto\sum_n\lambda^na_n\)。因此该同态满射，且其核为
\[
\theta_\lambda^{-1}\bigl((z-1)\mathcal R_F(A)\bigr)
= (z-\lambda)\mathcal R_F(A),
\]
得到每个非零纤维的同构。

\(A[z]\) 作为 \(k[z]\) 模无挠：一个非零标量多项式与非零向量系数多项式相乘，其最高次系数仍非零。子模 \(\mathcal R_F(A)\) 因而无挠。\(k[z]\) 是主理想整环，由\cref{b2:thm:pid-flat-torsionfree}得到平坦性，无需假设 Rees 代数有限生成。
\end{proof}

平坦性是 \(k[z]\) 模的性质，并不保证各个纤维作为代数同构，也不保证其乘法相同。这里的非零纤维都恢复 \(A\)，零纤维却只保留相应滤过次数上的乘积；\cref{b2:prop:filtered-domain-converse-failure}将给出原代数没有零因子而零纤维含非零幂零元的例子。

\subsection{滤过上的有限性问题}
\label{b2:subsec:2001:03}

\begin{theorem}\label{b2:thm:filtered-generator-lifting}
设 \(k\) 为域，\(A\) 为带有非负穷尽滤过 \((F_nA)\) 的含幺结合 \(k\) 代数，\(M\) 为带有相容、穷尽且负次数为零的滤过 \((F_nM)\) 的左 \(A\) 模。若 \(\operatorname{gr}_F M\) 由有限个次数为 \(d_i\) 的齐次元素 \(\bar m_i\) 作为 \(\operatorname{gr}_F A\) 模生成，则任意代表 \(m_i\in F_{d_i}M\) 生成 \(M\)，并有
\[
F_nM=\sum_i(F_{n-d_i}A)m_i.
\]
\end{theorem}
\begin{proof}
右侧包含于左侧来自滤过相容性。反向要把分次中的表达提升到原模：先匹配最高的一层，再相减，使剩余部分降到较低层。对 \(n\) 归纳，给定 \(m\in F_nM\)，其在 \(F_nM/F_{n-1}M\) 中的类可写为 \(\sum_i\bar a_i\bar m_i\)，其中取齐次分量后可令 \(\deg\bar a_i=n-d_i\)，负次数项为零。选择 \(a_i\in F_{n-d_i}A\) 提升这些系数，便有
\[
m-\sum_i a_im_i\in F_{n-1}M.
\]
对差值用归纳假设；所得较低次系数仍包含在 \(F_{n-d_i}A\) 中。\(n=0\) 时差值在 \(F_{-1}M=0\)，归纳开始。穷尽性再说明这些元素生成全部 \(M\)。
\end{proof}

\begin{theorem}\label{b2:thm:filtered-noetherian-transfer}
设 \(k\) 为域，\(A\) 为带有非负穷尽滤过 \((F_nA)\) 的含幺结合 \(k\) 代数，\(M\) 为带有相容、穷尽且负次数为零的滤过 \((F_nM)\) 的左 \(A\) 模。若 \(\operatorname{gr}_F A\) 是左 Noether 环，则 \(A\) 也是左 Noether 环。若 \(\operatorname{gr}_F M\) 是 Noether 左 \(\operatorname{gr}_F A\) 模，则 \(M\) 是 Noether 左 \(A\) 模。将左模及左 Noether 条件同时换为右侧条件，结论也成立。
\end{theorem}
\begin{proof}
对任意子模 \(N\le M\)，取诱导滤过 \(F_nN=N\cap F_nM\)。这样，\(F_nN\) 映入 \(F_nM/F_{n-1}M\) 的核恰为 \(N\cap F_{n-1}M=F_{n-1}N\)，所以每个商片单射，\(\operatorname{gr}_F N\) 才能成为 \(\operatorname{gr}_F M\) 的齐次子模；任意另选的滤过未必有这一性质。由 Noether 假设它有限生成。可以取齐次生成元：将一个有限生成组各自拆成齐次分量，所得有限族仍在该齐次子模内并生成它。由\cref{b2:thm:filtered-generator-lifting}，\(N\) 有限生成。每个子模都有限生成，故 \(M\) Noether。取 \(M=A\) 为左正则模即得到环的断言。对于右模，所有系数写在生成元右侧，同一个降次归纳与诱导滤过论证成立。
\end{proof}

\begin{proposition}\label{b2:prop:filtered-domain}
设 \(k\) 为域，\(A\) 为带有非负穷尽滤过 \((F_nA)\) 的含幺结合 \(k\) 代数。若伴随分次代数 \(\operatorname{gr}_F A\) 没有零因子，则 \(A\) 没有零因子，且任意非零 \(a,b\in A\) 的乘积滤过次数等于两者次数之和。
\end{proposition}
\begin{proof}
非零 \(a,b\) 都有最低的滤过次数，主符号非零。它们的符号乘积不为零，故 \(ab\) 的相应商类不为零。这同时说明 \(ab\ne0\) 及乘积不降次。
\end{proof}

穷尽性和下界不能略去。若在非零模上取全部 \(F_nM=0\)，则伴随分次为零，却看不到原模；这违反穷尽性。若允许所有整数指标并令 \(F_nM=M\)，则滤过虽穷尽，伴随分次仍为零，因为没有开始降次归纳的下界。本节后面使用的字长滤过都满足已经列出的条件。

\begin{proposition}\label{b2:prop:filtered-domain-converse-failure}
设 \(k\) 为任意域，给含幺 \(k\) 代数 \(A=k[x]\) 赋予滤过
\[
F_nA=\{f\in k[x]:\deg f\le 2n\}\quad(n\ge0),
\qquad F_nA=0\quad(n<0),
\]
其中零多项式属于每个非负滤过项。则 \(A\) 是整环，而
\[
\operatorname{gr}_F A\cong k[u,v]/(u^2),
\qquad \deg u=\deg v=1.
\]
因此\cref{b2:prop:filtered-domain}的逆命题不成立。
\end{proposition}
\begin{proof}
多项式次数的乘法公式保证 \(F_iA\,F_jA\subseteq F_{i+j}A\)，而这个滤过显然穷尽且 \(1\in F_0A\)。令 \(u=\operatorname{symb}(x)\)、\(v=\operatorname{symb}(x^2)\)。两者都在一次分量中；但 \(x^2\in F_1A\)，所以把 \(u\) 与自身相乘，所得在 \(F_2A/F_1A\) 中的类为零。于是有分次代数同态 \(k[u,v]/(u^2)\to\operatorname{gr}_F A\)。

零次分量以一为基。对每个 \(n\ge1\)，\(F_nA/F_{n-1}A\) 的一组基是 \(x^{2n-1}\)、\(x^{2n}\) 的类，它们分别是 \(uv^{n-1}\)、\(v^n\) 的像。另一方面，\(k[u,v]/(u^2)\) 的第 \(n\) 次分量恰以 \(uv^{n-1},v^n\) 为基，所以此同态逐次为同构。特别地，\(u\ne0\) 而 \(u^2=0\)，尽管原代数 \(k[x]\) 没有零因子。
\end{proof}

% !TeX root = ../../Book2.tex
% 纲要标题：### 20.2 Weyl 代数
% 纲要位置：计划纲要.md，第 1704 行。

\section{Weyl 代数}
\label{b2:sec:2002}

在多项式上，先乘不定元再求导，与先求导再乘不定元，恰好相差恒等算子。Weyl 代数只保留这条交换关系。我们先证明抽象代数中的正规单项式确实独立，再判断它对通常多项式空间的作用是否忠实；这两项问题在正特征中有不同答案。

% 纲要要点（供目录核对）：
%
% * 微分与乘法的交换关系；区分任意特征下的抽象 Weyl 代数、通常多项式作用与辅助变量忠实作用
% * 正规单项式及作用；下降重排给生成性，辅助算子在一上的取值记录两个指数而给独立性
% * 特征零与正特征的差别；中心秩 p²、有限维表示迹障碍和中心零纤维的矩阵商在正文完整证明；原题改为商模分解、平移纤维及中心分式除代数
%

\subsection{微分与乘法的交换关系}
\label{b2:subsec:2002:01}

在 \(k[X]\) 上令 \(X\) 表示乘以不定元的算子，令 \(D_0\) 为形式求导。Leibniz 公式\footnote{莱布尼茨（Gottfried Wilhelm Leibniz，1646-1716），德国数学家及哲学家，创立微积分，并研究符号演算。}给出 \(D_0X-XD_0=1\)，这在任意特征下都成立。

\begin{definition}\label{b2:def:first-weyl-algebra}
设 \(k\) 为域。将自由结合代数 \(k\langle x,\mathrm{d}\rangle\) 对关系 \(\mathrm{d}\cdot x-x\cdot\mathrm{d}=1\) 生成的双边理想取商，所得代数称为 \(k\) 上的第一\term[Weyl-daishu]{Weyl 代数}[first Weyl algebra]，记为
\[
A_1(k)=k\langle x,\mathrm{d}\rangle/(\mathrm{d}\cdot x-x\cdot\mathrm{d}-1).
\]
其代数同态由生成元的像决定，并要求这些像满足 \(\mathrm{d}\cdot x-x\cdot\mathrm{d}=1\)。
\end{definition}
\symidx{\(A_1(k)\)：第一 Weyl 代数}

由\cref{b2:prop:algebra-relations-universal}，\(x\mapsto X,\mathrm{d}\mapsto D_0\) 给出 \(A_1(k)\) 在 \(k[X]\) 上的作用。但关系成立只证明作用存在，尚未证明这个作用单射。

\subsection{正规单项式及其作用}
\label{b2:subsec:2002:02}

\begin{theorem}\label{b2:thm:weyl-normal-basis}
设 \(k\) 为域，\(A_1(k)=k\langle x,\mathrm{d}\rangle/(\mathrm{d}\cdot x-x\cdot\mathrm{d}-1)\)，则全部 \(x^i\mathrm{d}^j\)、\(i,j\ge0\)，构成 \(A_1(k)\) 的一组 \(k\) 基。因此每个元素唯一写为有限和 \(\sum_j f_j(x)\mathrm{d}^j\)，其中 \(f_j\in k[x]\)。
\end{theorem}
\begin{proof}
生成性来自 \(\mathrm{d}\cdot x=x\cdot\mathrm{d}+1\)。对一个字，按字长及其中位于某个 \(x\) 左边的 \(\mathrm{d}\) 的配对数作字典式归纳：将相邻 \(\mathrm{d}\cdot x\) 换成 \(x\cdot\mathrm{d}+1\)，第一项的字长不变、错序数减少一，第二项的字长减少二。这两个量都是非负整数，故不能无限下降，最终只剩全部 \(x\) 在前、全部 \(\mathrm{d}\) 在后的正规字。不过归约终止这里只保证正规字生成，尚不能排除它们之间仍有线性关系。

为证独立性，我们要找到一个作用，使 \(x^i\mathrm{d}^j\) 的像保留两个指数的信息。通常的求导算子在常数一上反复作用都为零，不能用一来检测 \(j\)。因此在两个交换变量的多项式空间 \(k[X,Y]\) 上令 \(X\) 仍表示乘法，并令
\[
D=\frac{\partial}{\partial X}+Y,
\]
其中 \(Y\) 也表示乘法算子。增加的乘 \(Y\) 算子与乘 \(X\) 交换，故仍有 \(DX-XD=1\)，给出一个表示。同时 \(\partial(Y^j)/\partial X=0\)，所以 \(D(Y^j)=Y^{j+1}\)，从常数一出发便有 \(D^j(1)=Y^j\)，因而
\[
X^iD^j(1)=X^iY^j.
\]
若抽象代数中有有限关系 \(\sum_{i,j}c_{ij}x^i\mathrm{d}^j=0\)，令该算子作用在一上，得到 \(\sum_{i,j}c_{ij}X^iY^j=0\)。交换多项式的单项式独立，故所有 \(c_{ij}=0\)。整个论证没有使用阶乘可逆性。
\end{proof}

这个双变量表示因而在任意特征下都忠实。在特征 \(p\) 中，虽然通常的第 \(p\) 次求导在 \(k[X]\) 上为零，这里的 \(D^p(1)=Y^p\) 却不为零。它证明的是抽象 Weyl 代数中正规字的独立性，并不能据此断言通常的多项式作用也忠实；后者将在\cref{b2:prop:weyl-polynomial-action}中单独判断。后面的\cref{b2:thm:diamond-lemma}还会把改写的下降条件与歧义检查结合起来，证明正规式唯一。Weyl 代数的正规单项式基还可参阅 Etingof 等人\cite[Proposition 2.7.1(i), pp. 17--18]{EtingofEtAl2011RepresentationTheory}；那里使用另一种带辅助变量的表示。

\ProseParagraphBreak
用字长滤过，\(F_nA_1(k)\) 恰由 \(i+j\le n\) 的正规单项式生成。交换关系中的修正项降低次数二，故最高次部分变得交换。

\begin{proposition}\label{b2:prop:weyl-associated-graded}
设 \(k\) 为域，\(A_1(k)=k\langle x,\mathrm{d}\rangle/(\mathrm{d}\cdot x-x\cdot\mathrm{d}-1)\)，按 \(x,\mathrm{d}\) 的字长赋予滤过 \(F_nA_1(k)\)，则 \(\operatorname{gr}_F A_1(k)\cong k[X,D]\)。因此 \(A_1(k)\) 没有零因子，且左右 Noether。
\end{proposition}
\begin{proof}
将交换变量 \(X,D\) 分别送到一次主符号 \(\operatorname{symb}(x),\operatorname{symb}(\mathrm{d})\)。它们交换且生成伴随分次，故得到满同态。\cref{b2:thm:weyl-normal-basis}说明第 \(n\) 个商片恰以 \(i+j=n\) 的单项式符号为基，因此同态逐次单射。

交换多项式环 \(k[X,D]\) 是整环；由第一卷定理12.4.10（Hilbert 基定理），它还为 Noether 环。应用\cref{b2:prop:filtered-domain}及\cref{b2:thm:filtered-noetherian-transfer}，得到无零因子性和两侧 Noether 性。
\end{proof}

例如 \(\mathrm{d}^2x^3=x^3\mathrm{d}^2+6x^2\mathrm{d}+6x\)，由两次使用 \(\mathrm{d}\cdot x^n=x^n\mathrm{d}+nx^{n-1}\) 得到。等式在正特征中仍成立，只是整数系数按该特征取值。

\subsection{特征零与正特征}
\label{b2:subsec:2002:03}

为避免与群交换子混淆，本节的 \([a,b]=ab-ba\) 表示代数中的加法交换子。直接归纳给出
\begin{equation}\label{b2:eq:weyl-commutators}
[\mathrm{d},x^i\mathrm{d}^j]=i x^{i-1}\mathrm{d}^j,\qquad
[x,x^i\mathrm{d}^j]=-j x^i\mathrm{d}^{j-1}.
\end{equation}
零指数时右侧解释为零。

\begin{theorem}\label{b2:thm:weyl-center-simplicity}
设 \(k\) 为域，\(A_1(k)=k\langle x,\mathrm{d}\rangle/(\mathrm{d}\cdot x-x\cdot\mathrm{d}-1)\)。若 \(\operatorname{char}k=0\)，则 \(Z\bigl(A_1(k)\bigr)=k\)，且 \(A_1(k)\) 是单环。若 \(\operatorname{char}k=p>0\)，则
\[
Z\bigl(A_1(k)\bigr)=k[x^p,\mathrm{d}^p],
\]
这是二元多项式环；\(A_1(k)\) 是其中心上的自由模，一组基为
\[
\{x^i\mathrm{d}^j:0\le i,j<p\},
\]
故在中心上秩为 \(p^2\)。此时 \(A_1(k)\) 不是单环。
\end{theorem}
\begin{proof}
由\eqref{b2:eq:weyl-commutators}及正规单项式基，元素 \(\sum c_{ij}x^i\mathrm{d}^j\) 同时与 \(x,\mathrm{d}\) 交换，当且仅当每个非零系数满足 \(i c_{ij}=j c_{ij}=0\)。在特征零中只剩 \(i=j=0\)；在特征 \(p\) 中要求 \(p\mid i,j\)。这证明两个中心公式，且正规单项式基保证 \(x^p,\mathrm{d}^p\) 之间没有多项式关系。

在特征零中，要证明非零双边理想 \(I\) 等于整个代数，只需在其中得到非零标量。双边性保证与 \(x\) 或 \(\mathrm{d}\) 取交换子后仍留在 \(I\)，而\eqref{b2:eq:weyl-commutators}表明这分别降低 \(\mathrm{d}\) 或 \(x\) 的指数。选其中非零元素 \(P=\sum_{j=0}^m f_j(x)\mathrm{d}^j\)，使 \(\mathrm{d}\) 次数 \(m\) 最小。若 \(m>0\)，则 \([P,x]\in I\) 的 \(\mathrm{d}\) 次数为 \(m-1\)，最高系数为 \(m\cdot f_m(x)\ne0\)，矛盾。因此 \(I\) 含非零 \(f(x)=a_rx^r+\cdots+a_0\)，其中 \(a_r\ne0\)。记 \(\operatorname{ad}_{\mathrm{d}}(u)=[\mathrm{d},u]\)，则
\[
\operatorname{ad}_{\mathrm{d}}^{\,r}(f)=r!a_r\in k^\times.
\]
这个交换子仍属于 \(I\)，故 \(1\in I\)，\(I=A_1(k)\)。这里 \(r!\ne0\) 正是特征零假设的作用。

在特征 \(p\) 中，将每个指数唯一写为 \(i=pa+r\)、\(j=pb+s\)，其中 \(0\le r,s<p\)。利用 \(x^p,\mathrm{d}^p\) 居中，便有
\[
x^i\mathrm{d}^j=(x^p)^a(\mathrm{d}^p)^b x^r\mathrm{d}^s,
\]
所以这些余数指数的单项式在中心上生成 \(A_1(k)\)。若它们之间有中心系数的线性关系，将各系数展开为 \(x^p,\mathrm{d}^p\) 的多项式，不同商和余数给出不同的正规单项式。\cref{b2:thm:weyl-normal-basis}迫使所有系数为零，证明中心上的自由性及秩公式。中心元 \(x^p\) 生成的双边理想就是 \(x^pA_1(k)\)，按正规单项式基它由 \(x\) 指数至少为 \(p\) 的单项式生成。因此它非零且不含一，是一个真双边理想。
\end{proof}

\begin{proposition}\label{b2:prop:weyl-polynomial-action}
设 \(k\) 为域，\(A_1(k)=k\langle x,\mathrm{d}\rangle/(\mathrm{d}\cdot x-x\cdot\mathrm{d}-1)\)。令 \(x\) 在 \(k[X]\) 上作用为乘 \(X\)，\(\mathrm{d}\) 作用为形式求导，则这个 \(A_1(k)\) 模作用在 \(\operatorname{char}k=0\) 时忠实，在 \(\operatorname{char}k>0\) 时不忠实。
\end{proposition}
\begin{proof}
特征零时，若非零 \(P=\sum_{j=0}^m f_j(x)\mathrm{d}^j\) 作用为零，记 \(M_X\) 为乘 \(X\) 的算子。由\eqref{b2:eq:weyl-commutators}，\([f_j(X)\mathrm{d}^j,M_X]=j f_j(X)\mathrm{d}^{j-1}\)。因此从 \(P\) 开始反复按 \([P,M_X]\) 的顺序取 \(m\) 次交换子，仍作用为零；低于 \(m\) 次的项全部消失，所得算子是乘 \(m!f_m(X)\)，令其作用在一上便得到矛盾。\(m=0\) 时直接如此。

特征 \(p\) 时，第 \(p\) 次形式求导在每个 \(X^n\) 上都为零，因为连续 \(p\) 个整数的乘积含因子 \(p\)。所以非零的抽象元素 \(\mathrm{d}^p\) 作用为零，非零性则由\cref{b2:thm:weyl-normal-basis}保证。
\end{proof}

\begin{proposition}\label{b2:prop:weyl-finite-dimensional-representations}
设 \(k\) 为任意域，\(A=A_1(k)=k\langle x,\mathrm{d}\rangle/(\mathrm{d}\cdot x-x\cdot\mathrm{d}-1)\)。若 \(\operatorname{char}k=0\)，则每个有限维幺性左 \(A\) 模都是零模。若 \(\operatorname{char}k=p>0\)，则每个有限维幺性左 \(A\) 模的 \(k\) 维数都被 \(p\) 整除；此时
\[
V_p=k[X]/(X^p)
\]
在 \(x\) 作用为乘 \(X\)、\(\mathrm{d}\) 作用为形式求导时，是一个 \(p\) 维简单左 \(A\) 模，该作用诱导 \(k\) 代数同构
\[
A/(x^p,\mathrm{d}^p)\cong\operatorname{End}_k(V_p)\cong M_p(k),
\]
其中 \((x^p,\mathrm{d}^p)\) 是这两个中心元生成的双边理想。
\end{proposition}
\begin{proof}
在一个有限维幺性左 \(A\) 模 \(V\) 上，记 \(x,\mathrm{d}\) 的作用为 \(X,D\)，则 \(DX-XD=1_V\)。矩阵迹满足 \(\operatorname{Tr}(DX)=\operatorname{Tr}(XD)\)，故
\[
0=\operatorname{Tr}(DX-XD)=(\dim_kV)\,1_k.
\]
在特征零中这迫使 \(\dim_kV=0\)；在特征 \(p\) 中则迫使 \(p\mid\dim_kV\)。这一迹障碍也见 Etingof 等人\cite[Problem 2.7.4, p. 18]{EtingofEtAl2011RepresentationTheory}。

以下设 \(\operatorname{char}k=p>0\)。形式求导保持理想 \((X^p)\)，因为 \((X^pf)'=X^pf'\)，所以乘法与求导都下降到 \(V_p\)，仍满足 Weyl 关系。若一个非零子模含有非零 \(f\)，取次数小于 \(p\) 的代表，设其次数为 \(r\)、最高系数为 \(a_r\ne0\)。反复求导 \(r\) 次得到非零常数 \(r!a_r\)；这里 \(r<p\)，所以阶乘不为零。因此子模含一，再乘 \(X\) 就得到 \(1,X,\ldots,X^{p-1}\)，故 \(V_p\) 简单。

在 \(V_p\) 上，\(X^p\) 与 \(D^p\) 都为零，故作用经过所写商代数。要证明其像是全部线性算子，可以在像中逐一构造矩阵单位。令 \(H=XD\)，则 \(H(X^r)=rX^r\) 对 \(0\le r<p\)。因而
\[
E_0=\prod_{r=1}^{p-1}\frac{H-r\,1_{V_p}}{-r}
\]
恰为到常数项的投影：它在一上为恒等，在其余基向量上为零。所有分母都是 \(k\) 中的非零元素。对 \(0\le i,j<p\)，令
\[
E_{ij}=\frac{1}{j!}X^iE_0D^j.
\]
当 \(r<j\) 时，\(D^j(X^r)=0\)；当 \(r=j\) 时，它等于 \(j!\)；当 \(r>j\) 时，它是正次数单项式的标量倍，随后被 \(E_0\) 消去。因此 \(E_{ij}(X^r)=\delta_{jr}X^i\)，正是该基下的矩阵单位。它们全在商代数的像中，故诱导映射满射。另一方面，商代数由 \(x^i\mathrm{d}^j\)、\(0\le i,j<p\) 张成，维数至多 \(p^2\)；目标维数恰为 \(p^2\)，所以这个满射也是单射。
\end{proof}

% !TeX root = ../../Book2.tex
% 纲要标题：### 20.3 量子平面与扭曲代数
% 纲要位置：计划纲要.md，第 1710 行。

\section{量子平面与扭曲代数}
\label{b2:sec:2003}

把交换关系中的修正换成一个非零参数，可以得到另一类可计算的代数。下面始终固定 \(q\in k^\times\)。参数为一时恢复交换多项式环；其他参数是否为单位根，将直接影响中心的大小。

% 纲要要点（供目录核对）：
%
% * 量子平面的定义；从移动 jr 对字母产生扭曲因子，用预定基上的结合乘法证明独立性
% * 扭曲多项式环的实例；量子关系齐次而 Weyl 修正低次，取伴随分次时保留的信息不同
% * 中心与参数的关系；有限中心秩要求参数的精确阶，区别于 Weyl 中心随特征变化的机制
%

\subsection{量子平面的定义}
\label{b2:subsec:2003:01}

\begin{definition}\label{b2:def:quantum-plane}
设 \(k\) 为域，\(q\in k^\times\)。将自由结合代数 \(k\langle x,y\rangle\) 对关系 \(yx=qxy\) 生成的双边理想取商，所得 \(k\) 代数称为参数 \(q\) 的\term[liangzi-pingmian]{量子平面}[quantum plane]，记为
\[
k_q[x,y]=k\langle x,y\rangle/(yx-qxy).
\]
这里不要求 \(x,y\) 可逆。将它们也取逆得到的是另一种代数，不能与当前定义混同。
\end{definition}
\symidx{\(k_q[x,y]\)：参数 \(q\ne0\) 的量子平面}

\begin{theorem}\label{b2:thm:quantum-plane-normal-basis}
设 \(k\) 为域，\(q\in k^\times\)，\(k_q[x,y]=k\langle x,y\rangle/(yx-qxy)\) 为量子平面，则全部 \(x^iy^j\)、\(i,j\ge0\)，构成它的一组 \(k\) 基，且对所有非负整数 \(i,j,r,s\)，有
\begin{equation}\label{b2:eq:quantum-plane-product}
(x^iy^j)(x^ry^s)=q^{jr}x^{i+r}y^{j+s}.
\end{equation}
\end{theorem}
\begin{proof}
用 \(yx=qxy\) 移动所有错序字母，得到正规字的生成性。在乘积 \(x^iy^jx^ry^s\) 中，中间的每个 \(y\) 都要越过 \(r\) 个 \(x\)，总共交换 \(jr\) 次，每次产生一个因子 \(q\)。因而
\[
y^jx^r=q^{jr}x^ry^j,
\]
这给出了\eqref{b2:eq:quantum-plane-product}，但尚不能保证正规字独立。按照这个乘法取以 \(e_{ij}\)、\(i,j\ge0\) 为基的向量空间，令 \(e_{ij}\) 对应预期的正规字 \(x^iy^j\)，定义
\[
e_{ij}e_{rs}=q^{jr}e_{i+r,j+s}.
\]
乘三个基向量 \(e_{ij},e_{rs},e_{uv}\) 时，两种结合方式的系数指数分别为 \(jr+(j+s)u\) 与 \(su+j(r+u)\)，它们相同，所以乘法结合，单位为 \(e_{00}\)。元素 \(e_{10},e_{01}\) 满足量子平面的关系，故由\cref{b2:prop:algebra-relations-universal}得到到该代数的同态。正规字 \(x^iy^j\) 恰映到 \(e_{ij}\)，其独立性由后者为基得到。
\end{proof}

关系本身齐次，所以取 \(\deg x=\deg y=1\)，量子平面就是分次代数，第 \(n\) 次分量维数为 \(n+1\)。它与二元多项式环具有相同的分次维数，但乘法不同；例如 \(q\ne1\) 时，\(yx-xy=(q-1)xy\ne0\)。对这个分次所给的字长滤过取伴随分次，仍得到量子平面自身。Weyl 关系 \(\mathrm{d}\cdot x-x\cdot\mathrm{d}=1\) 则含低次修正，常数在最高次层消失，所以它的伴随分次是交换多项式环。量子关系中的 \(q\) 保留在同一次数内，取伴随分次不会消去这个扭曲。

\subsection{扭曲多项式环的实例}
\label{b2:subsec:2003:02}

令 \(\sigma:k[x]\to k[x]\)、\(\sigma\bigl(f(x)\bigr)=f(qx)\)。因为 \(q\ne0\)，它是自同构，逆为代入 \(q^{-1}x\)。由关系逐项计算，得到
\[
y\cdot f(x)=\sigma\bigl(f(x)\bigr)y,
\qquad y^r\cdot f(x)=\sigma^r\bigl(f(x)\bigr)y^r.
\]
所以每个元素可以唯一写为 \(\sum_j f_j(x)y^j\)，而乘法为
\[
\bigl(f(x)y^i\bigr)\bigl(g(x)y^j\bigr)
=f(x)\sigma^i\bigl(g(x)\bigr)y^{i+j}.
\]
这称为一个\term[niuqu-duoxiangshi-huan]{扭曲多项式环}[skew polynomial ring]，记为 \(k[x][y;\sigma]\)。在\secref{b2:sec:2005}中，会把它推广为允许导子修正的 Ore 扩张。

\begin{proposition}\label{b2:prop:quantum-plane-domain}
设 \(k\) 为域，\(q\in k^\times\)，\(k_q[x,y]=k\langle x,y\rangle/(yx-qxy)\) 为量子平面，则它没有零因子。若以正规式中 \(y\) 的最高指数为次数，任意两个非零元素乘积的次数等于两因子次数之和。
\end{proposition}
\begin{proof}
设两因子的最高项为 \(f_m(x)y^m\)、\(g_n(x)y^n\)。乘积的 \(y^{m+n}\) 系数为 \(f_m(x)\sigma^m\bigl(g_n(x)\bigr)\)。\(\sigma\) 是自同构，且 \(k[x]\) 是整环，所以这个系数非零。其他项的 \(y\) 次数更低，不能消去它，得到结论。
\end{proof}

例如 \(y(x^2+x+1)=(q^2x^2+qx+1)y\)。这里移动系数改变的是整段多项式的输入；不能只给整个多项式乘上同一个 \(q\)。若 \(q=-1\) 且特征不为二，偶次幂与 \(y\) 交换，奇次幂则改变符号。

\subsection{中心与参数的关系}
\label{b2:subsec:2003:03}

\begin{proposition}\label{b2:prop:quantum-plane-center}
设 \(k\) 为域，\(q\in k^\times\)，\(k_q[x,y]=k\langle x,y\rangle/(yx-qxy)\) 为量子平面。若 \(q\) 在 \(k^\times\) 中具有无限阶，则 \(Z\bigl(k_q[x,y]\bigr)=k\)。若 \(q\) 的阶为有限整数 \(\ell\ge1\)，则
\[
Z\bigl(k_q[x,y]\bigr)=k[x^\ell,y^\ell].
\]
后一情形中，量子平面作为中心上的模自由，基为 \(x^iy^j\)、\(0\le i,j<\ell\)，秩为 \(\ell^2\)。
\end{proposition}
\begin{proof}
由\eqref{b2:eq:quantum-plane-product}，正规式 \(P=\sum c_{ij}x^iy^j\) 与 \(x\) 交换恰好要求 \((q^j-1)c_{ij}=0\)，与 \(y\) 交换恰好要求 \((q^i-1)c_{ij}=0\)。正规单项式基使这些条件逐项成立而不能相互抵消。无限阶时只剩 \(i=j=0\)；有限阶时，\(q^r=1\) 恰好等价于 \(\ell\mid r\)，故要求 \(\ell\mid i,j\)，得到中心公式。

对一般指数写 \(i=a\ell+r,j=b\ell+s\)，其中 \(0\le r,s<\ell\)，再把中心因子提出，得到所列有限生成组。若它们有中心系数关系，展开这些系数后便得到不同的正规单项式之间的 \(k\) 线性关系，所有系数必须为零。因此它们确为中心上的基。
\end{proof}

这里的 \(\ell\) 必须是 \(q\) 的确切阶。仅有 \(q^\ell=1\) 可以保证 \(x^\ell,y^\ell\) 属于中心，却不能保证它们生成整个中心；例如 \(q=1\) 时虽然 \(q^2=1\)，中心仍为 \(k[x,y]\)，不是 \(k[x^2,y^2]\)。特别地，\(q=1\) 对应 \(\ell=1\)；\(q=-1\) 且特征不为二时，确切阶为二，中心为 \(k[x^2,y^2]\)，代数在中心上的秩为四。在特征二中，\(-1=1\)，同一写法则恢复交换多项式环。

\ProseParagraphBreak
量子平面中心的上述变化由 \(q\) 的乘法阶决定。\cref{b2:thm:weyl-center-simplicity}中 Weyl 代数的中心变化则由域特征决定：特征零时中心只有 \(k\)，特征 \(p>0\) 时中心为 \(k[x^p,\mathrm{d}^p]\)，在中心上的秩为 \(p^2\)。例如在特征零中取 \(q=-1\)，量子平面已在中心上有限，而 Weyl 代数仍有标量中心。两个有限秩公式都来自将正规单项式的指数按中心幂的指数取余，但使这些幂成为中心元的条件不同。

% !TeX root = ../../Book2.tex
% 纲要标题：### 20.4 Lie 代数与 PBW 型结构
% 纲要位置：计划纲要.md，第 1716 行。

% 纲要要点（供目录核对）：
%
% * 由交换的伴随分次控制非交换代数；含常数的一次滤过项、重排与独立性分开证明；Rees 的 Z² 来自齐次化最高次数，习题推广加权参数
% * Lie 代数的交替括号、Jacobi 恒等式、交换子与矩阵例子；在特征二保留交替性，区别 Lie 理想与结合代数双边理想
% * Lie 代数的表示及伴随表示；括号关系转为包络代数上的幺性左模，Heisenberg 的全参数中心商与正规基在原命题内证明
% * 与第五卷包络代数的 PBW 定理的联系；一般结论完整证明仍归第五卷第15章，j 的单射性只由 PBW 后取得，具体模型不借一般结论
% * 低次修正的 Jacobi 约束与 PBW 相容性；zyx 歧义与 Jacobi 差产生同一个隐藏关系，正文命题完整辨认实际二生成元商及正规基，原题改为加权 Ore 实例
%

\section{Lie 代数与 PBW 型结构}
\label{b2:sec:2004}

Weyl 代数的交换关系含常数修正，量子平面的扭曲关系则是齐次的。前两节分别借具体模型证明了有序单项式独立；若只给定一组带低次修正的关系，最高次交换关系只能给出从多项式环到伴随分次的满射，额外关系可能使它并非单射。结合代数的交换子满足 Jacobi 恒等式，Lie 代数将这一约束抽取为公理，泛包络代数再把括号写回结合乘法。有序单项式能否组成基，仍须由伴随分次检验。以下 \(k\) 为任意域。


\subsection{伴随分次与非交换结构}
\label{b2:subsec:2004:01}

\begin{proposition}\label{b2:prop:pbw-filtered-criterion}
设 \(k\) 为域，\(A\) 为由 \(x_1,\ldots,x_n\) 生成的含幺结合 \(k\) 代数，\(n\geq1\)。以这些生成元的字长不超过 \(m\) 的乘积张成 \(F_mA\)，其中 \(F_0A=k\cdot1\)、\(F_1A=k\cdot1+kx_1+\cdots+kx_n\)。若全部交换子 \([x_i,x_j]=x_ix_j-x_jx_i\) 均属于 \(F_1A\)，则有满同态
\[
k[X_1,\ldots,X_n]\longrightarrow\operatorname{gr}_F A,
\qquad X_i\longmapsto x_i+F_0A\in F_1A/F_0A.
\]
它是同构，当且仅当全部有序单项式 \(x_1^{a_1}\cdots x_n^{a_n}\) 是 \(A\) 的一组基。
\end{proposition}
\begin{proof}
交换子属于 \(F_1A\)，所以一次符号交换；字长滤过又使这些符号生成整个伴随分次，得到满同态。关系 \(x_jx_i=x_ix_j+[x_j,x_i]\) 的修正项是常数与生成元的线性组合。交换相邻错序字母时，二次项的字长不变而错序数减少；修正项的字长则严格减少。因此按字长及错序数的字典序归纳，可以把任意长度不超过 \(m\) 的字整理为全次数不超过 \(m\) 的有序单项式的线性组合。

若这些单项式已经是一组基，就连各个滤过项的基也被上述整理确定。相应次数的商片以相同次数的单项式符号为基，故满同态逐次单射。

反过来，若伴随分次同构于多项式环，任意非零的有限有序单项式关系中，取最高全次数的部分，其符号就是交换多项式单项式之间的非平凡关系，矛盾。因此有序单项式线性无关；生成性已证，得到基。
\end{proof}

满足这种有序基性质时，常称所讨论的构造具有 PBW 型性质。这里所指的是一个已经说明的基与滤过结论，而不是把任意非交换关系都统称为 PBW 定理。

\ProseParagraphBreak
Weyl 代数给出完全可算的形变。它的 Rees 代数由 \(X=xz,D=\mathrm{d}\cdot z,Z=z\) 生成。两个原生成元各乘了一次 \(z\)，所以
\[
DX-XD=(\mathrm{d}\cdot x-x\cdot\mathrm{d})z^2=z^2=Z^2.
\]
原关系的常数项必须补上两个 \(Z\)，才能与左侧同为二次。\(z\) 的中心性又给出
\[
ZX=XZ,\qquad ZD=DZ.
\]
这些关系给出整个代数。先将 \(Z\) 移到右侧，再用 \(DX=XD+Z^2\) 移动错序的 \(D,X\)；修正项使 \(X,D\) 的总个数减少二，因此按这个数和错序数归纳，可以把任何字写成 \(X^iD^jZ^r\) 的线性组合。它们在 Rees 代数中的像为 \(x^i\mathrm{d}^jz^{i+j+r}\)，由\cref{b2:thm:weyl-normal-basis}及不同 \(z\) 次数的独立性，它们线性无关。因此该关系代数与 Rees 代数同构。

\ProseParagraphBreak
它作为 \(k[Z]\) 模以 \(X^iD^j\) 为基。在 \(Z=0\) 处得到交换平面，在 \(Z=1\) 处得到 Weyl 代数；这是\cref{b2:thm:rees-fibres}的一个具体实现。

\subsection{Lie 代数与基本例子}
\label{b2:subsec:2004:lie-basics}

在 Weyl 代数中，交换子 \(\mathrm{d}\cdot x-x\cdot\mathrm{d}\) 比乘积 \(\mathrm{d}\cdot x\) 的次数低。现在暂时不保留两个乘积，只把它们的差作为一项双线性运算。结合律会使这种运算满足额外的恒等式；这些恒等式正是下面定义的依据。

\begin{definition}\label{b2:def:lie-algebra}
设 \(k\) 为域，\(\mathfrak g\) 为 \(k\) 向量空间，并给定双线性映射
\[
[\ ,\ ]:\mathfrak g\times\mathfrak g\longrightarrow\mathfrak g,
\]
若它满足交替性 \([u,u]=0\) 对所有 \(u\in\mathfrak g\) 成立，以及对任意 \(u,v,w\in\mathfrak g\) 成立的\term[Jacobi-hengdengshi]{Jacobi 恒等式}[Jacobi identity]：\footnote{雅可比（Carl Gustav Jacob Jacobi，1804-1851），德国数学家，研究椭圆函数、行列式和力学。}
\begin{equation}\label{b2:eq:lie-jacobi}
\bigl[u,[v,w]\bigr]+\bigl[v,[w,u]\bigr]+\bigl[w,[u,v]\bigr]=0.
\end{equation}
则称 \((\mathfrak g,[\ ,\ ])\) 为 \(k\) 上的\term[lie-daishu]{Lie 代数}[Lie algebra]，称 \([\ ,\ ]\) 为 Lie 括号。这里不要求结合律，也不指定单位元。
\end{definition}
\symidx{\([u,v]\)：Lie 代数的括号}

展开 \([u+v,u+v]=0\)，得 \([u,v]=-[v,u]\)。反过来，在特征不等于二时，反对称性推出 \([u,u]=0\)；特征二时这一推论失效，所以定义中保留交替性。这与\cref{b2:prop:alternating-implies-skew,b2:def:symmetric-alternating-form}中交替映射和双线性型的区别相同。例如一维空间 \(ke\) 上的双线性运算 \([e,e]=e\) 在特征二时反对称，却不交替。将任意向量空间上的括号恒取为零，则满足全部公理，称为\term[jiaohuan-lie-daishu]{交换 Lie 代数}[abelian Lie algebra]。

\begin{proposition}\label{b2:prop:commutator-lie-algebra}
设 \(k\) 为域，\(A\) 为结合 \(k\) 代数。在其底向量空间上定义 \([a,b]=ab-ba\) 对任意 \(a,b\in A\) 成立，则所得结构为 \(k\) 上的 Lie 代数，记为 \(A^-\)。
\end{proposition}
\begin{proof}
乘法的双线性给出括号的双线性，且 \([a,a]=0\)。对于 Jacobi 恒等式，逐项展开为
\[
\begin{aligned}
\bigl[a,[b,c]\bigr]&=abc-acb-bca+cba,\\
\bigl[b,[c,a]\bigr]&=bca-bac-cab+acb,\\
\bigl[c,[a,b]\bigr]&=cab-cba-abc+bac.
\end{aligned}
\]
结合律使这里的三重乘积无需另加括号；相加后每个乘积的两个系数互为相反数，所得和为零。这个消去在任意特征下成立。
\end{proof}
\symidx{\(A^-\)：结合代数的交换子 Lie 代数}

特别地，\(\mathfrak{gl}(V)=\operatorname{End}_k(V)^-\)，有限维且选定基时也记为 \(\mathfrak{gl}_n(k)=M_n(k)^-\)。其括号由算子复合或矩阵乘法之差给出，与群中的乘法交换子 \(aba^{-1}b^{-1}\) 是不同运算。
\symidx{\(\mathfrak{gl}(V)\)：向量空间的一般线性 Lie 代数}

\ProseParagraphBreak
设 \(\mathfrak g\) 为 \(k\) 上的 Lie 代数。若线性子空间 \(\mathfrak h\subseteq\mathfrak g\) 对括号封闭，则称 \(\mathfrak h\) 为\term[lie-zidaishu]{Lie 子代数}[Lie subalgebra]；公理可从 \(\mathfrak g\) 直接限制下来。对 \(k\) 上的 Lie 代数 \(\mathfrak g,\mathfrak h\)，若 \(k\) 线性映射 \(f:\mathfrak g\to\mathfrak h\) 保持括号，即
\[
f\bigl([u,v]\bigr)=\bigl[f(u),f(v)\bigr].
\]
则称 \(f\) 为\term[lie-daishu-tongtai]{Lie 代数同态}[Lie algebra homomorphism]。

\ProseParagraphBreak
若线性子空间 \(I\subseteq\mathfrak g\) 满足 \([u,a]\in I\) 对所有 \(u\in\mathfrak g\)、\(a\in I\) 成立，则称 \(I\) 为 \(\mathfrak g\) 的\term[lie-lixiang]{Lie 理想}[ideal of a Lie algebra]。交替性使 \([a,u]\) 也属于 \(I\)。这里的封闭条件只涉及括号，不能换成结合乘法的左右封闭性。例如在 \(M_2(k)^-\) 中，\(kI_2\) 是 Lie 理想，因为与任意矩阵的括号为零；它却不是 \(M_2(k)\) 的双边理想，因为 \(E_{11}I_2=E_{11}\notin kI_2\)。若 \(I\) 为 Lie 理想，则
\[
[u+I,v+I]=[u,v]+I
\]
在 \(\mathfrak g/I\) 上良定义：替换两个代表元所增加的 \([a,v]+[u,b]+[a,b]\) 均属于 \(I\)。双线性、交替性和 Jacobi 恒等式由商映射传递，因而得到商 Lie 代数。特别地，Lie 同态的核是这样的理想，其像是 Lie 子代数。

\ProseParagraphBreak
设 \(\mathfrak g\) 为 \(k\) 上的 Lie 代数。若 \(u\in\mathfrak g\) 与 \(\mathfrak g\) 的全部元素括号为零，则称 \(u\) 为中心元素。所有中心元素组成 \(\mathfrak g\) 的\term[lie-daishu-zhongxin]{中心}[center of a Lie algebra]，记为
\[
Z(\mathfrak g)=\bigl\{u\in\mathfrak g\mid [u,v]=0\;\text{对所有}\;v\in\mathfrak g\bigr\}.
\]
它是一个 Lie 理想，因为对任意 \(u\in Z(\mathfrak g)\)、\(v\in\mathfrak g\)，都有 \([v,u]=0\in Z(\mathfrak g)\)。
\symidx{\(Z(\mathfrak g)\)：Lie 代数的中心}

\ProseParagraphBreak
最小的非交换例子已能在二阶矩阵中看到。令 \(h=E_{11}\)、\(e=E_{12}\)，则
\[
[h,e]=e,\qquad [h,h]=[e,e]=0.
\]
因此 \(kh\oplus ke\) 对括号封闭，且
\[
[ah+be,ch+de]=(ad-bc)e.
\]
它是二维非交换 Lie 代数。在这里，封闭性和\cref{b2:prop:commutator-lie-algebra}已经保证 Jacobi 恒等式，不需要对每一组三元组重新检验。

\begin{example}\label{b2:exm:heisenberg-lie-algebra}
设 \(k\) 为域。\(M_3(k)^-\) 中三阶严格上三角矩阵组成的 Lie 代数记为 \(\mathfrak h\)，称为三维\term[heisenberg-lie-daishu]{Heisenberg Lie 代数}[Heisenberg Lie algebra]。\footnote{海森堡（Werner Karl Heisenberg，1901-1976），德国物理学家，创立矩阵力学，提出不确定性原理。}取
\[
x=E_{23},\qquad d=E_{12},\qquad z=E_{13},
\]
则 \(x,d,z\) 是它的一组基，括号满足
\begin{equation}\label{b2:eq:heisenberg-brackets}
[d,x]=z,\qquad [z,x]=[z,d]=0.
\end{equation}
确实，\(E_{ij}E_{rs}=\delta_{jr}E_{is}\) 给出 \(dx=z\)、\(xd=0\)，而 \(z\) 与 \(x,d\) 的两个方向乘积均为零。因而这些矩阵对括号封闭，\(z\) 是一个非零的中心元素。
\end{example}

上述交换子与矩阵例子也见 Etingof 等人的定义和例表\cite[Definition 2.9.1; Example 2.9.2, pp. 22--23]{EtingofEtAl2011RepresentationTheory}。这里先用矩阵检验抽象公理；稍后将通过同一个 Heisenberg 例子回到 Weyl 代数。

\subsection{Lie 代数的表示}
\label{b2:subsec:2004:lie-representations}

矩阵例子说明，Lie 代数可以通过线性算子的交换子实现。对于一个给定的抽象 Lie 代数，不必要求这种实现单射；它只是指定各元素怎样作用在某个向量空间上。

\begin{definition}\label{b2:def:lie-representation}
设 \(k\) 为域，\(\mathfrak g\) 为 \(k\) 上的 Lie 代数，\(V\) 为 \(k\) 向量空间。若 \(k\) 线性映射
\[
\rho:\mathfrak g\longrightarrow\operatorname{End}_k(V)
\]
对任意 \(u,v\in\mathfrak g\) 满足
\begin{equation}\label{b2:eq:lie-representation}
\rho\bigl([u,v]\bigr)=\rho(u)\rho(v)-\rho(v)\rho(u).
\end{equation}
则称 \(\rho\) 为 \(\mathfrak g\) 在 \(V\) 上的\term[lie-daishu-biaoshi]{表示}[representation of a Lie algebra]；换言之，\(\rho:\mathfrak g\to\mathfrak{gl}(V)\) 是 Lie 代数同态。记 \(u\cdot w=\rho(u)(w)\)，其条件就是
\[
[u,v]\cdot w=u\cdot(v\cdot w)-v\cdot(u\cdot w).
\]
带有这种作用的 \(V\) 也称为 \(\mathfrak g\) 模。
\termidx[lie-mo]{Lie 代数模}
\foreignidx{Lie algebra module}
\end{definition}

若对所有 \(u\) 取 \(\rho(u)=0\)，得到平凡表示。若 \(\mathfrak g\) 是矩阵 Lie 子代数，其矩阵对列向量的自然作用就是一个表示。若表示映射 \(\rho\) 单射，则称它为\term[zhongshi-biaoshi]{忠实表示}[faithful representation]。两个表示间的同态是满足 \(F(u\cdot v)=u\cdot F(v)\) 的线性映射 \(F:V\to W\)；它要求作用相容，不要求 \(F\) 可逆。

\begin{proposition}\label{b2:prop:adjoint-lie-representation}
设 \(k\) 为域，\(\mathfrak g\) 为 \(k\) 上的 Lie 代数。令
\[
\operatorname{ad}(u)(v)=[u,v].
\]
则 \(\operatorname{ad}:\mathfrak g\to\mathfrak{gl}(\mathfrak g)\) 是一个表示，称为\term[bansui-biaoshi]{伴随表示}[adjoint representation]。它的核是
\[
Z(\mathfrak g)=\bigl\{u\in\mathfrak g\mid [u,v]=0\;\text{对所有}\;v\in\mathfrak g\bigr\},
\]
即 \(\mathfrak g\) 的中心。
\end{proposition}
\begin{proof}
括号的双线性使 \(\operatorname{ad}(u)\) 为线性算子，并使 \(u\mapsto\operatorname{ad}(u)\) 线性。把\eqref{b2:eq:lie-jacobi}中的后两项用反对称性改写，得到
\[
\bigl[[u,v],w\bigr]
=\bigl[u,[v,w]\bigr]-\bigl[v,[u,w]\bigr].
\]
因此，对任意 \(w\)，
\[
\operatorname{ad}\bigl([u,v]\bigr)(w)
=\bigl(\operatorname{ad}(u)\operatorname{ad}(v)
-\operatorname{ad}(v)\operatorname{ad}(u)\bigr)(w),
\]
正是\eqref{b2:eq:lie-representation}。最后，\(\operatorname{ad}(u)=0\) 等价于 \([u,v]=0\) 对全部 \(v\) 成立，故其核恰为所写的中心。
\end{proof}
\symidx{\(\operatorname{ad}\)：Lie 代数的伴随表示}

这说明 Jacobi 恒等式控制着一个实际作用：先取括号 \([u,v]\) 再作用，与两个作用算子的交换子相同。它不保证伴随表示忠实。例如\cref{b2:exm:heisenberg-lie-algebra}中的 \(z\ne0\)，却有 \(\operatorname{ad}(z)=0\)。表示的定义与伴随表示也可对照 Etingof 等人\cite[Definition 2.9.7; Example 2.9.8, p. 24]{EtingofEtAl2011RepresentationTheory}。

\subsection{包络代数与 PBW 定理}
\label{b2:subsec:2004:02}

表示条件\eqref{b2:eq:lie-representation}混合了 Lie 括号与算子复合。若把每个 \(u\in\mathfrak g\) 当成一个结合代数的生成元，再加入恰好保证该条件的关系，就能用已经建立的模论讨论 Lie 表示。

\begin{definition}\label{b2:def:universal-enveloping-algebra}
设 \(k\) 为域，\(\mathfrak g\) 为 \(k\) 上的 Lie 代数，\(T_k(\mathfrak g)\) 为其底向量空间的张量代数。令 \(J\) 为全部元素
\[
u\otimes v-v\otimes u-[u,v],\qquad u,v\in\mathfrak g,
\]
生成的双边理想，其中最后一项属于一次张量分量。定义商代数
\[
U(\mathfrak g)=T_k(\mathfrak g)/J
\]
并称它为 \(\mathfrak g\) 的\term[fan-baoluo-daishu]{泛包络代数}[universal enveloping algebra]。以 \(j:\mathfrak g\to U(\mathfrak g)\) 表示一次分量包含与商映射的复合。
\end{definition}
\symidx{\(U(\mathfrak g)\)：Lie 代数的泛包络代数}

关系 \(u\otimes v-v\otimes u-[u,v]\) 的前两项属于张量次数二，最后一项属于张量次数一。因此张量字长仍给出商代数的滤过，原来的张量次数分次通常却不能下降到商；随后考察 \(\operatorname{gr}U(\mathfrak g)\)，正是为了分离这些关系的最高次部分。

\ProseParagraphBreak
定义只给出线性映射 \(j\)，尚未证明它单射。在证明单射以前，只能使用其像 \(j(\mathfrak g)\)，不能把 \(\mathfrak g\) 直接视为 \(U(\mathfrak g)\) 的线性子空间。商关系保证
\begin{equation}\label{b2:eq:enveloping-canonical-lie-map}
j\bigl([u,v]\bigr)=j(u)j(v)-j(v)j(u),
\end{equation}
所以 \(j:\mathfrak g\to U(\mathfrak g)^-\) 是 Lie 同态。这个构造与\cref{b2:def:enveloping-algebra}的 \(A^e=A\otimes A^{\mathrm{op}}\) 不同：\(A^e\) 用于将结合代数的双模视为左模，而 \(U(\mathfrak g)\) 用于将 Lie 代数的表示视为结合代数的模。

\begin{theorem}[泛包络代数的泛性质]\label{b2:thm:enveloping-universal}
设 \(k\) 为域，\(\mathfrak g\) 为 \(k\) 上的 Lie 代数，\(U(\mathfrak g)\) 为其泛包络代数，\(j:\mathfrak g\rightarrow U(\mathfrak g)\) 为自然 Lie 同态。对任意含幺结合 \(k\) 代数 \(B\) 和 Lie 同态 \(f:\mathfrak g\to B^-\)，其中 \(B^-\) 使用交换子括号，存在唯一保幺的 \(k\) 代数同态
\[
\widetilde f:U(\mathfrak g)\longrightarrow B,
\qquad \widetilde f\circ j=f.
\]
反过来，每个这样的代数同态与 \(j\) 复合，都给出一个 Lie 同态。
\end{theorem}
\begin{proof}
先将 \(f\) 视为线性映射。\cref{b2:thm:tensor-algebra-universal}给出唯一代数同态 \(\widehat f:T_k(\mathfrak g)\to B\)，在一次分量上等于 \(f\)。对关系理想的每个指定生成元，有
\[
\widehat f\bigl(u\otimes v-v\otimes u-[u,v]\bigr)
=f(u)f(v)-f(v)f(u)-f\bigl([u,v]\bigr)=0,
\]
最后的等式正是 \(f\) 保持括号的条件。\(\widehat f\) 是代数同态，所以也零化这些生成元的任意左右倍数之和，即整个 \(J\)。于是 \(\widetilde f(t+J)=\widehat f(t)\) 良定义，给出所需代数同态。\(U(\mathfrak g)\) 由 \(j(\mathfrak g)\) 及单位元生成，故其像一经指定，同态唯一。

反向若给定 \(\widetilde f\)，将它作用于\eqref{b2:eq:enveloping-canonical-lie-map}，便得 \(\widetilde f\circ j\) 保持括号。由刚证的唯一性，两个过程互逆。
\end{proof}

取零 Lie 同态 \(\mathfrak g\to k^-\)，得到代数同态 \(U(\mathfrak g)\to k\)，它将每个 \(j(u)\) 送到零，将单位元送到一。因此 \(U(\mathfrak g)\) 的单位元非零。

\begin{corollary}\label{b2:cor:lie-representations-enveloping-modules}
设 \(k\) 为域，\(\mathfrak g\) 为 \(k\) 上的 Lie 代数，\(V\) 为 \(k\) 向量空间，\(U(\mathfrak g)\) 为泛包络代数，则 \(\mathfrak g\) 在 \(V\) 上的表示与 \(V\) 上的幺性左 \(U(\mathfrak g)\) 模结构一一对应；对任意两个这样的向量空间，Lie 表示同态恰为相应的 \(U(\mathfrak g)\) 模同态。
\end{corollary}
\begin{proof}
将\cref{b2:thm:enveloping-universal}用于 \(B=\operatorname{End}_k(V)\)，Lie 表示 \(\rho\) 唯一延拓为代数同态 \(\widetilde\rho:U(\mathfrak g)\to\operatorname{End}_k(V)\)。规定 \(a\cdot v=\widetilde\rho(a)(v)\)，乘法与单位元的保持性恰给出幺性左模。反向把一个幺性左模的作用限制到 \(j(\mathfrak g)\)，由\eqref{b2:eq:enveloping-canonical-lie-map}得到 Lie 表示；唯一性说明两个过程互逆。

若线性映射 \(F:V\to W\) 与每个 \(u\in\mathfrak g\) 的作用相容，对乘积的长度归纳便知它与每个 \(j(u_1)\cdots j(u_r)\) 的作用相容；线性性再使它与这些乘积及单位元的任意线性组合相容，即为 \(U(\mathfrak g)\) 模同态。反向的限制立即给出 Lie 表示同态。
\end{proof}

表示中的括号相容条件被写成了 \(U(\mathfrak g)\) 的代数关系，而连续施加 Lie 元素的作用成为这个代数中的乘法。由于 \(U(\mathfrak g)\) 由 \(j(\mathfrak g)\) 和单位元生成，一个子空间对所有 Lie 元素的作用稳定，恰好等价于它是 \(U(\mathfrak g)\) 子模。于是前面的子模、商模和自同态环理论都可以用于 Lie 表示。

\ProseParagraphBreak
在交换 Lie 代数的情形，关系只有 \(uv-vu=0\)，故由\cref{b2:def:symmetric-algebra}，\(U(\mathfrak g)=\operatorname{Sym}_k(\mathfrak g)\)。若 \(\mathfrak g\) 有有限基 \(x_1,\ldots,x_n\)，\cref{b2:thm:symmetric-polynomial-algebra}将它识别为 \(k[x_1,\ldots,x_n]\)。因而交换 Lie 代数的表示就是一组两两交换的线性算子。

\begin{proposition}\label{b2:prop:heisenberg-enveloping-weyl}
设 \(k\) 为域，\(\mathfrak h\) 为以 \(x,d,z\) 为基、满足 \([d,x]=z\)、\([z,x]=[z,d]=0\) 的三维 Heisenberg Lie 代数，\(j:\mathfrak h\rightarrow U(\mathfrak h)\) 为自然映射。对任意 \(\lambda\in k\)，记
\[
B_\lambda=U(\mathfrak h)/\bigl(j(z)-\lambda1\bigr),
\]
其中括号表示所列元素生成的双边理想，\(\bar x,\bar d\) 分别为 \(j(x),j(d)\) 在 \(B_\lambda\) 中的像。\(B_\lambda\) 的 \(k\) 基为全部 \(\bar x^i\bar d^j\)、\(i,j\ge0\)。当 \(\lambda=0\) 时，\(B_0\cong k[X,D]\)，同构将 \(\bar x,\bar d\) 送到 \(X,D\)；当 \(\lambda\ne0\) 时，\(B_\lambda\cong A_1(k)\)，同构将 \(\bar x,\bar d\) 分别送到 Weyl 生成元 \(x,\lambda\cdot\mathrm{d}\)。对任意 \(k\) 向量空间 \(V\)，其上的幺性左 \(B_\lambda\) 模结构恰好对应中心元素 \(z\) 作用为 \(\lambda\operatorname{id}_V\) 的 \(\mathfrak h\) 表示。
\end{proposition}
\begin{proof}
商关系给出 \(\bar d\bar x-\bar x\bar d=\lambda1\)，并使 \(j(z)\) 的像为 \(\lambda1\)，所以 \(B_\lambda\) 由 \(\bar x,\bar d\) 及单位元生成。

若 \(\lambda=0\)，这两个生成元交换，故 \(X\mapsto\bar x,D\mapsto\bar d\) 给出代数同态 \(k[X,D]\to B_0\)。反向指定 Lie 同态 \(\mathfrak h\to k[X,D]^-\)，将 \(x,d,z\) 分别送到 \(X,D,0\)。因为目标的括号为零，这些像满足\eqref{b2:eq:heisenberg-brackets}，由\cref{b2:thm:enveloping-universal}得到 \(U(\mathfrak h)\to k[X,D]\)，再由 \(j(z)\) 映为零下降到 \(B_0\)。两个复合均固定生成元和单位元，因而互为逆映射。交换多项式环的基给出 \(B_0\) 的所列基。

若 \(\lambda\ne0\)，在 \(A_1(k)^-\) 中指定 \(x\mapsto x,d\mapsto\lambda\cdot\mathrm{d},z\mapsto\lambda1\)。等式 \([\lambda\cdot\mathrm{d},x]=\lambda1\) 和标量的中心性保证这是 Lie 同态；同样的泛性质和取商过程给出 \(B_\lambda\to A_1(k)\)。反向，\(\lambda^{-1}\bar d\) 与 \(\bar x\) 满足 Weyl 关系，故由\cref{b2:prop:algebra-relations-universal}得到 \(A_1(k)\to B_\lambda\)，将 \(x,\mathrm{d}\) 分别送到 \(\bar x,\lambda^{-1}\bar d\)。两个复合再次固定生成元和单位元，所以互为逆映射。所列单项式在前一个同构下成为 \(\lambda^jx^i\mathrm{d}^j\)，其中 \(\lambda^j\ne0\)，由\cref{b2:thm:weyl-normal-basis}得到基。

最后，\cref{b2:cor:lie-representations-enveloping-modules}将 \(\mathfrak h\) 表示视为 \(U(\mathfrak h)\) 模，而这种作用下降到 \(B_\lambda\) 当且仅当 \(j(z)-\lambda1\) 的作用为零，即 \(z\) 作用为 \(\lambda\operatorname{id}_V\)。
\end{proof}

取 \(\lambda=1\) 就得到原来的 Weyl 代数；这个辨认可与 Etingof 等人的 Heisenberg 例子对照\cite[Example 2.9.13, p. 24]{EtingofEtAl2011RepresentationTheory}。取 \(\lambda=0\) 则得到交换多项式环。同一个 Lie 代数的中心元素在表示中取不同的标量作用，所对应的结合代数商可以由非交换变为交换。

\ProseParagraphBreak
对有限维 \(\mathfrak g\) 选有序基 \(x_1,\ldots,x_n\)，令 \(F_mU(\mathfrak g)\) 为长度不超过 \(m\) 的 \(j(x_i)\) 乘积所生成的线性空间。\eqref{b2:eq:enveloping-canonical-lie-map}的右侧交换子属于一次滤过项，所以\cref{b2:prop:pbw-filtered-criterion}给出自然满射
\[
\operatorname{Sym}_k(\mathfrak g)\longrightarrow\operatorname{gr}U(\mathfrak g),
\qquad x_i\longmapsto j(x_i)+F_0U(\mathfrak g).
\]
这里还使用\cref{b2:thm:symmetric-polynomial-algebra}把给定基上的对称代数与多项式环识别。满射的单射性是下面具名结果的内容。

\begin{claim}[PBW]\label{b2:claim:pbw-enveloping-preview}
设 \(k\) 为域，\(\mathfrak g\) 为\(k\) 上的有限维 Lie 代数，\(x_1,\ldots,x_n\) 为其有序基，\(U(\mathfrak g)\) 为泛包络代数，\(j:\mathfrak g\rightarrow U(\mathfrak g)\) 为自然映射。按 \(j(x_i)\) 的乘积长度滤过 \(U(\mathfrak g)\)，则自然映射 \(\operatorname{Sym}_k(\mathfrak g)\rightarrow\operatorname{gr}U(\mathfrak g)\) 为同构；等价地，全部 \(j(x_1)^{a_1}\cdots j(x_n)^{a_n}\)、\(a_i\geq0\)，构成 \(U(\mathfrak g)\) 的 \(k\) 基。特别地，\(j\) 单射。
\end{claim}

这称为\term[PBW-dingli]{Poincaré--Birkhoff--Witt 定理}[Poincaré--Birkhoff--Witt theorem]。\footnote{庞加莱（Henri Poincaré，1854-1912），法国数学家，研究拓扑、动力系统与微分方程。伯克霍夫（Garrett Birkhoff，1911-1996），美国数学家，研究格论与普遍代数。Witt 已在本卷前文介绍。}完整证明见第五卷第15章。前面的 Weyl 与量子平面正规基，以及上述 Heisenberg 商的基，已经由具体模型或互逆同态证明。较低次的括号关系必须满足相容条件，Jacobi 恒等式正是其中的关键。

\subsection{低次修正与 PBW 相容性}
\label{b2:subsec:2004:03}

考虑三个生成元 \(x,y,z\)，随意指定
\[
yx=xy+x,\qquad zy=yz+y,\qquad zx=xz.
\]
这些关系的最高次部分都是交换关系，但不保证基仍是 \(x^ay^bz^c\)。

\begin{proposition}\label{b2:exm:incompatible-pbw-relations}
设 \(k\) 为域，并令
\[
A=k\langle x,y,z\rangle/(yx-xy-x,\ zy-yz-y,\ zx-xz).
\]
在 \(A\) 中必有 \(x=0\)，且将 \(y,z\) 的同名商类对应起来给出代数同构
\[
A\cong k\langle y,z\rangle/(zy-yz-y).
\]
全部 \(y^iz^j\)、\(i,j\ge0\)，构成 \(A\) 的 \(k\) 基。因此预期的三变量 PBW 基虽然失败，余下两个生成元仍有唯一的正规式。
\end{proposition}
\begin{proof}
对同一个字 \(zyx\) 先整理左边的 \(zy\)，得到
\[
(zy)x=(yz+y)x=yxz+yx=xyz+xz+xy+x.
\]
先整理右边的 \(yx\)，则得到
\[
z(yx)=z(xy+x)=xzy+xz=xyz+xy+xz.
\]
结合律使两者相等，故 \(x=0\)。这里出现的是一次额外关系，单看最高次的三条交换关系无法预见。

记 \(B=k\langle y,z\rangle/(zy-yz-y)\)。将 \(x,y,z\) 分别送到 \(0,y,z\)，\(A\) 的三条定义关系均成立，得到同态 \(A\to B\)。反向，\(A\) 中的 \(y,z\) 满足 \(B\) 的定义关系，给出同态 \(B\to A\)。它们的复合固定 \(y,z\) 和单位元；在 \(A\) 中又有 \(x=0\)，所以两个复合处处为恒等，得到所述同构。

在 \(B\) 中，用 \(zy=yz+y\) 整理错序字母，二次项的错序数下降，一次项的字长下降，故按字长及错序数归纳，\(y^iz^j\) 张成 \(B\)。为了检验独立性，取交换多项式空间 \(k[Y,Z]\)，以 \(M_Y\) 表示乘以 \(Y\) 的算子。用 \(Y\partial_Y\) 产生关系中的 \(y\)，再加上与 \(M_Y\) 交换的乘 \(Z\) 算子，以便作用于 \(1\) 时仍能记录 \(z\) 的指数。具体令
\[
H(f)=Y\frac{\partial f}{\partial Y}+Zf.
\]
形式求导的 Leibniz 公式在任意特征下都给出
\[
(HM_Y-M_YH)(f)=Yf=M_Y(f).
\]
因此 \(y\mapsto M_Y,z\mapsto H\) 满足 \(B\) 的关系，得到代数同态 \(B\to\operatorname{End}_k(k[Y,Z])\)。由于 \(H(Z^j)=Z^{j+1}\)，有 \(H^j(1)=Z^j\)，从而
\[
M_Y^iH^j(1)=Y^iZ^j.
\]
若正规单项式间存在有限线性关系，将相应算子作用于 \(1\)，便得到不同交换单项式 \(Y^iZ^j\) 间的同一线性关系，所有系数只能为零。这证明了 \(B\) 以及 \(A\) 的所列基。
\end{proof}

这种失败可以用括号解释：指定的规则给出 \([y,x]=x\)、\([z,y]=y\)、\([z,x]=0\)。在以 \(x,y,z\) 为基的三维空间上，它们使
\[
\bigl[z,[y,x]\bigr]+\bigl[y,[x,z]\bigr]+\bigl[x,[z,y]\bigr]=-x\ne0,
\]
所以不满足 Jacobi 恒等式。这个 Jacobi 差为 \(-x\)，正是前面两种整理结果之差 \(x\) 的相反数，二者暴露的是同一个额外关系。\cref{b2:prop:commutator-lie-algebra}已证明，结合代数中的加法交换子必满足该恒等式，因此关系只能通过使 \(x\) 消失来相容。\secref{b2:sec:2006}将把这种同一字的两种整理方式作为重叠歧义处理。

% !TeX root = ../../Book2.tex
% 纲要标题：### 20.5* Ore 扩张与微分算子环
% 纲要位置：计划纲要.md，第 1724 行。

\OptionalSection{Ore 扩张与微分算子环}{b2:sec:2005}

Weyl 关系在移动系数时增加一个导数项，量子平面则对系数施加自同构。Ore 扩张把这两种操作合在一起，添加一个新的形式变量，并规定它与原环系数怎样相乘。第十九章的 Ore 局部化则把原环中指定的元素变成可逆元；这里的新变量并不因构造而可逆。

% 纲要要点（供目录核对）：
%
% * 环自同态 \(\sigma\) 与 \(\sigma\)-导子 \(\delta\)：\(\delta(ab)=\sigma(a)\delta(b)+\delta(a)b\)；结合律强迫扭曲 Leibniz 条件
% * Ore 扩张 \(R[t;\sigma,\delta]\) 的关系 \(ta=\sigma(a)t+\delta(a)\) 与正规表达；自由左模上的加法算子在一上读取系数，正文命题证明自同构时的右正规式及非单射边界
% * 微分域上的 \(K[\partial;\delta]\)：\(\partial a=a\partial+\delta(a)\)；区别形式生成元与其作用，Weyl 嵌入由正规式而非关系成立保证
% * 左模与线性微分方程；左理想、首一左除法和解同态就地证明，伴随矩阵的模坐标与取值转置写为正式命题；原题改为含系数导数的换基，第三卷附录 E 和第五卷附录 E 作为延伸
%

\subsection{环自同态与扭曲导子}
\label{b2:subsec:2005:01}

\begin{definition}\label{b2:def:sigma-derivation}
设 \(R\) 为含幺环，\(\sigma:R\to R\) 为保幺环自同态（不要求可逆），\(\delta:R\to R\) 为加法映射。若对任意 \(a,b\in R\)，有
\[
\delta(ab)=\sigma(a)\delta(b)+\delta(a)b,
\]
则称 \(\delta\) 为\term[sigma-daozi]{左 \(\sigma\) 导子}[left sigma-derivation]。这里的“左”指 \(\sigma\) 作用于乘积的左因子；\(\sigma=\operatorname{id}\) 时就是通常的导子。
\end{definition}

由 \(\delta(1)=\delta(1)+\delta(1)\) 得 \(\delta(1)=0\)。例如 \(k[x]\) 上的形式求导是恒等自同态的导子；对于任意自同态，零映射总是一个 \(\sigma\) 导子。

\ProseParagraphBreak
扭曲导子公式是结合律 \((ta)b=t(ab)\) 强迫的相容条件。若希望 \(ta=\sigma(a)t+\delta(a)\)，则先移动 \(a\) 再移动 \(b\) 会得到
\[
(ta)b=\sigma(a)\sigma(b)t+\sigma(a)\delta(b)+\delta(a)b.
\]
它须与 \(t(ab)\) 的展开相同，因而需要上述 \(\sigma\) 导子等式。不过检查这一处相容还不是整个代数存在性的证明，下面给出一个实际模型。

\subsection{Ore 扩张的关系与正规表达}
\label{b2:subsec:2005:02}

\begin{theorem}\label{b2:thm:ore-extension-construction}
设 \(R\) 为含幺环，\(\sigma:R\rightarrow R\) 为保幺环自同态，\(\delta:R\rightarrow R\) 为满足 \(\delta(ab)=\sigma(a)\delta(b)+\delta(a)b\) 的加法映射，则存在含幺环 \(R[t;\sigma,\delta]\)，含 \(R\) 为子环，且每个元素唯一写成有限和
\[
a_0+a_1t+\cdots+a_nt^n,\qquad a_i\in R,
\]
满足 \(ta=\sigma(a)t+\delta(a)\)。它作为左 \(R\) 模以 \(1,t,t^2,\ldots\) 为基。

对任意含幺环 \(B\)、保幺环同态 \(f:R\to B\) 及满足下式对所有 \(a\in R\) 成立的 \(u\in B\)，
\[
u\cdot f(a)=f\bigl(\sigma(a)\bigr)u+f\bigl(\delta(a)\bigr)
\]
存在唯一保幺环同态 \(R[t;\sigma,\delta]\to B\)，延伸 \(f\) 并把 \(t\) 送到 \(u\)。
\end{theorem}
\begin{proof}
为证明各个 \(t^i\) 左线性无关，取形式左自由模 \(P=\bigoplus_{n\ge0}Rz^n\)，并寻找一个算子 \(T\)，使 \(T^i(1)=z^i\)。这样，把候选关系作用于一，就能读出每个幂的系数。同时，\(T\) 必须实现题设的扭曲系数关系。因此在 \(P\) 的加法群上定义
\[
L_a\left(\sum_n b_nz^n\right)=\sum_n ab_nz^n,
\]
\[
T\left(\sum_n b_nz^n\right)
=\sum_n\bigl(\sigma(b_n)z^{n+1}+\delta(b_n)z^n\bigr).
\]
它们是加法群自同态。对 \(b\in R,p\in P\)，导子公式给出
\[
T(bp)=\sigma(b)T(p)+\delta(b)p,
\]
所以 \(T\) 通常不是左 \(R\) 线性的；非交换环中，\(L_a\) 也不必左 \(R\) 线性。这就是在 \(\operatorname{End}_{\mathbb Z}(P)\) 中构造模型，而不使用 \(\operatorname{End}_R(P)\) 的原因。直接在 \(bz^n\) 上计算得
\[
TL_a=L_{\sigma(a)}T+L_{\delta(a)}.
\]
令 \(S\) 为 \(\operatorname{End}_{\mathbb Z}(P)\) 中由全部 \(L_a\) 和 \(T\) 生成的子环。利用刚才的等式，可以在每个有限算子字中把所有 \(L_a\) 移到左侧；对 \(T^iL_a\) 按 \(i\) 归纳即可看到，每次展开仍为有限和。因此 \(S\) 中元素都具有 \(\sum_iL_{a_i}T^i\) 的形式。

由于 \(\sigma(1)=1\)、\(\delta(1)=0\)，有 \(T^i(1)=z^i\)。若 \(\sum_iL_{a_i}T^i=0\)，作用在一上便得 \(\sum_i a_iz^i=0\)，所有系数为零。这同时证明正规式唯一，以及 \(a\mapsto L_a\) 单射。令 \(R[t;\sigma,\delta]=S\)、\(t=T\)，结合律由算子复合继承。

对于泛性质，正规式唯一使公式 \(\sum_i a_it^i\mapsto\sum_i f(a_i)u^i\) 良定义。计算两个正规式的乘积时，逐次用 \(ta=\sigma(a)t+\delta(a)\) 移动系数；将每一步的 \(a,t\) 换成 \(f(a),u\)，题设等式保证同一步展开仍成立。因此该公式保持乘法，也保持加法与单位。任何延伸同态的值都由 \(R\) 及 \(t\) 确定，故唯一。
\end{proof}

这个环称为\term[Ore-kuozhang]{Ore 扩张}[Ore extension]。\(\delta=0\) 时简记为 \(R[t;\sigma]\)；\(\sigma=\operatorname{id}\) 时简记为 \(R[t;\delta]\)。上述定理给出的是系数在左侧的正规式；右侧的表达 \(\sum_i t^ib_i\) 是否存在、是否唯一，还须另证。\cref{b2:prop:ore-right-normal-form}将证明自同构情形的右正规式，并给出一般自同态情形的失败例子。
\symidx{\(R[t;\sigma,\delta]\)：关系 \(ta=\sigma(a)t+\delta(a)\) 的 Ore 扩张}

\begin{proposition}\label{b2:prop:ore-leading-degree}
设 \(R\) 为含幺环，\(\sigma:R\rightarrow R\) 为保幺环自同态，\(\delta\) 为 \(\sigma\) 导子，\(R[t;\sigma,\delta]\) 为相应 Ore 扩张。若非零元素 \(f,g\) 的最高项分别为 \(a_mt^m,b_nt^n\)，则乘积的 \(t^{m+n}\) 系数为 \(a_m\sigma^m(b_n)\)。因此，若 \(R\) 无零因子且 \(\sigma\) 单射，该 Ore 扩张也无零因子，非零乘积的次数相加。
\end{proposition}
\begin{proof}
由基本关系归纳得 \(t^m b=\sigma^m(b)t^m+\text{较低次项}\)。所以唯一可能贡献最高次项的是两个最高项的乘积，系数正如所述。单射性保证 \(\sigma^m(b_n)\ne0\)，再由 \(R\) 无零因子得到结论。
\end{proof}

不假设 \(\sigma\) 单射时，左正规式仍唯一，但环可以有零因子。例如 \(R=k[x]\)、\(\sigma(f)=f(0)\)、\(\delta=0\) 满足所有构造条件，却有 \(tx=0\)，而 \(t,x\) 都由正规式定理保证非零；另一方面，左正规单项式 \(xt\ne0\)。因此左、右乘法的行为不能从这条左正规式定理推成对称的结论。

\subsection{微分算子环与 Weyl 代数}
\label{b2:subsec:2005:03}

设 \(K\) 为带导子 \(\delta\) 的交换域。定义形式\term[weifen-suanzi-huan]{微分算子环}[differential operator ring]
\[
D=K[\partial;\delta],\qquad
\partial a=a\partial+\delta(a).
\]
元素为 \(\sum_i a_i\partial^i\)，乘法表示先作用右侧算子，再作用左侧算子。其在 \(K\) 上的自然作用为
\[
\left(\sum_i a_i\partial^i\right)(f)=\sum_i a_i\delta^i(f).
\]
这里 \(\partial\) 是 \(D\) 中的形式生成元，\(\delta\) 则是它在 \(K\) 上的作用算子。Ore 泛性质给出保幺同态 \(D\to\operatorname{End}_{\mathbb Z}(K)\)，保证这个作用与乘法相容，但不保证忠实。例如 \(\delta=0\) 时，形式变量 \(\partial\ne0\)，却在 \(K\) 上作用为零。

\ProseParagraphBreak
取 \(K=k(x)\)、\(\delta=\mathrm{d}/\mathrm{d}x\)，便有嵌入 \(A_1(k)\to K[\partial;\delta]\)，把 \(x,\mathrm{d}\) 送到 \(x,\partial\)。关系成立保证同态存在，Ore 正规式又保证 \(\sum_j f_j(x)\partial^j=0\) 时每个多项式系数都为零，所以它单射。这把多项式系数扩大为有理函数系数。正特征中仍应区别抽象形式算子与它在函数上的具体作用。

\ProseParagraphBreak
例如当 \(\delta=\mathrm{d}/\mathrm{d}x\) 时，
\[
\partial^2 a=a\partial^2+2\delta(a)\partial+\delta^2(a).
\]
这由连续两次移动系数得到。若把 \(a\) 直接与 \(\partial\) 交换，就会同时漏掉中间项和最后一项。

\subsection{左模与线性微分方程}
\label{b2:subsec:2005:04}

在 \(K\) 向量空间 \(M\) 上指定左 \(D\) 模结构，等价于指定一个加法算子 \(\nabla:M\to M\)，满足
\[
\nabla(am)=a\nabla(m)+\delta(a)m.
\]
正向取 \(\nabla(m)=\partial m\)，反向把标量作用及 \(\nabla\) 代入\cref{b2:thm:ore-extension-construction}的泛性质。配有这个算子的空间称为相对于 \(\delta\) 的\term[weifen-mo]{微分模}[differential module]。模同态就是与 \(\nabla\) 交换的 \(K\) 线性映射。\(\nabla\) 通常不是 \(K\) 线性的。

\ProseParagraphBreak
若要用循环生成元 \(e\) 记录方程 \(Le=0\)，左模作用还要求每个 \(QL\) 都消去 \(e\)。因此应模掉 \(L\) 生成的左理想 \(DL\)，从而构造左 \(D\) 模；不能无条件改用右理想 \(LD\)。

\begin{proposition}\label{b2:prop:differential-cyclic-module}
设 \(K\) 为带导子 \(\delta:K\rightarrow K\) 的交换域，\(D=K[\partial;\delta]\) 为满足 \(\partial a=a\partial+\delta(a)\) 的微分算子环，并让 \(D\) 按 \(\partial(y)=\delta(y)\) 作用于 \(K\)。取整数 \(m\geq1\) 及 \(a_0,\ldots,a_{m-1}\in K\)，令 \(L=\partial^m+a_{m-1}\partial^{m-1}+\cdots+a_0\)、\(M_L=D/DL\)，其中 \(DL\) 是左理想，则 \(M_L\) 有左 \(K\) 基 \(1,\partial,\ldots,\partial^{m-1}\) 的陪集，且有自然双射
\[
\operatorname{Hom}_D(M_L,K)\longrightarrow\bigl\{\,y\in K\mid L(y)=0\,\bigr\},
\qquad f\longmapsto f(1+DL).
\]
\end{proposition}
\begin{proof}
对任意算子 \(P\)，若最高项为 \(b_n\partial^n\)、\(n\ge m\)，就减去左倍数 \(b_n\partial^{n-m}L\)。由于 \(L\) 首一，这个倍数的最高项正是 \(b_n\partial^n\)，不必再求逆某个最高系数。每步严格降低次数，所以得到 \(P=QL+R\)，其中 \(R=0\) 或 \(\deg R<m\)。若有两个这样的余项，其差属于 \(DL\)；由\cref{b2:prop:ore-leading-degree}，非零 \(QL\) 的次数至少为 \(m\)，故余项之差只能为零。这证明所列商类为基。

左模同态由 \(1+DL\) 的像 \(y\) 决定，并必须满足 \(L(y)=0\)。反过来，若该方程成立，公式 \(P+DL\mapsto P(y)\) 与代表元无关，因为 \((QL)(y)=Q\bigl(L(y)\bigr)=0\)；它直接保持左 \(D\) 作用。两种构造互逆。
\end{proof}

\(M_L\) 记录的是方程对一个循环生成元施加的关系；解对应从 \(M_L\) 到函数模 \(K\) 的 \(D\) 模同态，并不是 \(M_L\) 的所有元素。解空间自然是常数域 \(C=\bigl\{\,a\in K\mid\delta(a)=0\,\bigr\}\) 上的向量空间：对 \(c\in C\)，有 \(\delta^i(cy)=c\delta^i(y)\)，故 \(L(cy)=cL(y)\)。一般 \(K\) 系数却会产生额外的导数项。例如方程 \(\partial y=0\) 的解包含一，但 \(\delta(a)\ne0\) 时 \(a\cdot1\) 不是解，所以不能一般地把解空间视为 \(K\) 子空间。\(C\) 确为域：导子公式保证对和、积封闭，并由 \(\delta(aa^{-1})=0\) 得到非零常数的逆仍为常数。

\ProseParagraphBreak
\(M_L\) 的 \(K\) 维数为 \(m\)，并不保证在指定的微分域中能找到 \(m\) 个独立解。例如取 \(K=\mathbb C(x)\)、\(\delta=\mathrm{d}/\mathrm{d}x\)、\(L=\partial-1\)。上述命题给出 \(\dim_KM_L=1\)，但 \(y'=y\) 没有非零有理函数解。若 \(y=p/q\ne0\)，其中 \(p,q\in\mathbb C[x]\) 均非零，则方程要求
\[
p'q-pq'=pq.
\]
左边为零或次数至多为 \(\deg p+\deg q-1\)，右边的次数却是 \(\deg p+\deg q\)，矛盾。因此这里 \(\operatorname{Hom}_D(M_L,K)=0\)。方程所对应的模已构造出来，解空间是否足够大还取决于允许解落在哪个微分域中。

\begin{proposition}\label{b2:prop:differential-companion-coordinates}
设 \(K\) 为带导子 \(\delta:K\to K\) 的交换域，\(D=K[\partial;\delta]\) 满足 \(\partial a=a\partial+\delta(a)\)，并按 \(\partial(y)=\delta(y)\) 左作用于 \(K\)。取整数 \(m\ge1\) 及 \(a_0,\ldots,a_{m-1}\in K\)，令
\[
L=\partial^m+\sum_{i=0}^{m-1}a_i\partial^i,\qquad
M_L=D/DL,\qquad e_i=\partial^i+DL\quad(0\le i<m).
\]
以 \(e_0,\ldots,e_{m-1}\) 为左 \(K\) 基，将元素写成坐标列 \(v=(v_0,\ldots,v_{m-1})^{\mathsf T}\)。令 \(A\in M_m(K)\) 的行、列按 \(0,\ldots,m-1\) 编号，条目为
\[
A_{ij}=
\begin{cases}
1,&i=j+1,\quad 0\le j<m-1,\\
-a_i,&j=m-1,\\
0,&\text{其余情形}.
\end{cases}
\]
则左乘 \(\partial\) 的坐标为 \(\nabla(v)=\delta(v)+Av\)，其中 \(\delta\) 逐坐标作用。另一方面，\(K\) 线性映射 \(f:M_L\to K\) 为 \(D\) 模同态，当且仅当其取值列 \(Y=(f(e_0),\ldots,f(e_{m-1}))^{\mathsf T}\) 满足
\[
\delta(Y)=A^{\mathsf T}Y.
\]
这些取值列恰为 \(Y=(y,\delta y,\ldots,\delta^{m-1}y)^{\mathsf T}\)，其中 \(L(y)=0\)。特别地，\(m=1\) 时 \(A=(-a_0)\)，系统为 \(\delta y=-a_0y\)。
\end{proposition}
\begin{proof}
由\cref{b2:prop:differential-cyclic-module}，所列 \(e_i\) 确为基。左乘 \(\partial\) 使 \(\partial e_j=e_{j+1}\)、\(j<m-1\)，而 \(Le_0=0\) 给出
\[
\partial e_{m-1}=-\sum_{i=0}^{m-1}a_ie_i.
\]
所以 \(A\) 的第 \(j\) 列正是 \(\partial e_j\) 的坐标。对 \(v=\sum_jv_je_j\)，Leibniz 公式给出
\[
\partial v=\sum_j\delta(v_j)e_j+\sum_jv_j\partial e_j,
\]
即坐标式 \(\delta(v)+Av\)。

若 \(f\) 为 \(D\) 模同态，则
\[
\delta\bigl(f(e_i)\bigr)=f(\partial e_i)
=\sum_j A_{ji}f(e_j),
\]
这就是 \(\delta(Y)=A^{\mathsf T}Y\)：模中基元素的作用记在 \(A\) 的各列，而同态对这些基元素的取值将同一列系数收集成系统的一行。反过来，若 \(Y\) 满足该系统，令 \(f(\sum_jv_je_j)=\sum_jv_jY_j\)，则
\[
f(\partial v)
=\sum_j\delta(v_j)Y_j+\sum_jv_j\delta(Y_j)
=\delta\bigl(f(v)\bigr).
\]
它已经保持 \(K\) 作用，又与 \(\partial\) 作用交换，而 \(D\) 由 \(K\) 和 \(\partial\) 生成，故为 \(D\) 模同态。

系统的前 \(m-1\) 行为 \(\delta(Y_i)=Y_{i+1}\)，逐次得到 \(Y_i=\delta^i(Y_0)\)；最后一行则是 \(\delta^m(Y_0)=-\sum_i a_i\delta^i(Y_0)\)，即 \(L(Y_0)=0\)。反向把任意这样的解及其连续导数代入，就满足全部系统。\(m=1\) 时没有前面的递推行，最后一行仍给出所述一阶方程。
\end{proof}

上面的转置来自模坐标与同态取值的不同约定，并不把 \(\nabla\) 变成 \(K\) 线性算子。回到一般 Ore 扩张，能否把系数放在右侧，同样取决于系数关系是否能够反向求解。

\begin{proposition}\label{b2:prop:ore-right-normal-form}
设 \(R\) 为含幺环，\(\sigma:R\to R\) 为保幺环自同构，\(\delta:R\to R\) 为加法映射，且对任意 \(a,b\in R\) 满足 \(\delta(ab)=\sigma(a)\delta(b)+\delta(a)b\)。则 Ore 扩张 \(R[t;\sigma,\delta]\) 作为右 \(R\) 模也以 \(1,t,t^2,\ldots\) 为基，每个元素唯一写成有限和 \(\sum_i t^ib_i\)、\(b_i\in R\)。

若只要求 \(\sigma\) 为保幺自同态，这个结论可能失败。具体地，设 \(k\) 为域，\(R=k\times k\)、\(\sigma(a,b)=(a,a)\)、\(\delta=0\)，令 \(e=(0,1)\)。则在 \(R[t;\sigma]\) 中 \(te=0\)、\(et\ne0\)，而 \(et\) 没有右系数表达 \(\sum_i t^ib_i\)。
\end{proposition}
\begin{proof}
将基本关系用于 \(\sigma^{-1}(a)\)，便得
\[
at=t\sigma^{-1}(a)-\delta\bigl(\sigma^{-1}(a)\bigr).
\]
这使一个系数能够从 \(t\) 的左侧移到右侧。进一步归纳可得
\[
at^n=t^n\sigma^{-n}(a)+\sum_{i<n}t^ib_i
\]
对适当的 \(b_i\in R\) 成立：将 \(at^{n+1}=(at^n)t\) 中每个右侧系数 \(b\) 按上述 \(bt\) 公式再移一次即可。因此每个左正规式都可改成右系数表达。

若有非平凡的有限关系 \(\sum_i t^ib_i=0\)，取最大的 \(n\) 使 \(b_n\ne0\)。由基本关系归纳，\(t^ib_i\) 的左系数表达最高项为 \(\sigma^i(b_i)t^i\)。故所给关系的 \(t^n\) 系数为 \(\sigma^n(b_n)\ne0\)，与\cref{b2:thm:ore-extension-construction}的左正规式唯一性矛盾。这证明右线性独立。

在第二种情形，\(\sigma\) 确为保幺环自同态，零映射满足扭曲导子公式。由 \(\sigma(e)=0\) 得 \(te=0\)，而左正规式定理给出 \(et\ne0\)。对 \(i\ge1\)，有
\[
t^i(a,b)=\sigma^i(a,b)t^i=(a,a)t^i.
\]
若 \(et=\sum_i t^ib_i\)，左正规式比较一次项就要求某个 \((a,a)\) 等于 \((0,1)\)，不可能。因此右侧表达的存在性失败；\(te=0\) 还直接说明这些幂不是右线性无关的。
\end{proof}

第三卷附录E继续介绍微分代数与消去；第五卷附录E讨论微分模及微分 Galois 理论，附录D则简介 D-模的几何观点。代数上的微分作用与Hilbert空间上的有界作用还有区别。\appref{b2:app:D}的\cref{b2:prop:bounded-weyl-impossible}证明，非零含幺 Banach 代数\footnote{巴拿赫（Stefan Banach，1892-1945），波兰数学家，建立赋范线性空间理论，著有《线性算子理论》。}中不能出现交换子等于一的两个元素。因此研究这些微分算子的分析表示时，必须另外考察无界算子及其定义域。

% !TeX root = ../../Book2.tex
% 纲要标题：### 20.6* 非交换 Gröbner 基与 Diamond 引理
% 纲要位置：计划纲要.md，第 1731 行。

% 纲要要点（供目录核对）：
%
% * 自由结合代数中的可容许字序、重写规则及包含和重叠歧义；下降次序须与两侧连接相容，两个最小失败例的实际商在正文命题完整证明
% * Diamond 引理与正规字；分别证明终止、首步选择无关及正规映射核等于关系理想，不把交换 Buchberger 判据原样照搬
% * PBW 型代数及可解型多项式代数中的有效约化；区别单次约化终止、全部歧义可解、规则有效匹配与补全终止，有限展示不自动给出有限 Gröbner 基
% * 建议先读第一卷附录 F.1—F.3 的交换 Gröbner 基；第三卷第8章发展计算应用、附录 F 处理局部标准基；正文用同一理想的两种字序证明有限基依赖字序，原题改为乘法、中心与单性计算，不把维数数据当作乘法结构
% ---
%

\OptionalSection{非交换 Gröbner 基与 Diamond 引理}{b2:sec:2006}

第一卷附录F.1—F.3讨论交换多项式的除法、首项理想和Buchberger 算法\footnote{布赫贝格尔（Bruno Buchberger，1942-），奥地利数学家，提出 Gröbner 基算法。}。本节在这些概念的背景下重新定义非交换字的次序与重写规则，说明字母顺序带来的差别。

\ProseParagraphBreak
正规式不仅要能够算出，还须证明不同计算次序给出相同结果。本节在自由结合代数的字上建立一个首一重写版本的 Diamond 引理。规则下降保证每次约化过程终止，歧义可解才保证不同选择得到相同正规式。要把它用于算法，还须说明规则如何匹配、系数如何计算；有限条原始关系也不保证补全后仍只有有限条规则。这些是不同的要求。


\subsection{可容许字序、重写规则与歧义}
\label{b2:subsec:2006:01}

令 \(X\) 为有限字母集，\(X^*\) 为全部有限字，空字记为一。自由结合代数 \(k\langle X\rangle\) 以这些字为基，乘法为连接。

\begin{definition}\label{b2:def:admissible-word-order}
设 \(X\) 为有限字母集，\(X^*\) 为由 \(X\) 构成的全部有限字的集合，空字记为 \(1\)。若 \(X^*\) 上的全序 \(<\) 为良序，以空字为最小元，且对任意字 \(a,b,u,v\in X^*\)，\(u<v\) 蕴含 \(aub<avb\)，则称 \(<\) 为\term[kerongxu-zixu]{可容许字序}[admissible word order]。例如先比较字长，再按固定的字母次序作字典比较，就满足这些条件。
\end{definition}

对一个非零多项式 \(f\)，其最大的字称为\term[shou-danxiangshi]{首单项式}[leading monomial]，记为 \(\operatorname{LM}(f)\)。若其系数为一，称 \(f\) 首一。规定一族规则 \(\ell_\alpha\to r_\alpha\)，其中 \(\ell_\alpha\ne1\)，右侧为有限线性组合且其中每个字都严格小于 \(\ell_\alpha\)。这相当于使用首一关系 \(g_\alpha=\ell_\alpha-r_\alpha\)。

\ProseParagraphBreak
一次重写将多项式中一个完整项 \(c\,a\ell_\alpha b\) 换成 \(c\,ar_\alpha b\)，然后合并同类项。字若不含任何 \(\ell_\alpha\) 作为连续子字，称为\term[zhenggui-zi]{正规字}[normal word]。右侧的字小于 \(\ell_\alpha\)，还必须在加入任意前缀 \(a\)、后缀 \(b\) 后继续小于 \(a\ell_\alpha b\)。可容许字序的双侧相容性正保证这件事；仅有良序而没有这项相容性，局部的下降不能保证大字内部的重写也下降。

\ProseParagraphBreak
同一字中若有两处可替换，可能出现两类\term[chongxie-qiyi]{歧义}[rewriting ambiguity]。一类是包含：\(\ell_\alpha=a\ell_\beta b\)，可先改整个左端，也可先改其中一段；另一类是重叠：\(\ell_\alpha=au\)、\(\ell_\beta=ub\)，其中 \(a,u,b\) 都非空，字 \(aub\) 有两种第一步。相同左端的不同规则也算包含歧义。若两种第一步结果能继续重写为同一个多项式，称这项歧义可解。

\ProseParagraphBreak
例如规则 \(x^2\to1\)、\(yx\to0\) 在 \(yxx\) 上重叠：改末尾 \(xx\) 得到 \(y\)，改开头 \(yx\) 却得到零。规则 \(xyx\to x\)、\(yx\to0\) 则在 \(xyx\) 上包含：改整个字得到 \(x\)，改其中的末尾 \(yx\) 得到零。两组原始规则都下降，但各自的两支已不能继续合到一起。\cref{b2:prop:diamond-ambiguity-failures}将计算这两组关系实际定义的商代数。

\ProseParagraphBreak
两处彼此分离时不形成新的歧义：在字 \(a\ell_\alpha b\ell_\beta c\) 中，先改哪一处，再改另一处，都会得到 \(ar_\alpha br_\beta c\)。这里利用有限和的分配律，不交换任何字母。

\subsection{Diamond 引理与正规字}
\label{b2:subsec:2006:02}

\begin{theorem}[Diamond 引理，首一字重写版本]\label{b2:thm:diamond-lemma}
设 \(k\) 为域，\(X\) 为有限字母集，\(<\) 为 \(X^*\) 上的可容许字序。给定一族首一字重写规则 \(\ell_\alpha\rightarrow r_\alpha\)，其中每个 \(\ell_\alpha\) 都是非空字，\(r_\alpha\) 是严格小于 \(\ell_\alpha\) 的字的有限 \(k\) 线性组合。每次重写在多项式中将一个完整项 \(c\,a\ell_\alpha b\) 换为 \(c\,ar_\alpha b\)，其中 \(c\in k\)，\(a,b\in X^*\)，然后合并同类项，则从任意多项式出发的每个重写过程都在有限步后停止。若全部包含与重叠歧义可解，则每个多项式有唯一的正规字线性组合形式；这里正规字指不含任何 \(\ell_\alpha\) 为连续子字的字。令 \(I\) 为全部 \(\ell_\alpha-r_\alpha\) 生成的双边理想，则正规字的商类构成 \(k\langle X\rangle/I\) 的一组 \(k\) 基。
\end{theorem}

这项结果称为\term[Diamond-yinli]{Diamond 引理}[Diamond lemma]。这里证明的是自由结合代数上采用首一规则与可容许全序的专门形式；原论文还讨论更一般的约化系统\cite{Bergman1978DiamondLemma}。只写若干等式而不给下降次序，还不能使用它。

\begin{proof}
每次重写将一个完整项换为有限多个更小的字，下降依据是良序，右侧的有限性则保证仍得到有限多项式。当前最大字未必每步都下降，因为也可以先改较小项；要证明任意选择都终止，对有限多项式的最大字 \(w\) 作良序归纳。只要 \(w\) 尚未被替换，对较小项进行的操作都是从一个有限的较小字多项式出发，按归纳假设不可能无限进行。较小字也不能产生 \(w\)。若 \(w\) 可替换，把它的整个系数项换掉后，余下所有字都小于 \(w\)，再次由归纳假设终止；若 \(w\) 不可替换，只需处理下面各项。最小字一没有可用规则，提供归纳起点。合并同类项造成的消去只会删除项，不会破坏这个下降论证。

再对单个字 \(w\) 作良序归纳证明正规结果唯一，并记为 \(N(w)\)。假设所有更小的字已经有唯一正规结果，将 \(N\) 线性延伸到它们的有限和上。若 \(w\) 正规，规定 \(N(w)=w\)。否则，选择一次替换得到多项式 \(f\)，其各字小于 \(w\)，候选结果为 \(N(f)\)。须证明它不依赖第一步。

比较两种第一步。如果两处互不相交，继续改另一处会得到上一小节写出的同一个表达式，因此按归纳假设两者的正规结果相同。如果两处包含或重叠，它们在某个连续子字中构成一项题设歧义。该歧义的两支能化为共同结果；在两侧加上原字的相同前缀和后缀，仍是合法的重写序列。所有这些字在第一步以后都严格小于 \(w\)，所以归纳假设再次保证两支的正规结果相同。这就排除了第一步的全部选择，并完成归纳。

由线性延伸，得到到正规字线性生成空间的映射 \(N\)。每次重写都保持 \(N\)，所以每个终止结果等于 \(N(f)\)。又每次重写之差属于 \(I\)，因而 \(f-N(f)\in I\)，这已经保证正规字的商类生成商代数。为证明独立，还要排除正规字之间藏在 \(I\) 中的非零关系。对任意字 \(a,b\) 和规则，有 \(N(a\ell_\alpha b)=N(ar_\alpha b)\)，所以 \(N\) 零化 \(a(\ell_\alpha-r_\alpha)b\)。这些元素线性生成 \(I\)，故 \(N(I)=0\)，即 \(I\subseteq\ker N\)；而 \(N(f)=0\) 时，\(f=f-N(f)\in I\)，给出反向包含。因此 \(\ker N=I\)。正规字上的 \(N\) 是恒等映射，所以它们的非零线性组合不可能属于 \(I\)，得到独立性。
\end{proof}

例如选择 \(y>x\) 的长度字典序，规则 \(yx\to xy+1\) 没有自身的包含或重叠歧义：左端的非空真后缀为 \(x\)，非空真前缀为 \(y\)，它们不同。正规字恰为 \(x^iy^j\)，所以引理给出 Weyl 代数的正规单项式基的另一证明。规则 \(yx\to qxy\) 同样处理量子平面；这里依赖的是字序和歧义检查。

\subsection{PBW 型代数的有效约化}
\label{b2:subsec:2006:03}

当规则有限且系数可精确计算时，可以固定一种约化程序：每步取最大的可约字，再选取其中最左的可约位置；若仍有多条规则，按事先编号选择。上一定理保证停止。若歧义条件也成立，输出就与这个具体选择无关，并且
\[
f\in I\quad\Longleftrightarrow\quad N(f)=0.
\]
因此正规式可以用于相等判定与理想成员判定。终止性本身没有给出统一的时间复杂度界。

\ProseParagraphBreak
例如 \(yx\to xy+x^2\) 在 \(y>x\) 的长度字典序下严格下降，且仍没有自身歧义，因此 \(x^iy^j\) 也构成正规单项式基。计算得
\[
yx^2=x^2y+2x^3,\qquad
y^2x=xy^2+2x^2y+2x^3.
\]
这个代数也可用\cref{b2:thm:ore-extension-construction}构造：取 \(R=k[x]\)、\(\sigma=\operatorname{id}\)、\(\delta(f)=x^2f'\)，再把扩张变量记为 \(y\)。两种方法一个检查重写歧义，一个通过导子建立乘法，得到同一条关系及同一组基。

\begin{definition}\label{b2:def:noncommutative-groebner-basis}
设 \(k\) 为域，\(X\) 为有限字母集，\(<\) 为 \(X^*\) 上的可容许字序，\(I\subseteq k\langle X\rangle\) 为双边理想，\(G\subseteq I\setminus\{0\}\) 为 \(I\) 的一组生成元。若每个非零 \(f\in I\) 的首单项式都含某个 \(g\in G\) 的首单项式为连续子字，则称 \(G\) 为 \(I\) 关于 \(<\) 的\term[feijiaohuan-Groebner-ji]{非交换 Gröbner 基}[noncommutative Gröbner basis]。
\end{definition}

满足 Diamond 引理条件的首一关系族就是这样的基。事实上，若 \(f\in I\) 的最大字不可约，它在任何对较小字的重写中都不可能被消去，所以 \(N(f)\ne0\)，与 \(N(I)=0\) 矛盾。

\ProseParagraphBreak
若某项歧义产生两个不同的正规结果，它们之差是理想中的新关系；可以将差的首项归一化，加入新的下降规则，再检查新出现的歧义。这种补全同样利用尚未消去的关系来扩大首项条件，但非交换字的歧义来自连续包含与连续重叠，不能用交换单项式的最小公倍数替代。即使原规则有限，每次新增的左端仍可能产生下一项歧义，不能由单次约化会停止推断整个补全过程也会停止。

\subsection{非交换 Gröbner 基的有限性问题}
\label{b2:subsec:2006:04}

有限条原始关系不保证只需补充有限多条规则，甚至不保证在给定字序下存在任何有限 Gröbner 基。下面对一个仅有两个生成元、一条关系的代数证明这一点：不仅所写的补全过程无穷延续，其首字反链还将排除所有有限替代。不过这个结论依赖指定的字序，换序后同一理想可以有有限 Gröbner 基。

\ProseParagraphBreak
在 \(k\langle x,y\rangle\) 中，取 \(x>y\) 的长度字典序，从原关系 \(x^2-xy\) 得到规则 \(x^2\to xy\)。对同一个字 \(x^3\)，两种第一步产生分叉：
\[
(x^2)x\longrightarrow xyx,
\qquad
x(x^2)\longrightarrow x^2y\longrightarrow xy^2.
\]
两支的末项都已不能用原规则继续约化，却应代表同一个商类。为使这项歧义可解，须补上 \(xyx\to xy^2\)。

\ProseParagraphBreak
这种补充会继续下去。若已补到规则 \(xy^nx\to xy^{n+1}\)，它与原规则在 \(xy^nx^2\) 上又有两种约化：
\[
\begin{aligned}
(xy^nx)x&\longrightarrow xy^{n+1}x,\\
xy^n(x^2)&\longrightarrow xy^nxy\longrightarrow xy^{n+2}.
\end{aligned}
\]
于是下一条规则应是 \(xy^{n+1}x\to xy^{n+2}\)。这些左端中的两个 \(x\) 相隔越来越多的 \(y\)，以前的规则不能直接改写新的左端。下面证明这一过程产生的全部关系确实属于原理想，并形成无限 Gröbner 基。

\begin{proposition}\label{b2:exm:infinite-noncommutative-groebner}
设 \(k\) 为域，在 \(k\langle x,y\rangle\) 中令 \(I=(x^2-xy)\)。采用 \(x>y\) 的长度字典序时，该理想没有有限非交换 Gröbner 基。一组无限首一 Gröbner 基为
\[
g_n=xy^nx-xy^{n+1},\qquad n\ge0.
\]
商代数的正规字为 \(y^a\) 及 \(y^axy^b\)、\(a,b\ge0\)。改用 \(y>x\) 的长度字典序时，单条关系 \(xy-x^2\) 就是有限首一 Gröbner 基，相应正规字为 \(y^ax^b\)、\(a,b\ge0\)。
\end{proposition}
\begin{proof}
\(g_0=x^2-xy\) 是原关系，而直接展开给出
\[
g_{n+1}=g_ny-g_nx+xy^n g_0.
\]
由归纳，全部 \(g_n\in I\)，所以它们生成的双边理想仍恰为 \(I\)。规则 \(xy^nx\to xy^{n+1}\) 保持字长，最后一个 \(x\) 改成较小的 \(y\)，因此下降。

不同左端不互相包含，因为它们各有且只有两个 \(x\)。非平凡重叠只能共享这个末端与另一个左端的首个 \(x\)，得到字 \(xy^axy^bx\)。先改左边，得到 \(xy^{a+b+1}x\)，再改为 \(xy^{a+b+2}\)；先改右边，得到 \(xy^axy^{b+1}\)，再改左边，也得到 \(xy^{a+b+2}\)。因此所有歧义可解。\cref{b2:thm:diamond-lemma}给出正规字基，而不可约字恰好是至多含一个 \(x\) 的字，即所列形式。

这也确定了理想的全部首单项式：它们恰为包含某个 \(xy^nx\) 的字。一个方向由每个 \(g_n\in I\) 及两边乘字得到；另一个方向由正规式判定得到，因为理想中的非零元素不可能有不可约的最大字。

各个字 \(xy^nx\) 对连续子字包含关系构成无限反链，因而理想的这些最小首字不能由有限多个其他首字统一覆盖。具体地，若存在有限 Gröbner 基，对每个 \(n\)，其中某个首单项式必须是 \(xy^nx\) 的子字；而它又必须包含某个 \(xy^mx\)。因为两个 \(x\) 必须同时保留，只能有 \(m=n\)，且该首单项式就是整个 \(xy^nx\)。于是基中必须为每个 \(n\) 提供一个不同的首单项式，与有限性矛盾。

改用 \(y>x\) 的长度字典序后，\(xy-x^2\) 的首字为 \(xy\)，规则 \(xy\to x^2\) 严格下降。它没有非平凡的自身包含歧义，也没有重叠歧义，因为唯一的非空真后缀 \(y\) 与唯一的非空真前缀 \(x\) 不同。因此\cref{b2:thm:diamond-lemma}给出正规字基，并由前面首项判定说明这条关系构成 Gröbner 基。不含连续子字 \(xy\) 的字只能先写若干个 \(y\)，再写若干个 \(x\)，恰为 \(y^ax^b\)。这证明同一理想的有限 Gröbner 基存在性确实可以随字序改变。
\end{proof}

上述理想在 \(x>y\) 字序下的无限规则仍可用于计算：给定长度为 \(N\) 的字，只可能匹配 \(0\le n\le N-2\) 的左端 \(xy^nx\)。沿字从左到右寻找相邻两个 \(x\) 之间全由 \(y\) 组成的连续子字，就能判定可约位置并执行下一步；对有限多项式逐项检查即可。每步都在可容许字序下下降，所以在系数运算有效时得到终止的正规式计算。“没有有限 Gröbner 基”不等于“任何计算都不可能”。这里证明有限展示不保证给定字序下有有限 Gröbner 基，同时另行证明了这组无限规则可以有效匹配；不能仅由有限展示推得后一性质。第一卷附录F已经给出交换 Gröbner 基的有限性及基本算法；第三卷第8章进一步发展交换情形的计算与几何应用，第三卷附录F则处理局部标准基。使用时须分别核对对象、字序或项序，以及终止性的实际依据。

\begin{proposition}\label{b2:prop:diamond-ambiguity-failures}
设 \(k\) 为域。在 \(k\langle x,y\rangle\) 中，规则 \(x^2\to1\)、\(yx\to0\) 都按长度字典序下降，却在 \(yxx\) 上有不可解的重叠歧义，且
\[
k\langle x,y\rangle/(x^2-1,yx)\cong k[s]/(s^2-1).
\]
在这个商中 \(y=0\)，基为 \(1,x\) 的商类。规则 \(xyx\to x\)、\(yx\to0\) 也都下降，却在 \(xyx\) 上有不可解的包含歧义，且
\[
k\langle x,y\rangle/(xyx-x,yx)\cong k[s].
\]
在这个商中 \(x=0\)，基为 \(1,y,y^2,\ldots\) 的商类。
\end{proposition}
\begin{proof}
第一组规则在 \(yxx\) 上，一条约化先改 \(xx\) 而得到 \(y\)，另一条先改 \(yx\) 而得到零；\(y\) 与零对于原规则都已正规，所以歧义不可解。令 \(A=k\langle x,y\rangle/(x^2-1,yx)\)，暂将生成元的商类仍记为 \(x,y\)。结合律与两条关系给出
\[
y=yx^2=(yx)x=0.
\]
向 \(k[s]/(s^2-1)\) 发送 \(x\mapsto s\)、\(y\mapsto0\)，两条关系均变成零，故得到同态 \(A\to k[s]/(s^2-1)\)。反向以 \(s\mapsto x\) 定义同态，因 \(x^2=1\) 而良定义。两个复合保持生成元；在 \(A\) 中，已经证明 \(y=0\)，所以这两个同态互逆。首一多项式 \(s^2-1\) 的带余除法给出 \(1,s\) 的商类为基，得到 \(A\) 中所列的基。这个论证在特征二时也成立。

第二组规则中，改整个 \(xyx\) 得到 \(x\)，改末尾 \(yx\) 却得到零；\(x\) 与零对原规则仍都正规。令 \(B=k\langle x,y\rangle/(xyx-x,yx)\)，其关系给出
\[
x=xyx=x(yx)=0.
\]
发送 \(x\mapsto0\)、\(y\mapsto s\) 得到同态 \(B\to k[s]\)，反向以 \(s\mapsto y\) 得到同态 \(k[s]\to B\)。两个复合都保持生成元，其中 \(B\) 内的 \(x=0\)，所以它们互逆。\(k[s]\) 的幂基给出所列商类为基。两组失败都产生了原先终止结果未显示的额外关系，因此原规则的终止性没有保证正规字的商类线性独立。
\end{proof}


\begin{chaptersummary}
滤过保留“不超过给定次数”的信息，伴随分次取每一层的新部分。穷尽性保证元素都有有限滤过次数，下界使逐次消去主符号的过程停止；由伴随分次提升 Noether 性时，子模须用诱导滤过。Rees 代数的零纤维为伴随分次，非零纤维恢复原代数，但平坦性不要求纤维具有相同乘法或相同零因子性质。粗化的多项式滤过产生平方零主符号，已经表明原代数为整环不保证伴随分次也是整环。

\ProseParagraphBreak
有序单项式的生成性与独立性分别需要证明。Weyl 关系的重排给出生成性，辅助变量作用 \(x^i\mathrm d^j(1)=X^iY^j\) 给出独立性；量子平面的扭曲乘法直接构造具有独立基的代数；Ore 构造用 \(T^i(1)=z^i\) 读取左正规式的系数。右正规式还须另查 \(\sigma\) 的可逆性。Weyl 代数的中心在特征 \(p\) 时扩大，中心上的秩为 \(p^2\)，而量子平面的中心由参数的精确阶决定，有限阶 \(\ell\) 时相应秩为 \(\ell^2\)。特征零的 Weyl 代数无非零有限维幺性表示；正特征则可由固定中心值的商得到矩阵代数。

\ProseParagraphBreak
Lie 括号的 Jacobi 恒等式保证伴随作用是表示，泛包络代数把括号关系转为结合乘法，使 Lie 表示对应于幺性左模。Heisenberg 的中心作用为零时，参数商是交换多项式环；作用为非零标量时，商经缩放成为 Weyl 代数。一般 PBW 定理确定 \(\operatorname{gr}U(\mathfrak g)\cong\operatorname{Sym}(\mathfrak g)\)，从而保证有序单项式的独立性。仅把写出的关系取最高次部分，不能代替这个同构的证明：坏 PBW 关系的 \(zyx\) 分叉与 Jacobi 失效产生同一个隐藏关系 \(x=0\)，实际伴随分次也会记录这一坍缩；剩余代数仍可构造正规基。

\ProseParagraphBreak
线性微分方程由左理想商 \(M_L=D/DL\) 编码，函数解是 \(\operatorname{Hom}_D(M_L,K)\) 的元素。模元素的坐标作用要保留系数求导项，模同态的取值系统则使用转置矩阵。Diamond 引理中，下降字序保证约化终止，歧义可解保证正规结果唯一；补充规则是否终止是另一问题。理想 \((x^2-xy)\) 在 \(x>y\) 的字序下无有限 Gröbner 基，换成 \(y>x\) 时一条规则已足够。商的每次分次维数虽与二元多项式环相同，仍有 \(x(x-y)=0\)：维数数据无法单独恢复乘法结构。
\end{chaptersummary}

\BookExercises

\begin{optionalexercise}\label{b2:ex:20-coarse-filtration}
在 \(A=k[x]\) 上取 \(F_nA=\{f\mid\deg f\le2n\}\)，负次滤过项为零。证明其 Rees 代数有展示
\[
\mathcal R_F(A)\cong k[Z,U,V]/(U^2-ZV),
\qquad Z\mapsto z,\quad U\mapsto xz,\quad V\mapsto x^2z.
\]
直接计算 \(Z=0\) 与 \(Z=1\) 的两个商，并解释平方零元为何只出现在第一个商中。
\end{optionalexercise}
\begin{solution}
每个 Rees 单项式 \(x^iz^n\) 满足 \(i\le2n\)。写 \(i=2b+a\)、\(a\in\{0,1\}\)，则 \(n\ge b+a\)，所以它等于 \(U^aV^bZ^{n-b-a}\) 的像。关系 \(U^2=ZV\) 把每个交换单项式化为 \(Z^cU^aV^b\)、\(a=0,1\)。这些元素的像是 \(x^{a+2b}z^{a+b+c}\)；不同三元组给出不同单项式，因为 \(x\) 次数先确定 \(a,b\)，再由 \(z\) 次数确定 \(c\)。故所给满同态也单射。

模 \(Z\) 后得到 \(k[U,V]/(U^2)\)，其中 \(U\ne0\) 而 \(U^2=0\)；它就是\cref{b2:prop:filtered-domain-converse-failure}中的伴随分次。模 \(Z-1\) 后关系成为 \(V=U^2\)，得到 \(k[U]\cong k[x]\)。原 Rees 代数中 \(U^2=ZV\ne0\)，取零纤维才把右侧杀掉；因此平坦族的不同纤维可以有不同的零因子性质。
\end{solution}

\begin{optionalexercise}\label{b2:ex:20-polynomial-rees}
在 \(k[x,y]\) 中赋予 \(x,y\) 权重一、二，以权重不超过 \(n\) 的单项式张成 \(F_nA\)。证明
\[
\mathcal R_F(A)=k[z,xz,yz^2]\cong k[Z,U,V],
\qquad \deg Z=\deg U=1,\quad\deg V=2.
\]
辨认零纤维与一纤维，并说明零纤维的分次如何记录原来两个变量的不同权重。
\end{optionalexercise}
\begin{solution}
Rees 单项式为 \(x^iy^jz^n\)、\(i+2j\le n\)，故等于 \((xz)^i(yz^2)^jz^{n-i-2j}\)。不同的 \(Z^aU^iV^j\) 映到不同的 \(x^iy^jz^{a+i+2j}\)，这给出所述多项式环同构和次数。模 \(Z\) 后得到分次环 \(k[U,V]\)，其中 \(U,V\) 分别对应 \(x\) 的一次主符号与 \(y\) 的二次主符号；模 \(Z-1\) 后得到原环 \(k[x,y]\)。两纤维作为未分次的环都是二元多项式环，但零纤维保留的权重分次并不是两个变量都为一次的通常分次。
\end{solution}

\begin{optionalexercise}\label{b2:ex:20-filtered-isomorphism}
设 \(k\) 为域，\(M,N\) 为具有非负、穷尽滤过的 \(k\) 向量空间，\(k\) 线性映射 \(f:M\to N\) 满足 \(f(F_nM)\subseteq F_nN\)。若 \(\operatorname{gr}f\) 是同构，证明 \(f\) 是同构，且 \(f(F_nM)=F_nN\) 对每个 \(n\) 成立。
\end{optionalexercise}
\begin{solution}
若 \(0\ne m\in\ker f\)，取其最低滤过次数，则它的非零主符号被 \(\operatorname{gr}f\) 送到零，矛盾。故 \(f\) 单射。对 \(y\in F_nN\)，商片上的满射性给出 \(x\in F_nM\)，使 \(y-f(x)\in F_{n-1}N\)。归纳提升这个差值，在 \(F_{-1}N=0\) 处结束，就得 \(y\in f(F_nM)\)。穷尽性再使 \(f\) 满射。整个论证同时证明每一滤过项上的满射性。
\end{solution}

\begin{optionalexercise}\label{b2:ex:20-weyl-normal-ordering}
在任意特征的 Weyl 代数中证明
\[
\mathrm{d}^mx^n=\sum_{r=0}^{\min(m,n)}\binom mr n(n-1)\cdots(n-r+1)x^{n-r}\mathrm{d}^{m-r},
\]
其中 \(r=0\) 的下降乘积为一，整数系数取其在 \(k\) 中的像。
\end{optionalexercise}
\begin{solution}
\(m=0\) 时成立。由 \(\mathrm{d}f(x)=f(x)\mathrm{d}+f'(x)\)，对第 \(m\) 步的每项再左乘 \(\mathrm{d}\)，得到保持导数次数的一项及增加一次导数的一项。相同 \(r\) 的系数由 \(\binom mr+\binom m{r-1}=\binom{m+1}r\) 合并，下降乘积的新增因子为 \(n-r+1\)。这正给出第 \(m+1\) 步。证明只用整数恒等式及形式导数，没有把可能为零的阶乘作除数，故对任意特征成立。
\end{solution}

\begin{optionalexercise}\label{b2:ex:20-weyl-finite-representations}
设 \(\operatorname{char}k=p>0\)，\(V\ne0\) 为有限维幺性左 \(A_1(k)\) 模，且 \(x^pV=\mathrm d^pV=0\)。证明 \(V\) 是 \(k[X]/(X^p)\) 的若干个直和；说明直和项数如何由 \(\dim_kV\) 决定，以及 \(V\) 何时简单。
\end{optionalexercise}
\begin{solution}
由\cref{b2:prop:weyl-finite-dimensional-representations}，作用经过 \(A_1(k)/(x^p,\mathrm d^p)\cong M_p(k)\)。记这个矩阵环的矩阵单位为 \(E_{ij}\)，其中 \(0\le i,j<p\)，令 \(W=E_{00}V\)。恒等式 \(1=\sum_iE_{ii}\) 给出 \(V=\bigoplus_iE_{ii}V\)，而 \(E_{i0}:W\to E_{ii}V\) 与 \(E_{0i}\) 互逆。因此
\[
k^p\otimes_kW\longrightarrow V,\qquad e_i\otimes w\longmapsto E_{i0}w
\]
是 \(M_p(k)\) 模同构，矩阵环只作用于第一个因子。正文的截断多项式表示正是标准列模 \(k^p\)。取 \(W\) 的一组基，得到 \(V\cong\bigl(k[X]/(X^p)\bigr)^{\oplus r}\)，其中 \(r=\dim_kW=\dim_kV/p\)。标准列模简单；若 \(r>1\)，一个直和项给出真非零子模，所以 \(V\) 简单当且仅当 \(r=1\)。
\end{solution}

\begin{optionalexercise}\label{b2:ex:20-weyl-characteristic-p-matrix}
设 \(\operatorname{char}k=p>0\)，\(a,b\in k\)。证明
\[
A_1(k)/(x^p-a^p,\mathrm d^p-b^p)\cong M_p(k).
\]
在 \(k[X]/(X^p)\) 上写出这个商的简单表示，并求其作为 \(A_1(k)\) 表示的核。
\end{optionalexercise}
\begin{solution}
置 \(u=x-a\)、\(v=\mathrm d-b\)，则 \(vu-uv=1\)。标量与生成元交换，故特征 \(p\) 的二项式公式给出 \(u^p=x^p-a^p\)、\(v^p=\mathrm d^p-b^p\)。代换与逆代换 \(x=u+a,\mathrm d=v+b\) 因而把所给商同构到 \(A_1(k)/(u^p,v^p)\)，由\cref{b2:prop:weyl-finite-dimensional-representations}它是 \(M_p(k)\)。

相应作用为 \(x\cdot f=(X+a)f\)、\(\mathrm d\cdot f=f'+bf\)。减去标量后恢复正文的截断多项式表示，因此这个表示简单。商矩阵环在标准列模上的作用忠实，故原代数的作用核恰为双边理想 \((x^p-a^p,\mathrm d^p-b^p)\)。这里只处理中心值已经是 \(k\) 中 \(p\) 次幂的纤维，未假定任意中心值都能作这样的变量平移。
\end{solution}

\begin{optionalexercise}\label{b2:ex:20-weyl-center-module-rank}
设 \(\operatorname{char}k=p>0\)，\(C=Z(A_1(k))\)、\(F=\operatorname{Frac}(C)\)。证明 \(F\otimes_CA_1(k)\) 是中心为 \(F\)、维数为 \(p^2\) 的除代数，因而不是 \(M_p(F)\)。
\end{optionalexercise}
\begin{solution}
由\cref{b2:thm:weyl-center-simplicity}，\(C=k[x^p,\mathrm d^p]\)，且 \(A_1(k)\) 在 \(C\) 上以 \(x^i\mathrm d^j\)、\(0\le i,j<p\)，为自由基。原代数因而嵌入中心局部化；该局部化以同一组元素为 \(F\) 基，维数为 \(p^2\)。原 Weyl 代数是整环；中心非零元作分母后仍是整环，因为两个非零分式的乘积有非零分子。

有限维域代数中，非零元素的左乘在整环条件下单射，故满射，给出右逆；整环条件又使该右逆同时为左逆。因此这个局部化是除代数。若 \(c=\sum_{i,j<p}c_{ij}x^i\mathrm d^j\) 居中，所有系数属于 \(F\)，与 \(\mathrm d\) 取交换子使所有 \(i>0\) 的系数为零，与 \(x\) 取交换子再使所有 \(j>0\) 的系数为零；这里 \(1,\ldots,p-1\) 在 \(k\) 中都非零。故中心恰为 \(F\)。而 \(p\ge2\)，\(M_p(F)\) 有非零零因子，所以不能与这个除代数同构。有限秩和相同维数本身不决定中心上的代数是否分裂。
\end{solution}

\begin{optionalexercise}\label{b2:ex:20-quantum-units-ideals}
求量子平面 \(k_q[x,y]\) 的单位，并证明它对任意 \(q\ne0\) 都不是单环。特别解释中心等于 \(k\) 为什么不足以推出单性。
\end{optionalexercise}
\begin{solution}
由\cref{b2:prop:quantum-plane-domain}，非零乘积的 \(y\) 次数相加。若两个元素互逆，它们都不含 \(y\)，于是属于 \(k[x]\)，只能是非零常数。因此单位群为 \(k^\times\)。关系使 \(Ax=xA\)，双边理想 \((x)\) 由 \(x\) 指数为正的正规单项式生成，非零且不含一；商为 \(k[y]\)。所以环不单。参数无限阶时中心虽为 \(k\)，这仍不改变该真双边理想的存在。
\end{solution}

\begin{optionalexercise}\label{b2:ex:20-quantum-binomial}
在量子平面中把 \((x+y)^3\) 写成正规式。若 \(q\ne1\) 且 \(1+q+q^2=0\)，证明 \((x+y)^3\) 属于中心。
\end{optionalexercise}
\begin{solution}
先有 \((x+y)^2=x^2+(1+q)xy+y^2\)。再乘 \(x+y\)，用 \(yx=qxy\)、\(y^2x=q^2xy^2\)，得到
\[
(x+y)^3=x^3+(1+q+q^2)x^2y+(1+q+q^2)xy^2+y^3.
\]
题设使 \(q^3=1\) 且 \(q\ne1\)，所以参数的阶为三，式子化为 \(x^3+y^3\)。由\cref{b2:prop:quantum-plane-center}，两项均居中。
\end{solution}

\begin{optionalexercise}\label{b2:ex:20-quantum-inverse-parameter}
证明交换两个生成元给出 \(k_q[x,y]\cong k_{q^{-1}}[x,y]\)。说明它与中心阶数公式相容，但不要据此断言同阶的任意两个参数都给出同构代数。
\end{optionalexercise}
\begin{solution}
在目标代数中记生成元为 \(X,Y\)，有 \(YX=q^{-1}XY\)。赋值 \(x\mapsto Y,y\mapsto X\) 将原关系送到 \(XY-qYX=0\)，所以诱导同态。反向仍交换两生成元，两复合固定生成元，因而互逆。\(q\) 与 \(q^{-1}\) 的阶相同，故中心中的幂指数一致；这个构造只比较互逆的参数，没有处理其他同阶参数。
\end{solution}

\begin{optionalexercise}\label{b2:ex:20-weyl-rees-fibres}
固定正整数 \(r,s\)，在 Weyl 代数中令 \(x,\mathrm d\) 的权重分别为 \(r,s\)，以正规单项式 \(x^i\mathrm d^j\) 的权重 \(ri+sj\) 定义滤过。证明其 Rees 代数由 \(Z=z,X=xz^r,D=\mathrm dz^s\) 生成，且有展示
\[
k\langle X,D,Z\rangle/(ZX-XZ,ZD-DZ,DX-XD-Z^{r+s}).
\]
证明 \(X^iD^j\) 构成它在 \(k[Z]\) 上的一组基，并辨认各个 \(Z=\lambda\) 的纤维。
\end{optionalexercise}
\begin{solution}
Weyl 重排中的每个交换子项都将权重降低 \(r+s\)，故所给空间确为乘法相容的滤过。Rees 单项式 \(x^i\mathrm d^jz^n\)、\(ri+sj\le n\)，是 \(X^iD^jZ^{n-ri-sj}\) 的像。关系的像为零，反复移动 \(Z\) 并使用 \(DX=XD+Z^{r+s}\)，使任意字成为 \(X^iD^jZ^a\) 的线性组合：按 \(X,D\) 的字长及逆序数作下降归纳即可。它们的像 \(x^i\mathrm d^jz^{ri+sj+a}\) 由 Weyl 正规基及 \(z\) 次数线性无关，故没有其他关系，并得 \(k[Z]\) 自由基。

取纤维后仍以 \(X^iD^j\) 为 \(k\) 基，关系变为 \(DX-XD=\lambda^{r+s}\)。\(\lambda=0\) 时是交换多项式环；\(\lambda\ne0\) 时把 Weyl 的 \(x,\mathrm d\) 分别送到 \(X,\lambda^{-(r+s)}D\)，逆向将 \(X,D\) 送到 \(x,\lambda^{r+s}\mathrm d\)，得到互逆同构。指数 \(r+s\) 来自原关系最高权重，而不是参数的一次幂。
\end{solution}

\begin{optionalexercise}\label{b2:ex:20-bad-pbw-quotient}
固定整数 \(r\ge2\)，令 \(B_r=k\langle y,z\rangle/(zy-yz-y^r)\)。证明 \(B_r\cong k[y][z;\delta]\)，其中 \(\delta(f)=y^rf'\)，并给出正规单项式基。再令 \(y,z\) 的权重分别为一、\(r\)，证明相应滤过的伴随分次是加权多项式环。
\end{optionalexercise}
\begin{solution}
\(\delta\) 是 \(k[y]\) 的导子；Ore 构造给出 \(zy=yz+y^r\) 和基 \(y^iz^j\)。反向，从这个生成元关系归纳得到 \(zf=fz+y^rf'\)，再用 Ore 泛性质得到回到 \(B_r\) 的同态。两复合固定生成元，所以互逆，基结论随之成立。

赋权后，关系中的 \(zy,yz\) 权重为 \(r+1\)，\(y^r\) 权重为 \(r\)。移动每个逆序对时，交换项保持权重，导子项降低权重。因此权重不超过 \(n\) 的正规单项式恰好张成商上的第 \(n\) 个滤过项，且乘法相容。正规基的独立性使每个商片仍由权重恰为 \(n\) 的单项式组成；最高项关系只要求两个主符号交换，得到 \(\operatorname{gr}B_r\cong k[Y,Z]\)，其中 \(\deg Y=1,\deg Z=r\)。选择权重可将高次修正项变为低次项，但正规基的独立性仍须先由 Ore 构造保证。
\end{solution}

\begin{optionalexercise}\label{b2:ex:20-affine-lie-ideals}
设二维 Lie 代数 \(\mathfrak a=kh\oplus ke\) 满足 \([h,e]=e\)。求其全部 Lie 理想，并求全部一维表示。说明任何一维表示为什么都不能忠实。
\end{optionalexercise}
\begin{solution}
\(0,ke,\mathfrak a\) 显然都是理想。若理想 \(I\) 含有 \(ah+be\) 且 \(a\ne0\)，则 \([ah+be,e]=ae\in I\)，故 \(e\in I\)，继而 \(h\in I\)，于是 \(I=\mathfrak a\)。否则 \(I\subseteq ke\)，只有零空间与 \(ke\) 两种可能。

在一维空间上，每个算子为标量，算子交换子恒为零。表示条件因此给出 \(\rho(e)=\rho\bigl([h,e]\bigr)=0\)，而 \(\rho(h)=\lambda\) 可为任意 \(\lambda\in k\)。反过来，这些赋值满足唯一非零基本括号的关系，双线性保证在全部向量上成立，所以恰为全部一维表示。它们的核均包含 \(ke\ne0\)，因而没有一个忠实。
\end{solution}

\begin{optionalexercise}\label{b2:ex:20-heisenberg-adjoint-matrices}
对\cref{b2:exm:heisenberg-lie-algebra}中的 \(\mathfrak h\)，在有序基 \((x,d,z)\) 下求 \(\operatorname{ad}(x),\operatorname{ad}(d),\operatorname{ad}(z)\) 的矩阵及伴随表示的核。与该例中 \(\mathfrak h\subseteq M_3(k)^-\) 的自然表示比较。
\end{optionalexercise}
\begin{solution}
由 \([x,d]=-z\)、\([d,x]=z\) 及 \(z\) 居中，三矩阵依次为
\[
\begin{pmatrix}0&0&0\\0&0&0\\0&-1&0\end{pmatrix},\qquad
\begin{pmatrix}0&0&0\\0&0&0\\1&0&0\end{pmatrix},\qquad 0.
\]
前两个矩阵线性无关，所以核恰为 \(kz\)。这也由\cref{b2:prop:adjoint-lie-representation}等于中心：若 \(ax+bd+cz\) 居中，与 \(x\) 取括号得 \(bz=0\)，与 \(d\) 取括号得 \(-az=0\)，故 \(a=b=0\)。自然表示把 \(x,d,z\) 分别送到线性无关的 \(E_{23},E_{12},E_{13}\)，所以忠实。伴随表示不忠实，并不妨碍同一个 Lie 代数有别的忠实表示；这些计算在特征二时仍然成立。
\end{solution}

\begin{optionalexercise}\label{b2:ex:20-abelian-lie-jordan-modules}
令 \(\mathfrak a=ku\) 的括号为零，在 \(V=k^3\) 上规定 \(u\) 作用为
\[
J=\begin{pmatrix}0&1&0\\0&0&1\\0&0&0\end{pmatrix}.
\]
将这个表示对应的 \(U(\mathfrak a)\) 模写成一个循环模商，求 \(U(\mathfrak a)\) 对 \(V\) 的作用核，并求与 \(J\) 交换的全部线性算子。
\end{optionalexercise}
\begin{solution}
括号为零，且只有一个给定算子，所以确为表示。由\cref{b2:def:symmetric-algebra,b2:thm:symmetric-polynomial-algebra}及泛包络的关系，\(U(\mathfrak a)=k[t]\)，其中 \(t=j(u)\)。向量 \(v=e_3\) 满足 \(Jv=e_2,J^2v=e_1,J^3v=0\)，故映射 \(k[t]\to V\)、\(f\mapsto f(J)v\) 满射，其核为 \((t^3)\)：除以 \(t^3\) 所得次数小于三的余式，作用于 \(v\) 后的系数正是独立向量 \(e_3,e_2,e_1\) 的系数。因此 \(V\cong k[t]/(t^3)\)，作用核同样为 \((t^3)\)。

若 \(FJ=JF\)，则 \(F\) 由 \(F(v)\) 决定，因为 \(F(J^rv)=J^rF(v)\)。唯一写成 \(F(v)=(aI+bJ+cJ^2)v\) 后，这个等式在整组循环基上成立，故 \(F=aI+bJ+cJ^2\)。反过来，这些算子都与 \(J\) 交换。特别地，Lie 表示 \(ku\to\operatorname{End}_k(V)\) 本身单射，而其泛包络的作用不单射，二者的忠实性不可混同。
\end{solution}

\begin{optionalexercise}\label{b2:ex:20-enveloping-ideal-quotient}
设 \(I\) 是 Lie 代数 \(\mathfrak g\) 的理想，\(\bigl(j(I)\bigr)\) 表示 \(U(\mathfrak g)\) 中由 \(j(I)\) 生成的双边理想。不使用 PBW 定理，证明
\[
U(\mathfrak g/I)\cong U(\mathfrak g)/\bigl(j(I)\bigr).
\]
说明两边的模对应于哪些 \(\mathfrak g\) 表示。
\end{optionalexercise}
\begin{solution}
记 \(\pi:\mathfrak g\to\mathfrak g/I\) 为商 Lie 同态，\(j'\) 为商 Lie 代数到其泛包络的典范映射。\(j'\pi\) 由\cref{b2:thm:enveloping-universal}延拓为 \(U(\mathfrak g)\to U(\mathfrak g/I)\)，并零化 \(j(I)\)，故下降为同态 \(f:U(\mathfrak g)/\bigl(j(I)\bigr)\to U(\mathfrak g/I)\)。

反向令 \(q:U(\mathfrak g)\to U(\mathfrak g)/\bigl(j(I)\bigr)\) 为商同态。因 \(qj\) 零化 \(I\)，公式 \(u+I\mapsto qj(u)\) 不依赖代表元；\eqref{b2:eq:enveloping-canonical-lie-map}又保证它保持括号。再次应用泛性质，得到 \(g:U(\mathfrak g/I)\to U(\mathfrak g)/\bigl(j(I)\bigr)\)。两复合分别固定全部 \(qj(u)\) 与 \(j'(u+I)\)，这些元素连同单位元生成各自的代数，所以两复合为恒等，得到同构。

由\cref{b2:cor:lie-representations-enveloping-modules}，模掉 \(\bigl(j(I)\bigr)\) 正是要求 \(I\) 的每个元素在表示空间上作用为零。这也恰是 Lie 表示下降到 \(\mathfrak g/I\) 的条件。
\end{solution}

\begin{optionalexercise}\label{b2:ex:20-heisenberg-central-character}
对三维 Heisenberg Lie 代数 \(\mathfrak h\)，在特征零时构造一个三维忠实表示，使中心元素 \(z\) 作用非零；说明它为何不违背正文关于非零标量中心作用的有限维表示结论。特征 \(p>0\) 时，对任意 \(\lambda\in k^\times\) 构造一个 \(p\) 维简单表示，使 \(z\) 作用为 \(\lambda I\)，并证明这个 Lie 表示也忠实。
\end{optionalexercise}
\begin{solution}
在 \(k^3\) 上取 \(\rho(x)=E_{23},\rho(d)=E_{12},\rho(z)=E_{13}\)，则 \([E_{12},E_{23}]=E_{13}\)，且 \(E_{13}\) 与另两个矩阵交换。三个矩阵线性无关，故表示忠实；中心作用是非零平方零矩阵，并非非零标量矩阵。\cref{b2:prop:heisenberg-enveloping-weyl,b2:prop:weyl-finite-dimensional-representations}所排除的是特征零中中心作用为指定非零标量的情形，不能由此排除所有有限维 Heisenberg 表示。

特征 \(p\) 时，在 \(k[X]/(X^p)\) 上取 \(\rho(x)f=Xf\)、\(\rho(d)f=\lambda f'\)、\(\rho(z)f=\lambda f\)。关系成立，非零 \(\lambda\) 的缩放又不改变截断 Weyl 表示的子模，故它简单。若 \(a\rho(x)+b\rho(d)+c\rho(z)=0\)，作用于一得到 \(aX+c\lambda=0\)，从而 \(a=c=0\)；再作用于 \(X\) 得 \(b\lambda=0\)，故 \(b=0\)。因此三个作用算子独立，Lie 表示忠实。这里用到 \(p\ge2\)，所以一和 \(X\) 在商中确实独立。
\end{solution}

\begin{optionalexercise}\label{b2:ex:20-inner-sigma-derivation}
设 \(R\) 为含幺环，\(\sigma:R\to R\) 为保幺环自同态，固定 \(c\in R\)，令 \(\delta(a)=ca-\sigma(a)c\)。证明 \(\delta\) 是 \(\sigma\) 导子，并通过变量代换证明 \(R[t;\sigma,\delta]\cong R[u;\sigma]\)。
\end{optionalexercise}
\begin{solution}
直接展开得到
\[
\sigma(a)\delta(b)+\delta(a)b
=\sigma(a)cb-\sigma(a)\sigma(b)c+cab-\sigma(a)cb
=\delta(ab).
\]
在原扩张中置 \(u=t-c\)，则
\[
ua=\sigma(a)t+ca-\sigma(a)c-ca=\sigma(a)u.
\]
所以 Ore 泛性质给出从 \(R[u;\sigma]\) 到原扩张的同态。反向将 \(t\) 送到 \(u+c\)，同一个等式反向验证所需关系。两复合固定 \(R\) 和扩张变量，故互逆。
\end{solution}

\begin{optionalexercise}\label{b2:ex:20-ore-noninjective-sidedness}
令 \(R=k[x]\)、\(\sigma(f)=f(0)\)、\(S=R[t;\sigma]\)。对每个 \(m\ge1\)，分别求
\[
\{q\in S\mid t^mq=0\},\qquad \{q\in S\mid qt^m=0\}.
\]
证明第一个集合是双边理想，辨认相应商环，并说明这两个集合为何不能混同。
\end{optionalexercise}
\begin{solution}
把 \(q\) 唯一写成 \(\sum_i f_i(x)t^i\)。由于 \(\sigma^m(f)=f(0)\)，有 \(t^mq=\sum_i f_i(0)t^{m+i}\)，故它为零当且仅当每个 \(f_i\in xk[x]\)。另一方面，\(qt^m=\sum_i f_i(x)t^{m+i}\)，左正规式的独立性使其为零当且仅当 \(q=0\)。

第一个集合为 \(J=\bigoplus_{i\ge0}xk[x]t^i=xS\)。它显然对右乘封闭；左乘 \(k[x]\) 保持系数被 \(x\) 整除，左乘 \(t\) 则将全部元素杀掉，所以也对左乘封闭。将系数在零处求值、保留 \(t\)，得到满同态 \(S\to k[t]\)，核恰为 \(J\)，故 \(S/J\cong k[t]\)。非单射的 \(\sigma\) 能造成完全不同的左右零化行为；两侧集合的计算依靠左正规式，不可交换乘积次序。
\end{solution}

\begin{optionalexercise}\label{b2:ex:20-q-derivation}
固定 \(q\in k^\times\)，在 \(k[x]\) 上令 \(\sigma\bigl(f(x)\bigr)=f(qx)\)，并按线性定义
\[
\delta(x^n)=(1+q+\cdots+q^{n-1})x^{n-1}\quad(n\ge1),
\qquad\delta(1)=0.
\]
证明这是左 \(\sigma\) 导子，并写出相应 Ore 扩张的 \(tx\)、\(tx^2\) 及正规单项式基。
\end{optionalexercise}
\begin{solution}
记 \([n]_q=1+q+\cdots+q^{n-1}\)，则 \([m+n]_q=[m]_q+q^m[n]_q\)。因此在两个单项式上，\(\delta(x^mx^n)=\sigma(x^m)\delta(x^n)+\delta(x^m)x^n\)。两边对多项式分别线性，故等式对任意乘积成立。由构造定理，
\[
tx=qxt+1,\qquad tx^2=q^2x^2t+(1+q)x,
\]
正规单项式基为 \(x^it^j\)。当 \(q=1\) 时，\([n]_q=n\)，恢复通常求导及 Weyl 关系；所有定义都没有使用除以 \(q-1\)。
\end{solution}

\begin{optionalexercise}\label{b2:ex:20-differential-line}
在微分域 \((K,\delta)\) 上取 \(L=\partial-a\)，写出 \(M_L\) 上的作用。若将基向量 \(e\) 改为 \(f=be\)、\(b\ne0\)，新作用系数是什么？最后在 \(K=k(x)\)、\(\delta=\mathrm{d}/\mathrm{d}x\)、\(a=1/x\) 时，用常数域 \(C=\ker\delta\) 表示全部 \(K\) 内的解。
\end{optionalexercise}
\begin{solution}
模以 \(e=1+DL\) 为 \(K\) 基，满足 \(\partial e=ae\)。因此 \(\partial(ue)=\bigl(\delta(u)+au\bigr)e\)。换基后
\[
\partial f=\bigl(\delta(b)+ab\bigr)e
=\bigl(a+b^{-1}\cdot\delta(b)\bigr)f.
\]
对于最后的方程，\(\delta(y)=y/x\) 等价于 \(\delta(y/x)=0\)，所以所有解为 \(y=cx\)、\(c\in C\)。这同时表明解空间自然只需对常数域封闭；在正特征中，\(C\) 可以大于给定的 \(k\)。
\end{solution}

\begin{optionalexercise}\label{b2:ex:20-differential-companion}
设 \((K,\delta)\) 为微分域，左微分模在一组基 \(e=(e_0,\ldots,e_{m-1})\) 中的坐标作用为 \(\nabla\mathbf v=\delta\mathbf v+A\mathbf v\)。令 \(e'=eP\)、\(P\in\operatorname{GL}_m(K)\)，求新作用矩阵和模同态取值列 \(Y'\) 的变换式。在 \(K=k(x)\)、\(L=\partial^2-x\) 时，取
\[
P=\begin{pmatrix}1&x\\0&1\end{pmatrix},
\]
实际计算这两个变换。
\end{optionalexercise}
\begin{solution}
同一元素满足 \(e\mathbf v=e'\mathbf v'\)，即 \(\mathbf v=P\mathbf v'\)。由 Leibniz 公式，
\[
\delta(P\mathbf v')+AP\mathbf v'
=P\delta\mathbf v'+(\delta P+AP)\mathbf v'.
\]
改为新基坐标，得 \(B=P^{-1}AP+P^{-1}\delta P\)。系数求导项不可丢弃。若 \(f\) 是到函数模的同态，\(Y_i=f(e_i)\)，则 \(Y'=P^{\mathsf T}Y\)。正文的系统 \(\delta Y=A^{\mathsf T}Y\) 因而变为 \(\delta Y'=B^{\mathsf T}Y'\)，这也可直接对 \(P^{\mathsf T}Y\) 求导验证。

对于 \(L=\partial^2-x\)，\cref{b2:prop:differential-companion-coordinates}给出 \(A=\begin{psmallmatrix}0&x\\1&0\end{psmallmatrix}\)。代入可得
\[
B=\begin{pmatrix}-x&x+1-x^2\\1&x\end{pmatrix},
\qquad Y'=\begin{pmatrix}y\\xy+\delta y\end{pmatrix}.
\]
对第二个分量求导，用 \(\delta^2y=xy\)，即得 \(\delta Y'=B^{\mathsf T}Y'\)。这里两个基之间的 \(K\) 线性换坐标并不将 \(\nabla\) 变成 \(K\) 线性算子。
\end{solution}

\begin{optionalexercise}\label{b2:ex:20-diamond-idempotent-rules}
用规则 \(x^2\to x\)、\(yx\to y\) 检查全部非平凡歧义，证明代数 \(k\langle x,y\rangle/(x^2-x,yx-y)\) 有基 \(y^n,xy^n\)、\(n\ge0\)。给出其中两个非零元素的乘积为零的例子。
\end{optionalexercise}
\begin{solution}
两条规则都降低字长。只有 \(xxx\) 与 \(yxx\) 存在重叠：\(xxx\) 无论先改哪一对都化为 \(x\)，\(yxx\) 两种改法都化为 \(y\)。没有非平凡包含歧义，所以\cref{b2:thm:diamond-lemma}适用。不可约字中不能有连续的两个 \(x\)，也不能有 \(y\) 在 \(x\) 左边，故恰为所列形式。正规字基保证 \(y\ne0\)、\(1-x\ne0\)，但关系给出 \(y(1-x)=0\)。
\end{solution}

\begin{optionalexercise}\label{b2:ex:20-diamond-two-failures}
固定 \(a,b,c\in k\)，在 \(k\langle x,y\rangle\) 中取规则 \(x^2\to ax+b\)、\(yx\to cy\)。证明这组规则的全部歧义可解当且仅当 \(c^2-ac-b=0\)。在该条件成立时给出商代数的一组基；条件不成立时，辨认隐藏关系及实际商代数。

再对 \(\alpha,\beta\in k\) 考察规则 \(xyx\to\alpha x\)、\(yx\to\beta\) 的包含歧义，辨认相应商代数，并在非零商的情形给出一组基。
\end{optionalexercise}
\begin{solution}
两条规则降低字长，非平凡重叠只有 \(xxx\) 与 \(yxx\)，没有非平凡包含。前者两支都成为 \((a^2+b)x+ab\)，后者两支分别为 \(c^2y\) 与 \((ac+b)y\)。它们已是正规式，所以歧义可解恰好要求 \(c^2-ac-b=0\)。在该条件下，Diamond 引理给出不可约字基 \(y^n,xy^n\)、\(n\ge0\)：每个 \(x\) 只能位于字首，且至多有一个。

若 \(c^2-ac-b\ne0\)，重叠差强迫 \(y=0\)。将 \(y\) 送到零、\(x\) 送到自身，得到商到 \(k[x]/(x^2-ax-b)\) 的同态；反向保留 \(x\) 的类，利用 \(y=0\)，两复合为恒等。因此实际商正是这个二次多项式商，基为 \(1,x\)。这既计算出额外关系，也表明规则不相容时不能沿用原先预期的无限正规字基。

第二组规则在 \(xyx\) 处，整体改写得到 \(\alpha x\)，改写其中的 \(yx\) 得到 \(\beta x\)。若 \(\alpha=\beta\)，第一条关系由第二条左乘 \(x\) 得到，故商只需关系 \(yx=\beta\)。规则 \(yx\to\beta\) 没有自身包含或重叠，Diamond 引理给出基 \(x^iy^j\)、\(i,j\ge0\)。若 \(\alpha\ne\beta\)，包含差强迫 \(x=0\)，第二条关系随即要求 \(\beta1=0\)。当 \(\beta=0\) 时，正反赋值 \(x\mapsto0\)、保留 \(y\) 给出商 \(k[y]\)，其基为 \(y^j\)；当 \(\beta\ne0\) 时，\(1=0\)，商为零代数。这一包含歧义可以删掉一个生成元，也可以使整个商为零。
\end{solution}

\begin{optionalexercise}\label{b2:ex:20-infinite-groebner-growth}
对 \(B=k\langle x,y\rangle/(x^2-xy)\) 取字长分次。利用\cref{b2:exm:infinite-noncommutative-groebner}求 \(\dim_kB_n\) 及形式 Hilbert 级数 \(\sum_{n\ge0}\dim_k(B_n)t^n\)，并证明 \(B\) 虽与二元多项式环具有相同的分次维数，却有零因子。
\end{optionalexercise}
\begin{solution}
次数 \(n\ge1\) 的正规字中，有一个 \(y^n\)，以及满足 \(a+b=n-1\) 的 \(n\) 个 \(y^axy^b\)，所以维数为 \(n+1\)；零次维数为一。因此级数为
\[
\sum_{n\ge0}(n+1)t^n=\frac1{(1-t)^2},
\]
这里是形式幂级数恒等式。正规字基使 \(x\) 与 \(x-y\) 均非零，但 \(x(x-y)=0\)。所以分次维数不能决定乘法结构，也不能单凭这一数列判断无零因子性。
\end{solution}

\begin{optionalexercise}\label{b2:ex:20-groebner-order-dependence}
在 \(B=k\langle x,y\rangle/(x^2-xy)\) 中，使用\cref{b2:exm:infinite-noncommutative-groebner}关于 \(y>x\) 的正规基，计算 \((y^ax^b)(y^cx^d)\)，其中 \(a,b,c,d\ge0\)。据此求中心，并证明中心为 \(k\) 并不使 \(B\) 成为单环。
\end{optionalexercise}
\begin{solution}
由 \(xy=x^2\)，归纳得 \(x^by^c=x^{b+c}\) 对 \(b\ge1\) 成立。因此
\[
(y^ax^b)(y^cx^d)=
\begin{cases}
y^{a+c}x^d,&b=0,\\
y^ax^{b+c+d},&b\ge1.
\end{cases}
\]
若 \(z=\sum_{b=0}^Nf_b(y)x^b\) 与 \(y\) 交换，且 \(N\ge1,f_N\ne0\)，则 \(zy\) 的 \(x^{N+1}\) 系数为 \(f_N(y)\)，\(yz\) 却没有这一项，违反正规基的独立性。因此 \(z=f_0(y)\)。再要求它与 \(x\) 交换：若 \(f_0\) 的次数 \(n\ge1\)，\(xf_0(y)\) 含非零 \(x^{n+1}\) 项，而 \(f_0(y)x\) 只有 \(x\) 次数一，仍矛盾。故 \(Z(B)=k\)。

双边理想 \((x)\) 非零，因为 \(x\) 是正规基元素；它又是真理想，因为模掉 \(x\) 得到非零商 \(k[y]\)。所以 \(B\) 不是单环。不同字序下的正规式允许不同计算，但所得中心和理想仍是同一个代数的内在对象。
\end{solution}

\begin{optionalexercise}\label{b2:ex:20-differential-associated-module}
在 \(D=K[\partial;\delta]\) 中按 \(\partial\) 次数滤过，\(K\) 在零次。若 \(L\) 为 \(m\) 阶首一算子，证明 \(\operatorname{gr}D\cong K[T]\)，并对商的诱导滤过证明
\[
\operatorname{gr}(D/DL)\cong K[T]/(T^m)
\]
为分次左模。
\end{optionalexercise}
\begin{solution}
关系 \(\partial a-a\partial=\delta(a)\) 降低一次 \(\partial\) 次数，故最高次符号与 \(K\) 交换；Ore 正规式使每个次数商片一维，得到 \(\operatorname{gr}D\cong K[T]\)。对于左理想 \(DL\)，非零 \(QL\) 的次数为 \(\deg Q+m\)，其主符号为 \(\operatorname{symb}(Q)T^m\)，所以该理想的伴随分次恰为 \((T^m)\)。也可直接使用唯一左余式：商模在次数小于 \(m\) 时各新增一个基向量，随后不再新增，且乘 \(T\) 逐次移位，最后为零。于是得到所写的分次模同构。
\end{solution}

\begin{optionalexercise}\label{b2:ex:20-quantum-matrix-fibre}
设 \(\operatorname{char}k\ne2\)，\(a=s^2\ne0\)、\(b\ne0\) 为 \(k\) 中元素。证明
\[
k_{-1}[x,y]/(x^2-a,y^2-b)\cong M_2(k).
\]
\end{optionalexercise}
\begin{solution}
由量子平面的正规单项式基及平方关系，商由 \(1,x,y,xy\) 线性生成。取
\[
X=\begin{pmatrix}s&0\\0&-s\end{pmatrix},\qquad
Y=\begin{pmatrix}0&b\\1&0\end{pmatrix}.
\]
它们满足 \(YX=-XY\)、\(X^2=aI\)、\(Y^2=bI\)，所以给出到矩阵代数的同态。\(I,X\) 线性独立并生成对角矩阵；\(Y,XY\) 的两个非对角位置组成的系数矩阵行列式为 \(-2sb\ne0\)，所以生成全部非对角矩阵。因此像为四维矩阵代数，来源又至多四维，所得满同态必为同构。
\end{solution}

\begin{optionalexercise}\label{b2:ex:20-ore-right-normal-form}
设 \(R\) 为含幺环，\(\sigma\) 为环自同构，\(\delta\) 为左 \(\sigma\) 导子。置 \(\tau=\sigma^{-1}\)、\(\varepsilon=-\delta\sigma^{-1}\)。把 \(\tau,\varepsilon\) 视为 \(R^{\mathrm{op}}\) 上的映射，证明 \(\varepsilon\) 是左 \(\tau\) 导子，并证明
\[
R[t;\sigma,\delta]^{\mathrm{op}}\cong R^{\mathrm{op}}[u;\tau,\varepsilon].
\]
同构在底环上为恒等，将 \(u\) 送到反环中的 \(t\)。
\end{optionalexercise}
\begin{solution}
在原环记乘积，扭曲 Leibniz 公式给出
\[
\varepsilon(ba)=b\varepsilon(a)+\varepsilon(b)\tau(a).
\]
在反环中它恰好是 \(\varepsilon(a\cdot_{\mathrm{op}}b)=\tau(a)\cdot_{\mathrm{op}}\varepsilon(b)+\varepsilon(a)\cdot_{\mathrm{op}}b\)，故是所需导子。\(\tau\) 也是反环自同构。

在原 Ore 扩张中，\(at=t\tau(a)+\varepsilon(a)\)。反过来读乘积即为 \(t\cdot_{\mathrm{op}}a=\tau(a)\cdot_{\mathrm{op}}t+\varepsilon(a)\)。Ore 泛性质给出题设同态；它的像包含底环和 \(t\)，故满射。由\cref{b2:prop:ore-right-normal-form}，原扩张以 \(t^i\) 为右 \(R\) 模基，因而其反环以同一组元素为左 \(R^{\mathrm{op}}\) 模基。这个同态将 Ore 左正规基逐项送到该基，所以单射。扭曲导子中的负号和逆自同构由反转乘法次序产生。
\end{solution}

